\def\jobname{lansoprazole}
\documentclass[11pt,a4paper,notitlepage]{report}
\usepackage[margin=2.2cm]{geometry}
\usepackage{iftex}
\ifPDFTeX
  \usepackage[T2A]{fontenc}
  \usepackage[utf8]{inputenc}
  \usepackage[english,russian]{babel}
\else
  \usepackage{fontspec}
  \setmainfont{DejaVu Serif}
  \setsansfont{DejaVu Sans}
  \setmonofont{DejaVu Sans Mono}
  \usepackage[english,russian]{babel}
\fi
\usepackage{amsmath,amssymb}
\usepackage{array,booktabs,longtable,tabularx}
\usepackage{etoolbox}
\makeatletter
\IfPackageAtLeastTF{longtable}{2025-12-23}{}{
  \patchcmd{\LT@output}
    {\unvbox\z@\copy\LT@foot\vss}
    {\unvbox\z@\copy\LT@foot\vskip\z@\@plus1fil\@minus\normalbaselineskip}
    {}{\PackageError{lansoprazole-review}{Longtable final-page fix failed}{Check the installed longtable version.}}
  \patchcmd{\LT@output}
    {\unvbox\@cclv\copy\LT@foot\vss}
    {\unvbox\@cclv\copy\LT@foot\vskip\z@\@plus1fil\@minus\normalbaselineskip}
    {}{\PackageError{lansoprazole-review}{Longtable page-break fix failed}{Check the installed longtable version.}}
}
\renewcommand{\@tocrmarg}{2.55em plus 1fil}
\makeatother
\patchcmd{\subsubsection}{\bfseries}{\bfseries\raggedright}
  {}{\PackageError{lansoprazole-review}{Subsubsection alignment fix failed}{Check the document class.}}
\usepackage{enumitem}
\usepackage{xcolor}
\usepackage{tikz}
\usetikzlibrary{arrows.meta}
\usepackage{microtype}
\usepackage{xurl}
\usepackage{hyperref}
\hypersetup{unicode=true,colorlinks=true,linkcolor=blue!45!black,citecolor=blue!45!black,urlcolor=blue!45!black,pdftitle={Лансопразол: Междисциплинарная монография},pdfauthor={}}
\setcounter{tocdepth}{2}
\setcounter{secnumdepth}{3}
\setlength{\parindent}{0pt}
\setlength{\parskip}{5pt plus 1pt minus 1pt}
\setlength{\emergencystretch}{3em}
\setlist{nosep,leftmargin=*,topsep=0.4\baselineskip}
\renewcommand{\arraystretch}{1.18}
\newcolumntype{P}[1]{>{\raggedright\arraybackslash}p{#1}}
\newcommand{\Hp}{\textit{H. pylori}}
\newcommand{\src}[1]{\unskip\space\cite{#1}}
\title{\textbf{Лансопразол: молекулярные механизмы действия и стратегии междисциплинарной гастропротекции}\\[0.4em]\large Междисциплинарная монография\\}
\author{Вербицкий С.В., https://orcid.org/0009-0008-5706-0174}
\date{}

\begin{document}
\pagenumbering{roman}
\maketitle

\begin{abstract}
Научное издание представляет собой междисциплинарную монографию, посвященную системному клинико-фармакологическому анализу лансопразола. В работе систематизированы данные регуляторных документов, профильных клинических рекомендаций уровня 2024--2026 гг. и результаты рандомизированных клинических исследований. В монографии подробно описаны физико-химические параметры молекулы, константы диссоциации ($pK_a$), фазы протонирования и ковалентной блокады H$^+$/K$^+$-АТФазы, а также механизмы системной pH-селективности пролекарства. На основе принципов доказательной медицины разработаны алгоритмы гастропротекции для коморбидных пациентов в смежных областях интернальной медицины: кардиологии (при сочетанной антитромботической терапии), ревматологии (при НПВП-гастропатиях), аллергологии (при эозинофильном эзофагите по критериям ACG 2025), пульмонологии и оториноларингологии. Заключительный раздел монографии содержит структурированную программу перспективных научно-практических исследований для верификации молекулярно-генетических парадоксов молекулы. Издание предназначено для клинических фармакологов, гастроэнтерологов, терапевтов, научных сотрудников и врачей смежных специальностей.
\end{abstract}

\clearpage
\tableofcontents
\clearpage
\pagenumbering{arabic}

\chapter{Молекулярная биология, фармакокинетика и кинетика мишени}
\label{chap:molecular}

\section*{Введение к главе}
\addcontentsline{toc}{section}{Введение к главе}
Фармакологическое действие лансопразола определяется последовательностью доставки, кислотной активации и взаимодействия с мишенью, а не одним свойством молекулы. Липофильность исходного соединения необходимо рассматривать совместно со степенью ионизации: мембранопроницаемая фракция достигает секреторного канальца, тогда как протонирование способствует локальному удержанию и последующему превращению. Значения $pK_a$ характеризуют соответствующие кислотно-основные равновесия, но не задают универсальный порог клинической эффективности. Разграничение накопления и химической активации объясняет, почему преждевременное превращение препарата в просвете желудка неэквивалентно образованию реакционноспособного продукта у люминальной поверхности помпы.\src{molchem,molnagaya90,pharmalabel}

Ковалентная блокада желудочной H$^+$/K$^+$-АТФазы связана с образованием дисульфидного аддукта с доступными цистеинами и должна рассматриваться в контексте каталитического цикла, мембранной локализации и функционального состояния фермента. Продолжительность подавления секреции поэтому не совпадает непосредственно с плазменным периодом полувыведения; вместе с тем ковалентное связывание не означает одновременной блокады всех помп или химической невосстановимости аддукта. Для клинического фармаколога это различие принципиально при интерпретации отсроченного эффекта, остаточной кислотопродукции и восстановления секреции после прекращения приема.\src{molstructure,molcys,molrestore,uslabel}

Практические задачи первой главы состоят в объяснении требований к гастрорезистентной форме, связи приема с секреторной активностью и ограничений переноса плазменной экспозиции на эффект в ткани. Нарушение доставки, недостаточная экспозиция и недоступность части мишеней представляют разные причины неполной кислотосупрессии; их нельзя объединять в недифференцированное понятие «слабого ИПП». Поэтому параметры $pK_a$, липофильность и реакционная способность рассматриваются как взаимосвязанные механистические характеристики, а не самостоятельные основания для выбора дозы или заявления о превосходстве над другой молекулой. Этот подход задает основу последующего анализа клинических исходов и персонализации лечения.\src{molchem,molnagaya90,molp450,acggerd,cpic}

\section{Область охвата, метод и навигация}
\label{sec:reading-guide}

Объектом настоящего исследования является молекула ингибитора протонной помпы (ИПП) лансопразола. Предметом исследования выступают ее физико-химические свойства, молекулярно-генетические механизмы действия и профили междисциплинарной безопасности. Методологическая база монографии включает системный анализ нормативно-правовых и регистрационных документов (FDA, EMA), национальных и международных клинических рекомендаций (ACG, AGA, ESGE, DGVS, NICE, JSGE, JSHR), а также данных репрезентативных контролируемых исследований уровня 2024--2026 годов.\src{uslabel,methodema2010,acghp,agapcab,esge2026,dgvs,nice,jgerd,jhp} В издании реализован принцип разграничения класс-специфических эффектов ИПП, изомерных характеристик пролекарства и индивидуальных стратегий дозирования с учетом генетического полиморфизма.\src{molchem,molcys,cpic}

\textbf{Навигационная карта издания для специалистов смежных медицинских направлений:}
\begin{itemize}[itemsep=0.45\baselineskip]
\item \textbf{Кардиология и антитромботическая терапия:} особенности сочетанного применения лансопразола на фоне двойной антитромбоцитарной терапии (ДАТТ: клопидогрел + аспирин), оценка и минимизация рисков MACE интегрированы в главу~\ref{chap:interdisciplinary} (раздел~\ref{sec:clopidogrel} и подраздел~\ref{sec:dapt-protocol}).\src{cardioacc2025,cardioplavix2025,cardiocpic2022}

\item \textbf{Ревматология и неврология:} стратегии длительной гастропротекции при приеме ульцерогенных агентов (НПВП, аспирин) и фармакодинамический синергизм комбинаций с фиксированными дозами разобраны в главе~\ref{chap:molecular} (подразделы~\ref{sec:fixed-combinations} и~\ref{sec:ph-retention-synergy}) и главе~\ref{chap:interdisciplinary} (раздел~\ref{sec:nsaid}).\src{aspirintrial,aspirin15,nsaid15}

\item \textbf{Аллергология и иммунология:} терапия эозинофильного эзофагита по критериям ACG 2025, ее молекулярный базис и экспериментальные данные о сигнальной оси Keap1--Nrf2--HO-1 изложены в главе~\ref{chap:interdisciplinary} (раздел~\ref{sec:eoe}).\src{eoe,eoehighlights2025,molnrf2}

\item \textbf{Оториноларингология и пульмонологическая практика:} клиническая оценка эффективности лансопразола при внепищеводных рефлюкс-ассоциированных масках (хронический ларингит, кашель, бронхиальная астма) приведена в главе~\ref{chap:gastroenterology} (раздел~\ref{sec:gerd}) и главе~\ref{chap:interdisciplinary} (раздел~\ref{sec:extraesophageal}).\src{acggerd,toppits,asthma}

\item \textbf{Клиническая фармакогенетика и безопасность длительного лечения:} влияние полиморфизмов CYP2C19, риски феноконверсии, компенсаторная ECL-клеточная гиперплазия и пошаговый алгоритм безопасного депрескриптинга на основе принципов AGA систематизированы в главе~\ref{chap:personalization} (разделы~\ref{sec:pharmacogenetics}, \ref{sec:safety}, \ref{sec:withdrawal} и~\ref{sec:deprescribing-algorithm}).\src{cpic,withdrawsheng2020,agadepres,withdrawfarrell2017}

\item \textbf{Программа перспективных исследований (фтизиатрия):} экспериментальный анализ активности метаболита лансопразол-сульфида в отношении QcrB микобактерий и дорожная карта для будущих ученых представлены в заключении (глава~\ref{chap:conclusion}, раздел~\ref{sec:research-agenda}).\src{moltb,moltbstructure}
\end{itemize}

\subsection{Версии документов и региональная применимость}
\begin{longtable}{P{0.20\textwidth}P{0.73\textwidth}}
\toprule
Регион / тема & Документальная основа монографии \\
\midrule
\endhead
США & ACG по ГЭРБ, 2022; ASGE по ГЭРБ, 2025; AGA по персонализированному ведению и отмене ИПП, 2022, и пищеводу Баррета, 2025; AGA по P-CAB, 2024; ACG по \Hp, 2024, и эозинофильному эзофагиту, 2025; SCCM/ASHP, 2024.\src{acggerd,asge2025,agagerd,agadepres,barrett2025,agapcab,acghp,eoe,sccm} \\
Европа & Maastricht VI/Florence, 2022; DGVS S2k по ГЭРБ, март 2023; NICE CG184, обновление 18.10.2019; BSG по функциональной диспепсии, 2022; Lyon 2.0, онлайн 2023 / выпуск 2024; ESGE по язвенному кровотечению, обновление 2026. Великобритания рассматривается отдельно от ЕС.\src{maastricht,dgvs,nice,bsgfd,lyon,esge2026} \\
Япония & JSGE по ГЭРБ: редакция 2021, англоязычная публикация 2022; JSHR по \Hp: редакция 2024, англоязычная публикация 28.03.2026; новое руководство JSGE по язвенной болезни, опубликованное 24.08.2026.\src{jgerd,jhp,jpud2026} \\
Республика Корея & Фокусированное обновление Seoul Consensus 2025 по кислотосупрессии, опубликованное также в выпуске 2026; руководство по \Hp «2025 Revised Edition», публикация 2026.\src{seoul,khp} \\
Китай & Полнотекстовое национальное руководство по эрадикации, публикация 20.12.2022; новое комплексное руководство от 01.07.2026: подтверждены выход и аннотация, но не все конкретные схемы полного текста.\src{china2022,china2026} \\
Индия & Bhubaneswar Consensus ISG по \Hp, 2021; экспертный консенсус Society for Clinical Gastroenterology по кислотозависимым заболеваниям, 05.03.2026. Различаются разработчики и методология.\src{indiahp,india2026} \\
Юго-Восточная Азия & Региональный консенсус по рациональному назначению и отмене ИПП, 2026.\src{seadepres2026} \\
РФ & Рекомендации российских обществ по ГЭРБ и язвенной болезни, 2024; КР «Язвенная болезнь», 277\_2, и «Гастрит и дуоденит», 708\_2, 2024, в доступных републикациях утвержденного текста.\src{rfgerd,rfpud,rfpudmin,rfgastrit} \\
\bottomrule
\end{longtable}

\textbf{Применение в Российской Федерации.} Перед использованием положения в клиническом решении необходимо сверить редакцию рекомендаций в Рубрикаторе Минздрава и инструкцию конкретного торгового наименования в ГРЛС. Публикация профессионального общества и иностранная инструкция не заменяют этой проверки; юридическую актуальность каждой регистрации устанавливают по официальной карточке. Возрастные допуски и дозовые режимы зарубежных продуктов на российские препараты автоматически не переносят.\src{registry,grls}

\textbf{Применение региональных документов Азии.} Рекомендации Японии, Республики Корея, Китая и Индии, а также региональный консенсус Юго-Восточной Азии анализируют раздельно по разработчику, версии и клинической задаче. Единого обязательного для всех стран режима не предполагается. При использовании китайских схем следует различать полнотекстовый документ 2022 года и руководство июля 2026 года: конкретные обновленные режимы устанавливают по полному тексту последнего, а не по его аннотации.\src{china2026}

\subsection{Ключевые выводы и принципы интерпретации}
\begin{enumerate}
\item Лансопразол --- эффективный ИПП для подавления кислотопродукции; он не устраняет все механизмы рефлюкса и не является универсальным средством от боли в животе.\src{uslabel,acggerd}
\item В европейских дозовых классификациях \textbf{30 мг/сут} --- стандартная доза, \textbf{15 мг/сут} --- низкая. Однако американская инструкция предусматривает 15 мг/сут и для ряда лечебных показаний; «стандартная доза класса» не заменяет таблицу конкретной инструкции.\src{nice,uslabel}
\item Прием до еды и соблюдение режима принципиальны. Симптомы на фоне ИПП сначала требуют проверки диагноза и приема, а не бесконечного повышения дозы.\src{acggerd,agagerd}
\item При \Hp\ лансопразол --- кислотосупрессивный компонент комбинации, а не самостоятельная эрадикационная терапия. Современный выбор антибиотиков определяется устойчивостью и предшествующим лечением.\src{acghp,maastricht}
\item В США эмпирическая кларитромициновая тройная терапия больше не является стандартным выбором. Наличие старой схемы в инструкции не означает, что она остается предпочтительной.\src{acghp}
\item При тяжелом эрозивном эзофагите калий-конкурентные блокаторы кислотопродукции могут иметь преимущество; это не делает лансопразол неэффективным при всех формах ГЭРБ.\src{vontrial,seoul}
\item Длительное лечение оправдано при сохраняющемся показании. Наблюдательные ассоциации с осложнениями нельзя автоматически считать доказанным вредом, но реальные лекарственные реакции требуют действий.\src{agadepres,uslabel,kidney}
\item Детям, беременным и кормящим препарат выбирают отдельно; применение у младенцев для обычного плача и срыгивания не обосновано.\src{infant,uktis,lactmed}
\item При НПВП и аспирине гастропротекция определяется риском. Доказательства профилактики язвы или язвенного осложнения не равнозначны доказательству защиты всего кишечника либо снижению смертности.\src{nice,aspirintrial,aspirin15,nsaid15}
\item Эмпирический курс лансопразола не показал преимущества при персистирующих глоточных симптомах в TOPPITS; добавление препарата не улучшило контроль астмы у исследованных детей. Эти результаты не подменяют отдельную оценку показаний к лечению ГЭРБ.\src{toppits,asthma}
\item ИПП входят в варианты лечения ЭоЭ по ACG 2025, но клинический ответ не заменяет эндоскопической и гистологической оценки. Индукция HO-1 и изменение IL-8 в других моделях не доказывают механизм ремиссии ЭоЭ.\src{eoe,molnrf2,molil8}
\item В ОРИТ стресс-профилактика зависит от риска и не тождественна лечению состоявшегося язвенного кровотечения; универсальное преимущество лансопразола в этой нише не установлено.\src{sccm,esge2026}
\item Данные о LPZS и микобактериальном дыхательном комплексе относятся к доклиническому направлению. Они не обосновывают перенос обычных доз лансопразола в лечение туберкулеза.\src{moltb,moltbstructure}
\end{enumerate}

\section{Что представляет собой препарат}
\label{sec:drug-profile}

\subsection{Молекула, механизм и место в классе}
Международное непатентованное наименование --- \textit{lansoprazole}; код ATХ --- A02BC03. Это замещенный бензимидазол, ингибитор протонной помпы (ИПП; также употребляется «ингибитор протонного насоса»). После всасывания и накопления в кислой среде секреторных канальцев париетальных клеток активируется и блокирует конечный этап секреции кислоты --- H$^+$/K$^+$-АТФазу. Подавляет базальную и стимулированную секрецию. Восстановление кислотопродукции зависит от восстановления функциональных помп, поэтому действие продолжительнее присутствия вещества в плазме.\src{pharmalabel,hpra}

Антациды нейтрализуют уже имеющуюся кислоту, H$_2$-блокаторы воздействуют на гистаминовую стимуляцию, а ИПП --- на конечный секреторный механизм. Из этого следует практическое различие: лансопразол лучше рассматривать как курсовое или поддерживающее лечение по показанию, а не как средство немедленной помощи при единичном приступе.\src{pharmalabel,otc}

Подробная химическая и клеточная последовательность действия, включая границы экспериментальных данных, разобрана в разделе~\ref{sec:molecular}.

\subsection{Лансопразол и декслансопразол --- не одно и то же}
Лансопразол является рацематом. Декслансопразол --- его R-энантиомер, применяемый в препаратах с иной системой высвобождения и собственными дозировками. Декслансопразол может приниматься независимо от еды согласно американской инструкции; это правило нельзя переносить на обычный лансопразол. Замена «миллиграмм на миллиграмм» некорректна.\src{dexlabel,pharmalabel}

\subsection{Фармакокинетика и ее клинический смысл}
\begin{longtable}{P{0.29\textwidth}P{0.64\textwidth}}
\toprule
Характеристика & Практическое значение \\
\midrule
\endhead
Кислотолабильность & Необходима гастрорезистентная защита гранул. Разжевывание или растирание разрушает механизм доставки. \\
Всасывание & Пик концентрации обычно примерно через 1,5--2 часа; биодоступность высокая, около 80--90\% и более в зависимости от условий исследования. Еда может существенно уменьшать экспозицию. \\
Связывание с белками & Около 97\%; обычный гемодиализ не является эффективным способом удаления препарата. \\
Плазменный период полувыведения & Обычно около 1--2 часов. Это не интервал дозирования: кислотосупрессия сохраняется значительно дольше. \\
Метаболизм & Преимущественно печеночный, с участием CYP2C19 и CYP3A4. Генотип и взаимодействия влияют на экспозицию. \\
Выведение & Преимущественно в виде метаболитов. Почечная недостаточность сама по себе обычно не требует изменения дозы; при печеночной недостаточности ситуация иная. \\
\bottomrule
\end{longtable}
Источники фармакокинетических сведений: инструкции и руководство CPIC. Приведенные диапазоны не являются индивидуальным прогнозом концентрации или эффективности.\src{pharmalabel,hpra,cpic}

\section{Молекулярные механизмы действия: от пролекарства к клеточному ответу}
\label{sec:molecular}

\subsection{Как разграничены механизмы и доказательства}
В этом разделе разделены четыре предмета: химическое превращение лансопразола; устройство и блокада его основной мишени; вторичные последствия изменения кислотности; экспериментальные эффекты вне этой мишени. Результат на очищенном белке показывает биохимическую возможность, опыт на клетках --- действие в конкретной модели, а исследование у человека --- клиническую применимость лишь в пределах своего дизайна. Ни молекулярный докинг, ни изменение экспрессии отдельного гена сами по себе не устанавливают нового показания. Перечень ниже охватывает основные механистические направления, но не претендует на каталог всех опубликованных скрининговых взаимодействий.

\subsection{Химическое устройство и три разных значения слова «активация»}
Лансопразол имеет формулу $\mathrm{C}_{16}\mathrm{H}_{14}\mathrm{F}_{3}\mathrm{N}_{3}\mathrm{O}_{2}\mathrm{S}$ и молекулярную массу около 369,36 г/моль. В молекуле соединены бензимидазольный фрагмент, сульфоксидная связь и замещенное пиридиновое кольцо с метильной и трифторэтоксигруппой. Атом серы сульфоксида является стереогенным центром; обычный препарат содержит R- и S-формы.\src{uslabel}

\textbf{Активация клетки} означает включение секреции и перемещение помп. \textbf{Активация пролекарства} --- кислотокатализируемую химическую перестройку. \textbf{Метаболизм в печени} --- отдельные ферментативные превращения, определяющие экспозицию. Эти процессы происходят в разных компартментах: CYP2C19 не нужен для образования антисекреторного ингибитора на поверхности помпы. Это принципиальное отличие от пролекарств, которым необходима печеночная биоактивация.\src{molnagaya90,molp450}

\subsection{Мишень: строение желудочной H+/K+-АТФазы}
Желудочная H$^+$/K$^+$-АТФаза --- мембранный комплекс P-типа, состоящий из каталитической $\alpha$-субъединицы ATP4A и вспомогательной $\beta$-субъединицы ATP4B. У $\alpha$-субъединицы десять трансмембранных спиралей и цитоплазматические домены: N связывает нуклеотид, P содержит участок фосфорилирования, A обеспечивает сопряженные конформационные перестройки. Гликозилированная $\beta$-субъединица участвует в стабилизации и организации комплекса. Обозначение P-типа связано с образованием фосфорилированного промежуточного состояния белка, а не с названием препарата.\src{molstructure}

Помпа переносит протоны из цитозоля в секреторный каналец в обмен на калий, используя энергию ATP. \textbf{Каналец --- впячивание апикальной мембраны, сообщающееся с просветом железы, а не закисленный цитозоль клетки.} Лансопразол действует с люминальной стороны мембраны; его основное действие не состоит в конкуренции с ATP за цитоплазматический нуклеотидный карман.\src{molstructure,molcys}

\subsection{Каталитический цикл: что именно прерывает ингибитор}
Схематически цикл можно представить так:
\[
E_1 \;\longrightarrow\; E_1P
\;\longrightarrow\; E_2P
\;\longrightarrow\; E_2(\mathrm{K}^+)
\;\longrightarrow\; E_1.
\]
В состоянии $E_1$ транспортные участки доступны из цитозоля. Фосфорилирование консервативного аспартата за счет MgATP сопрягается с переходом к $E_2P$, открытием люминального пути и отдачей протона. Связывание K$^+$ снаружи способствует дефосфорилированию и окклюзии катиона --- его временному запиранию внутри белка; последующий переход возвращает доступ к цитозолю. Обозначения здесь упрощены и не включают все промежуточные состояния.\src{molk}

Структурное исследование 2019 года выявило один окклюдированный K$^+$ и поддержало обмен $1\mathrm{H}^+:1\mathrm{K}^+$ на гидролиз одной ATP. Такой электронейтральный цикл согласуется с поддержанием градиента порядка шести единиц pH. Поэтому некорректно механически переносить на эту помпу стехиометрию Na$^+$/K$^+$-АТФазы или утверждать, что она всегда переносит два протона на ATP.\src{molk}

\subsection{Ион-связывающий карман: лизин, глутаматы и микросреда}
На атомном уровне важны не только «наличие канала» и ATP, но и протонирование боковых цепей. Lys791 и Glu820 участвуют в конформационно-зависимых электростатических взаимодействиях; вместе с другими остатками они определяют геометрию катион-связывающего участка. В молекулярном моделировании и мутационных опытах исследованы роли Glu343, Glu795, Glu820 и Asp824. Одновременное протонирование Glu343 и Glu795 поддерживало селективность к K$^+$; протонирование Glu936, напротив, повышало чувствительность к Na$^+$. Мутант Glu936Gln демонстрировал слабую Na$^+$-активируемую ATPазную активность. Это показывает, что важны конкретные состояния ионизации, а не просто суммарная кислотность среды.\src{molselect}

Это описание \emph{физиологии мишени}, а не доказательство ковалентного присоединения лансопразола к перечисленным глутаматам. Мишени тиоловой реакции препарата --- цистеины; ион-координирующие остатки и участки лекарственного связывания нельзя смешивать. Нумерация остатков далее соответствует используемой в цитируемых структурных и биохимических работах; перед переносом на конкретную видовую последовательность требуется выравнивание.\src{molselect,molcys}

\subsection{Рецепторные каскады: разные сигналы сходятся на одной помпе}
\textbf{Гистаминовое звено.} Базолатеральный H$_2$-рецептор через G$_s$ стимулирует аденилатциклазу: из ATP образуется cAMP, активирующий PKA. Последующие фосфорилирования меняют организацию секреторного аппарата. Но формула «H$_2$ действует исключительно через cAMP» чрезмерно категорична: в мембранах париетальных клеток собак и трансфицированных клетках Wang и соавт. обнаружили также H$_2$-зависимую активацию PLC через отдельный GTP-зависимый путь. Блокирование G$_s\alpha$ подавляло аденилатциклазную, но не PLC-ветвь. Это уточнение экспериментальной модели, а не отмена центральной роли cAMP/PKA.\src{molh2signal,molezrin}

\textbf{Холинергическое звено.} Упрощенная последовательность M$_3$--G$_q$--PLC включает гидролиз мембранного фосфатидилинозитол-4,5-бисфосфата (PIP$_2$):
\[
\mathrm{PIP}_2\xrightarrow{\mathrm{PLC}}\mathrm{IP}_3+\mathrm{DAG}.
\]
IP$_3$ способствует высвобождению Ca$^{2+}$ из внутриклеточного депо; DAG участвует в регуляции PKC. Не следует считать все PKC-зависимые события стимулирующими и сводить сложную секреторную реакцию к одной линейной стрелке. Эксперименты с фармакологическими блокаторами требуют контроля их неспецифических эффектов.\src{molprobes}

Удаление M$_3$-рецептора у мышей уменьшало базальную секрецию и ответы на несколько стимулов, включая гистамин и гастрин. Это показывает взаимозависимость рецепторных систем, а не прямое связывание лансопразола с M$_3$. Пожизненный нокаут также меняет развитие и адаптацию ткани и не воспроизводит краткий прием антагониста.\src{molm3}

\subsection{Гастрин, ECL-клетки и соматостатин: межклеточный контур}
Гастрин действует через CCK$_2$-рецепторы; существенная часть его кислотостимулирующего действия проходит через ECL-клетки, выделяющие гистамин. Гистамин затем активирует H$_2$-рецепторы париетальных клеток. CCK$_2$-рецепторы описаны и на самих париетальных клетках, однако вклад прямого гастринового сигнала зависит от модели; его нельзя без оговорки считать единственным или главным путем у человека. Опыты с CCK-рецепторными нокаутами подтверждают также влияние CCK$_2$ на максимальный секреторный ответ и чувствительность к карбахолу.\src{molcck}

Соматостатин через SSTR$_2$ подавляет гастриновые клетки, ECL-клетки и париетальные клетки, создавая несколько уровней отрицательной регуляции. В SSTR$_2$-дефицитных мышах Zhao и соавт. наблюдали компенсаторное усиление галанинового тормозного пути и перестройку стимуляции париетальных клеток. Поэтому отсутствие ожидаемой гиперсекреции у нокаута не доказывает неважность удаленного рецептора.\src{molsst}

\textbf{Связь с препаратом.} Гипергастринемия на фоне подавления кислотности не означает, что лансопразол является агонистом гастринового рецептора. Это ответ регулируемой системы на изменение конечного продукта секреции. Межклеточный контур объясняет, почему концентрация одного гормона не позволяет непосредственно вычислить число заблокированных помп.\src{uslabel,molsst}

\subsection{Почему помпа должна быть доступна в секретирующей клетке}
Гистамин запускает cAMP/PKA-зависимую стимуляцию париетальной клетки. Существенный этап --- перемещение H$^+$/K$^+$-АТФазы из тубуловезикул к апикальной мембране. В первичных культурах париетальных клеток показано, что фосфорилирование эзрина по Ser66 позволяет ему взаимодействовать с ACAP4; система ACAP4--ARF6 и мембранно-цитоскелетные перестройки участвуют в организации секреторной поверхности. Нефосфорилируемый вариант эзрина нарушал локализацию ACAP4 и секреторный ответ.\src{molezrin}

Лансопразол не требуется считать прямым ингибитором PKA, эзрина или ACAP4: эти белки объясняют \emph{доступность его мишени}. Более ранние опыты с изолированными париетальными клетками собак показали подавление секреции, вызванной гистамином, карбахолом и дибутирил-cAMP. Действие после стимуляции разными путями согласуется с блокадой общего конечного звена, а не одного рецептора. Значения активности в таких опытах не являются сравнением клинических доз у человека.\src{molnagaya90}

\subsection{Калиевый рециклинг: помпа не работает изолированно}
Выведенный в цитозоль K$^+$ должен снова становиться доступным на люминальной стороне. Канал KCNQ1 с регуляторной субъединицей KCNE2 обеспечивает важную часть такого апикального калиевого рециклинга. В работе Lambrecht и соавт. анализ экспрессии в изолированных клетках, подавление секреции блокаторами KCNQ и совместная локализация канала с H$^+$/K$^+$-АТФазой в канальце поддержали его функциональную связь с помпой. Канал поставляет ион для обмена, тогда как помпа осуществляет ATP-зависимый транспорт.\src{molkchannel}

Следовательно, наличие каталитически исправной ATPазы само по себе недостаточно для полноценной секреции при нарушенном ионном обеспечении. Но этот факт не доказывает, что антисекреторное действие лансопразола опосредовано прямой блокадой KCNQ1--KCNE2. Необходимо различать компоненты одной физиологической системы и непосредственную лекарственную мишень.\src{molkchannel,molcys}

\subsection{Откуда берутся протоны: карбоангидраза и обмен бикарбоната}
Желудочная помпа не синтезирует водород и не переносит готовую молекулу HCl через мембрану. Один из основных кислотно-основных процессов в клетке описывается суммарным равновесием:
\[
\mathrm{CO}_2+\mathrm{H_2O}
\overset{\text{карбоангидраза}}{\rightleftharpoons}
\mathrm{H}^{+}+\mathrm{HCO}_3^{-}.
\]
Карбоангидраза ускоряет установление равновесия, а направленный транспорт продуктов поддерживает секреторный процесс. Иммунозолотое исследование париетальных клеток крысы выявило карбоангидразу II преимущественно в цитоплазме и нуклеоплазме, а не исключительно в апикальной мембране. Это отдельный фермент, не идентичный ATP4A.\src{molca}

Обменник \textbf{AE2/SLC4A2} на базолатеральной стороне обеспечивает выход HCO$_3^{-}$ и вход Cl$^{-}$. Так согласуются удаление основания, поддержание внутриклеточного pH и поступление хлорида для секреции. Полное удаление Ae2 у мышей сопровождалось ахлоргидрией и нарушением развития канальцев и тубуловезикул. Данный опыт устанавливает значение AE2 для системы, но не доказывает его прямого ингибирования лансопразолом.\src{molae2}

Выход бикарбоната к крови и выделение протона в просвет --- разные стороны векторного транспорта. Из этой схемы нельзя выводить, что обычный прием ИПП обязательно вызывает системный ацидоз или алкалоз: она не учитывает буферные системы, другие ткани и компенсаторные потоки. Здесь представлена локальная модель секреции, а не расчет кислотно-основного состояния пациента.\src{molae2}

\subsection{Путь хлорида: почему нельзя назначить один канал «виновником всего»}
Для образования кислого хлоридсодержащего секрета Cl$^{-}$ должен пройти апикальный путь отдельно от обмена H$^+$/K$^+$. Удаление Slc26a9 у мышей приводило к нарушениям секреторной функции и утрате тубуловезикул. Это поддерживает участие SLC26A9 в организации секреции, но не доказывает, что весь апикальный хлоридный поток человека проходит через один белок.\src{molslc26}

Для ClC-2 также получены нокаутные данные: снижались стимулированная секреция, число париетальных клеток, количество H$^+$/K$^+$-АТФазы и тубуловезикул. Авторы прямо отмечают, что совокупность этих изменений не позволяет из такого опыта отдельно вычислить вклад ClC-2 в непосредственный выход Cl$^{-}$. Снижение секреции может быть следствием и изменения транспорта, и перестройки клетки. Поэтому на схеме ниже указан \emph{апикальный путь Cl$^{-}$}, а не объявлен единственный окончательно установленный канал.\src{molclc2}

\subsection{Интегральная схема париетальной клетки}
На рисунке~\ref{fig:parietal-mechanism} показано взаимное расположение основных процессов. Отдельный препарат, отдельный сигнальный путь и вся секреторная система --- разные уровни описания. Красным отмечено место основного лекарственного воздействия, а не все процессы, меняющиеся после лечения.

\begin{figure}[htbp]
\centering
\begin{tikzpicture}[
  x=1cm,y=1cm,
  every node/.style={font=\scriptsize,align=center},
  process/.style={draw=blue!45!black,fill=white,rounded corners=2pt,inner sep=4pt},
  flux/.style={-{Stealth[length=2mm]},thick,draw=blue!55!black},
  signal/.style={-{Stealth[length=2mm]},dashed,draw=black!65},
  drug/.style={draw=red!65!black,fill=red!5,rounded corners=2pt,inner sep=4pt}
]
\fill[black!3] (0,0) rectangle (3.1,10.5);
\fill[blue!3] (3.1,0) rectangle (11.3,10.5);
\fill[yellow!12] (11.3,0) rectangle (15.8,10.5);
\draw[thick,blue!45!black] (3.1,0) -- (3.1,10.5);
\draw[thick,blue!45!black] (11.3,0) -- (11.3,10.5);
\node[font=\small\bfseries] at (1.55,10) {Интерстиций};
\node[font=\small\bfseries] at (7.2,10) {Париетальная клетка};
\node[font=\small\bfseries] at (13.55,10) {Секреторный\\каналец};
\node[text width=2.3cm] (histamine) at (1.55,8.75) {Гистамин\\из ECL-клеток};
\node[process,text width=1.25cm] (h2) at (4.15,8.75) {H$_2$};
\node[process,text width=2.75cm] (camp) at (7.7,8.75) {G$_s$--аденилатциклаза\\cAMP--PKA};
\draw[signal] (histamine) -- (h2);
\draw[signal] (h2) -- (camp);
\node[text width=2.3cm] (acetylcholine) at (1.55,7.3) {Ацетилхолин};
\node[process,text width=1.25cm] (m3) at (4.15,7.3) {M$_3$};
\node[process,text width=2.75cm] (calcium) at (7.7,7.3) {G$_q$--PLC\\IP$_3$--Ca$^{2+}$};
\draw[signal] (acetylcholine) -- (m3);
\draw[signal] (m3) -- (calcium);
\node[process,text width=3.6cm] (traffic) at (7.6,5.9) {Организация апикальной мембраны\\и доступность помп};
\draw[signal] (camp.south east) -- (9.6,8.2) -- (9.6,6.3) -- (traffic.east);
\draw[signal] (calcium.south) -- (traffic.north);
\node[drug,text width=2.8cm] (activated) at (13.65,8.05) {Доставленный лансопразол\\$\downarrow$ кислая среда\\активное производное};
\node[drug,text width=1.45cm] (pump) at (11.3,5.8) {H$^+$/K$^+$-\\АТФаза\\ATP4A/B};
\draw[signal] (traffic.east) -- (pump.west);
\draw[red!70!black,line width=1pt,-{Bar[width=3mm]}] (activated.south) -- (13.65,5.8) -- (pump.east);
\draw[flux] (9.9,6.6) -- node[above] {H$^+$} (12.4,6.6);
\draw[flux] (12.4,5.05) -- node[below] {K$^+$} (9.9,5.05);
\node[text width=1.8cm] at (9.65,4.55) {MgATP $\rightarrow$\\ADP + P$_i$};
\node[process,text width=4.2cm] (carbonic) at (6.7,3.65) {Карбоангидраза II\\CO$_2$ + H$_2$O $\rightleftharpoons$ H$^+$ + HCO$_3^{-}$};
\node[process,text width=1.45cm] at (3.1,3.65) {AE2\\SLC4A2};
\draw[flux] (1.8,4.45) -- node[above] {Cl$^{-}$} (4.45,4.45);
\draw[flux] (4.45,2.85) -- node[below] {HCO$_3^{-}$} (1.8,2.85);
\node[process,text width=1.5cm] at (11.3,2.7) {KCNQ1\\KCNE2};
\draw[flux] (10,3.4) -- node[above] {K$^+$} (12.7,3.4);
\node[text width=2.3cm] at (13.9,2.7) {Калиевый\\рециклинг};
\node[process,text width=1.6cm] at (11.3,0.8) {Выход\\Cl$^{-}$};
\draw[flux] (9.9,1.55) -- node[above] {Cl$^{-}$} (12.7,1.55);
\node[text width=5.2cm,align=left] at (6.7,1.5) {Сплошные стрелки: ионные потоки.\\Пунктир: регуляторные связи.\\Красная линия: блокада помпы.};
\end{tikzpicture}
\caption{Авторская схема, не в масштабе. Каналец сообщается с просветом железы; цитозоль не изображен как кислая полость. Положение белков условно, показаны основные направления стимулированной секреции, а не все переносчики и сигнальные ветви. Гастриновый и соматостатиновый контуры описаны в тексте. Прямое ингибирование AE2, карбоангидразы, KCNQ1 или хлоридных каналов лансопразолом из рисунка не следует.\src{molh2signal,molprobes,molezrin,molca,molae2,molkchannel,molslc26,molcys}}
\label{fig:parietal-mechanism}
\end{figure}

\subsection{Доставка к мишени: почему нужна кишечнорастворимая защита}
Неактивированный препарат должен пережить прохождение желудочного содержимого, высвободиться из гастрорезистентных гранул, всосаться и попасть к париетальной клетке со стороны кровотока. Далее мембранопроницаемая фракция достигает кислого канальца, где задерживается и превращается в реакционноспособную форму. Поэтому образование активного продукта \emph{до всасывания в просвете желудка} не эквивалентно его образованию рядом с помпой.\src{pharmalabel,molnagaya90}

В опытах на изолированных клетках наружное добавление заранее образованных активных производных AG-2000 и AG-1812 не воспроизводило действие исходного AG-1749. Это иллюстрирует зависимость эффекта от доставки и локального превращения, а не только от реакционной способности вещества в пробирке. Разжевывание защищенных гранул потому и недопустимо, что нарушает последовательность доставки, а не потому, что лекарству вообще нельзя контактировать с кислотой.\src{molnagaya90,uslabel}

\subsection{Протонирование и ионное удержание: роль двух pKa}
В принятой химической модели для лансопразола приводятся приблизительные значения $pK_{a1}=3{,}83$ для пиридинового центра и $pK_{a2}=0{,}62$ для второй стадии протонирования бензимидазольного компонента. Это физико-химические параметры конкретной модели и условий измерения, а не значения pH, при которых эффект внезапно включается. Первая стадия способствует накоплению; вторая --- химической активации.\src{molchem}

Для \emph{упрощенного} однокислотного сопряженного равновесия слабого основания:
\[
\frac{[\mathrm{BH}^{+}]}{[\mathrm{B}]}=10^{pK_a-pH}.
\]
Если свободная нейтральная форма условно уравновешена между двумя средами, то
\[
\frac{C_{\text{каналец}}}{C_{\text{плазма}}}
\approx
\frac{1+10^{pK_a-pH_{\text{каналец}}}}
     {1+10^{pK_a-pH_{\text{плазма}}}}.
\]
При подстановке 3,83; 1,0 и 7,4 получается около $6{,}8\times10^2$. Это \textbf{авторский учебный расчет}, не измеренная тканевая концентрация: он игнорирует вторую протонизацию, связывание с белками, транспорт, химическое потребление и неоднородность pH. Именно поэтому его нельзя использовать для расчета дозы.\src{molchem}

\addtocontents{toc}{\protect\setcounter{tocdepth}{3}}
\subsubsection{Кинетика активации, константы \texorpdfstring{$pK_a$}{pKa} и системная pH-селективность лансопразола}
\label{sec:activation-ph-selectivity}
\addtocontents{toc}{\protect\setcounter{tocdepth}{2}}

\textbf{Сопоставимость физико-химических параметров.} Для корректного межмолекулярного сравнения необходимо использовать константы, относящиеся к одному центру протонирования и сопоставимому методу определения. В исследовании Shin и соавт. значения $pK_{a1}$ пиридинового центра составляют приблизительно 3,83 для лансопразола, 4,06 для омепразола и 3,83 для пантопразола. Значение 2,96 для пантопразола этим источником не подтверждается и в данное сопоставление не включается. Константа диссоциации описывает равновесие ионизации, а не скорость образования активного производного: близость $pK_{a1}$ лансопразола и пантопразола не означает одинаковой кинетики их превращения.\src{molchem}

\textbf{Разделение накопления и активации.} В резко кислой среде секреторного канальца, при pH ниже 2, протонирование способствует удержанию препарата, а кислотокатализируемая перестройка --- быстрому образованию реакционноспособных производных. Первая стадия не тождественна второй: протонирование бензимидазольного компонента создает условия для внутримолекулярной реакции с образованием сульфеновой кислоты и циклического сульфенамида. Поэтому pH-селективность определяется сочетанием кислотно-основных равновесий, химических превращений и доступности цистеинов помпы; ее нельзя свести к единственному значению $pK_{a1}$ или представить как абсолютный порог включения при pH 2,0.\src{molchem}

\textbf{Прямое кинетическое сопоставление.} В исследовании Al-Matar и соавт. кислотные превращения трех ИПП изучали полярографически в присутствии 2-меркаптоэтанола при pH 2,0--5,0. Общая последовательность скоростей деградации и последующих реакций с тиолом была следующей: лансопразол $>$ омепразол $>$ пантопразол. Следовательно, эти данные не поддерживают обобщение о значительно большей химической стабильности лансопразола относительно омепразола в диапазоне pH 3,0--5,0. Необходимо различать исчезновение исходного сульфоксида, образование реакционноспособного промежуточного продукта и накопление дисульфида: измерение одного этапа не заменяет определения другого. Порядок скоростей в тиолсодержащем растворе не является непосредственным ранжированием скорости блокады помпы в организме.\src{molkinetics}

\textbf{Умеренно кислая среда и время экспозиции.} В работе Kromer и соавт. при pH 5,1 период полуактивации составлял 4,7 часа для пантопразола и 1,4 часа для омепразола. Таким образом, количественно показанное в этом сравнении замедление активации относится к пантопразолу, а не устанавливает аналогичное преимущество лансопразола. Авторы сопоставляли химическую активацию с длительностью системного присутствия, свободной экспозицией, связыванием с белками и липофильностью; различия внежелудочных эффектов исследовались в моделях \textit{in vitro}. Эти величины нельзя объединять с результатами другого экспериментального протокола в универсальный ряд клинической безопасности.\src{molphselect1998}

\textbf{Органеллярная селективность и предел клинического вывода.} Ионное удержание слабого основания не доказывает образования достаточного количества активного продукта, а замедленная активация не доказывает отсутствия накопления. Поэтому для лизосом макрофагов, почечного эпителия или сосудистого эндотелия необходимы данные о локальной свободной концентрации, времени пребывания и конкретной мишени; один показатель pH не характеризует всю клетку или орган. Исследование лизосомального закисления либо почечной Na$^+$/K$^+$-АТФазы \textit{in vitro} не равнозначно оценке частоты клинического повреждения соответствующих тканей.\src{molphselect1998}

\textbf{Итоговая интерпретация.} Кислотозависимая активация служит механистической предпосылкой избирательного действия ИПП на желудочную помпу. Однако из представленных исследований нельзя заключить, что лансопразол благодаря особой устойчивости при pH 3,0--5,0 минимизирует нецелевую активацию или обладает доказанным преимуществом системной безопасности перед омепразолом. Это разграничение химической гипотезы и клинического результата сохраняется при чтении раздела~\ref{sec:safety}.\src{molchem,molkinetics,molphselect1998}

\addtocontents{toc}{\protect\setcounter{tocdepth}{3}}
\subsubsection{Фармакокинетическое обоснование фиксированных комбинаций лансопразола с антитромботическими и противовоспалительными агентами}
\label{sec:fixed-combinations}
\addtocontents{toc}{\protect\setcounter{tocdepth}{2}}

\textbf{Регистрационная конкретизация: фиксированная таблетка и совместная упаковка.} В Японии применяется Takelda --- фиксированная комбинация ацетилсалициловой кислоты 100 мг и лансопразола 15 мг; о начале продаж объявлено 12 июня 2014 года. Это объединение антиагреганта и ИПП в одной лекарственной форме, а не комбинация двух противовоспалительных средств.\src{fdctakeldalaunch} Японская инструкция предусматривает одну таблетку в сутки для определенных сердечно-сосудистых и цереброваскулярных показаний либо после коронарной реваскуляризации у пациентов с язвой желудка или двенадцатиперстной кишки в анамнезе. Регистрация не означает назначения такой комбинации всем получающим аспирин или ее использования для лечения активной язвы.\src{fdctakeldalabel}

Для напроксена необходимо различать другие продукты. Историческая инструкция FDA для Prevacid NapraPAC 500 описывает \emph{совместную упаковку} капсул лансопразола 15 мг с отсроченным высвобождением и отдельных таблеток напроксена 500 мг, а не единую фиксированную таблетку; документ 2009 года не подтверждает текущую доступность продукта.\src{fdcnaprapac} В препарате Vimovo фиксированная таблетка содержит напроксен и \emph{эзомепразол}, а не лансопразол. Поэтому ее конструкцию и результаты нельзя приписывать лансопразолу; наличие зарегистрированной комбинации ИПП с НПВП не устанавливает взаимозаменяемости молекул и форм.\src{fdcvimovo}

\textbf{Биохимическая комплементарность вместо недоказанного усиления антитромботического действия.} Аспирин необратимо подавляет тромбоцитарную COX-1 и образование тромбоксана A$_2$; лансопразол после кислотной активации ингибирует желудочную H$^+$/K$^+$-АТФазу. Это воздействие на разные мишени с разными терапевтическими целями.\src{fdctakeldalabel} Ульцерогенное действие НПВП включает системное подавление синтеза защитных простагландинов и местное повреждение. Лансопразол снижает кислотную нагрузку на уязвимую слизистую, но не восстанавливает синтез простагландинов и не является аналогом мизопростола. Термин «синергизм» здесь допустим лишь как клиническая взаимодополняемость, а не как доказанное сверхаддитивное усиление антиагрегации, обезболивания или противовоспалительного эффекта.\src{fdcnaprapac,misotrial}

\textbf{Липофильность и кинетика мишени: разные условия гастропротекции.} Липофильность нейтральной фракции способствует мембранному переносу, однако скорость появления препарата в крови определяется также гастрорезистентной доставкой, опорожнением желудка и всасыванием. После системной абсорбции препарат достигает париетальной клетки, удерживается в кислой среде канальца и превращается в реакционноспособное производное. Продолжительность подавления секреции связана с модификацией помпы и восстановлением ее активности, а не с длительным сохранением высокой плазменной концентрации. Поэтому одна липофильность не доказывает ни преимущественно быстрого системного распределения относительно других ИПП, ни постоянного внутрижелудочного pH при одновременном приеме ульцерогенного средства.\src{molchem,molcys,molrestore,uslabel}

\textbf{Архитектура Takelda и разнесение экспозиции во времени.} Согласно отчету PMDA, наружная часть таблетки содержит гастрорезистентные гранулы лансопразола, а внутреннее ядро --- кишечнорастворимый аспирин. В исследовании влияния пищи медиана $t_{\max}$ лансопразола составляла 1,0 часа натощак и 4,0 часа после еды; для аспирина --- 3,5 и 5,5 часа соответственно. Эта последовательность плазменных максимумов не доказывает достижения защитного pH до начала действия аспирина: $t_{\max}$ не измеряет начало кислотосупрессии или длительность пребывания pH выше заданного порога.\src{fdctakeldareview} Гастрорезистентная доставка защищает кислотолабильный ИПП от преждевременного разрушения; фармакодинамический результат оценивают отдельно, с учетом повторного приема, CYP2C19 и доступности активных помп.\src{uslabel,cpic,molrestore}

\textbf{Биоэквивалентность не равна усилению эффекта.} В основном исследовании PMDA у 227 участников отношения геометрических средних для лансопразола из Takelda относительно совместно принятых монопрепаратов составляли 1,121 для $C_{\max}$ (90\% ДИ 1,068--1,177) и 1,075 для AUC$_{0-24}$ (1,046--1,104), удовлетворяя критериям биоэквивалентности. Для аспирина предварительное исследование не подтвердило полный набор фармакокинетических критериев; регуляторное обоснование дополнительно опиралось на сопоставление профилей растворения. Следовательно, регистрация не доказывает полной фармакокинетической идентичности обоих компонентов или превосходства фиксированной таблетки над раздельным приемом.\src{fdctakeldareview}

\textbf{Влияние пищи на экспозицию.} В японской инструкции приведены данные 12 здоровых участников с генотипом CYP2C19 EM: прием комбинации через 30 минут от начала завтрака уменьшал $C_{\max}$ лансопразола на 27,9\%, AUC --- на 10,7\% относительно приема натощак. Эти измерения характеризуют конкретную форму, но сами по себе не доказывают недостаточной гастропротекции после еды. Они показывают, почему время приема нельзя определять только по липофильности или механически переносить с другого комбинированного препарата.\src{fdctakeldalabel}

\textbf{Профиль высвобождения не устраняет системную токсичность.} В Vimovo немедленное высвобождение эзомепразола сочетается с кишечнорастворимым ядром напроксена; эту последовательность нельзя переносить на Takelda, где гастрорезистентными являются оба компонента. Кишечнорастворимая оболочка меняет место высвобождения, но не отменяет системное ингибирование циклооксигеназы. Даже при добавлении ИПП сохраняются риски кровотечения и внежелудочных осложнений НПВП, включая почечные и сердечно-сосудистые.\src{fdcvimovo,fdcnaprapac} Для аспириновой комбинации сохраняются антиагрегантное действие и обусловленный им геморрагический риск; защиту верхних отделов ЖКТ нельзя переименовывать в нейтрализацию всех токсических эффектов аспирина.\src{fdctakeldalabel}

\textbf{Клинический результат и приверженность требуют раздельной оценки.} Рандомизированное исследование Sugano и соавт. подтверждает эффективность лансопразола 15 мг в предотвращении рецидивов язв при продолжающемся приеме низких доз аспирина по сравнению с гефарнатом. Оно обосновывает сочетанную гастропротекцию, но не доказывает превосходства фиксированной таблетки над теми же компонентами, принимаемыми раздельно.\src{aspirin15} Уменьшение числа лекарственных единиц может облегчать соблюдение режима у коморбидных пациентов; производитель Takelda обозначает это как ожидаемое преимущество. Однако ожидание лучшей приверженности не равнозначно ее измеренному увеличению или доказанному снижению кровотечений вследствие самой фиксации доз.\src{fdctakeldalaunch}

\textbf{Практическая граница фиксирования доз.} В авторском клинико-фармакологическом синтезе такая форма рациональна, когда одновременно необходимы оба компонента именно в предусмотренных дозах. При необходимости отдельной коррекции ИПП или антиагреганта удобство одной таблетки не должно препятствовать персонализации. Прекращение комбинированного препарата означает прекращение обоих компонентов, поэтому самостоятельная отмена «гастропротектора» вместе с необходимым аспирином недопустима. Сравнение преимуществ должно учитывать не только таблеточную нагрузку, но и совместимость лечения, риск язвенного кровотечения и сохранение показания к антиагреганту; клинические основания подробнее разобраны в разделах~\ref{sec:clopidogrel} и~\ref{sec:nsaid}.\src{fdctakeldalabel,cardioesc2024,agadepres}

\addtocontents{toc}{\protect\setcounter{tocdepth}{3}}
\subsubsection{Фармакодинамический синергизм удержания внутрижелудочного pH}
\label{sec:ph-retention-synergy}
\addtocontents{toc}{\protect\setcounter{tocdepth}{2}}

\textbf{Фармакодинамическая комплементарность и предмет анализа.} Сочетание лансопразола с нестероидным противовоспалительным препаратом (НПВП) или низкодозной ацетилсалициловой кислотой объединяет подавление кислотной секреции с сохранением необходимого противовоспалительного либо антитромботического действия. В данном подразделе синергизм означает функциональную комплементарность воздействий на разные звенья патогенеза, а не установленное сверхаддитивное взаимодействие на одной молекулярной мишени. Для его оценки необходимо последовательно сопоставлять лекарственную форму, свободную системную экспозицию, доступность секретирующих помп, продолжительность повышения внутрижелудочного pH и частоту гастродуоденальных осложнений. Фармакокинетическая совместимость компонентов сама по себе не равнозначна доказанной клинической гастропротекции.\src{fdcnaprapac,fdctakeldareview,misotrial,aspirin15}

\textbf{Липофильность и пространственная последовательность доставки.} Липофильность неионизированной фракции лансопразола создает предпосылку для распределения молекулы в фосфолипидном бислое и трансцеллюлярного переноса. После растворения гастрорезистентного покрытия в кишечнике всасывание предполагает преодоление апикального барьера энтероцита и последующий выход в кровоток. Париетальная клетка получает препарат уже из крови через базолатеральную поверхность; затем мембранопроницаемая фракция достигает апикальной канальцевой мембраны и кислого секреторного пространства. Следовательно, кишечную абсорбцию и поступление к желудочной мишени нельзя представлять как единый перенос через две апикальные мембраны. Скорость доставки определяется не только липофильностью, но и растворением формы, опорожнением желудка, площадью всасывания, кровоснабжением и свободной концентрацией препарата. Указанные в инструкции высокая степень связывания с белками плазмы и зависимость экспозиции от приема пищи дополнительно ограничивают возможность выводить скорость тканевого поступления из одного физико-химического параметра.\src{pharmalabel,molchem}

\textbf{Диффузионный поток и кислотозависимое удержание.} В приближении пассивного переноса поток нейтральной молекулы определяется мембранной проницаемостью и разностью ее свободных концентраций по обе стороны мембраны. В кислой среде функционирующего канальца протонирование пиридинового центра уменьшает долю вещества, способного к обратной диффузии; последующее химическое превращение также уменьшает локальный пул неизмененного пролекарства. Такое сопряжение переноса с протонированием и реакцией поддерживает поступление новых молекул, пока сохраняется доступная системная экспозиция. Быстрая доставка важна для совпадения экспозиции с периодом секреторной активности, однако повышенная суммарная канальцевая концентрация не означает столь же высокой концентрации свободной нейтральной формы. Канальцевое пространство представляет апикальный люминальный компартмент, а не цитозоль. Поэтому корректно различать локальное накопление протонированного пролекарства, образование активного производного и количество уже модифицированной помпы; их объединение в понятие «внутриклеточного запаса сульфенамида» искажает кинетическую модель.\src{molchem,molnagaya90}

\textbf{Кислотная активация и длительность блокады мишени.} Протонирование бензимидазольного компонента создает условия для кислотокатализируемой перестройки сульфоксида с образованием тиофильных производных, включая сульфеновую кислоту и циклический сульфенамид. Реакционноспособный продукт взаимодействует с доступными люминальными цистеинами H$^+$/K$^+$-АТФазы и формирует дисульфидный аддукт. Локальная реакция с белком, а не длительная циркуляция готового сульфенамида, объясняет сохранение эффекта после уменьшения плазменной концентрации исходного препарата. Активированное производное не следует рассматривать как свободно диффундирующий резерв, равномерно распределенный между клетками. Степень кислотосупрессии определяется долей доступных и модифицированных помп; повторное поступление лансопразола позволяет воздействовать на мишени, включающиеся в последующие секреторные циклы. Восстановление функционального пула помп и восстановительные реакции ограничивают длительность подавления секреции, поэтому первая доза не обеспечивает одномоментной блокады всей кислотопродукции.\src{molchem,molnagaya89,molrestore,pharmalabel}

\textbf{Перекрест с местным и системным ульцерогенным действием.} НПВП нарушают простагландин-зависимую защиту слизистой вследствие ингибирования циклооксигеназ; вклад ЦОГ-1 и ЦОГ-2 зависит от селективности препарата, дозы и состояния ткани. Уменьшение синтеза защитных простагландинов ослабляет поддержание слизисто-бикарбонатного барьера, микроциркуляции и репаративных процессов. Прямое контактное действие перорального агента и системное подавление простагландинового синтеза являются взаимодополняющими, но не тождественными механизмами повреждения. Аспирин дополнительно подавляет агрегацию тромбоцитов, что имеет самостоятельное значение для кровотечения из уже поврежденной слизистой. Лансопразол воздействует на кислотный компонент агрессии: уменьшение поступления протонов и кислотно-пептической нагрузки снижает дополнительное повреждение при ослабленном барьере. Это частичная патофизиологическая компенсация, а не восстановление активности ЦОГ, синтеза простагландинов, микроциркуляции или функции тромбоцитов.\src{fdcnaprapac,fdctakeldalabel,misotrial}

\textbf{Содержание показателя pH и пределы его интерпретации.} Достижение pH выше 4,0 либо 4,5 характеризует отдельный уровень кислотосупрессии, но эти пороги необходимо задавать раздельно. Для фармакодинамического анализа определяют время первого превышения выбранного порога, долю суток выше него и продолжительность непрерывного удержания; однократное пересечение порога не заменяет двух последних показателей. Повышение pH уменьшает кислотную агрессию и создает более благоприятные условия для восстановления слизистой, но ни pH $>4{,}0$, ни pH $>4{,}5$ в приведенных исследованиях не валидированы как универсальная гарантия отсутствия НПВП-эрозий, язв или кровотечений. Средний внутрижелудочный pH также не воспроизводит непосредственно условия у поверхности эпителия. Поэтому фармакодинамический показатель следует сопоставлять с эндоскопическими и клиническими исходами, а не использовать вместо них.\src{phthoring1999,phbell1996,misotrial,aspirin15}

\textbf{Прямые измерения скорости и устойчивости кислотосупрессии.} В открытом рандомизированном перекрестном исследовании Thoring и соавт. у 15 здоровых добровольцев после однократного приема лансопразола 30 мг pH впервые превышал 4 через 130 минут, после омепразола 20 мг --- через 250 минут; средний pH за восемь часов составлял соответственно 2,9 и 2,0. Эти результаты относятся к сопоставленным дозам и условиям опыта; независимый вклад липофильности в скорость эффекта не измеряли.\src{phthoring1999} В исследовании Bell и Hunt у девяти здоровых мужчин, получавших лансопразол 30 мг/сут семь дней, средний суточный pH увеличился с 2,11 исходно до 3,57 в первый день, а доля времени с pH $>3$ составляла 54,7\% в первый и 67,4\% в седьмой день. Измеренный порог pH 3 не эквивалентен порогу 4 или 4,5. Совокупность этих данных подтверждает раннее начало и увеличение продолжительности кислотосупрессии при повторных дозах, но не доказывает непрерывного целевого pH у каждого пациента, получающего НПВП.\src{phbell1996}

\textbf{Фиксированная форма и совместная упаковка: различия профилей высвобождения.} Фиксированная таблетка Takelda содержит лансопразол 15 мг и ацетилсалициловую кислоту 100 мг; ее регистрационное обоснование необходимо анализировать по данным именно этой лекарственной формы. Отчет PMDA рассматривает сопоставление экспозиций компонентов комбинированной таблетки и раздельно назначаемых препаратов, а не доказательство универсального временного опережения блокады помпы относительно действия аспирина.\src{fdctakeldareview,fdctakeldalabel} Для напроксена исторический Prevacid NapraPAC представляет совместную упаковку отдельной гастрорезистентной капсулы лансопразола и таблетки напроксена, а не единую таблетку с программируемой последовательностью высвобождения. В фиксированной форме Vimovo компонентом ИПП является эзомепразол; переносить ее конструкцию на лансопразол нельзя. Различия между фиксированной комбинацией, совместной упаковкой и свободным сочетанием препаратов определяют допустимость фармакокинетических и клинических выводов.\src{fdcnaprapac,fdcvimovo}

\textbf{Временное соотношение высвобождения, абсорбции и эффекта.} В архивной инструкции NapraPAC среднее время достижения максимальной плазменной концентрации лансопразола составляет около 1,7 часа, напроксена --- 2--4 часа. Даже такое соотношение не устанавливает, что блокада помп завершена до первого системного действия НПВП: абсорбция начинается до достижения максимума концентрации, а поступление лансопразола в кровь предшествует его доставке, активации и связыванию с мишенью.\src{fdcnaprapac} Для доказательства фармакодинамического опережения необходимы согласованные во времени измерения высвобождения обоих компонентов, их свободной экспозиции, внутрижелудочного pH и повреждения слизистой при начальном и повторном приеме. Биоэквивалентность не устанавливает равенства этих временных характеристик и не доказывает полного устранения контактной токсичности.\src{fdctakeldareview,phthoring1999} Формулировка о предотвращении «отсроченного системного некробиоза делящихся клеток желудочных желез» также не следует из фармакокинетического профиля: она требует отдельного морфологического подтверждения и не является установленной конечной точкой приведенных испытаний гастропротекции.\src{misotrial,aspirin15,nsaid15}

\textbf{Клиническая верификация и итоговое положение.} В рандомизированном исследовании у 461 пациента с язвенным анамнезом, продолжавшего низкодозный аспирин, кумулятивная частота язв желудка или двенадцатиперстной кишки к 361-му дню составляла 3,7\% при лансопразоле 15 мг/сут и 31,7\% при гефарнате 50 мг дважды в сутки. Этот результат подтверждает снижение риска повторных язв в изученной популяции; он не устанавливает полного устранения повреждения и не доказывает преимущества фиксированной таблетки над раздельным приемом тех же компонентов.\src{aspirin15} Клинически обоснованное положение состоит в том, что адекватная кислотосупрессия может поддерживать гастродуоденальную защиту при необходимом продолжении ульцерогенного лечения, уменьшая язвенный риск без устранения всех местных и системных механизмов токсичности. Частоту эрозий, язв и кровотечений оценивают раздельно, с учетом компаратора и исходного риска; прямые сравнения представлены в разделе~\ref{sec:gastroprotection-trials}.\src{misotrial,aspirintrial,nsaid15}

\textbf{Приверженность и индивидуализация длительного режима.} Объединение показанных компонентов в одной форме может уменьшать таблеточную нагрузку и вероятность пропуска гастропротектора, однако ожидаемое удобство не равнозначно измеренному повышению приверженности. Этот исход требует собственного сравнительного исследования с учетом прекращения приема, переносимости и сохранения необходимых доз. Фиксирование доз рационально только при одновременном показании к обоим компонентам; необходимость изменить дозу ИПП либо антитромботического средства может потребовать раздельного назначения. Таким образом, продолжительная гастропротекция определяется сочетанием обоснованного показания, подходящей формы, достаточного режима кислотосупрессии и регулярной оценки безопасности, а не одним свойством липофильности или предполагаемой синхронностью высвобождения.\src{fdctakeldalaunch,fdctakeldareview,fdctakeldalabel}

\subsection{Кислотокатализируемая перестройка сульфоксида}
Протонирование бензимидазольного компонента повышает электрофильность его атома C2. Внутримолекулярное взаимодействие с непротонированной долей пиридинового азота запускает перегруппировку с образованием тиофильного серосодержащего продукта. «Тиофильный» здесь означает способный реагировать с серой тиольных групп белка. Это не простое обратимое присоединение протона и не реакция, катализируемая CYP.\src{molchem,molkinetics}

В механистических схемах различают \textbf{сульфеновую кислоту} и \textbf{циклический сульфенамид}; последнему для лансопразола соответствует историческое обозначение AG-2000. Сульфеновая кислота может циклизоваться с отщеплением воды. В зависимости от условий среды и реакции с белком нельзя сводить все превращения к единственной свободно циркулирующей «активной молекуле». В ранних опытах изучали также дисульфидное производное AG-1812.\src{molnagaya89,molkinetics}

Химическая модель в присутствии 2-меркаптоэтанола при pH 2--5 показала различия скоростей кислотных превращений и образования дисульфидов между лансопразолом, омепразолом и пантопразолом. Но порядок скоростей в модельном растворе не ранжирует препараты по клинической эффективности: необходимо учитывать экспозицию, доставку и доступные помпы.\src{molkinetics}

\subsection{Ковалентная блокада: дисульфид, а не обычное рецепторное связывание}
Тиольная группа цистеина помпы вступает в реакцию с активированным производным. Условная схема для сульфеновой кислоты:
\[
\mathrm{Pump{-}Cys{-}SH}+\mathrm{Drug{-}SOH}
\longrightarrow
\mathrm{Pump{-}Cys{-}S{-}S{-}Drug}+\mathrm{H_2O}.
\]
Обозначение Drug здесь относится к \emph{перегруппированному} фрагменту, а не к неизмененному лансопразолу. При участии сульфенамида исходная реакционная последовательность иная, но принципиальный результат --- образование дисульфидного аддукта с белком. Такая схема не означает алкилирования цистеина с образованием связи C--S.\src{molnagaya89,molcys}

В ранних опытах с желудочными микросомами связывание меченого AG-1749 сопровождалось уменьшением свободных SH-групп и угнетением ATPазной активности; дитиотреитол противодействовал этим изменениям. Это биохимическая аргументация участия тиолов и дисульфидов, более содержательная, чем само по себе совпадение эффектов препарата и повышения pH.\src{molnagaya89}

\subsection{Какие цистеины важны и почему списки в источниках различаются}
Для блокады общего для ИПП люминального участка особенно важен \textbf{Cys813} в области между пятым и шестым трансмембранными сегментами. Для лансопразола описано также присоединение в области \textbf{Cys321} у наружного конца третьего сегмента. При картировании реакции в 1997 году выявлялось более широкое связывание с цистеин-содержащими доменами. Поэтому краткую учебную пару «813/321» не следует превращать в утверждение об исключительности этих двух остатков при любых условиях.\src{molcys}

Ключевое различие --- между \emph{обнаруженным присоединением} и \emph{присоединением, определяющим потерю функции}. В той работе ингибирование связывали прежде всего с модификацией Cys813. Для пантопразола важна также реакция с более глубоко расположенным Cys822, что помогает объяснить иную чувствительность аддукта к восстановителям. Нельзя автоматически приписывать лансопразолу «необратимую блокаду Cys822» по аналогии с другим ИПП.\src{molcys,molrestore}

\subsection{«Необратимость» в фармакологии не означает химическую вечность}
Ковалентный аддукт не обязан быстро распадаться при снижении свободной концентрации препарата; этим он отличается от обычного обратимого комплекса лиганд--рецептор. Однако дисульфидная связь потенциально восстанавливаема. В исследовании Shin и Sachs 2002 года активность ATPазы, угнетенной лансопразолом, восстанавливалась под действием дитиотреитола или глутатиона. Аналогичный вывод нельзя без уточнения переносить на все ИПП.\src{molrestore}

Восстановление кислотопродукции поэтому включает не только синтез нового белка, но и возможную реактивацию части старых помп, а также включение ранее не заблокированного пула. Эксперимент с восстановителем не является рекомендацией принимать глутатион или антиоксиданты для отмены действия препарата. Клиническая скорость восстановления определяется целой системой, а не одним временем химического распада.\src{molrestore,agadepres}

\subsection{Связь молекулярного механизма с режимом приема}
\begin{itemize}
\item \textbf{Прием до еды.} Временное присутствие пролекарства в крови желательно совместить с активацией секреторной поверхности; отсюда рекомендация принимать обычный лансопразол до приема пищи.\src{acggerd,molezrin}
\item \textbf{Нарастание эффекта в первые дни.} Одна доза не охватывает все помпы одновременно. Повторный прием влияет на вновь доступную часть пула; это не обязательно означает накопление исходной молекулы в плазме.\src{pharmalabel}
\item \textbf{Длительность после элиминации.} Важен срок сохранения модифицированной мишени, а не только плазменный период полувыведения.\src{molrestore}
\item \textbf{Неполный ночной контроль.} Короткое окно экспозиции и неодновременная активность помп ограничивают подавление секреции. Это не доказательство молекулярной устойчивости помпы к препарату у каждого пациента с ночной изжогой.\src{acggerd}
\end{itemize}

Разделение суточной терапии иногда меняет охват активируемых помп, но биологическое объяснение не заменяет показания к двукратному режиму. Высокая доза не исправляет функциональную изжогу, регургитацию некислого содержимого или неверный диагноз.\src{acggerd,agagerd}

\subsection{Молекулярное отличие от H2-блокаторов и P-CAB}
H$_2$-блокатор уменьшает рецепторную гистаминовую стимуляцию, тогда как лансопразол блокирует конечный фермент после разных секреторных стимулов. P-CAB также действуют на желудочную H$^+$/K$^+$-АТФазу, но \textbf{не требуют превращения сульфоксида в сульфенамид и не используют тот же дисульфидный механизм}. В структурной работе Abe и соавт. исследованы комплексы помпы с вонопразаном и SCH28080 в люминально доступной области. Их калий-конкурентное ингибирование следует отличать от ковалентной модификации ИПП.\src{molstructure,molnagaya90}

Обратимость связывания не обязательно означает короткий клинический эффект: значение имеют аффинность, время удержания на мишени и экспозиция. Напротив, ковалентный механизм не гарантирует полного подавления всех помп. Отсюда следует, что одной характеристики «сильнее связывается» недостаточно для выбора препарата или объяснения результата сравнительного исследования.\src{agapcab,vontrial}

\subsection{Как измерять молекулярный эффект: pH не равен кислотопродукции}
По определению $pH=-\log_{10}a_{\mathrm{H}^{+}}$, где $a_{\mathrm{H}^{+}}$ --- активность протонов. Ее изменение при переходе от pH 1 к pH 4 составляет $10^3$ раз. Но это не означает автоматически уменьшения \emph{скорости секреции} на 99,9\%: концентрация зависит также от объема, буферов, пищи, перемешивания и удаления содержимого. Это математическое различие концентрационной характеристики и потока, а не новая оценка эффективности препарата.

В механистических опытах необходимо различать активность ATPазы, накопление кислоты в препарате железы, титруемое количество секретированных H$^+$ за время и внутрижелудочный pH. Например, исследования хлоридных переносчиков использовали и измерения pH, и оценку кислотно-основных эквивалентов; структурные изменения клетки оценивались отдельно.\src{molslc26,molclc2}

В клинической интерпретации доля времени с pH выше заданного порога, облегчение изжоги, заживление эрозий и эрадикация \Hp\ также не взаимозаменяемы. Нельзя вычислить процент модифицированного Cys813 по одному измерению pH или объявить отсутствие молекулярного эффекта по сохранению симптома, который не обусловлен кислотой.\src{acggerd,agagerd,acghp}

\subsection{Как доказать прямую мишень, а не только сопутствующий эффект}
Для исследовательской оценки полезна следующая последовательность вопросов:
\begin{enumerate}
\item Измерено ли взаимодействие с очищенной мишенью или только изменение функции клетки?
\item Установлены ли химическая форма вещества, время инкубации, pH и свободная концентрация?
\item Подтверждается ли функциональный эффект независимым методом, например анализом аддукта или целевой мутацией?
\item Исключены ли токсичность, изменение числа клеток и неспецифическая блокада других белков?
\item Совместимы ли экспериментальные условия с экспозицией в ткани человека?
\end{enumerate}
Это авторская схема критического чтения, не формальная шкала качества исследования.

Ее необходимость показана непосредственно на париетальных клетках: Sugita и соавт. исследовали неспецифические действия распространенных фармакологических зондов. Поэтому угнетение секреции после добавления «ингибитора киназы» само по себе не устанавливает, что именно эта киназа является прямой мишенью лансопразола. Для основной тиоловой реакции препарата существенна совокупность данных о связывании, доступных SH-группах, восстановлении активности и локализации участков реакции, а не один функциональный тест.\src{molprobes,molnagaya89,molcys,molrestore}

\subsection{Метаболизм, стереохимия и экспозиция: не путать с мишенью}
\label{sec:metabolism}
В опытах с человеческими печеночными ферментами основными направлениями были образование 5-гидроксилансопразола с преимущественным участием CYP2C19 и сульфона с участием CYP3A. Вклад ферментов зависит от концентрации субстрата; результат микросомального опыта при высокой концентрации не обязательно описывает метаболизм при обычной экспозиции.\src{molp450}

Плохой метаболизатор CYP2C19 обычно дольше экспонирован к исходному пролекарству, но это не означает иную аминокислотную последовательность его протонной помпы. Быстрый клиренс, наоборот, сокращает доступность вещества для местной активации. Поэтому фармакогенетика лансопразола преимущественно описывает фармакокинетическую, а не мишень-зависимую вариабельность ответа.\src{cpic,molp450}

R- и S-энантиомеры могут по-разному метаболизироваться. При химическом превращении исходная сульфоксидная стереохимия не должна автоматически переноситься на конечный реакционноспособный продукт. Особенности декслансопразола нельзя объяснять только «более сильной связью R-формы с Cys813»: значение имеют экспозиция и технология высвобождения.\src{molchem,dexlabel}

\textbf{Межмолекулярные различия.} В обзоре Леоновой метаболическая зависимость от CYP2C19 выраженнее для омепразола, лансопразола и пантопразола, чем для рабепразола и эзомепразола. Для рабепразола показан также неферментативный путь превращения в тиоэфир. Это объясняет возможность меньшей генотипической вариабельности, но не означает полного отсутствия ферментативного метаболизма. Снижение клиренса исходного ИПП увеличивает его доступность для кислотной активации, а не активирует препарат в печени. Употребленное в статье слово «сульфонамиды» не переносится на описание активной формы: в настоящем обзоре сохраняется химически отличное понятие сульфенамида.\src{leonova2015,molp450,molchem}

Сравнение зависимости разных ИПП от CYP2C19 и границы переноса дозовых ориентиров CPIC приведены в разделе~\ref{sec:pharmacogenetics}.

\subsection{Ферментативная кинетика: почему R- и S-формы исчезают с разной скоростью}
В работе Kim и соавт. на человеческих микросомах печени суммарный внутренний клиренс образования гидрокси- и сульфонового метаболитов составлял 13,5 для R-формы и 57,3 мкл/мин/мг белка для S-формы. Это параметры экспериментальной системы, а не общий клиренс пациента в мл/мин. Для CYP2C19-зависимого гидроксилирования кажущиеся $K_m$ составляли 13,1 и 2,3 мкмоль/л соответственно. В превращениях участвовали также CYP3A4 и, в отдельных условиях, CYP2C9.\src{molstereokin}

Для простейшей модели Михаэлиса--Ментен:
\[
v=\frac{V_{\max}[S]}{K_m+[S]},
\qquad [S]\ll K_m:\quad
\frac{v}{[S]}\approx\frac{V_{\max}}{K_m}.
\]
Здесь $[S]$ обозначает концентрацию \emph{субстрата}, а не S-энантиомер. При малой концентрации важным становится отношение $V_{\max}/K_m$, поэтому ранжирование ферментов в опыте с избытком субстрата может отличаться от ранжирования при низкой экспозиции. $K_m$ не следует без дополнительных условий отождествлять с константой диссоциации $K_d$ или объявлять прямым измерением «силы связывания». Приведенная формула --- учебное объяснение кинетики, не индивидуальная модель дозирования.\src{molstereokin}

\subsection{Стереоселективность у человека: состав капсулы не равен составу плазмы}
В исследовании Miura и соавт. 2004 года участвовали 18 здоровых японских добровольцев --- по шесть в трех группах CYP2C19. Отношение $\mathrm{AUC}_R/\mathrm{AUC}_S$ составляло 12,7 у гомозиготных extensive metabolizers, 8,5 у гетерозиготных extensive metabolizers и 5,8 у poor metabolizers. Исторические названия этих групп не следует механически приравнивать ко всем современным категориям CPIC: необходим конкретный диплотип.\src{molhumanstereo,cpic}

Следовательно, равные количества двух энантиомеров в рацемате не создают одинаковую плазменную экспозицию. Но эти отношения не являются универсальными константами: результат получен в небольшой выборке и при определенном протоколе. Преобладание R-формы в плазме не доказывает, что S-форма не участвует в подавлении секреции, а AUC не является прямым измерением количества ковалентно модифицированных помп.\src{molhumanstereo,molcys}

\subsection{Генотип и ингибитор CYP2C19: пример феноконверсии}
\label{sec:phenoconversion}
Генотип задает исходную способность к метаболизму, но сопутствующее средство может изменить фактическую активность фермента без изменения ДНК. В исследовании 2005 года после шести дней флувоксамина по 25 мг дважды в сутки изучали разовый прием 60 мг рацемического лансопразола. Экспозиция увеличивалась, причем S-форма была чувствительнее к взаимодействию. У гомозиготных extensive metabolizers среднее отношение $\mathrm{AUC}_R/\mathrm{AUC}_S$ уменьшалось с 12,7 до 6,4.\src{molfluvox}

Величина взаимодействия зависела от исходной группы CYP2C19: она была выраженнее при исходно сохраненной активности, чем при низкой. Нельзя считать плохого метаболизатора полностью защищенным от взаимодействий, поскольку остаются другие пути элиминации и фармакодинамические эффекты комбинации. Это иллюстрация функционального изменения метаболического фенотипа, а не основание назначать флувоксамин для усиления кислотосупрессии. Изменять дозу по указанному отношению AUC без клинического контекста нельзя.\src{molfluvox,cpic}

\textbf{Разграничение для клинической главы.} Исследованное лекарственное торможение не тождественно наследственному ультрабыстрому фенотипу *17/*17, а увеличение лечебной дозы не является феноконверсией. Границы переноса этого эксперимента на UM и совместная оценка генотипа, взаимодействий и дозы приведены в разделе~\ref{sec:pharmacogenetics}.\src{molfluvox,cpic}

\subsection{Лансопразол как субстрат и как ингибитор CYP: разные эксперименты}
Факт метаболизма ферментом не означает обязательного сильного торможения этого фермента. Напротив, ингибирование в микросомах возможно и тогда, когда клиническое взаимодействие при стандартной экспозиции невелико. В работе Liu и соавт. конкурентное торможение CYP2C19-зависимого гидроксилирования S-мефенитоина характеризовалось $K_i$ около 0,6 мкмоль/л для S-лансопразола и 6,1 мкмоль/л для R-лансопразола. Исследовались и другие CYP-реакции, поэтому S-форму нельзя считать строго селективным лабораторным ингибитором CYP2C19.\src{molcypinhib}

$K_i$ относится к модели взаимодействия фермента с ингибитором; $K_m$ --- к кинетике субстрата; IC$_{50}$ --- к концентрации, уменьшающей измеренный ответ в конкретном тесте наполовину. Эти величины не взаимозаменяемы. Для переноса результата на пациента важны свободная экспозиция, доля метаболизма сопутствующего препарата данным ферментом и альтернативные пути его выведения. Клиническую совместимость определяют данные о конкретной комбинации, а не одно микросомальное $K_i$.\src{molcypinhib,uslabel}

\subsection{AhR--CYP1A: активация транскрипционного пути не равна прямому лигандному связыванию}
В исследовании Novotna и соавт. 2014 года изучали репортерную активность AhR, уровни мРНК и белка CYP1A1/1A2 и ферментативный EROD-ответ в клеточных линиях и первичных человеческих гепатоцитах. В диапазоне 1--10 мкмоль/л S-лансопразол в целом сильнее активировал путь AhR--CYP1A, чем R-форма; при 100 мкмоль/л соотношение могло меняться. Профили первичных гепатоцитов различались между культурами.\src{molahr}

Поэтому выражение «S-форма всегда активнее R-формы» неверно без указания мишени, концентрации и модели. Усиление репортерного сигнала AhR не само по себе доказывает прямое высокоаффинное связывание с его лигандным карманом: авторы обсуждали непрямые регуляторные механизмы, транспорт и метаболизм. Результат также не доказывает канцерогенность, противоопухолевую эффективность или клинически значимую индукцию CYP1A у каждого пациента.\src{molahr}

\subsection{PXR/GR--CYP3A4: индукция и ингибирование могут сосуществовать}
В другой работе 2014 года R-, S- и рацемическая формы лансопразола повышали экспрессию мРНК и белка CYP3A4 в человеческих гепатоцитах. Репортерные опыты выявляли дозозависимую активацию pregnane X receptor (PXR), а также небольшую активацию глюкокортикоидного рецептора (GR) и усиление его дексаметазон-индуцированного сигнала. В исследованиях использовали диапазон 1--250 мкмоль/л и экспозицию 24--48 часов; выраженность ответа зависела от модели и концентрации.\src{molpxr}

Индукция означает изменение количества фермента через регуляцию экспрессии и оборот белка; конкурентное ингибирование --- изменение активности уже присутствующего фермента. Поэтому результаты двух типов опыта не обязаны противоречить друг другу. Однако итоговое изменение клиренса у пациента нельзя получить простым вычитанием процентов из разных клеточных тестов. Даже если авторы считают часть концентраций терапевтически релевантными, клиническая величина взаимодействия требует отдельной проверки. Эти опыты не обосновывают использование лансопразола как глюкокортикоидного препарата.\src{molpxr,molcypinhib}

\subsection{Количественный мост между клеточным опытом и экспозицией}
При молекулярной массе 369,36 г/моль пересчет номинальной концентрации исходного лансопразола дает:
\begin{center}
\begin{tabular}{rr}
\toprule
Концентрация, мкмоль/л & Массовая концентрация, мкг/мл \\
\midrule
1 & 0,36936 \\
10 & 3,6936 \\
100 & 36,936 \\
250 & 92,34 \\
\bottomrule
\end{tabular}
\end{center}
Это авторский арифметический пересчет по формуле $C_{\text{мкг/мл}}=0{,}36936\,C_{\text{мкмоль/л}}$, а не измеренные концентрации у пациента. Таблица относится к исходной молекуле; для сульфида, сульфона или другого производного нужна соответствующая молекулярная масса.\src{uslabel}

В инструкции для лансопразола указано около 97\% связывания с белками плазмы. Для \emph{условного} примера $C_{\mathrm{total}}=1$ мкг/мл при свободной доле $f_u=0{,}03$:
\[
C_{\mathrm{total}}\approx2{,}71\ \text{мкмоль/л},
\qquad
C_{\mathrm{free}}=f_u C_{\mathrm{total}}
\approx0{,}081\ \text{мкмоль/л}.
\]
Значение 1 мкг/мл выбрано для иллюстрации, не обозначает обязательный $C_{\max}$ после какой-либо дозы. Свободная фракция в культуральной среде зависит от ее состава и не обязана совпадать с плазменной.\src{pharmalabel}

Для оценки переносимости клеточного результата нужно знать свободную и внутриклеточную концентрации, продолжительность воздействия, стабильность вещества и pH. Краткий плазменный пик нельзя приравнивать к постоянной инкубации на 24--48 часов. И наоборот, низкая свободная плазменная концентрация сама по себе не исключает действия в кислом канальце, где происходят удержание, локальная активация и реакция с мишенью. Эти оговорки особенно важны для AhR/PXR- и лизосомальных моделей; они не служат расчетом безопасной противоопухолевой дозы.\src{molahr,molpxr,mollysosome,molchem}

\subsection{Вторичные последствия: pH, гастрин и репарация слизистой}
Уменьшение выхода H$^+$ изменяет химическую среду, но не нейтрализует уже присутствующую кислоту так, как антацид. Меняются условия кислотно-пептического повреждения, стабильность и растворение некоторых веществ. Поэтому заживление язвы не требует предполагать прямое стимулирование лансопразолом пролиферации эпителия; снижение агрессивного воздействия само по себе является важным звеном лечения.\src{pharmalabel,nice}

Подавление кислотности ослабляет кислотозависимую отрицательную обратную связь секреции гастрина. Гипергастринемия может сопровождаться стимуляцией ECL-клеток и повышением хромогранина A; это учитывается при лабораторной диагностике. Трофический ответ не равнозначен доказанному развитию рака у каждого длительно леченного человека.\src{uslabel}

\subsection{Helicobacter pylori: средовой эффект и прямые биохимические наблюдения}
Основное клиническое значение ИПП в эрадикации --- создание более благоприятных условий для работы антибактериальной комбинации. Повышение pH влияет на физиологическое состояние бактерии и действие кислоточувствительных антибиотиков. Это не означает преодоления любой генетической антибиотикорезистентности одной только кислотосупрессией.\src{maastricht,acghp}

\textbf{Уреаза.} В исследовании 1993 года лансопразол и его производные подавляли активность уреазы \Hp\ в бесклеточных препаратах и интактных бактериях. Уменьшение доступных SH-групп и устранение ингибирования глутатионом/дитиотреитолом поддерживали гипотезу тиол-зависимого взаимодействия. Данные не устанавливают для уреазы ту же карту цистеинов, что для ATP4A, и не доказывают эквивалентного эффекта после обычной дозы у человека.\src{molurease}

\textbf{Не только уреаза.} В исследовании 1995 года подавление роста сохранялось и у уреаза-дефицитных вариантов; выраженность антибактериального эффекта не объяснялась одной уреазной активностью. Работа 2001 года показала угнетение дыхания и снижение клеточного ATP в культуре \Hp. Следовательно, нельзя подменять наблюдаемую антибактериальную активность единственным механизмом «блокада уреазы». Эти эксперименты не обосновывают монотерапию лансопразолом для эрадикации.\src{molgrowth,molresp,acghp}

\subsection{Экспериментальный путь Keap1--Nrf2--HO-1}
\label{sec:nrf2}
В клетках слизистой желудка крысы RGM-1 лансопразол повышал экспрессию гемоксигеназы-1 (HO-1). В работе Takagi и соавт. 2009 года были выявлены окисление Keap1, фосфорилирование ERK и Nrf2, ядерное перемещение Nrf2 и усиление его связывания с регуляторными элементами. Уменьшение Nrf2 с помощью siRNA и мутация регуляторного участка репортера ослабляли ответ; подавление каскада MEK/ERK с помощью U0126 также снижало индукцию HO-1. Такая совокупность вмешательств поддерживает участие пути, а не только корреляцию двух маркеров.\src{molnrf2}

Функциональный смысл --- переключение клеточного стрессового ответа и снижение образования отдельных хемокинов в этой модели. Но исследование не установило, что лансопразол является прямым лигандом Nrf2 или ковалентно модифицирует конкретный цистеин Keap1 у человека. Нельзя переносить доказанную реакцию с цистеинами помпы на любой тиол-чувствительный белок без отдельного картирования.\src{molnrf2}

\subsection{HO-1, ферритин и ROS: почему нет единой универсальной цепи}
\label{sec:ho1}
HO-1 участвует в распаде гема с образованием биливердина, CO и железа; сопутствующая индукция ферритина важна для связывания высвобожденного железа. В опытах Schulz-Geske и соавт. на эндотелиальных клетках и макрофагах лансопразол увеличивал HO-1 и ферритин, что сопровождалось снижением NADPH-зависимого образования реактивных форм кислорода (ROS). Эффект ослаблялся ингибитором HO.\src{molho1}

В отличие от RGM-1-модели, авторы не связали индукцию HO-1 с окислительным стрессом или активацией MAPK; ее подавлял LY294002, ингибитор PI3K. Поэтому корректно говорить о клеточно-зависимых путях, а не о доказанной для всех тканей цепи «лансопразол всегда активирует ERK--Nrf2». Фармакологический блокатор пути также не доказывает, что лансопразол непосредственно связывает соответствующую киназу.\src{molho1,molnrf2}

\subsection{Воспалительные маркеры у человека: предел причинного вывода}
\label{sec:inflammatory-markers}
В исследовании слизистой пищевода пациентов с эндоскопически негативной ГЭРБ наблюдались усиление экспрессии IL-8 и признаки активации NF-$\kappa$B. В небольшой подгруппе из девяти пациентов после восьми недель лансопразола экспрессия IL-8 уменьшилась. Это ближе к клинической ткани, чем клеточная линия, но дизайн не отделяет прямое действие препарата на транскрипцию от вторичного уменьшения кислотного повреждения.\src{molil8}

Следовательно, фразы «лансопразол непосредственно блокирует NF-$\kappa$B у пациента» или «противовоспалительный эффект не зависит от кислотосупрессии» требуют отдельного доказательства. Изменение цитокинового маркера не равно установленной эффективности при системном воспалительном заболевании.\src{molil8}

\subsection{Лизосомы, аутофагия и противоопухолевые модели}
Желудочная P-типовая H$^+$/K$^+$-АТФаза и лизосомальная V-АТФаза --- разные молекулярные машины. Упоминание «протонной помпы» в опухолевой работе не дает права переносить туда карту Cys813/Cys321. В клеточных экспериментах Takeda и соавт. лансопразол в сочетании с макролидами усиливал гибель ряда опухолевых линий и изменения проницаемости лизосомальных мембран. Изучались концентрации вплоть до 250 мкмоль/л, а не обычный клинический режим в пересчете на пациента.\src{mollysosome}

Накопление аутофагосом/аутолизосом может означать нарушение их утилизации, а не усиление полноценного аутофагического потока. В указанной работе выключение ATG5 не устраняло цитотоксичность комбинации, поэтому результат нельзя безусловно назвать аутофагия-зависимой гибелью. Ни эти данные, ни увеличение лизосомального pH сами по себе не доказывают прямое специфическое связывание препарата с V-АТФазой и не обосновывают лечение рака такой комбинацией.\src{mollysosome}

\subsection{Сводная карта: что можно утверждать с разной уверенностью}
\begin{longtable}{P{0.25\textwidth}P{0.36\textwidth}P{0.30\textwidth}}
\toprule
Уровень & Установленное содержание & Недопустимая подмена \\
\midrule
\endhead
Основная мишень & Локальная активация и дисульфидная модификация люминально доступных цистеинов желудочной помпы.\src{molnagaya89,molcys} & «Любая протонная помпа организма блокируется одинаково». \\
Фармакодинамика & Эффект переживает плазменную экспозицию; возврат функции связан с восстановлением доступного пула помп.\src{molrestore} & «Восстановление возможно исключительно после синтеза нового белка». \\
Клеточная сигнализация & HO-1/Nrf2, ROS и лизосомальные процессы изменяются в отдельных экспериментальных моделях.\src{molnrf2,molho1,mollysosome} & «Для этих путей доказана самостоятельная клиническая польза». \\
Антимикробные эффекты & Уреазные, дыхательные и QcrB-связанные эффекты исследованы в разных бактериальных системах.\src{molurease,molresp,moltb} & «Обычная капсула заменяет антибактериальную комбинацию». \\
\bottomrule
\end{longtable}

\textbf{Итоговая причинная последовательность основного действия:}
\begin{center}
\fbox{\parbox{0.91\textwidth}{\centering
Защищенное пролекарство $\longrightarrow$ всасывание и системная доставка\\[3pt]
$\longrightarrow$ кислая секреторная микросреда и протонирование\\[3pt]
$\longrightarrow$ реакционноспособное серосодержащее производное\\[3pt]
$\longrightarrow$ дисульфидная модификация H$^+$/K$^+$-АТФазы\\[3pt]
$\longrightarrow$ уменьшение секреции H$^+$ и вторичные pH-зависимые эффекты.}}
\end{center}
Эта схема объединяет фармакокинетические, химические и биохимические данные; дополнительные клеточные механизмы не подменяют ее и не расширяют автоматически перечень показаний.\src{pharmalabel,molnagaya89,molnagaya90}

\section{Лекарственные формы и правильный прием}
\label{sec:formulations}

Основные пероральные формы --- гастрорезистентные капсулы 15 и 30 мг; в отдельных странах доступны диспергируемые во рту таблетки с защищенными микрогранулами. «Растворяется во рту» не означает, что гранулы следует разжевать, и не означает сублингвальное всасывание. Наличие формы в другой стране не подтверждает ее доступность или регистрацию в РФ.\src{uslabel,emc}

\subsection{Практические правила}
\begin{itemize}
\item При однократном режиме принимать обычно за 30--60 минут до завтрака; при двукратном --- до завтрака и вечернего приема пищи. При нестандартном режиме питания время согласуют с врачом.\src{acggerd}
\item Капсулу или микрогранулы не разжевывать. Возможность вскрытия капсулы и смешивания содержимого с пищей определяется инструкцией именно этой формы.\src{uslabel}
\item В американских инструкциях для некоторых капсул разрешены определенные мягкие продукты или соки; смесь проглатывают сразу. Самостоятельно переносить эту технологию на любой генерик нельзя.\src{uslabel}
\item При введении через зонд проверяют его диаметр, растворитель, объем промывания и допустимость метода. Диспергируемая таблетка и содержимое капсулы могут иметь разные инструкции.\src{emc,uslabel}
\item Пропущенную дозу не компенсируют удвоением следующей. Полный симптоматический эффект может формироваться несколько дней; американская безрецептурная маркировка указывает 1--4 дня.\src{otc,uslabel}
\end{itemize}

\subsection{Безрецептурный режим в США}
Американские OTC-препараты 15 мг предназначены для частой изжоги у взрослых, а не для самостоятельного лечения язвы или \Hp. Типовой курс --- 15 мг один раз в день перед едой в течение 14 дней; повторять чаще одного раза в четыре месяца без врачебной оценки не следует. Дисфагия, кровотечение, похудение, длительные или нарастающие симптомы требуют обследования. Этот режим не является правилом отпуска препарата в РФ.\src{otc}

\subsection{Торговые наименования и цены: примеры по регионам}
\label{sec:brands-prices}
\textbf{Дата проверки и границы выборки.} Ниже приведены конкретные препараты лансопразола и опубликованные ценовые ориентиры, проверенные при дополнении обзора \textbf{14 сентября 2026 года}. Для Европы выбран пример Франции; для Азии --- Японии и Индии; российские предложения относятся к московской витрине. Это не исчерпывающий перечень торговых наименований и не единая цена для США, Европы, Азии или РФ. Использованы официальные ценовые базы и страницы продавцов, в том числе доступные индексированные копии; дата обращения не гарантирует, что каждая торговая страница обновлена в тот же день. Цены не подтверждались оформлением заказа.\src{priceusprev,priceusequate,pricefrbio,pricejptake,priceinlanzol,pricerulancid}

\begin{longtable}{P{0.13\textwidth}P{0.30\textwidth}P{0.18\textwidth}P{0.28\textwidth}}
\toprule
Страна & Наименование, дозировка и количество & Цена & Основание цены \\
\midrule
\endhead
США & Prevacid 24HR; капсулы с отсроченным высвобождением, 15 мг, 42 шт. & 23,42 USD за упаковку & Онлайн-предложение Walmart; не страховая доплата.\src{priceusprev} \\
США & Equate Lansoprazole; капсулы с отсроченным высвобождением, 15 мг, 42 шт. (3 по 14) & 9,97 USD за упаковку & Предложение Walmart в доступной индексированной карточке; не цена за три упаковки по 42 шт.\src{priceusequate} \\
Франция & LANSOPRAZOLE BIOGARAN; таблетки, диспергируемые в полости рта, 30 мг, 28 шт. & 5,12 EUR за упаковку & Официальная база: 4,10 EUR за препарат и 1,02 EUR сбора за отпуск. Не сумма после возмещения.\src{pricefrbio} \\
Япония & Takepron OD; таблетки, диспергируемые в полости рта, 15 мг & 17,90 JPY за 1 таблетку & Официальная лекарственная цена MHLW; не розничная стоимость упаковки.\src{pricejptake} \\
Япония & Takepron OD; таблетки, диспергируемые в полости рта, 30 мг & 29,90 JPY за 1 таблетку & Официальная лекарственная цена MHLW; не индивидуальная доплата пациента.\src{pricejptake} \\
Индия & Lanzol 30 (Cipla); капсулы 30 мг, 10 шт. & 69,80 INR за упаковку; MRP 112,29 INR & Tata 1mg, витрина Gurgaon: показанная цена со скидкой, без дополнительных купонов; MRP --- указанная максимальная розничная цена.\src{priceinlanzol} \\
РФ, Москва & Ланцид (Микро Лабс Лимитед); капсулы 15 мг, 30 шт. & 426 руб. за упаковку & Apteka.ru: цена заказа, а не единая установленная цена для РФ.\src{pricerulancid} \\
РФ, Москва & Ланцид (Микро Лабс Лимитед); капсулы 30 мг, 30 шт. & 518 руб. за упаковку & Apteka.ru: цена заказа; наличие зависит от выбранной аптеки.\src{pricerulancid} \\
\bottomrule
\end{longtable}

\textbf{Как интерпретировать стоимость.} Розничное предложение, официальная лекарственная цена и собственные расходы пациента --- разные величины. Французская сумма включает указанный сбор за отпуск, но не является расчетом доплаты после страхового возмещения. Японская цена приведена за одну таблетку в официальном перечне; умножение на число таблеток не определяет итоговый счет пациента. Для торговых площадок доставка, дополнительные сборы, место получения и условия скидки проверяются отдельно; в индийской карточке применимые налоги включены в показанную цену. Пересчет валют и сравнение стоимости полного курса здесь не выполнялись.\src{pricefrbio,pricejptake,priceusprev,priceusequate,priceinlanzol,pricerulancid}

\textbf{Наименование и размер упаковки не заменяют назначения.} Таблица описывает монопрепараты лансопразола, а не сравнительную эффективность торговых марок. Количество 42 капсулы в американской упаковке не означает рекомендованный непрерывный 42-дневный курс самолечения: ограничения OTC-режима приведены выше. При выборе проверяют дозировку, форму, разрешенный способ введения и инструкцию конкретного продукта. Наличие российской аптечной карточки не заменяет проверку регистрационного статуса и действующей инструкции в ГРЛС; ранее сформулированное ограничение российского раздела сохраняется.\src{otc,uslabel,grls}

\chapter{Клиническая гастроэнтерология и эрадикационные протоколы}
\label{chap:gastroenterology}

\section*{Введение к главе}
\addcontentsline{toc}{section}{Введение к главе}
При ГЭРБ и язвенной болезни биохимическая последовательность «доставка --- кислотная активация --- ковалентная блокада» приобретает значение через конкретные клинические исходы: контроль симптомов, заживление повреждения и предупреждение рецидива. Липофильность исходной молекулы и ее ионизация участвуют в условиях доступа к секреторному канальцу, а связанные с $pK_a$ процессы протонирования --- в локальном накоплении и активации. Однако эти свойства не устраняют другие механизмы рефлюкса и повреждения слизистой. Поэтому сохранение жалоб нельзя автоматически интерпретировать как недостаточную силу связывания с помпой или основание для последовательного увеличения дозы.\src{molchem,molnagaya90,acggerd,agagerd}

Ковалентный характер ингибирования объясняет необходимость разграничивать экспозицию препарата и длительность подавления кислотообразования, но не заменяет проверки времени приема, соблюдения режима и диагноза. Для гастроэнтеролога практическая проблема заключается в том, чтобы соотнести требуемую кислотосупрессию с характером поражения и объективным ответом, не приравнивая уменьшение симптомов к эндоскопическому заживлению. Сравнение ИПП и P-CAB также требует клинических конечных точек: различия $pK_a$ или механизма связывания сами по себе не устанавливают превосходства по частоте заживления, поддержанию ремиссии либо исходам язвенного кровотечения.\src{molrestore,acggerd,agagerd,vontrial,esge2026}

В эрадикационных протоколах изменение кислотности служит компонентом комбинированного воздействия, а не самостоятельным устранением \Hp. Ни липофильность лансопразола, ни образование длительно действующего аддукта с желудочной помпой не доказывают преодоления антибиотикорезистентности. Клинический выбор должен учитывать предшествующее лечение, чувствительность возбудителя, достаточность кислотосупрессии и контроль эрадикации. Поэтому глава связывает биохимические предпосылки с дозовыми режимами и региональными протоколами, сохраняя различие между зарегистрированной комбинацией, историческим исследованием и рекомендацией общества; фармакогенетические основания персонализации рассматриваются далее в главе~\ref{chap:personalization}.\src{acghp,maastricht,leonova2015,cpic}

\section{Дозы у взрослых: инструкция и клинический контекст}
\label{sec:doses}
\textbf{Таблица справочная, не рецепт.} Длительность и доза зависят от эндоскопической картины, причины поражения, сопутствующих препаратов и локальной инструкции. «До 8 недель» не означает разрешение бесконтрольно продлевать лечение при отсутствии эффекта.

\begin{longtable}{P{0.31\textwidth}P{0.29\textwidth}P{0.31\textwidth}}
\toprule
Ситуация & Типовой режим лансопразола & Срок / оговорка \\
\midrule
\endhead
Симптоматическая ГЭРБ, инструкция США & 15 мг 1 раз/сут & До 8 недель \\
Эрозивный эзофагит & 30 мг 1 раз/сут & Обычно 4--8 недель; США: до 8 недель \\
Поддержание заживления эзофагита & 15 мг 1 раз/сут & Дозу и необходимость пересматривают \\
Язва ДПК, инструкция США & 15 мг 1 раз/сут & 4 недели \\
Язва ДПК, европейская SmPC & 30 мг 1 раз/сут & 2 недели, при необходимости еще 2 \\
Доброкачественная язва желудка & 30 мг 1 раз/сут & 4--8 недель \\
НПВП-ассоциированная язва & 30 мг 1 раз/сут & Обычно до 8 недель \\
Профилактика НПВП-язв & 15 мг 1 раз/сут & В ряде SmPC при недостаточности 30 мг/сут \\
Кислотосупрессия в эрадикации & Обычно 30 мг 2 раза/сут & Срок определяется всей схемой \\
Синдром Золлингера--Эллисона & Старт обычно 60 мг/сут & Дальнейшая индивидуальная титрация \\
\bottomrule
\end{longtable}
Регистрационные режимы сопоставлены по американской инструкции и европейской SmPC; для эрадикации приоритетно учитываются современные рекомендации, а не только исторические схемы вкладыша.\src{uslabel,hpra,acghp}

\textbf{Двукратный режим при ГЭРБ.} 30 мг дважды в сутки используется при отдельных клинических сценариях после оптимизации приема. NICE прямо обозначает этот режим для ГЭРБ как off-label. Наличие эрадикационного показания к той же дозе не делает ее автоматически зарегистрированной для любого заболевания.\src{nice}

\textbf{Высокие дозы при гиперсекреции.} В инструкциях описаны дозы до 180 мг/сут; при потребности более 120 мг/сут дозу делят. Это специализированное лечение, а не допустимый предел самостоятельного повышения дозы при изжоге.\src{emc}

\section{ГЭРБ: когда лансопразол уместен}
\label{sec:gerd}

\subsection{Начальная терапия и оценка ответа}
ACG рекомендует при типичной изжоге и регургитации без тревожных признаков пробный восьминедельный курс ИПП один раз в день до еды. AGA допускает начальный интервал 4--8 недель с последующей оценкой результата и снижением до минимально эффективной дозы при хорошем ответе. ИПП предпочтительнее H$_2$-блокаторов для заживления эрозивного эзофагита и поддержания заживления.\src{acggerd,agagerd}

Лечение включает не только таблетку: коррекцию избыточной массы тела, избегание поздней еды, при ночных симптомах --- приподнятое изголовье. Пищевые ограничения разумнее подбирать по индивидуальной воспроизводимости симптомов, чем назначать всем длинный список запретов.\src{acggerd,agagerd}

Ответ на ИПП сам по себе не доказывает ГЭРБ. DGVS отдельно подчеркивает различие между лечением рефлюксоподобных жалоб и подтвержденной болезни. При отсутствии ответа на правильно проведенный курс необходима дальнейшая диагностика.\src{dgvs}

\subsection{Неэрозивная болезнь, тяжелый эзофагит и поддержание}
При неэрозивной форме и легком эзофагите после достижения контроля возможны снижение дозы и, у подходящих пациентов, прием по требованию. При тяжелом эзофагите LA C/D, пептической стриктуре или других осложнениях обычно необходима непрерывная поддерживающая стратегия. Для пищевода Баррета кислотосупрессия не заменяет эндоскопическое наблюдение и лечение дисплазии.\src{agadepres,dgvs,nice}

Корейское обновление 2025 года допускает стандартную дозу ИПП или P-CAB в течение 4--8 недель при неэрозивной болезни / легком эзофагите и 8 недель при тяжелом эзофагите. Для легких фенотипов обсуждается режим по требованию; для тяжелых --- непрерывное поддержание. Авторы отмечают возможное преимущество P-CAB в заживлении тяжелых поражений.\src{seoul}

В японских рекомендациях 2021 года вонопразан занимает более заметное место при тяжелом эзофагите, тогда как при легком заболевании допустимы и ИПП. Это региональное различие терапевтической стратегии, а не доказательство непригодности лансопразола для японских или других азиатских пациентов.\src{jgerd}

\textbf{Генотип и заживление: исторические данные по лансопразолу.} Леонова приводит частоту заживления при ГЭРБ 45,8\% у экстенсивных и 84,6\% у медленных метаболизаторов CYP2C19 на фоне лансопразола; различие описано как статистически значимое. Это сравнение фенотипических подгрупп при одном препарате, а не прямое доказательство превосходства рабепразола над лансопразолом. Доза и срок оценки именно этого результата в изложении статьи не уточнены, поэтому проценты нельзя автоматически относить к стандартному восьминедельному курсу 30 мг/сут. Приведенное там же отсутствие выявленной зависимости ответа на рабепразол от генотипа не доказывает эквивалентности всех ИПП. Термин «экстенсивный» в этих работах требует отдельного сопоставления с CPIC (раздел~\ref{sec:pharmacogenetics}).\src{leonova2015}

\subsection{Пищевод Баррета: обновление AGA 2025}
AGA в руководстве от 20 октября 2025 года условно рекомендует ежедневную терапию ИПП взрослым с пищеводом Баррета для снижения риска неопластического прогрессирования по сравнению с отсутствием ИПП. Это рекомендация класса, а не доказательство специального противоопухолевого эффекта лансопразола. Отсутствие изжоги не означает отсутствия показаний; ИПП не заменяет наблюдение и лечение дисплазии. Антирефлюксную операцию не следует выбирать вместо ИПП только как способ профилактики рака.\src{barrett2025}

\subsection{Что означает «рефрактерная ГЭРБ»}
\label{sec:refractory-gerd}
Практический алгоритм, синтезированный из AGA, DGVS и Lyon 2.0:\src{agagerd,dgvs,lyon}
\begin{enumerate}
\item Проверить фактическую дозу, время приема и пропуски, а также ведущую жалобу: изжога, регургитация, дисфагия и боль требуют разных решений.
\item Убедиться, что диагноз обоснован. Возможны функциональная изжога, рефлюксная гиперчувствительность, эозинофильный эзофагит, нарушение моторики, руминация и другие причины.
\item При неподтвержденной ГЭРБ рассмотреть эндоскопию и мониторирование рефлюкса вне кислотосупрессии; срок отмены зависит от исследования.
\item При доказанной ГЭРБ и сохраняющихся симптомах на оптимизированной терапии использовать pH-импедансометрию на лечении по показаниям.
\item Решать вопрос о двукратном режиме, смене препарата, P-CAB или антирефлюксном вмешательстве только после уточнения механизма жалоб.
\end{enumerate}

Lyon 2.0 рассматривает эзофагит LA B/C/D, подтвержденный пищевод Баррета и пептическую стриктуру как убедительные эндоскопические свидетельства ГЭРБ. Один лишь LA A менее специфичен. Отсутствие эрозий не исключает ГЭРБ; нормальная кислотная экспозиция с отсутствием связи симптомов и рефлюкса, напротив, заставляет искать иной механизм.\src{lyon}

\section{Язвенная болезнь и профилактика повреждения ЖКТ}
\label{sec:ulcer-disease}

\subsection{Лечение язвы и устранение причины}
Кислотосупрессия способствует заживлению, но необходимо установить причину: \Hp, НПВП/аспирин, другие лекарственные воздействия, реже гиперсекреторное состояние или иной диагноз. При язве желудка исключают злокачественность и планируют контроль заживления по клинической ситуации. При незаживающей язве проверяют приверженность, пропущенную инфекцию, продолжающийся прием НПВП и альтернативные причины.\src{nice}

После успешной эрадикации и устранения ульцерогенного фактора пожизненный ИПП нужен не каждому. При продолжающемся высоком риске кровотечения гастропротекция может оставаться необходимой даже после удаления \Hp.\src{aspirintrial,agadepres}

\subsection{Язвенное кровотечение}
\label{sec:ulcer-bleeding}
Это стационарная ситуация, а не повод самостоятельно принять увеличенную дозу капсул. ESGE в обновлении 2026 года рассматривает раннюю эндоскопическую оценку, эндоскопический гемостаз и интенсивную кислотосупрессию в единой стратегии. Предэндоскопический внутривенный ИПП не должен задерживать эндоскопию.\src{esge2026}

После гемостаза язвы высокого риска применяют высокодозный режим ИПП; среди подходов --- внутривенный болюс с инфузией, прерывистое внутривенное или подходящее пероральное введение. Конкретный препарат, дозу и путь определяет стационарный протокол; \textbf{обычный амбулаторный прием 15--30 мг/сут не равнозначен такой терапии}.\src{esge2026}

Рекомендации российских профессиональных обществ по язвенной болезни 2024 года отдельно описывают внутривенный лансопразол: 30 мг дважды в сутки в виде 30-минутной инфузии, краткосрочно до 7 дней, с переходом на пероральную форму. Это описание не подтверждает доступность любой соответствующей торговой формы сегодня; необходима проверка ее регистрации и инструкции.\src{rfpud,grls}

\section{Эрадикация Helicobacter pylori}
\label{sec:eradication}

\subsection{Роль лансопразола}
Сравнение с P-CAB, генотипическая оптимизация и экономические ограничения подробно рассмотрены в разделе~\ref{sec:hp-pcab-parity}.
Кислотосупрессия обеспечивает условия для работы антибиотиков, но не компенсирует устойчивость бактерии к неэффективному препарату. Типовая доза лансопразола в ИПП-содержащей комбинации --- 30 мг дважды в сутки; конкретная кратность и интенсивность могут отличаться в специальных схемах. Схему выбирают целиком, учитывая аллергию, предшествующие антибиотики, прошлые попытки лечения и локальные данные эффективности.\src{maastricht,cpic,nice}

\textbf{Фармакогенетически ориентированная эрадикация: что добавляет обзор 2015 года.} В изложении Леоновой японское исследование у 551 пациента сопоставляло стандартную тройную схему на основе лансопразола 30 мг дважды в сутки с подходом, учитывающим генотип и предусматривающим рабепразол 10 мг дважды в сутки у выделенных пациентов. Общая частота эрадикации составляла 80\% и 88,7\% соответственно; при приведенном $p=0{,}078$ статистически значимое преимущество по обычному порогу 0,05 не установлено. Это сравнение лечебных стратегий, а не изолированный эффект замены одной молекулы другой при полностью одинаковых условиях. Описанные в статье дробные режимы, включая рабепразол 10 мг четыре раза в сутки, относятся к историческим исследованиям и не являются универсальным назначением после любой неудачи эрадикации.\src{leonova2015}

Фармакогенетический вклад особенно заметен в приведенных исследованиях омепразола и лансопразола, однако результаты разных метаанализов неоднородны: в таблице~7 статьи различие между медленными и экстенсивными метаболизаторами для лансопразола отмечено как значимое только в одном из двух анализов. Отсутствие такого сигнала для рабепразола не означает, что препарат преодолевает любую антибиотикорезистентность. Генотипический выбор кислотосупрессии дополняет, но не заменяет учет чувствительности \Hp, предшествующих антибиотиков и проверку результата лечения. Исторические кларитромициновые комбинации не подменяют изложенную далее стратегию ACG 2024.\src{leonova2015,acghp,maastricht}

\subsection{США: ACG 2024}
Выбор эрадикационного протокола определяется чувствительностью \Hp, предшествующей антибактериальной терапией, переносимостью компонентов и доказанной эффективностью полной схемы. При неизвестной чувствительности ACG 2024 рекомендует для ранее не леченных пациентов \textbf{оптимизированную висмутовую квадротерапию продолжительностью 14 дней}: ИПП дважды в сутки, висмут четыре раза в сутки, тетрациклин 500 мг четыре раза в сутки и метронидазол 500 мг три или четыре раза в сутки. Конкретная доза висмута зависит от соли и препарата; доксициклин не рассматривается как автоматическая замена тетрациклина. Термин «оптимизированная» относится к совокупности компонентов, доз, кратности и длительности, а не к изолированному увеличению дозы лансопразола.\src{acghp}

Альтернативы для пациентов без аллергии на пенициллины включают предусмотренные рекомендациями рифабутиновые комбинации и двойную терапию вонопразан--амоксициллин. Назначение кларитромициновых и левофлоксациновых режимов требует подтвержденной чувствительности; усиление кислотосупрессии не заменяет этого условия. \textbf{Доказательства для двойной схемы с вонопразаном нельзя переносить на произвольное сочетание лансопразола с амоксициллином.} Исторические зарегистрированные режимы лансопразола и современные предпочтения ACG сопоставляют раздельно.\src{acghp,uslabel}

Фармакогенетическая оптимизация является следующим уровнем выбора внутри клинически обоснованной схемы. Она должна устранять потенциально недостаточную экспозицию ИПП, не изменяя доказательных требований к антибиотикам и контролю излечения. После завершения лечения результат подтверждают валидированным тестом не ранее четырех недель после антибиотиков; ИПП и P-CAB отменяют за две недели до исследования. Уменьшение симптомов либо повышение внутрижелудочного pH не являются критериями эрадикации.\src{acghp}

\addtocontents{toc}{\protect\setcounter{tocdepth}{3}}
\subsubsection{Молекулярно-генетические механизмы преодоления резистентности Helicobacter pylori при полиморфизмах CYP2C19}
\label{sec:hp-exposure-resistance}
\addtocontents{toc}{\protect\setcounter{tocdepth}{2}}

\textbf{Фармакогенетическая предпосылка и количественная модель компенсации экспозиции.} Полиморфизм CYP2C19 влияет на метаболический клиренс лансопразола и вероятность недостаточной кислотосупрессии, но не изменяет непосредственно бактериальные мишени антибиотиков. Поэтому термин «преодоление резистентности» в данном подразделе обозначает устранение фармакокинетического фактора неудачи, а не восстановление чувствительности генетически резистентного штамма. Фенотип быстрого метаболизма RM (*1/*17) отличают от ультрабыстрого UM (*17/*17), нормального NM и промежуточного IM; историческая категория экстенсивных метаболизаторов не может автоматически приравниваться к UM.\src{cpic,acghp} В линейном приближении при повторном приеме $\mathrm{AUC}_{24,ss}\approx F D_{24}/CL$, где $F$ --- биодоступность, $D_{24}$ --- суточная доза, $CL$ --- клиренс. Если клиренс возрастает в $k$ раз при неизменном $F$, сохранение прежней экспозиции требует увеличения суточной дозы приблизительно в $k$ раз. Таким образом, удвоение дозы теоретически компенсирует двукратное увеличение клиренса, но не любое снижение AUC. Это расчетная модель, а не индивидуальная гарантия: результат зависит от абсорбции, пресистемного метаболизма, функции печени, взаимодействий и отклонений от линейности.\src{molp450,uslabel} Снижение экспозиции нельзя смешивать с недостаточной приверженностью или устойчивостью микроорганизма; каждый из этих факторов требует отдельного вмешательства.\src{acghp,cpic}

\textbf{Кинетическая интенсификация дозы и кратности приема.} CPIC рекомендует для UM увеличение начальной суточной дозы на 100\%; для RM при эрадикации \Hp\ допускается рассмотрение увеличения на 50--100\% с разделением дозы и контролем эффективности. Режимы 60 мг дважды в сутки и 30 мг четыре раза в сутки дают одинаковые 120 мг/сут: это удвоение по отношению к 30 мг дважды в сутки, но четырехкратное увеличение относительно 30 мг один раз в сутки. Эти арифметические соотношения не превращают 120 мг/сут в обязательное назначение каждому носителю *17.\src{cpic,uslabel} Гастрорезистентная форма сохраняет кислотолабильное пролекарство до кишечного высвобождения; после всасывания препарат поступает к париетальной клетке и активируется в кислом секреторном компартменте. Увеличение общей дозы повышает доступную экспозицию, тогда как ее дробление создает повторные периоды поступления ингибитора к помпам, функционально доступным в разные моменты суток. При одинаковой суточной дозе более частый прием не обязан увеличивать AUC: его предполагаемое преимущество относится прежде всего к временному профилю взаимодействия с мишенью.\src{uslabel,molrestore,hpphblum1998,hpphfuruta2001} Исследования кислотосупрессии поддерживают обоснованность оптимизации, но не доказывают универсальную эквивалентность двух указанных режимов или полную нормализацию экспозиции у UM. Последовательность клинического решения включает проверку чувствительности и предшествующих курсов, соблюдения приема и питания, генотипа при наличии результата, функции печени и взаимодействий; затем врач выбирает антисекреторный режим внутри доказанной полной схемы. Высокодозовая интенсификация не заменяет выбора антибиотиков и требует сверки с инструкцией конкретного препарата.\src{acghp,cpic,uslabel}

\textbf{Средовой контекст репликации и антибактериального действия.} Уменьшение кислотного стресса создает условия для более активного роста \Hp\ и реализации действия антибиотиков, зависящего от физиологического состояния бактерии. Для амоксициллина существенна активность синтеза пептидогликана: блокада пенициллинсвязывающих белков эффективнее проявляется при формировании клеточной стенки растущего микроорганизма. При этом максимальную эффективность нельзя свести к единственному обязательному порогу pH. В эксперименте Marcus и соавт. ампициллин проявлял бактерицидность при pH 4,5 и 7,4, но не при pH 3,0; это механистическое подтверждение влияния среды, а не испытание амоксициллина в конкретном эрадикационном режиме и не доказательство универсального перехода всей популяции в деление при pH выше 5 или 6.\src{hpgrowth2012} Поддержание менее кислой среды также уменьшает кислотную деградацию антибиотиков, особенно кларитромицина; различия их химической стабильности необходимо учитывать раздельно. Кларитромицин действует на бактериальную рибосому, и повышение pH не устраняет устойчивость, связанную с изменением его мишени.\src{hpstability1997,acghp} Поэтому доли суток с pH выше 5,0 и выше 6,0 следует трактовать как разные фармакодинамические показатели, а не взаимозаменяемую запись одного порога. Среднесуточный pH не доказывает непрерывного удержания целевого уровня и не измеряет непосредственно pH микросреды слизистого слоя. Увеличение времени эффективной кислотосупрессии поддерживает антибиотическое действие, но результат дополнительно зависит от чувствительности штамма, экспозиции антибиотиков, длительности курса и соблюдения назначений.\src{hpphblum1998,hpphfuruta2001,hpgrowth2012,acghp}

\textbf{ACG 2024, JSHR 2026 и пределы доказанного терапевтического паритета.} ACG 2024 при неизвестной чувствительности рекомендует 14-дневную оптимизированную висмутовую квадротерапию; среди альтернатив предусматривается двойная терапия вонопразан--амоксициллин. Англоязычный документ JSHR, опубликованный 28 марта 2026 года, представляет редакцию японского руководства 2024 года и отдает предпочтение вонопразан-содержащим схемам первой линии. Ни один из этих документов не устанавливает универсальный генотип-специфичный паритет эрадикации выше 90\% между лансопразолом 120 мг/сут и вонопразаном у NM/IM.\src{acghp,jhp} Высокая результативность отдельных высокодозовых лансопразол-содержащих схем в условиях низкой устойчивости или подбора терапии по чувствительности подтверждает возможность эффективного применения, но не является прямым доказательством равенства P-CAB.\src{hphighdose2013} В PHALCON-HP при чувствительных к амоксициллину и кларитромицину штаммах частота эрадикации составляла 84,7\% для тройной схемы с вонопразаном, 78,5\% для двойной схемы с ним и 78,8\% для тройной схемы с лансопразолом 30 мг дважды в сутки; в общей популяции схемы с вонопразаном превосходили лансопразол-содержащую тройную терапию. Испытание не оценивало лансопразол 120 мг/сут и не доказывало паритет именно у NM/IM.\src{hpphalcon2022} Для утверждения о равной эффективности необходимы прямое сопоставление полных схем, заранее установленная граница клинически допустимого различия, учет чувствительности штаммов и CYP2C19-фенотипов, а также согласованная интерпретация анализов всех рандомизированных пациентов и завершивших протокол. Показатель выше 90\% в отдельной группе сам по себе этих требований не выполняет.\src{hpphalcon2022,acghp}

\textbf{Клинико-экономическое заключение и критерий завершенного лечения.} Нормальный или промежуточный CYP2C19-фенотип сам по себе не требует перехода с лансопразола на P-CAB, но также не доказывает достаточности любой лансопразол-содержащей схемы. Рациональным основанием выбора служит сочетание подтвержденной эффективности полного режима в соответствующей популяции, переносимости, доступности и совокупных затрат; сохранение доступного эффективного лечения обоснованно, если клинические условия его применения соблюдены.\src{cpic,acghp,jhp} Для экономической оценки необходимо заранее определить перспективу плательщика и временной горизонт, учесть стоимость всех компонентов курса, диагностики, контроля излечения, нежелательных явлений и повторной терапии после неудачи. Сопоставление цены упаковок или только стоимости ИПП не устанавливает ни экономического паритета, ни бюджетной экономии. Японский анализ страховых данных показывает, что более высокая первичная эффективность вонопразан-содержащих режимов могла сопровождаться меньшими суммарными затратами; его региональный наблюдательный дизайн не позволяет переносить этот вывод на все системы здравоохранения или специально на NM/IM.\src{hpcost2021} Следовательно, предотвращение необоснованного удорожания означает выбор эффективной схемы с приемлемой полной стоимостью, а не запрет перехода на P-CAB или объявление такого перехода заведомо нерациональным. Условия сравнительного анализа изложены в разделе~\ref{sec:hp-pcab-parity}. Эрадикацию подтверждают валидированным тестом не ранее четырех недель после завершения антибиотиков, с отменой ИПП и P-CAB за две недели до тестирования; улучшение симптомов и достижение целевого pH не заменяют подтверждения излечения.\src{acghp,hpcost2021}

\subsection{Европа, Азия и РФ: реальные различия}
\begin{longtable}{P{0.20\textwidth}P{0.73\textwidth}}
\toprule
Документ & Стратегия и значение для лансопразола \\
\midrule
\endhead
Maastricht VI, 2022 & При высокой ($>15\%$) или неизвестной устойчивости к кларитромицину предпочтительна висмутовая квадротерапия. Обычно 14 дней; сокращение возможно только при доказанной локальной эффективности соответствующего режима. Если висмутовая схема недоступна, рассматривают альтернативы с учетом двойной устойчивости.\src{maastricht} \\
Япония, JSHR 2024 / публикация 2026 & Более значимая роль вонопразана и семидневных комбинаций. Рекомендуется выбор по чувствительности; для второй линии после амоксициллина/кларитромицина используются семидневные комбинации амоксициллин--метронидазол с P-CAB или ИПП. Японские короткие схемы не следует автоматически переносить в США или РФ.\src{jhp} \\
Корея, редакция 2025 / публикация 2026 & Две основные стратегии: лечение по чувствительности и эмпирические четырехкомпонентные схемы. Рекомендована 10-дневная сочетанная терапия без висмута; висмутовая квадротерапия 10--14 дней --- условный вариант первой линии, часто сохраняемый для последующей. P-CAB допускаются как альтернатива ИПП.\src{khp} \\
Китай, руководство 2022 & Висмутовые четырехкомпонентные режимы предпочтительнее стандартной тройной терапии; антибиотики подбирают с учетом предыдущего приема и устойчивости. Применяется также концепция высокодозной двойной терапии. Это содержание документа 2022 года, а не подтвержденный пересказ новых схем 2026 года.\src{china2022,china2026} \\
РФ, язвенная болезнь 2024 & Описаны 14-дневные варианты: тройная схема, усиленная висмутом; классическая висмутовая квадротерапия; сочетанная схема без висмута. Простая тройная терапия допускается там, где подтверждена ее эффективность. Это отличается от более строгого американского требования доказанной чувствительности.\src{rfpud,acghp} \\
\bottomrule
\end{longtable}

Корейская редакция конкретизирует персонализированный подход: при чувствительности к кларитромицину возможна семидневная тройная схема, при устойчивости --- десятидневная висмутовая; эмпирическая тройная терапия ограничена особыми условиями. После двух неудач предлагается направление в специализированный центр с возможностью определения чувствительности.\src{khp}

\textbf{JSHR: область установленного преимущества и область индивидуализации.} В CQ2-2 приведены результаты метаанализа первой линии: эрадикация в ITT-анализе 87,6\% при вонопразан-содержащих против 70,8\% при ИПП-содержащих схемах; при устойчивости к кларитромицину --- 75\% против 38\%. Это не сравнение с лансопразолом 120 мг/сут у UM. Для второй линии с амоксициллином и метронидазолом результаты сопоставимы. CQ2-3/2-4 дополнительно поддерживают подбор антибиотиков по чувствительности, а не сохранение неэффективного кларитромицина посредством усиления кислотосупрессии. Руководство допускает экономическое преимущество вонопразана за счет меньшей потребности в повторном лечении; следовательно, оно не обосновывает безусловный запрет перехода на P-CAB.\src{jhp}

\subsection{РФ: уточнение по утвержденным КР, а не только публикациям обществ}
В доступной републикации КР «Язвенная болезнь» 2024 года, ID 277\_2, для базисного лечения язвы указана доза лансопразола 30 мг/сут. Это не то же самое, что доза кислотосупрессивного компонента эрадикации, для которого важна отдельная схема комбинации.\src{rfpudmin}

В том же документе стандартная тройная эрадикация на 14 дней допускается \textbf{в регионах с подтвержденной эффективностью}; приведены также варианты усиления висмутом и четырехкомпонентной терапии. Поэтому формулировка «российские рекомендации разрешают тройную терапию» без условий неполна. Эта оговорка отличается от позиции ACG 2024, требующей подтвержденной чувствительности для кларитромициновой тройной схемы.\src{rfpudmin,acghp}

Найдена прямая ссылка на PDF КР 277\_2 в официальном API Минздрава; клинические формулировки для этого дополнения сверены по доступной републикации. Наличие ссылки и года 2024 не подтверждает автоматически отсутствие последующей редакции на дату конкретного назначения. Проверка всего Рубрикатора и всех инструкций ГРЛС остается неполной.\src{rfpudmin,registry,grls}

\subsection{Индия: локальная устойчивость и документы разного уровня}
\textbf{Bhubaneswar Consensus ISG, 2021.} Выбор эрадикации связывается с местной антибиотикорезистентностью. Тройная комбинация ИПП, амоксициллина и кларитромицина предусмотрена при низкой устойчивости к кларитромицину; для тройной, сочетанной, гибридной и четырехкомпонентной терапии рекомендуется 14-дневная продолжительность. В документе отдельно предостерегают от имидазол-содержащей \emph{тройной} схемы в индийских условиях; это не равнозначно запрету метронидазола во всех висмутовых квадросхемах. После неудачи второй линии при планировании эндоскопии рекомендуется культуральное исследование с определением чувствительности. Старые региональные проценты устойчивости нельзя выдавать за актуальную чувствительность конкретного пациента в 2026 году.\src{indiahp}

\textbf{Экспертный консенсус от 5 марта 2026 года.} Society for Clinical Gastroenterology опубликовало 20 положений по кислотозависимым заболеваниям. При типичной ГЭРБ без тревожных признаков рассматривается ИПП один раз в сутки на 4--8 недель, с последующим снижением интенсивности при ответе и отсутствии показаний к поддержанию; P-CAB обсуждаются как альтернатива. Это экспертный консенсус другого общества, а не обозначенная авторами новая редакция Bhubaneswar Consensus. Раскрыты финансирование Alembic Pharmaceuticals и работа двух авторов в этой компании. Эти сведения важны для критической оценки, но сами по себе не опровергают выводы.\src{india2026}

\textbf{Вывод для лансопразола.} Индийские документы не создают единой «азиатской схемы» и не дают основания переносить местный выбор антибиотиков в США или РФ. Подходы к классу ИПП следует сочетать с инструкцией конкретного препарата и результатами локального лечения. Это сравнительный вывод настоящего обзора.\src{indiahp,india2026,acghp,rfpudmin}

\subsection{Что проверить до начала курса}
Практический контрольный список, вытекающий из региональных рекомендаций:\src{acghp,maastricht,khp}
\begin{itemize}
\item Подтверждена ли активная инфекция, а не только давняя положительная серология?
\item Были ли макролиды, фторхинолоны и предыдущие курсы эрадикации?
\item Истинна ли заявленная аллергия на пенициллины и возможно ли ее уточнение?
\item Нет ли противопоказаний к антибиотикам, беременности, значимой почечной/печеночной дисфункции или опасных взаимодействий?
\item Доступны ли все компоненты на полный срок и понятен ли пациенту график приема?
\end{itemize}

\subsection{Обязательный контроль результата}
Успех лечения оценивают не по исчезновению боли. Используют уреазный дыхательный тест, антиген \Hp\ в кале или соответствующий биопсийный метод. Обычно исследование выполняют не ранее четырех недель после антибиотиков; ИПП отменяют на две недели, антибиотики и висмут --- на четыре недели, чтобы уменьшить риск ложноотрицательного результата. Серология не подходит для подтверждения эрадикации.\src{acghighlights,maastricht}

Отмена кислотосупрессии перед тестом должна быть безопасной. При необходимости обсуждают временные средства для симптомов и точные требования лаборатории. После неудачи выбирают новую схему с учетом уже использованных антибиотиков, а не просто повторяют тот же набор.\src{acghighlights,khp}

\subsection{Эрадикация у детей: отдельный алгоритм}
\label{sec:hp-children}
Совместные ESPGHAN/NASPGHAN-рекомендации редакции 2023 года, опубликованные в 2024 году, не поддерживают перенос взрослой стратегии «тестировать и лечить» на детей с любыми неспецифическими болями в животе. Предпочтителен выбор антибиотиков по чувствительности; кларитромицин без ее определения не используют. Для соответствующих схем рекомендуются 14 дней, а результат оценивают через 6--8 недель дыхательным тестом или подходящим моноклональным анализом антигена в кале. Перед контролем выдерживают не менее двух недель без ИПП и четырех недель без антибиотиков и висмута. Дозу лансопразола и допустимость его применения определяют отдельно по массе, показанию и местной инструкции.\src{pedhp2024}

\subsection{Лансопразол в эпоху P-CAB: фармакогенетическая оптимизация и пределы экономического паритета}
\label{sec:hp-pcab-parity}
\textbf{Предмет сравнения и версии руководств.} Далее рассматриваются взрослые пациенты; предметом оценки является подтвержденная эрадикация, а не выраженность кислотосупрессии сама по себе. ACG 2024 сохраняет оптимизированную висмутовую квадротерапию на 14 дней как предпочтительный эмпирический режим при неизвестной чувствительности и рассматривает вонопразан--амоксициллин как альтернативу у пациентов без аллергии на пенициллины. Поэтому появление P-CAB не исключает лансопразол из ИПП-содержащих схем. Однако сохранение терапевтической позиции не означает доказанной эквивалентности любому режиму с вонопразаном; кларитромициновые комбинации в американском алгоритме требуют подтверждения чувствительности.\src{acghp,acghighlights}

Обозначение «JSHR 2026» здесь относится к англоязычной публикации \emph{редакции японского руководства 2024 года}: онлайн 28 марта 2026 года, журнальный выпуск --- июль 2026 года. В CQ2-2 японский документ отдает предпочтение первой линии вонопразан--амоксициллин--кларитромицин; в приведенном анализе ITT эрадикация составляла 87,6\% против 70,8\% для ИПП-содержащих комбинаций. Для второй линии с амоксициллином и метронидазолом допускаются как ИПП, так и вонопразан, а результаты представлены как сопоставимые. Это региональное разграничение линий терапии, но не отдельное доказательство паритета высокодозового лансопразола и вонопразана у носителей конкретных генотипов. Японские семидневные режимы нельзя автоматически переносить в алгоритм ACG.\src{jhp}

\textbf{NM и IM: разные основания для оптимизации.} Нормальные метаболизаторы (NM; типичный диплотип CYP2C19 *1/*1) не тождественны промежуточным (IM; например, *1/*2 или *1/*3). При лансопразоле снижение метаболической активности у IM обычно увеличивает экспозицию относительно NM; нормальная активность фермента сама по себе не гарантирует достаточной кислотосупрессии. Исторические категории homozygous extensive и heterozygous extensive необходимо соотносить с конкретными исследованными аллелями, а не объединять в группу с одинаковым ответом. Вонопразан существенно меньше зависит от CYP2C19, но это фармакологическое различие не определяет исход независимо от антибиотиков.\src{cpic,jhp}

CPIC рекомендует для NM стандартную начальную дозу ИПП с возможностью увеличения на 50--100\% при эрадикации \Hp; для IM исходно предусмотрена стандартная доза, а не обязательное генотипическое удвоение. Следовательно, единое правило «высокая доза всем NM и IM» из CPIC не следует. Режим лансопразола 30 мг дважды в сутки уже является обычным компонентом ряда эрадикационных протоколов: его нельзя одновременно считать исходным режимом и результатом повторного удвоения без указания базы расчета. Решение об интенсификации согласуют с полной схемой, взаимодействиями и применимой инструкцией; подробная интерпретация фенотипов приведена в разделе~\ref{sec:pharmacogenetics}.\src{cpic,hpphalcon2022}

\textbf{Тройная и квадротерапия: неодинаковая доказательная база повышения дозы.} В положении 8 лечебного раздела Maastricht VI высокодозовый ИПП дважды в сутки связывается с повышением эффективности тройной терапии; для четырехкомпонентных режимов дополнительная польза такой интенсификации остается менее определенной. Поэтому результат нельзя переносить с тройной схемы на любую висмутовую либо невисмутовую квадротерапию. Оптимизация включает также продолжительность, адекватные антибиотики и приверженность. Приведенное в консенсусе удержание внутрижелудочного pH выше заданного уровня более 90\% времени является фармакодинамическим показателем, а не частотой эрадикации более 90\%. Целевой уровень излечения должен подтверждаться для полной схемы в соответствующей популяции.\src{maastricht}

\textbf{Достижимость результата выше 90\%: положительные данные и их границы.} В рандомизированном исследовании Furuta и соавт. (2007; 300 пациентов) персонализированная стратегия учитывала одновременно CYP2C19 и бактериальную чувствительность к кларитромицину; режимы лансопразола предварительно оценивали по внутрижелудочному pH. В первой линии результат ITT составил 144/150, или 96,0\% (95\% ДИ 91,5--98,2), против 105/150, или 70,0\% (95\% ДИ 62,2--77,2), при стандартной недельной тройной схеме. Это подтверждает потенциал индивидуализированной стратегии, но не устанавливает, что любое повышение дозы обеспечивает более 90\% отдельно у NM и IM: одновременно изменялись несколько лечебных решений, а опубликованный общий показатель не является результатом каждой генотипической подгруппы. Вонопразан и современные оптимизированные квадросхемы в этом сравнении не изучались.\src{hptailored2007}

\textbf{Прямое сравнение с вонопразаном: PHALCON-HP.} В исследовании Chey и соавт. (2022; 1046 рандомизированных пациентов США и Европы) 14-дневная тройная терапия включала лансопразол 30 мг или вонопразан 20 мг дважды в сутки с одинаковыми дозами амоксициллина и кларитромицина; отдельно изучалась двойная схема вонопразан--амоксициллин. Среди пациентов без устойчивости к кларитромицину и амоксициллину эрадикация составила 78,8\% при лансопразоле, 84,7\% при тройной и 78,5\% при двойной схеме с вонопразаном. Для тройных схем разность в пользу вонопразана равнялась 5,9 процентного пункта (95\% ДИ от $-0{,}8$ до 12,6). В общей популяции соответствующие показатели составляли 68,5\%, 80,8\% и 77,2\%. Исследование показало не меньшую эффективность режимов вонопразана в первичном анализе, но не доказало двустороннюю эквивалентность препаратов, генотипический паритет NM/IM или результат выше 90\% при оптимизированном высокодозовом лансопразоле.\src{hpphalcon2022}

\textbf{Экономический анализ: стоимость успешного лечения, а не одной упаковки.} В исследовании Furuta расчетные конечные затраты на одну успешную эрадикацию составляли 669 USD для персонализированной и 657 USD для стандартной стратегии. Эти исторические величины отражают конкретную модель и не доказывают экономию относительно вонопразана или бюджета 2026 года. Они показывают, почему улучшение результата первой линии следует оценивать вместе с затратами на тестирование и повторное лечение.\src{hptailored2007}

В японском анализе страховых данных Tokunaga и соавт. (2021; данные 2015--2018 годов), после сопоставления по propensity score по 22\,002 пациента в каждой группе, расчетные суммарные затраты первой и второй линий были ниже при P-CAB, чем при ИПП: 12\,952 против 13\,146 JPY. Оцененная авторами успешность первой линии составляла 93,6\% против 79,7\%. Это наблюдательное сравнение, а не генотипически стратифицированное испытание высокодозового лансопразола; данные страховых требований и расчет затрат не заменяют проспективного подтверждения эрадикации. Кроме того, три соавтора являлись сотрудниками Takeda. Результат не переносится непосредственно на другие страны, но опровергает предположение, что более высокая цена антисекреторного компонента неизбежно означает большие итоговые расходы.\src{hpcost2021}

\textbf{Условный экономический выбор вместо недоказанного полного паритета.} Авторский вывод настоящей монографии состоит в следующем: лансопразол сохраняет обоснованную позицию в эффективных, доступных и правильно подобранных схемах; его генотипическая зависимость допускает оптимизацию, а не требует автоматического отказа от молекулы. Если для конкретного протокола подтвержден высокий локальный результат и меньшая полная стоимость, его использование может быть экономически рациональным. Но для заявления о равном результате с вонопразаном необходимы прямое сопоставление оптимизированных схем, достаточная численность заранее выделенных генотипических групп, включая UM *17/*17, заранее заданная граница не меньшей эффективности либо эквивалентности в зависимости от цели исследования и учет устойчивости. Без этих данных «полный терапевтический паритет» и гарантированное отсутствие дополнительного бюджетного бремени остаются гипотезами. Экономическая оценка должна учитывать весь курс, генотипирование при его использовании, контроль излечения, нежелательные явления и повторные попытки; цена упаковки и генотип по отдельности недостаточны.\src{acghp,maastricht,cpic,hptailored2007,hpphalcon2022,hpcost2021}

\section{Другие показания и границы применения}

\subsection{Гиперсекреторные состояния}
Синдром Золлингера--Эллисона требует индивидуальной титрации и часто длительного лечения. Лансопразол контролирует кислотную гиперсекрецию, но не лечит саму гастриному; необходимы специализированная диагностика и лечение опухоли. Самостоятельная отмена ИПП ради исследования гастрина при подозрении на выраженную гиперсекрецию небезопасна.\src{uslabel,emc}

\subsection{Функциональная диспепсия и «гастрит»}
\label{sec:dyspepsia}
Ограниченный курс ИПП допустим у части пациентов с диспепсией после подходящей оценки и решения вопроса о \Hp. При отсутствии ответа следует пересмотреть диагноз и лечебную стратегию, а не считать любую диспепсию проявлением «избытка кислоты».\src{nice}

Российские КР «Гастрит и дуоденит» 2024 рассматривают ИПП на 4--6 недель для заживления эрозий, в том числе связанных с НПВП. При \Hp-ассоциированном гастрите этиотропная задача --- эрадикация. Изолированная эндоскопическая формулировка «поверхностный гастрит» без клинического контекста не обосновывает бессрочную кислотосупрессию; при атрофии и аутоиммунном процессе важны иные диагностические и профилактические задачи.\src{rfgastrit}

\textbf{Уточнение BSG 2022.} Британское руководство по функциональной диспепсии поддерживает ИПП как эффективный вариант лечения (сильная рекомендация, высокое качество доказательств), но не находит убедительной зависимости эффекта от увеличения дозы. Поэтому выбирают минимальную дозу, контролирующую симптомы. При отрицательном тесте на \Hp\ или сохраняющихся симптомах после эрадикации пробный курс допустим; при неэффективности следует пересмотреть диагноз и стратегию, а не автоматически переходить к многомесячному лансопразолу дважды в сутки. Рекомендация относится к классу ИПП и не доказывает превосходства лансопразола. Функциональная диспепсия не означает обязательную гиперсекрецию кислоты.\src{bsgfd}

\subsection{Чего от препарата не следует ожидать}
Вывод из рассмотренных показаний: лансопразол не следует считать антибиотиком, прокинетиком, лечением всех вариантов желчного рефлюкса, универсальной защитой от любых лекарств или средством профилактики всех видов рака ЖКТ. Ни повышение внутрижелудочного pH, ни подавление изжоги не заменяют этиологический диагноз.\src{uslabel,agagerd,rfgastrit}

\section{Сравнение с другими кислотосупрессивными препаратами}

\subsection{Другие ИПП}
В дозовых таблицах NICE стандартным режимам соответствуют лансопразол 30 мг, омепразол 20 мг, пантопразол 40 мг и рабепразол 20 мг один раз в день; для эзомепразола классификация зависит также от показания и тяжести эзофагита. Это клинические категории доз, а не точная эквивалентность pH, экспозиции или взаимодействий.\src{nice}

Выбор внутри класса определяется ответом, взаимодействиями, функцией печени, удобством формы и доступностью. Более слабая зависимость рабепразола от CYP2C19 может быть полезна в отдельных ситуациях, но не доказывает его универсального превосходства по всем исходам. Генотипические рекомендации CPIC для разных ИПП имеют неодинаковую доказательную базу.\src{cpic,agagerd}

\subsection{P-CAB: прямое сравнение с вонопразаном}
В исследовании Laine и соавт. 1024 взрослых с эрозивным эзофагитом рандомизировали на вонопразан 20 мг/сут или лансопразол 30 мг/сут до восьми недель. Заживление составило 92,9\% и 84,6\%; разница 8,3 процентного пункта, 95\% ДИ 4,5--12,2. Первично подтверждена не меньшая эффективность; преимущество по общему заживлению оценивалось в предусмотренном исследовательском анализе. Более заметные различия наблюдались при тяжелом эзофагите.\src{vontrial}

\textbf{Интерпретация:} результат поддерживает обсуждение P-CAB при тяжелом или недостаточно контролируемом заболевании, но не обязательную замену успешно работающего лансопразола. Короткие исследования заживления не отвечают полностью на вопрос сравнительной безопасности многолетнего приема. Следует также учитывать регистрацию и доступность конкретного P-CAB.\src{vontrial,seoul}

\subsection{Почему более сильное подавление кислоты не всегда означает замену ИПП}
В Clinical Practice Update AGA 2024 по P-CAB различаются клиническое преимущество, стоимость, доступность и объем долгосрочных данных. Эксперты не предлагают повсеместную замену ИПП при неисследованной изжоге, неэрозивной ГЭРБ, эзофагите LA A/B или профилактике язв. P-CAB можно рассматривать при LA C/D и у отдельных пациентов с объективно подтвержденным кислотным рефлюксом, не ответивших на ИПП дважды в сутки. Это экспертные Best Practice Advice, а не положения с формальной оценкой GRADE.\src{agapcab}

Для \Hp\ тот же документ более благоприятно оценивает замену ИПП на P-CAB у большинства пациентов. Однако ACG 2024 при неизвестной чувствительности сохраняет приоритет оптимизированной висмутовой квадротерапии на 14 дней, а вонопразан--амоксициллин рассматривает как альтернативу у подходящих пациентов. Эти тексты нельзя объединять в вымышленную единую рекомендацию «лансопразол больше не применяется». Сравниваются \emph{полные режимы}, а не только кислотосупрессивные компоненты.\src{agapcab,acghp}

\textbf{Практическое следствие.} Хороший ответ на переносимую терапию лансопразолом при сохраняющемся показании сам по себе не требует смены класса. Преимущество вонопразана над стандартной дозой лансопразола в конкретном исследовании не равно доказанному преимуществу над любым оптимизированным режимом ИПП. Это интерпретация границ имеющихся сравнений, а не запрет на обоснованную смену препарата.\src{agapcab,vontrial}

\subsection{Формы и вспомогательные вещества}
Для выбора препарата важны способность проглотить капсулу, наличие зонда и состав формы. Некоторые диспергируемые таблетки содержат аспартам; некоторые капсулы --- сахарозу и другие значимые вспомогательные компоненты. Проверка требуется при фенилкетонурии и соответствующих наследственных нарушениях обмена. Нельзя ранжировать генерики по клинической эффективности только по бренду, цене или стране производства без сравнительных данных.\src{emc,hpra}

\section{Региональный итог: где согласие, где расхождения}

\subsection{Общее ядро}
Синтез рассмотренных документов: ИПП остаются важным компонентом лечения кислотозависимых заболеваний; решение должно быть привязано к показанию и ответу. При \Hp\ необходимо эффективное комбинированное лечение, при высоком язвенном риске --- защита, при необъясненной рефрактерности --- диагностика. Ни один рассмотренный подход не обосновывает бессрочный лансопразол каждому человеку с неспецифическим дискомфортом.\src{agagerd,dgvs,seoul,rfgerd}

\subsection{Различия, меняющие практику}
\begin{enumerate}
\item \textbf{Дозы:} американский режим для язвы ДПК 15 мг/сут отличается от европейского 30 мг/сут; дозы не следует выравнивать по памяти.\src{uslabel,hpra}
\item \textbf{Эрадикация:} США строже ограничивают кларитромициновые схемы; Европа и РФ учитывают также доказанную локальную эффективность; в Японии и Корее есть обоснованные местными данными более короткие режимы.\src{acghp,maastricht,rfpud,jhp,khp}
\item \textbf{P-CAB:} в японских и новых корейских алгоритмах их роль заметнее, особенно при тяжелом эзофагите. Новое японское руководство по язвенной болезни 2026 сохраняет этиологический подход и приоритет лечения осложнений, а не объявляет один ИПП универсально предпочтительным.\src{jgerd,seoul,jpud2026}
\item \textbf{Дополнительные препараты:} российские рекомендации по ГЭРБ рассматривают комбинацию ИПП с ребамипидом. Это нельзя представлять как обязательный международный стандарт; более широкий набор препаратов в региональном документе не означает одинаковую доказательность всех комбинаций.\src{rfgerd,acggerd}
\item \textbf{Дети и печень:} возрастные допуски и формулировки снижения дозы отличаются. Нельзя использовать американское разрешение как замену российской инструкции.\src{uslabel,emc,grls}
\end{enumerate}

\chapter{Междисциплинарная интеграция и сравнительная гастропротекция}
\label{chap:interdisciplinary}

\section*{Введение к главе}
\addcontentsline{toc}{section}{Введение к главе}
Междисциплинарное использование лансопразола опирается на желудочную мишень, но оценивается по исходам, существенным для конкретной специальности. Липофильность и степень ионизации исходной молекулы относятся к ее доставке, $pK_a$ --- к условиям протонирования и кислотной активации, ковалентная модификация помпы --- к подавлению секреции. Эти механизмы объясняют уменьшение кислотного компонента повреждения, но не устраняют продолжающееся действие НПВП или антитромботических препаратов. Для ревматолога и невролога отсюда следует необходимость раздельной оценки профилактики язв, лечения существующего поражения и сохраняющегося риска при длительной лекарственной нагрузке.\src{molchem,molnagaya90,molcys,nice,aspirintrial,nsaid15}

В кардиологии гастропротективный эффект необходимо отделять от безопасности сочетания с антиагрегантами. Характер ковалентного связывания с H$^+$/K$^+$-АТФазой, липофильность и параметры кислотной активации не позволяют предсказать величину взаимодействия через CYP2C19 или сравнительный риск MACE. Выбор ИПП на фоне клопидогрела требует самостоятельной оценки лекарственного взаимодействия, тогда как необходимость профилактики определяется желудочно-кишечным риском. Аналогичное разграничение необходимо в ОРИТ: наличие механизма кислотосупрессии не устанавливает преимущества конкретной молекулы, лекарственной формы или пути введения при шоке и коагулопатии.\src{sps,molcypinhib,sccm}

Для оториноларинголога, пульмонолога и аллерголога основная трудность состоит в установлении связи между мишенью препарата и причиной симптома. Липофильность не является доказательством клинически значимого действия в любой внежелудочной ткани, а кислотная активация и блокада помпы не подтверждают рефлюксное происхождение кашля, охриплости или бронхоспазма. При эозинофильном эзофагите молекулярные противовоспалительные наблюдения также не заменяют клинической, эндоскопической и гистологической оценки. Поэтому глава сопоставляет подтвержденную гастропротекцию, отрицательные результаты внепищеводных исследований и границы экспериментальных данных, не превращая общий биохимический механизм в универсальное междисциплинарное показание.\src{toppits,asthma,eoe,molnrf2,molil8}

\section{Смежные специальности: сводная клиническая интерпретация}
\label{sec:interdisciplinary}
Сопоставление ниже основано на уже приведенных материалах. Оно различает рекомендацию общества, клиническое исследование и экспериментальную модель. Если формальная категория GRADE в обзоре не указана, ее нельзя восстановить лишь по наличию ссылки на руководство или по числу участников исследования.

\begin{longtable}{P{0.22\textwidth}P{0.35\textwidth}P{0.34\textwidth}}
\toprule
Направление & Клиническое место и данные & Основное ограничение \\
\midrule
\endhead
Кардиология & Взаимодействия с клопидогрелом и клинические исходы разобраны в разделе~\ref{sec:clopidogrel}; алгоритм ДАТТ --- в подразделе~\ref{sec:dapt-protocol}.\src{cardioacc2025,cardioplavix2025,sps} & Выбор ИПП для гастропротекции не равнозначен доказанному преимуществу конкретной молекулы по MACE. \\
Нефрология & Сравнительные оценки OR, HR и ROR, лабораторные ориентиры и действия при подозрении на ОИН представлены в разделе~\ref{sec:nephrology}.\src{renalblank2014,renalbigan2025,renalfaers2026,renalniceaki} & ОПП, ОИН и ХБП --- разные исходы; ROR сообщений не измеряет клиническую частоту или вероятность перехода ОИН в ХБП. \\
Онкогематология и эндокринология & Гастропротекция при иммуносупрессии и проверка pH-зависимых взаимодействий рассмотрены в разделе~\ref{sec:onco-endocrine}.\src{oncoppiposition2016,oncobosulif2026,oncocalcaps2026,oncocaltablet2026} & Показание к ИПП и совместимость с конкретной противоопухолевой формой оцениваются раздельно. \\
Ревматология и неврология: НПВП/аспирин & Гастропротекция по риску; лечебные и профилактические дозы разграничены. Прямые исследования разобраны в разделах~\ref{sec:nsaid}--\ref{sec:gastroprotection-trials}.\src{nice,aspirintrial,misotrial,aspirin15,nsaid15} & Это не лечение основного заболевания и не универсальная защита всего кишечника. Эндоскопическая язва, кровотечение и смертность --- разные исходы. \\
Оторино\-ларингология & TOPPITS: преимущества эмпирического курса 30 мг дважды/сут на 16 недель не выявлено; раздел~\ref{sec:extraesophageal}.\src{toppits} & Изученный режим с отрицательным результатом не является стандартом лечения любого ларингита или охриплости. \\
Пульмонология & У 306 детей добавление лансопразола не улучшило контроль астмы; раздел~\ref{sec:extraesophageal}.\src{asthma} & В обзоре нет количественных данных о бронхоспазмах, ОФВ$_1$ и подгруппе с объективно подтвержденным рефлюксом. ГЭРБ и астма требуют раздельной оценки показаний. \\
Аллергология и иммунология & ИПП включены ACG 2025 в варианты терапии ЭоЭ; необходимы клиническая, эндоскопическая и гистологическая оценки; раздел~\ref{sec:eoe}.\src{eoe} & Ответ $<15$ eos/hpf не равнозначен полному восстановлению слизистой. Критерии биопсии и ограничения молекулярного сравнения приведены в разделе~\ref{sec:eoe}. \\
Реаниматология & SCCM/ASHP 2024: ИПП или H$_2$-блокатор при риске стрессового кровотечения; прекращение после устранения риска; раздел~\ref{sec:stress-prophylaxis}.\src{sccm} & Универсальное преимущество лансопразола не установлено. Режим лечения язвенного кровотечения нельзя автоматически переносить на профилактику. \\
Фтизиатрия и инфекционные болезни & LPZS и QcrB в составе микобактериального комплекса $bc_1$: доклиническое направление; раздел~\ref{sec:mycobacteria}.\src{moltb,moltbstructure} & Другая химическая форма и другая мишень. Представленные данные не устанавливают клиническую дозу или эффективность лечения туберкулеза у человека. \\
\bottomrule
\end{longtable}

\subsection{От показания к решению: шесть клинических маршрутов}
Следующие маршруты представляют собой авторскую организацию уже изложенных данных, а не новый консенсус или самостоятельный протокол назначения. Их назначение --- связать клинический вопрос, допустимую цель кислотосупрессии и точку пересмотра решения. Дозы, не приведенные в исходных подразделах, здесь не реконструируются.

\begin{enumerate}
\item \textbf{НПВП или аспирин: сначала цель, затем доза.} Различают профилактику у пациента с риском и лечение уже установленного язвенного либо эрозивного поражения. Для профилактики фиксируют факторы риска и необходимость продолжения ульцерогенного воздействия; для лечения --- вид поражения и его причину. Режимы 15 мг/сут для профилактики и 30 мг/сут для лечения язвы сопоставлены в разделах~\ref{sec:doses} и~\ref{sec:nsaid}; предусмотренное отдельными SmPC повышение профилактической дозы не отменяет этого разграничения. Точка пересмотра --- заживление поражения и изменение потребности в НПВП или антитромботическом сочетании. Окончание лечебного курса не равнозначно автоматическому прекращению гастропротекции при сохраняющемся высоком риске; эрадикация \Hp\ также не устраняет этот риск во всех случаях.\src{uslabel,hpra,nice,aspirintrial,agadepres}

\item \textbf{ЛОР-жалобы: сначала обоснование рефлюкса.} При кашле, охриплости или ощущении «кома в горле» наличие симптома не устанавливает механизм его возникновения. До решения о продлении или усилении кислотосупрессии уточняют диагноз ГЭРБ и выбирают соответствующее исследование (раздел~\ref{sec:refractory-gerd}). Точка пересмотра --- сопоставление сохраняющейся жалобы с подтвержденным механизмом, а не достижение условного количества недель приема. Отсутствие преимущества в TOPPITS при 30 мг дважды в сутки за 16 недель не дает основания считать еще более длительную эмпирическую терапию доказанным следующим шагом.\src{agagerd,dgvs,lyon,toppits}

\item \textbf{Астма: два независимых вопроса.} Отдельно определяют, имеется ли гастроэнтерологическое показание к лансопразолу, и отдельно оценивают контроль астмы. Улучшение пищеводных симптомов не следует записывать как доказанное уменьшение бронхоспазмов или улучшение ОФВ$_1$. Отрицательный результат исследования у детей не поддерживает назначение лансопразола ради контроля астмы; при этом сам по себе он не отменяет независимо установленного показания для лечения ГЭРБ. Дозовый алгоритм для пациентов с астмой и объективно подтвержденным рефлюксом из представленных данных вывести нельзя.\src{asthma}

\item \textbf{ЭоЭ: ответ требует нескольких уровней оценки.} Уменьшение дисфагии не заменяет эндоскопию и биопсию. Критерий успешного гистологического ответа ACG, протокол отбора биоптатов и отличие ответа от полной морфологической нормализации изложены в разделе~\ref{sec:eoe}. Молекулярные наблюдения не заменяют оценку результата лечения.\src{eoe,eoeguide}

\item \textbf{ОРИТ: риск, профилактика и кровотечение --- разные этапы.} Сначала разграничивают профилактику стрессового кровотечения и лечение уже возникшего язвенного кровотечения. В первом случае оценивают факторы риска, названные SCCM/ASHP, и прекращают профилактику после их исчезновения; во втором применяют отдельную стратегию эндоскопии, гемостаза и кислотосупрессии. Точка пересмотра профилактики определяется сохранением риска, а не автоматическим завершением семидневного внутривенного курса, описанного для другой ситуации. Преимущество конкретно лансопразола из принадлежности к классу ИПП не следует.\src{sccm,esge2026,rfpud}

\item \textbf{Туберкулез: граница между исследованием и назначением.} Для интерпретации результата необходимо назвать химическую форму (LPZS), бактериальную мишень и модель исследования. Клиническая применимость не следует из одного факта связывания с мишенью. В пределах представленных данных это исследовательское направление, не алгоритм назначения; ограничения экспозиции, дозы и комбинаций разобраны в разделе~\ref{sec:mycobacteria}.\src{moltb,moltbstructure}
\end{enumerate}

\subsection{Три разных смысла ограничения доказательств}
\textbf{Первый: в конкретном клиническом исследовании преимущество не выявлено.} К этой категории относятся представленные результаты TOPPITS и исследования астмы у детей. Корректное заключение сохраняет изученную популяцию, вмешательство и конечную точку. Его нельзя расширять до утверждения «лансопразол неэффективен при любом рефлюкс-ассоциированном симптоме», но нельзя и игнорировать при обосновании рутинной эмпирической терапии.\src{toppits,asthma}

\textbf{Второй: необходимые подробности не воспроизведены в обзоре.} Таковы числовой критерий ремиссии ЭоЭ, отдельные показатели легочной функции в исследовании астмы и специальный режим лансопразоловой стресс-профилактики. Из этого следует недостаточность настоящего текста для соответствующего количественного решения, а не отсутствие сведений в первичной статье или полном руководстве. Для восполнения пробела требуется отдельная проверка первоисточника; ее результат нельзя подменять аналогией с другим показанием.\src{eoe,eoeguide,asthma,sccm}

\textbf{Третий: показан механизм, но не клиническая польза при обсуждаемой болезни.} Индукция HO-1 в RGM-1 не доказывает гистологическую ремиссию ЭоЭ, а структурное связывание LPZS с микобактериальным дыхательным комплексом не доказывает эффективность стандартной дозы лансопразола у больного туберкулезом. Детализация механизма повышает точность биологического объяснения, но сама по себе не заменяет клиническую верификацию.\src{molnrf2,molho1,moltb,moltbstructure}

\subsection{Как оформить междисциплинарное заключение без подмены показаний}
Для практического использования этой главы предлагается фиксировать пять элементов: \textbf{самостоятельное показание; тип доказательства; цель и режим лечения; способ оценки результата; основание для пересмотра}. Это редакционный шаблон, а не валидированная шкала риска. В графе доказательств следует различать зарегистрированную инструкцию, рекомендацию общества, результат исследования именно лансопразола и экспериментальную модель; формальную категорию GRADE не присваивают, если она в обзоре не воспроизведена.

В заключении следует явно указать, оценивается ли профилактика, лечение установленного заболевания или экспериментальная гипотеза. Затем результат соотносят именно с выбранной задачей: симптоматический ответ, эндоскопическое заживление, гистологический ответ и профилактика клинического осложнения не взаимозаменяемы. Конкретные ограничения для шести направлений приведены в таблице и маршрутах этого подраздела.\src{nice,aspirintrial,asthma,eoe,sccm,moltb,moltbstructure}

\section{Кардиология: клопидогрел, аспирин и двойная антитромбоцитарная терапия}
\label{sec:clopidogrel}
\textbf{Две самостоятельные цели лечения.} При двойной антитромбоцитарной терапии (ДАТТ; здесь --- аспирин и клопидогрел) необходимо предупредить желудочно-кишечное кровотечение, не создавая устранимого препятствия активации антиагреганта. Клопидогрел является пролекарством; CYP2C19 участвует в образовании его активного метаболита. Поэтому взаимодействие с ИПП оценивают отдельно от эффективности кислотосупрессии. Показание к гастропротекции не отменяют только из-за необходимости клопидогрела; выбор ИПП и решение о составе ДАТТ представляют разные клинические задачи.\src{cardioplavix2025,cardioesc2023}

\textbf{Прямое конкурентное ингибирование: лансопразол не является самым слабым ингибитором in vitro.} В работе Li и соавт. (2004) использованы микросомы печени человека и рекомбинантный CYP2C19; маркерными реакциями служили гидроксилирование S-мефенитоина и R-омепразола. Получены следующие значения константы конкурентного ингибирования $K_i$.\src{cardioli2004}
\begin{longtable}{P{0.32\textwidth}P{0.59\textwidth}}
\toprule
ИПП & $K_i$ для CYP2C19, мкмоль/л \\
\midrule
\endhead
Лансопразол & 0,4--1,5 \\
Омепразол & 2--6 \\
Эзомепразол & Около 8 \\
Пантопразол & 14--69 \\
Рабепразол & 17--21 \\
\bottomrule
\end{longtable}
Меньшее $K_i$ соответствует более сильному прямому ингибированию в данной экспериментальной системе. Следовательно, формулировка «лансопразол предпочтителен, поскольку конкурентно ингибирует CYP2C19 слабее омепразола» противоречит этим результатам. Значения относятся к модельным субстратам, а не непосредственно к активации клопидогрела; их нельзя превращать в клинический рейтинг без учета экспозиции и других механизмов взаимодействия.\src{cardioli2004}

\textbf{Метаболизм-зависимое ингибирование объясняет различие между лабораторной активностью и клиническим взаимодействием.} В исследовании Ogilvie и соавт. (2011) 30-минутная преинкубация с NADPH снижала $IC_{50}$ омепразола в 4,2 раза, эзомепразола --- в 10 раз; для лансопразола и пантопразола сдвиг был менее 1,5 раза. Метаболизм-зависимое торможение омепразолом и эзомепразолом не устранялось ультрацентрифугированием, что согласуется с необратимой либо квазинеобратимой инактивацией CYP2C19 в этой модели. Для лансопразола и пантопразола такой эффект в исследованных условиях не выявлен. Таким образом, одного $K_i$ недостаточно: важны свободная концентрация, метаболиты и зависимость ингибирования от времени. Инактивацию печеночного CYP2C19 нельзя отождествлять с кислотной активацией ИПП и дисульфидной блокадой желудочной H$^+$/K$^+$-АТФазы.\src{cardioogilvie2011,molcys}

\textbf{Проверка у человека: фармакокинетика и функция тромбоцитов, но не MACE.} В перекрестном исследовании Frelinger и соавт. (2012) участвовали 160 здоровых добровольцев с генотипом гомозиготных extensive metabolizers CYP2C19. Клопидогрел 75 мг/сут изучали в течение 9 дней без ИПП и с лансопразолом 30 мг/сут, эзомепразолом 40 мг/сут, омепразолом 80 мг/сут либо декслансопразолом 60 мг/сут. Эзомепразол, в отличие от лансопразола, значимо уменьшал AUC активного метаболита и ухудшал показатель VASP-PRI. Однако отсутствие значимости по этим показателям не означает отсутствия взаимодействия по всем тестам. Доза омепразола 80 мг/сут использована как положительный контроль; это не сравнение равных стандартных доз. Исследование не оценивало клинические события на ДАТТ и не устанавливает снижение MACE при выборе лансопразола.\src{cardiofrelinger2012}

\textbf{FDA: регистрационное ограничение, а не доказательство превосходства по исходам.} Американская инструкция Plavix, редакция мая 2025 года (FDA/DailyMed, разделы 5.1 и 7.2), предписывает избегать сочетания клопидогрела с омепразолом и эзомепразолом. Для омепразола снижение антиагрегантной активности отмечено и при разнесении приема на 12 часов. В инструкции прямо указано, что лансопразол, пантопразол и декслансопразол влияли на антиагрегантную активность меньше, чем омепразол или эзомепразол. Это обосновывает выбор альтернативы при необходимой гастропротекции, но не равнозначно утверждению FDA о доказанном снижении MACE именно лансопразолом.\src{cardioplavix2025}

\textbf{NHS SPS: предпочтительные альтернативы и неодинаковый потенциал взаимодействия.} Документ Specialist Pharmacy Service, обновленный 14 сентября 2023 года, рекомендует избегать омепразола и эзомепразола и предпочитать лансопразол, пантопразол или рабепразол, если пациенту на клопидогреле необходим ИПП. В приведенной SPS оценке ESC вероятность взаимодействия наиболее высока для омепразола и эзомепразола, промежуточна для лансопразола и наиболее низка для пантопразола и рабепразола. Поэтому лансопразол --- предпочтительная альтернатива двум первым препаратам, но не доказанно лучший вариант среди всех ИПП. Эта градация касается взаимодействия, а не частоты инфаркта, инсульта или тромбоза стента.\src{sps}

\textbf{ESC 2023: кому необходима гастропротекция на ДАТТ.} Рекомендации ESC по острому коронарному синдрому (раздел 13.3.7 и таблица 7) рассматривают ИПП при повышенном желудочно-кишечном риске: язва или кровотечение в анамнезе, антикоагулянтная терапия, хроническое применение НПВП либо кортикостероидов; также основанием служит сочетание не менее двух факторов --- возраст $\geq65$ лет, диспепсия, ГЭРБ, инфекция \Hp, хроническое употребление алкоголя. Возраст 65 лет сам по себе не следует выдавать за единственный универсальный критерий этой таблицы. ESC признает фармакодинамическое взаимодействие омепразола и эзомепразола с клопидогрелом, но отмечает отсутствие убедительного подтверждения увеличения ишемических событий или тромбоза стента в соответствующих клинических исследованиях.\src{cardioesc2023}

\textbf{ESC 2024: продолжительность защиты и пределы молекулярного предпочтения.} Рекомендации ESC по хроническим коронарным синдромам (раздел 4.3.1.4) предусматривают гастропротекцию ИПП у пациентов повышенного желудочно-кишечного риска на протяжении антитромботической терапии. Совместное применение омепразола или эзомепразола с клопидогрелом не приветствуется из-за снижения экспозиции активного метаболита; одновременно документ не констатирует однозначного влияния этого сочетания на ишемические события или тромбоз стента. Данные положения поддерживают устранение потенциального взаимодействия, но не устанавливают отдельного преимущества лансопразола перед пантопразолом по MACE.\src{cardioesc2024}

\textbf{MACE: количественные данные COGENT и их ограничения.} В COGENT рандомизированы 3873 пациента, в анализ включен 3761. На фоне аспирина 75--325 мг/сут сравнивали клопидогрел 75 мг/сут с омепразолом 20 мг/сут и без него. Через 180 дней частота комбинированной желудочно-кишечной конечной точки составляла 1,1\% против 2,9\% (HR 0,34; 95\% ДИ 0,18--0,63). Сердечно-сосудистая конечная точка включала сердечно-сосудистую смерть, нефатальный инфаркт, реваскуляризацию или инсульт: 4,9\% против 5,7\%, HR 0,99 (95\% ДИ 0,68--1,44; $p=0{,}96$). Испытание досрочно прекращено из-за прекращения финансирования; широкий ДИ не исключает клинически значимого вреда. Отсутствие значимого различия не доказывает эквивалентность. Изучался омепразол, а не лансопразол; эти результаты нельзя использовать как прямое сравнение MACE между ИПП.\src{cardiocogent2010}

\textbf{Аспирин: доказанная гастропротекция не равна доказанной кардиобезопасности комбинации.} В исследовании Lai и соавт. 123 пациента после осложненной язвы и эрадикации \Hp\ продолжали аспирин 100 мг/сут. При лансопразоле 30 мг/сут повторные язвенные осложнения возникли у 1,6\% против 14,8\% при плацебо примерно за год. Это прямой результат по желудочно-кишечной конечной точке в высокорисковой популяции; он не является исследованием ДАТТ и не доказывает сохранность эффекта клопидогрела. Результаты японского исследования лансопразола 15 мг/сут при аспирин-ассоциированном риске также относятся к вторичной профилактике язв, а не к сравнительной оценке MACE на ДАТТ.\src{aspirintrial,aspirin15}

\textbf{Дозирование и длительность: разные основания для разных решений.} Лансопразол 15 мг один раз в сутки является профилактической дозой при НПВП-ассоциированных язвах и изученной дозой вторичной профилактики язв на аспирине; 30 мг один раз в сутки изучены после аспирин-ассоциированного язвенного осложнения и применяются для лечения НПВП-ассоциированной язвы, обычно до 8 недель. Эти режимы не следует переименовывать в доказанную молекулоспецифичную схему снижения MACE на ДАТТ. Продолжительность гастропротекции определяется сохраняющимся риском и антитромботической терапией, а не автоматически восьминедельным курсом заживления язвы.\src{uslabel,hpra,aspirin15,aspirintrial,cardioesc2024}

\textbf{Персонализация не ограничивается заменой ИПП.} Меньшая зависимость собственной экспозиции эзомепразола от генотипа CYP2C19 не означает меньшего воздействия на метаболизм клопидогрела. Генетически сниженная активация антиагреганта и лекарственное ингибирование --- разные причины недостаточного ответа; замена ИПП не исправляет генотип. Фармакогенетические вопросы рассмотрены в разделе~\ref{sec:pharmacogenetics}. Самостоятельное прекращение клопидогрела, изменение его дозы или отмена показанной гастропротекции из приведенных данных не следуют.\src{leonova2015,cardioplavix2025,agadepres}

\textbf{Сравнение ИПП при клопидогреле: преимущество выбора, а не доказанное снижение MACE.} Лансопразол и пантопразол предпочтительнее омепразола и эзомепразола с точки зрения уменьшения взаимодействия с клопидогрелом; NHS SPS также включает рабепразол в предпочтительные альтернативы. Регуляторное основание и клинические рекомендации поддерживают такой выбор при необходимой гастропротекции, но не доказывают уменьшение MACE именно вследствие назначения лансопразола. Его превосходство перед пантопразолом по сердечно-сосудистым исходам не установлено.\src{cardioplavix2025,sps,cardioesc2024}

\textbf{Предел прямого сравнения.} Конкурентное ингибирование CYP2C19, метаболизм-зависимая инактивация, изменение экспозиции активного метаболита и клинические события --- последовательные, но не взаимозаменяемые уровни доказательств. Представленные исследования позволяют количественно сравнить первые уровни; COGENT оценивает клинические исходы добавления омепразола, а не замену его лансопразолом. Поэтому формулировки «лансопразол снижает MACE по сравнению с омепразолом» и «все ИПП эквивалентны по MACE» выходят за пределы приведенной доказательной базы.\src{cardioli2004,cardioogilvie2011,cardiofrelinger2012,cardiocogent2010}

\subsection{Клинический протокол антитромботической гастропротекции при двойной антитромбоцитарной терапии}
\label{sec:dapt-protocol}
\textbf{Отбор пациента и определение показания.} Протокол применяют при назначенной ДАТТ клопидогрелом и ацетилсалициловой кислотой, раздельно оценивая ишемический риск и вероятность желудочно-кишечного кровотечения. Документируют показание и продолжительность ДАТТ, дату коронарного вмешательства, язвенное кровотечение в анамнезе и сочетание с антикоагулянтами, НПВП или глюкокортикостероидами. Высокий кардиоваскулярный риск не является самостоятельным доказательством необходимости определенной молекулы ИПП. ESC 2023/2024 и ACC/AHA 2025 обосновывают гастропротекцию при повышенном желудочно-кишечном риске; рекомендации РКО уточняют национальный контекст ведения ОКС. Цель выбора состоит в сохранении показанной антитромботической терапии одновременно с профилактикой кровотечения и уменьшением потенциального лекарственного взаимодействия. Эти положения относятся к классу гастропротекторов и не устанавливают исключительного преимущества лансопразола по сердечно-сосудистым исходам.\src{cardioesc2023,cardioesc2024,cardioacc2025,cardiorko2024}

\textbf{Фармакокинетическое основание сравнительного предпочтения.} Клопидогрел требует метаболической активации с участием CYP2C19; уменьшение образования активного метаболита способно ослабить подавление тромбоцитарного рецептора P2Y$_{12}$. Однако утверждение о минимальном конкурентном ингибировании CYP2C19 лансопразолом не соответствует результатам изолированных ферментных исследований: в работе Li и соавт. он проявлял выраженное прямое ингибирование. Для оценки взаимодействия необходимо совместно учитывать свободную концентрацию ингибитора, обратимый и метаболизм-зависимый компоненты, дозу и временной профиль экспозиции.\src{cardioli2004} В исследовании Ogilvie и соавт. метаболизм-зависимое ингибирование CYP2C19 выявлялось у омепразола, но не у лансопразола или пантопразола в изученных системах. В рандомизированном перекрестном исследовании здоровых добровольцев лансопразол меньше влиял на образование активного метаболита и подавление функции тромбоцитов, чем изученные режимы омепразола или эзомепразола. Именно это различие взаимодействий, отраженное в американской инструкции Plavix, обосновывает предпочтение при переключении; оно не равнозначно доказанному уменьшению инфарктов или тромбоза стента.\src{cardioogilvie2011,cardiofrelinger2012,cardioplavix2025}

\textbf{Пошаговый выбор гастропротектора.} Сначала подтверждают сохраняющееся показание к ИПП и проверяют все назначенные и безрецептурные препараты. Если пациент получает клопидогрел вместе с омепразолом или эзомепразолом, рассматривают замену гастропротектора: американская инструкция рекомендует избегать этих сочетаний, а разнесение омепразола и клопидогрела на 12 часов не устраняет взаимодействия. Затем оценивают допустимость лансопразола с учетом переносимости, функции печени, сопутствующих лекарств и предшествующих реакций на ИПП. Лансопразол является одной из предпочтительных альтернатив указанным двум молекулам, а не безусловно первым выбором перед пантопразолом или другими совместимыми вариантами.\src{cardioplavix2025,sps,uslabel} При подозрении на ОИН или предшествующей тяжелой гиперчувствительности стандартное переключение внутри класса не проводят; применяют алгоритм подраздела~\ref{sec:renal-interstitial-markers}. В остальных случаях фиксируют выбранную дозу по гастроэнтерологическому показанию, дату замены, отменяемый ИПП и срок пересмотра. Два ИПП одновременно не назначают; сама замена гастропротектора не требует прекращения ДАТТ. При известном неблагоприятном CYP2C19-фенотипе стратегию антиагрегантной терапии пересматривает кардиолог: замена ИПП не устраняет генетическое ограничение активации клопидогрела.\src{uslabel,cardiocpic2022}

\textbf{Мониторинг и предел клинического вывода.} При наблюдении раздельно регистрируют признаки желудочно-кишечного кровотечения, гемоглобин по клиническим показаниям, соблюдение ДАТТ, переносимость ИПП и симптомы ишемии. Кровавая рвота, мелена, обморок или подозрение на ОКС требуют неотложной оценки, а не плановой замены гастропротектора; изменение антитромботического режима согласуют с учетом риска тромбоза стента.\src{cardioesc2023,cardioacc2025} В COGENT уменьшение желудочно-кишечных событий получено при добавлении омепразола к клопидогрелу и аспирину; сердечно-сосудистый исход HR 0,99 (95\% ДИ 0,68--1,44) не доказывал эквивалентности и не исключал клинически значимого различия. Лансопразол в этом испытании не изучался.\src{cardiocogent2010} Поэтому итоговый протокол обосновывает приоритет лансопразола относительно омепразола или эзомепразола при необходимости уменьшить избегаемое взаимодействие, но не приписывает ему доказанного превосходства по MACE перед всеми ИПП. Классовые сведения о безопасности и результаты испытаний других молекул рассматривают отдельно; ограничения COMPASS в главе~\ref{chap:personalization} сохраняют самостоятельное значение и не заменяют прямого исследования лансопразола.\src{cardioesc2024,cardioplavix2025,compass}

\begin{enumerate}
\item \textbf{Документировать обе группы риска.} Уточнить показание и плановую длительность ДАТТ, дату коронарного вмешательства, язвенное кровотечение в анамнезе, сопутствующие антикоагулянты, НПВП и кортикостероиды. Применить критерии желудочно-кишечного риска ESC, приведенные выше. Не сокращать ДАТТ и не менять ингибитор P2Y$_{12}$ только ради выбора ИПП; соответствующее решение принимает кардиолог.\src{cardioesc2023,cardioacc2025}
\item \textbf{Устранить избегаемое взаимодействие, не отменяя показанную защиту.} Если необходим ИПП, избегать омепразола и эзомепразола; разнесение омепразола и клопидогрела на 12 часов не решает проблему. Лансопразол является предпочтительной альтернативой этим двум препаратам наряду с пантопразолом и рабепразолом. Если пациент уже получает эффективный и переносимый лансопразол, само назначение клопидогрела не требует его замены. Однако преимущество лансопразола перед пантопразолом или рабепразолом по MACE не установлено.\src{cardioplavix2025,sps}
\item \textbf{Выбрать дозу по желудочно-кишечной задаче.} Режимы лансопразола 15 мг один раз в сутки, изученные для профилактики язв, и 30 мг один раз в сутки, изученные после аспирин-ассоциированного язвенного осложнения, имеют разные основания. Они не образуют универсальной шкалы дозирования по сердечно-сосудистому риску. Использовать режим соответствующего показания и инструкции; перенесенное кровотечение требует отдельной оценки. При сохраняющемся высоком риске продолжать гастропротекцию в течение антитромботического лечения с периодическим пересмотром необходимости.\src{uslabel,aspirin15,aspirintrial,cardioesc2024}
\item \textbf{Проверить ограничения конкретного пациента.} Сверить весь список лекарств и предшествующие реакции на ИПП; учесть нефрологический алгоритм раздела~\ref{sec:nephrology}. Лекарственное уменьшение активации клопидогрела и генетически обусловленный недостаточный ответ не тождественны: замена ИПП не исправляет фенотип CYP2C19. При подозрении на недостаточную эффективность антиагреганта пересматривают антитромботическую стратегию, а не повышают эмпирически дозу лансопразола.\src{cardioplavix2025,uslabel}
\item \textbf{Контролировать исходы, а не только совместимость.} При повторном визите оценивать приверженность обоим направлениям лечения и признаки кровопотери; анемия или подозрение на кровотечение требуют гемограммы и поиска источника. Мелена, кровавая рвота, обморок либо новые ишемические симптомы требуют срочной оценки. Решение об изменении ДАТТ при кровотечении принимают совместно; пациент не должен самостоятельно отменять клопидогрел или аспирин.\src{cardioacc2025,cardioesc2024}
\end{enumerate}

\textbf{Практический протокол переключения: приоритет относительно конкретного исходного препарата.} Если показанная гастропротекция на ДАТТ проводится омепразолом или эзомепразолом, лансопразол следует рассматривать как один из предпочтительных препаратов переключения, а не как единственную обязательную альтернативу. При документированной эффективности и переносимости пантопразола или рабепразола переключение на лансопразол ради предполагаемого снижения MACE не обосновано. Такое разграничение соответствует рекомендациям SPS: приоритет относится к устранению избегаемого взаимодействия, а не к ранжированию клинических сердечно-сосудистых исходов.\src{sps,cardioplavix2025}
\begin{longtable}{P{0.29\textwidth}P{0.62\textwidth}}
\toprule
Ситуация & Обязательное действие в авторском алгоритме \\
\midrule
\endhead
Плановая замена омепразола/эзомепразола & Записать причину переключения, выбранный ИПП, дозу по желудочно-кишечному показанию и дату начала; исключить дублирование двух ИПП. Не создавать запланированного перерыва в клопидогреле или аспирине ради замены гастропротекции.\src{cardioplavix2025,sps} \\
Лансопразол уже эффективен и переносится & Сохранить обоснованную гастропротекцию; проверить безрецептурный омепразол/эзомепразол и другие взаимодействия. Не повышать дозу ИПП для «усиления» антиагрегантного эффекта.\src{sps,cardioplavix2025} \\
Подозрение на ОИН или предшествующая тяжелая реакция на ИПП & Не применять алгоритм простой замены молекулы. Приоритетны отмена подозреваемого препарата и согласованное кардиологическое, гастроэнтерологическое и нефрологическое решение; см. раздел~\ref{sec:nephrology}.\src{uslabel,renalniceaki} \\
\bottomrule
\end{longtable}

\textbf{Генотипическая развилка у пациента после ОКС/ЧКВ.} При уже известном фенотипе CYP2C19 IM или PM замена омепразола на лансопразол не восстанавливает генетически сниженную активацию клопидогрела. CPIC 2022 для острого коронарного синдрома и/или чрескожного коронарного вмешательства рекомендует избегать стандартного клопидогрела у таких пациентов и использовать прасугрел или тикагрелор при отсутствии противопоказаний. Решение принимает кардиолог с учетом кровотечений, анамнеза инсульта/ТИА и исходного показания; прасугрел противопоказан при перенесенном инсульте/ТИА. Это интерпретация имеющегося генотипа, а не требование обязательного тестирования всех пациентов и не разрешение самостоятельной смены ДАТТ. В отличие от уменьшения потенциального взаимодействия ИПП, данная развилка направлена непосредственно на выбор эффективного антиагреганта.\src{cardiocpic2022}

\textbf{Граница молекулярного обоснования.} Термин «минимальное подавление CYP2C19 лансопразолом» некорректен для прямого конкурентного ингибирования: приведенные выше $K_i$ показывают обратное. Практическая приемлемость сочетания опирается на различия метаболизм-зависимого ингибирования и исследований у человека. Снижение потенциального взаимодействия не является доказательством снижения MACE; гастропротективная эффективность и сердечно-сосудистое превосходство должны подтверждаться разными конечными точками.\src{cardioli2004,cardioogilvie2011,cardiofrelinger2012,cardiocogent2010}

\section{НПВП, аспирин и антитромботическая терапия}
\label{sec:nsaid}
\textbf{Клиническая цель и механизм гастропротекции.} В ревматологии и неврологии лансопразол рассматривается как средство профилактики и лечения кислотозависимого повреждения верхних отделов ЖКТ, а не как терапия основного воспалительного или неврологического заболевания. Блокада желудочной H$^+$/K$^+$-АТФазы уменьшает кислотный компонент повреждения, но не устраняет продолжающееся ульцерогенное воздействие и не заменяет пересмотр НПВП или антитромботической комбинации. Экспериментальная модуляция воспалительных путей не обосновывает самостоятельного системного противовоспалительного показания.\src{uslabel,nice,molnrf2}

\textbf{Критерии назначения профилактики.} Основание --- желудочно-кишечный риск, особенно перенесенная язва или кровотечение, старший возраст и сочетание НПВП с антитромботическими препаратами. Длительность приема и полипрагмазия сами по себе не заменяют индивидуальной оценки риска. NICE CG184 предусматривает сопутствующий ИПП у пациента с язвенным анамнезом, которому необходимо продолжать НПВП; одновременно пересматривают выбор анальгетика. Единая балльная шкала для всех получателей НПВП здесь не вводится; отдельные критерии ESC для пациентов на ДАТТ приведены в разделе~\ref{sec:clopidogrel}.\src{nice,hpra}

\textbf{Профилактику, лечение язвы и лечение эрозий необходимо разграничивать:}
\begin{itemize}
\item \textbf{Профилактика НПВП-ассоциированных язв:} лансопразол 15 мг один раз в сутки; в отдельных европейских SmPC при недостаточности допускаются 30 мг/сут. Это не универсальное требование начинать профилактику с 30 мг.\src{uslabel,hpra}
\item \textbf{Лечение НПВП-ассоциированной язвы:} 30 мг один раз в сутки, обычно до 8 недель. Этот лечебный режим не следует подменять профилактической дозой или автоматически превращать в бессрочное назначение.\src{uslabel,hpra}
\item \textbf{Эрозивное поражение:} российские КР «Гастрит и дуоденит» 2024 года предусматривают ИПП на 4--6 недель, в том числе при связи с НПВП. Это положение о классе препаратов и другом типе поражения; специальная доза лансопразола для данного российского показания в обзоре не воспроизведена.\src{rfgastrit}
\end{itemize}

\textbf{Прямые доказательства по дозам.} Доза 15 мг/сут исследована в японских рандомизированных исследованиях вторичной профилактики язв на фоне аспирина и НПВП у пациентов с язвенным анамнезом; препаратом сравнения был гефарнат 50 мг дважды в сутки. Эти результаты поддерживают соответствующее профилактическое применение, но не устанавливают преимущество лансопразола над другими ИПП. Сравнение доз 15 и 30 мг с мизопростолом и плацебо и различия конечных точек подробно представлены в разделе~\ref{sec:gastroprotection-trials}.\src{aspirin15,nsaid15,misotrial}

\textbf{Аспирин после язвенного осложнения.} В исследовании Lai и соавт. 123 пациента после осложненной язвы и эрадикации \Hp\ продолжали аспирин 100 мг/сут. При лансопразоле 30 мг/сут повторные язвенные осложнения отмечены у 1,6\% против 14,8\% при плацебо примерно за год. Абсолютная разница 13,2 процентного пункта соответствует расчетному NNT около 8. Результат относится к отобранной высокорисковой популяции, а не ко всем пользователям аспирина; эрадикация не исключает необходимости гастропротекции при сохраняющемся риске.\src{aspirintrial}

\textbf{Сопутствующая терапия и пересмотр назначения.} Отдельно оценивают предупреждение для высокодозного метотрексата, мониторирование при варфарине и выбор ИПП на фоне клопидогрела (раздел~\ref{sec:interactions}). При завершении лечебного курса повторно оценивают заживление и сохраняющийся риск: окончание лечения язвы не означает автоматической отмены необходимой профилактики. Решение об изменении антитромботической терапии принимают отдельно с учетом тромботического риска; самостоятельно отменять ее из-за назначения лансопразола нельзя.\src{uslabel,sps,nice,agadepres,esge2026}

\textbf{Границы эффекта.} Профилактика эндоскопической язвы, предотвращение клинического кровотечения и снижение смертности --- разные исходы. Представленные исследования не доказывают универсальную защиту всего кишечника или одинаковую абсолютную пользу при любом исходном риске. Лансопразол снижает определенный компонент риска, но не делает потенциально опасное лекарственное сочетание безусловно безопасным.\src{aspirintrial,misotrial,aspirin15,nsaid15}

\textbf{Кардиологический контекст.} При клопидогреле выбор альтернатив омепразолу и эзомепразолу обоснован уменьшением лекарственного взаимодействия, а не доказанным преимуществом лансопразола по MACE. Регистрационные положения FDA/DailyMed, рекомендации ESC и NHS SPS, критерии риска и результаты COGENT разобраны в разделе~\ref{sec:clopidogrel}. Дозы профилактики НПВП- и аспирин-ассоциированных язв не заменяют индивидуального выбора режима на ДАТТ.\src{cardioplavix2025,cardioesc2023,cardioesc2024,sps}

\textbf{Прямое сравнение с мизопростолом: профилактическая эффективность и переносимость.} В исследовании Graham и соавт. 537 пользователей НПВП без \Hp, с язвой желудка в анамнезе, получали лансопразол 15 или 30 мг/сут, мизопростол 200 мкг четыре раза в сутки либо плацебо. Через 12 недель язва желудка отсутствовала соответственно у 80\%, 82\%, 93\% и 51\%. По этому показателю результат лансопразола был численно ниже результата мизопростола; преимущества лансопразола по предотвращению эндоскопической язвы эти цифры не показывают. Однако мизопростол чаще вызывал нежелательные явления и досрочное прекращение участия. При учете выбывших как неудач успех составил 69\% во всех активных группах и 35\% при плацебо. Это иллюстрирует зависимость практического результата от переносимости и способа учета выбываний. Совпадение 69\% не заменяет формального доказательства эквивалентности; различие 80\% и 82\% само по себе не устанавливает превосходства дозы 30 мг над 15 мг. Мизопростол здесь является активным компаратором, а не другим представителем ИПП.\src{misotrial}

\textbf{Скорость эпителизации: прямое сравнение доз лансопразола.} В исследовании Agrawal и соавт. при продолжении НПВП сравнивали лансопразол 15 и 30 мг один раз в сутки с ранитидином 150 мг дважды в сутки. Конечной точкой было эндоскопическое заживление язвы желудка, а не гистологическое измерение скорости регенерации эпителия. Частоты на фиксированных сроках позволяют оценить динамику заживления, но не дают медианного времени эпителизации.\src{nsaidagrawal2000}
\begin{longtable}{P{0.44\textwidth}P{0.22\textwidth}P{0.22\textwidth}}
\toprule
Режим & Заживление к 4-й неделе & Заживление к 8-й неделе \\
\midrule
\endhead
Лансопразол 15 мг/сут & 47\% (56/118) & 69\% (81/118) \\
Лансопразол 30 мг/сут & 57\% (67/117) & 73\% (85/117) \\
Ранитидин 150 мг дважды/сут & 30\% (34/115) & 53\% (61/115) \\
\bottomrule
\end{longtable}
Обе дозы лансопразола превосходили ранитидин на обоих сроках. Численное преимущество 30 мг перед 15 мг не достигало статистической значимости; результат не доказывает ни ускорения заживления при удвоении дозы, ни формальной эквивалентности доз. Наличие исследовательской группы 15 мг не отменяет зарегистрированного лечебного режима 30 мг/сут.\src{nsaidagrawal2000,uslabel,hpra}

\textbf{Омепразол и мизопростол: отдельное прямое сравнение OMNIUM.} В этом исследовании при продолжающемся приеме НПВП сравнивали омепразол 20 или 40 мг один раз в сутки и мизопростол 200 мкг четыре раза в сутки. Через 8 недель язвы желудка зажили соответственно у 87\% (102/117), 80\% (105/132) и 73\% (91/125). Преимущество перед мизопростолом было значимым для омепразола 20 мг ($p=0{,}004$), но не 40 мг ($p=0{,}14$). При эрозивном поражении направление различия было иным: критерий заживления эрозий достигнут у 77\%, 79\% и 87\% соответственно, с преимуществом мизопростола. Лансопразол в OMNIUM не изучался. Сопоставление этих процентов с предыдущей таблицей является межисследовательским, а не рандомизированным сравнением ИПП; различия популяций и поражений не позволяют объявить один препарат быстрее другого.\src{nsaidomnium1998}

\textbf{Лансопразол против мизопростола: профилактика не измеряет эпителизацию.} В прямом исследовании Graham лансопразол 15 и 30 мг/сут сравнивали с мизопростолом 200 мкг четыре раза/сут для предотвращения рецидива язвы, а не заживления исходного дефекта. Через 12 недель без язвы оставались 80\%, 82\% и 93\% соответственно; преимущество мизопростола по эндоскопической профилактике сопровождалось худшей переносимостью. При учете выбывших как неудач успех составлял 69\% во всех активных группах. Это не доказательство равной скорости эпителизации и не основание считать 30 мг лансопразола быстрее 15 мг или мизопростола. Подробности профилактической конечной точки сохранены в разделе~\ref{sec:gastroprotection-trials}.\src{misotrial}

\textbf{Сравнение с омепразолом и другими ИПП: недостающие конечные точки.} Для профилактики НПВП-ассоциированных язв в обзоре приведены 15 мг лансопразола один раз в сутки с возможностью 30 мг/сут при недостаточности по отдельным SmPC; для лечения установленной язвы --- 30 мг один раз в сутки, обычно до 8 недель. Представленные прямые сравнения относятся к лансопразолу против ранитидина, лансопразолу против мизопростола при профилактике и омепразолу против мизопростола при лечении. Прямого сравнения лансопразола с омепразолом по скорости заживления НПВП-язв эти испытания не содержат. Исследование Graham посвящено профилактике и не устанавливает ускорения заживления существующего дефекта. Следовательно, ни терапевтическая эквивалентность другим ИПП, ни преимущество лансопразола по срокам заживления на основании представленных данных не доказаны.\src{uslabel,hpra,misotrial}

\textbf{Дозовые категории не являются результатом сравнительного исследования.} В таблицах NICE стандартным режимам соответствуют лансопразол 30 мг, омепразол 20 мг, пантопразол 40 мг и рабепразол 20 мг один раз в сутки; классификация эзомепразола зависит от показания и тяжести эзофагита. Эти категории не доказывают одинаковую профилактику НПВП-язв, скорость заживления, экспозицию или профиль взаимодействий. Результаты сравнения лансопразола с плацебо, мизопростолом или гефарнатом нельзя переносить на прямое сравнение с омепразолом.\src{nice,misotrial,nsaid15}

\section{Прямые исследования гастропротекции: величина эффекта}
\label{sec:gastroprotection-trials}
Ниже приведены отдельные рандомизированные исследования именно лансопразола, а не косвенное сравнение всех ИПП. Результаты относятся к пациентам с язвенным анамнезом и не обосновывают профилактический прием у любого человека, принимающего аспирин или НПВП. Дозы указаны как изученные режимы, а не как персональное назначение.\src{aspirintrial,misotrial,aspirin15,nsaid15}

\begin{longtable}{P{0.23\textwidth}P{0.34\textwidth}P{0.34\textwidth}}
\toprule
Исследование и популяция & Сравнение и результат & Что ограничивает вывод \\
\midrule
\endhead
Lai и соавт., 2002; 123 пациента после язвенного осложнения и эрадикации \Hp & Аспирин 100 мг/сут с лансопразолом 30 мг/сут или плацебо. Повторные язвенные осложнения: 1,6\% против 14,8\% примерно за год.\src{aspirintrial} & Клиническое осложнение, а не только эндоскопическая язва. Небольшая высокорисковая выборка; различия смертности не установлены. \\
Graham и соавт., 2002; 537 пользователей НПВП без \Hp, с язвой желудка в анамнезе & Через 12 недель без язвы желудка: 80\% при 15 мг/сут, 82\% при 30 мг/сут лансопразола, 93\% при мизопростоле 200 мкг четыре раза/сут, 51\% при плацебо.\src{misotrial} & Мизопростол чаще вызывал нежелательные явления и досрочный выход. При учете выбывших как неудач успех составил 69\% во всех активных группах и 35\% при плацебо. \\
Sugano и соавт., 2011; 461 пациент в Японии, аспирин и язвенный анамнез & Лансопразол 15 мг/сут против гефарната 50 мг дважды/сут. Кумулятивная частота язвы к 361-му дню: 3,7\% против 31,7\%; отношение рисков во времени 0,099, 95\% ДИ 0,042--0,230.\src{aspirin15} & Первичная конечная точка --- язва желудка/ДПК. Сравнение с гефарнатом, а не с другим ИПП или отсутствием гастропротекции. \\
Sugano и соавт., 2012; 366 пациентов в Японии, НПВП и язвенный анамнез & Лансопразол 15 мг/сут против гефарната 50 мг дважды/сут. Кумулятивная частота язвы к 361-му дню: 12,7\% против 36,9\%.\src{nsaid15} & Оценки Каплана--Майера, неодинаковая длительность наблюдения и выбывания. Не доказательство снижения смертности или защиты всего кишечника. \\
\bottomrule
\end{longtable}

\textbf{Как не переоценить цифры.} Отсутствие эндоскопической язвы, отсутствие кровотечения и выживаемость --- разные исходы. Процент пациентов без язвы нельзя подписывать как «снижение кровотечений». Кумулятивную вероятность события по Каплану--Майеру нельзя без оговорки заменять обычной долей событий от числа включенных. В частности, большое преимущество перед гефарнатом не доказывает такого же преимущества перед другим ИПП. Это методологические ограничения интерпретации приведенных исследований.\src{aspirin15,nsaid15,misotrial}

\textbf{Абсолютная польза зависит от исходного риска.} В исследовании Lai расчет по округленным опубликованным частотам дает $NNT=1/(0{,}148-0{,}016)\approx 7{,}6$, то есть около 8 пациентов за год. Это авторский расчет, а не универсальное NNT лансопразола: при меньшем исходном риске абсолютная польза обычно будет меньше даже при сходном относительном эффекте.\src{aspirintrial}

\section{Кашель, охриплость, «ком в горле» и астма}
\label{sec:extraesophageal}
\textbf{Диагноз ГЭРБ и причинность внепищеводного симптома.} Кашель, охриплость, ларингит и ощущение «кома в горле» нельзя автоматически считать кислотозависимыми проявлениями. В изложенной диагностической рамке AGA, DGVS и Lyon 2.0 сначала уточняют диагноз и режим приема: при неподтвержденной ГЭРБ рассматривают исследование рефлюкса вне кислотосупрессии; при доказанной ГЭРБ и сохраняющихся жалобах --- pH-импедансометрию на оптимизированном лечении по показаниям. Решение об усилении терапии должно следовать из уточнения механизма жалоб, а не из одного факта их сохранения (раздел~\ref{sec:refractory-gerd}).\src{agagerd,dgvs,lyon}

\textbf{Сопоставление позиций гастроэнтерологических и ЛОР-обществ.} ACG 2022 рекомендует сначала исключать нерефлюксные причины: при изолированных внепищеводных симптомах без изжоги и регургитации исследование рефлюкса проводят до назначения ИПП. При сочетании типичных и внепищеводных симптомов можно рассмотреть ИПП дважды в сутки на 8--12 недель; это условная рекомендация с низким качеством доказательств. Ларингоскопические признаки сами по себе не подтверждают рефлюксную этиологию. AGA 2023 также поддерживает предварительную объективизацию при отсутствии типичных симптомов; после одного неудачного курса до 12 недель последовательные эмпирические замены ИПП имеют низкую результативность. Ни один из этих документов не выбирает лансопразол по липофильности.\src{acggerd,agaextra2023}

Рекомендация AAO-HNSF по дисфонии (2018, ключевое положение 6) относится к иной исходной задаче: не назначать антирефлюксные препараты при изолированной дисфонии только по симптомам предполагаемой ГЭРБ/ларингофарингеального рефлюкса без визуализации гортани. Это не противоречит ACG: оториноларинголог исключает локальную патологию, а гастроэнтеролог проверяет предполагаемую рефлюксную причинность. Ларингоскопия необходима для оценки гортани, но сама по себе не доказывает, что ларингит ответит на ИПП. Специального преимущества лансопразола перед омепразолом, эзомепразолом, пантопразолом или рабепразолом данная позиция AAO-HNSF не устанавливает.\src{aao2018,acggerd}

\textbf{Пульмонологическая позиция: хронический кашель и астма не взаимозаменяемы.} CHEST 2016 при подозрении на рефлюкс-кашлевой синдром без изжоги и регургитации рекомендует против монотерапии ИПП, поскольку она маловероятно устранит кашель (Grade 1C). При типичных симптомах рефлюкса кислотосупрессия направлена на их контроль и сочетается с немедикаментозными мерами; возможный ответ кашля может запаздывать до трех месяцев. Это не рекомендация любого ИПП как универсального противокашлевого средства.\src{chestcough2016}

GINA 2026, раздел о сопутствующей ГЭРБ, рекомендует лечить симптоматический рефлюкс, но не назначать антирефлюксную терапию для контроля плохо контролируемой астмы без симптомов рефлюкса (Evidence A). Бессимптомный рефлюкс маловероятно объясняет недостаточный контроль астмы; наиболее вероятная ограниченная польза ИПП относится к сочетанию симптоматического рефлюкса с ночными респираторными симптомами. Это не установление предпочтительного ИПП и не замена базисной противоастматической терапии. В отличие от эмпирического лечения изолированного ларингита, здесь гастроэнтерологическое показание может быть самостоятельным, но улучшение ГЭРБ и улучшение астмы оценивают раздельно.\src{gina2026}

\textbf{ЛОР-патология: что показано непосредственно для лансопразола.} В TOPPITS 346 взрослых с персистирующими глоточными симптомами получали лансопразол 30 мг дважды в сутки либо плацебо в течение 16 недель. Клинического преимущества лансопразола не обнаружено. Следовательно, 30 мг дважды в сутки на 16 недель --- изученный режим с отрицательным результатом, а не доказанный стандарт лечения хронической охриплости, ларингита или ощущения «кома в горле». Продление такого эмпирического курса не становится доказанным следующим шагом только потому, что симптомы сохраняются.\src{toppits}

Отрицательный результат относится к изученной популяции и не исключает рефлюксного механизма у каждого отдельного пациента. Однако в обзоре не представлена подгруппа, в которой конкретная доза и длительность лансопразола достоверно уменьшали перечисленные ЛОР-симптомы; позиция AAO-HNSF, приведенная выше, ограничивает эмпирическое лечение изолированной дисфонии, но не превращает отрицательный результат TOPPITS в доказанную эффективность другого ИПП. Гастроэнтерологическую диагностическую рамку нельзя выдавать за самостоятельный ЛОР-протокол. Режимы лечения собственно ГЭРБ из раздела~\ref{sec:doses} также не являются автоматически доказанными режимами лечения ее предполагаемых внепищеводных проявлений.\src{toppits,agagerd,dgvs,lyon,nice}

\textbf{Астма: контроль заболевания и кислотосупрессия --- разные конечные точки.} В исследовании 306 детей с плохо контролируемой астмой добавление лансопразола не улучшило контроль заболевания. Это не обосновывает применение препарата в качестве противоастматического средства при отсутствии самостоятельного гастроэнтерологического показания. При наличии такого показания лечение ГЭРБ оценивается отдельно: уменьшение пищеводных симптомов нельзя приравнивать к снижению частоты бронхоспазмов или улучшению легочной функции.\src{asthma}

\textbf{Дозы и отрицательные результаты: сравнение без подмены популяций.} У 306 детей с плохо контролируемой астмой без явной симптоматической ГЭРБ лансопразол назначали на 24 недели: 15 мг/сут при массе менее 30 кг и 30 мг/сут при массе не менее 30 кг. Преимущества по контролю астмы не выявлено. В другом исследовании у 402 взрослых с плохо контролируемой астмой и минимальными либо отсутствующими симптомами ГЭРБ эзомепразол 40 мг дважды в сутки также не улучшал контроль заболевания по сравнению с плацебо. Эти результаты согласуются с ограничением GINA, но не являются прямым сравнением лансопразола с эзомепразолом: возраст, популяции и режимы различались. Ни экспериментальная индукция HO-1, ни изменение IL-8 при ГЭРБ не отменяют отрицательных клинических результатов.\src{asthma,molnrf2,molil8}\src{asthmaesome2009,gina2026}

\textbf{Липофильность: где заканчивается механизм и начинается недоказанное преимущество.} Липофильность исходной молекулы влияет на мембранный перенос только совместно с ионизацией, связыванием с белками, метаболизмом и кислотной активацией. Из этого не следует достижение противовоспалительной концентрации в гортани или бронхах при стандартном приеме. В рассмотренных исследованиях и рекомендациях не установлена ни одна ЛОР- или пульмонологическая ситуация, где преимущество лансопразола по клиническому исходу было бы доказано именно за счет его липофильности. Общим свойством ИПП остается подавление желудочной кислотопродукции; оно не равнозначно общему доказанному лечению ларингита или астмы. Поэтому «класс-эффект» допустим как характеристика антисекреторного механизма, но не как заявление о доказанной эквивалентности всех ИПП по кашлю, бронхоспазмам или восстановлению голоса. Это авторский вывод из сопоставления механизма, рекомендаций и клинических испытаний, а не отдельная рекомендация о тканевом распределении.\src{molchem,pharmalabel,acggerd,aao2018,gina2026,toppits,asthma,asthmaesome2009}

\textbf{Межмолекулярное сравнение при внепищеводных проявлениях.} TOPPITS сравнивало лансопразол 30 мг дважды в сутки в течение 16 недель с плацебо, а не с омепразолом, эзомепразолом, пантопразолом или рабепразолом. Отсутствие преимущества в этой популяции не доказывает ни превосходства другого ИПП, ни одинаковой неэффективности всего класса. Аналогично, отрицательный результат добавления лансопразола у 306 детей с плохо контролируемой астмой не является сравнительным испытанием разных ИПП. В представленном материале нет прямых данных, позволяющих выбрать один из перечисленных препаратов по доказанному преимуществу в уменьшении кашля, ларингита или бронхоспазмов.\src{toppits,asthma}

\textbf{Класс-эффект: механизм нельзя подменять клиническим исходом.} Общий механизм подавления кислотопродукции не устанавливает общего доказанного эффекта на внепищеводные симптомы. Для этих исходов обзор не подтверждает ни специфического преимущества лансопразола, ни клинической эквивалентности всех ИПП. Положения AGA, DGVS и Lyon 2.0 об уточнении рефлюксного механизма определяют диагностическую стратегию, а не рейтинг молекул по эффективности при хроническом ларингите или астме. Доказанный сравнительный режим для уменьшения бронхоспазмов здесь отсутствует; режим TOPPITS следует обозначать как изученный с отрицательным результатом, а не как рекомендуемую схему.\src{pharmalabel,agagerd,dgvs,lyon,toppits,asthma}

\section{Эозинофильный эзофагит}
\label{sec:eoe}
\textbf{Место ИПП по ACG 2025.} Эозинофильный эзофагит (ЭоЭ) рассматривается как иммунноопосредованное заболевание, а не разновидность кислотного рефлюкса. ACG 2025 включает ИПП в варианты первой линии наряду с проглатываемыми местными глюкокортикостероидами и эмпирической элиминационной диетой; выбор предполагает совместное решение врача и пациента. Ответ на ИПП не исключает ЭоЭ, а неудача пробного курса больше не является обязательным диагностическим условием. Дупилумаб рассматривают преимущественно при недостаточном ответе на первоначальное лечение; дилатация при стриктурах дополняет, но не заменяет противовоспалительную терапию.\src{eoe,eoeguide,eoehighlights2025}

\textbf{Режим терапии не выводится из молекулярного механизма.} В официальной памятке ACG для взрослых указан омепразол 20 мг дважды в сутки либо 40 мг один раз в сутки или эквивалентный режим другого ИПП. Это ориентир класса, а не доказательство одинаковой эффективности всех молекул при ЭоЭ. Памятка не выделяет специальной дозы лансопразола по активности Nrf2 или подавлению эотаксина-3; выбор его режима требует учета лекарственной формы, инструкции, возможного off-label применения и объективного ответа. Усиление HO-1 в культуре клеток не служит основанием самостоятельно повышать дозу.\src{eoehighlights2025,molnrf2}

\textbf{Keap1--Nrf2--HO-1: что непосредственно показано для лансопразола.} В желудочных эпителиальных клетках крысы RGM-1 лансопразол вызывал окисление Keap1, активацию и ядерную транслокацию Nrf2, увеличение активности регуляторных элементов и экспрессии HO-1. Ослабление эффекта при подавлении Nrf2 с помощью siRNA и при вмешательстве в MEK/ERK-сигнализацию поддерживает причинное участие пути в этой модели. Однако это не эпителий пищевода пациента с ЭоЭ и не исследование клинической ремиссии. Прямое связывание препарата с Nrf2, обязательность этого пути для ответа ЭоЭ на ИПП и превосходство лансопразола над омепразолом, эзомепразолом, пантопразолом или рабепразолом здесь не установлены; подробности приведены в разделе~\ref{sec:nrf2}.\src{molnrf2}

Клеточный контекст ограничивает обобщение механизма. В эндотелиальных клетках и макрофагах лансопразол повышал HO-1 и ферритин и уменьшал образование активных форм кислорода; индукцию подавлял ингибитор PI3K LY294002, тогда как зависимость от MAPK не подтверждалась. Эти данные не позволяют считать единую цепь Keap1--Nrf2--HO-1 универсальным, одинаково выраженным эффектом всех ИПП во всех тканях. Тем более они не доказывают последовательность «активация HO-1 --- подавление эотаксина-3 --- ремиссия ЭоЭ»: для такой связи требуется отдельная экспериментальная и клиническая проверка (раздел~\ref{sec:ho1}).\src{molho1,molnrf2}

\textbf{Эотаксин-3: противовоспалительное действие не является исключительным свойством лансопразола.} В работе Zhang и соавт. (2012) на культурах плоского эпителия пищевода пациентов с ЭоЭ кислотно-активированные омепразол и лансопразол подавляли IL-4-индуцированную секрецию эотаксина-3 (CCL26). Для омепразола при концентрациях от 5 мкмоль/л также показано снижение экспрессии мРНК; детальный анализ выявил уменьшение связывания STAT6 и РНК-полимеразы II с промотором CCL26 без подавления фосфорилирования STAT6 или его ядерной транслокации. Следовательно, эффект на CCL26 продемонстрирован по меньшей мере для двух ИПП и не уникален для лансопразола. Вместе с тем детали, установленные преимущественно для омепразола, нельзя автоматически приписывать всем представителям класса; работа не дает клинического сравнения пяти молекул и не подтверждает достижение таких эффектов в пищеводе при любой пероральной дозе.\src{eoestat6}

\textbf{Критерии оценки ответа: числовой порог и качество биопсии.} Контроль включает симптомы, эндоскопическую оценку EREFS и гистологию; уменьшение дисфагии само по себе недостаточно. Получают не менее шести прицельных биоптатов как минимум из двух уровней пищевода и количественно определяют эозинофилы. В официальном разъяснении ACG 2025 успешным гистологическим ответом считается \textbf{менее 15 эозинофилов в поле большого увеличения ($<15$ eos/hpf)}. В практическом комментарии к этому документу контрольная эндоскопия после начала ИПП проводится через 8--12 недель; это описанный подход к мониторингу, а не жесткий срок для каждого пациента и любой терапии.\src{eoeguide,eoehighlights2025}

\textbf{Полная гистологическая ремиссия: не подменять термин порогом ответа.} Порог $<15$ eos/hpf нельзя обозначать как доказанное полное восстановление слизистой. Более строгий критерий $\leq6$ eos/hpf использован, например, как гистологическая конечная точка исследования дупилумаба III фазы; это критерий конкретного испытания, а не установленный ACG универсальный критерий «полной ремиссии» на лансопразоле. Даже такой результат не тождествен нормализации всей тканевой архитектуры. Помимо количества эозинофилов учитывают гиперплазию базального слоя, расширение межклеточных пространств и фиброз собственной пластинки при ее наличии в биоптате. Поэтому в настоящей монографии следует указывать фактическое число эозинофилов, выбранный порог и остаточные морфологические изменения; термин «полная гистологическая ремиссия» без заранее определенных критериев не используется. Ни исчезновение симптомов, ни снижение CCL26 или повышение HO-1 этих критериев не заменяют.\src{eoeguide,eoedupilumab2022}

\textbf{Поддержание эффекта.} При объективно подтвержденном ответе ACG предусматривает продолжение эффективного лечения для предупреждения рецидива симптомов, гистологического воспаления и эндоскопических изменений. Сохраняющаяся дисфагия при уменьшении эозинофилии требует отдельной оценки стриктурного компонента, а симптоматическое улучшение при сохраняющемся воспалении не подтверждает достаточный контроль заболевания.\src{eoe,eoehighlights2025}

\textbf{Сравнение ИПП в рамках ACG 2025.} Рекомендация относится к классу ИПП; лансопразол не выделен как предпочтительный перед омепразолом, эзомепразолом, пантопразолом или рабепразолом по активности Keap1--Nrf2--HO-1 либо частоте гистологической ремиссии. Отсутствие такого предпочтения не доказывает клиническую эквивалентность всех молекул и доз. Решение об эффективности основывается на объективном контроле ЭоЭ, а не на предполагаемой выраженности отдельного внепомпового эффекта.\src{eoe,eoeguide}

\textbf{Keap1--Nrf2--HO-1 и эотаксин-3: границы механистического сравнения.} Для лансопразола описана индукция HO-1 в непищеводных моделях; подавление IL-4-индуцированной секреции CCL26 показано как для лансопразола, так и для омепразола в клетках пищевода. Это поддерживает общность части противовоспалительных свойств нескольких ИПП, но не универсальную идентичность механизмов всего класса. Причинная связь этих двух путей между собой и специфическое преимущество лансопразола в достижении полной гистологической ремиссии ЭоЭ не установлены.\src{molnrf2,molho1,eoestat6}

\section{Профилактика стрессовых кровотечений}
\label{sec:stress-prophylaxis}
\textbf{Критерий назначения --- риск, а не само пребывание в ОРИТ.} В изложении SCCM/ASHP 2024 профилактическая кислотосупрессия рассматривается у критически больных с риском стрессового кровотечения, а не у всех госпитализированных. Среди значимых обстоятельств названы коагулопатия, шок и хроническая болезнь печени. При наличии риска используются ИПП или H$_2$-блокаторы с учетом клинической ситуации; после исчезновения факторов риска профилактику прекращают. Эти факторы обосновывают оценку необходимости профилактики, но сами по себе не устанавливают предпочтительность лансопразола.\src{sccm}

\textbf{Имеются ли преимущества именно лансопразола?} Обзор не устанавливает его универсального преимущества перед другими ИПП для стресс-профилактики. Здесь не приведены количественное сравнительное преимущество, доказанный эффект на летальность, отдельные доза, путь введения и длительность именно лансопразолового режима. Поэтому этот подраздел нельзя использовать как самостоятельный протокол назначения при шоке или коагулопатии. Отсутствие доказанного превосходства не является доказательством эквивалентности всех препаратов и способов введения; недостаток подробностей в обзоре также не означает отсутствия соответствующих положений в полном руководстве.\src{sccm}

\textbf{Не переносить режим лечения кровотечения на профилактику.} В российских рекомендациях профессиональных обществ 2024 года описан внутривенный лансопразол 30 мг дважды в сутки в виде 30-минутной инфузии до 7 дней с последующим переходом на прием внутрь. Этот режим относится к контексту язвенного кровотечения (раздел~\ref{sec:ulcer-bleeding}), а не является универсальной схемой стресс-профилактики. Аналогично, амбулаторные 15--30 мг/сут не равнозначны интенсивной кислотосупрессии после эндоскопического гемостаза в стратегии ESGE 2026. Доступность и регистрационный статус парентеральной формы необходимо проверять отдельно.\src{rfpud,esge2026,grls}

\textbf{Форма и доставка препарата.} При рассмотрении введения через зонд проверяют допустимость метода для конкретного продукта, диаметр зонда, растворитель и объем промывания (раздел~\ref{sec:formulations}). Гастрорезистентные гранулы нельзя разжевывать или превращать в обычный порошок; диспергируемая таблетка и содержимое капсулы могут требовать разных способов подготовки. Техническая возможность введения не доказывает преимущества энтерального лансопразола при шоке и не заменяет оценки конкретной клинической ситуации.\src{uslabel,emc,sccm}

\textbf{Лансопразол против пантопразола и омепразола в ОРИТ.} В приведенном изложении SCCM/ASHP 2024 выбор профилактической кислотосупрессии основан на риске стрессового кровотечения; при наличии показаний рассматриваются ИПП или H$_2$-блокаторы. Отдельных прямых сравнений лансопразола с пантопразолом или омепразолом по клинически значимым кровотечениям, летальности либо безопасности в обзоре нет. Поэтому преимущество лансопразола в этой нише не установлено. Из этого, однако, не следует доказанная эквивалентность молекул или всех способов введения: класс-ориентированная рекомендация и сравнительное исследование решают разные вопросы.\src{sccm}

\textbf{Внутривенные формы: не конструировать эквивалентную схему по миллиграммам.} Приведенный режим внутривенного лансопразола 30 мг дважды в сутки, инфузия 30 минут до 7 дней, относится к язвенному кровотечению, а не к доказанному преимуществу в стресс-профилактике. Сопоставимые профилактические внутривенные режимы пантопразола и омепразола в этом изложении отсутствуют. Следовательно, нельзя ни объявить указанную схему лансопразола лучшей, ни пересчитать ее в эквивалентные дозы других ИПП. При шоке или коагулопатии наличие показания к профилактике не доказывает преимуществ конкретной молекулы, а техническая возможность энтерального введения не устанавливает его равнозначности внутривенному. Регистрационный статус и доступность формы проверяют отдельно.\src{rfpud,sccm,esge2026,uslabel,grls}

\section{Лансопразол-сульфид и микобактерии: другая молекула, другая мишень}
\label{sec:mycobacteria}
{\raggedright \textbf{Два разных превращения пролекарства.} Для желудочного действия исходного лансопразола-сульфоксида существенны накопление в кислой среде, химическая активация и дисульфидная модификация доступных цистеинов H$^+$/K$^+$-АТФазы. В противомикобактериальном направлении действующую форму связывают с восстановленным производным --- лансопразол-сульфидом (LPZS). Поэтому обозначение «пролекарство» в этих двух контекстах не описывает одну и ту же реакцию: образование LPZS нельзя отождествлять с кислотным образованием сульфенамида. Исходный сульфоксид, восстановленный сульфид, печеночный сульфон и кислотно-активированный сульфенамид необходимо различать.\src{moltb,molnagaya89,molcys}\par}

\textbf{Функционально-генетическое обоснование.} В работе Rybniker и соавт. 2015 года внутриклеточное образование LPZS связывали с антибактериальным действием лансопразола, исследованного как экспериментальное противотуберкулезное пролекарство. Устойчивые мутанты указывали на QcrB --- компонент микобактериального комплекса $bc_1$. Такое наблюдение связывает эффект с бактериальным дыхательным аппаратом, но не означает, что у микобактерии воспроизводится механизм блокады желудочной помпы. Генетическая связь устойчивости с мишенью поддерживает механистическую гипотезу; сама по себе она не устанавливает клиническую эффективность у человека.\src{moltb}

\textbf{Структурный механизм: участок окисления менахинола.} В работе 2024 года LPZS исследовали с дыхательным суперкомплексом $\mathrm{III}_2\mathrm{IV}_2$ \textit{Mycobacterium smegmatis} методами криоэлектронной микроскопии, функциональных опытов и моделирования. Описано связывание в области окисления менахинола $Q_o$, включая взаимодействие с His368 соответствующего участка, и нарушение дыхательного электронного переноса. Таким образом, другая мишень --- компонент бактериального дыхательного комплекса, а не ATP4A париетальной клетки. Эти результаты не являются доказательством ковалентной модификации QcrB по механизму дисульфидной блокады желудочной H$^+$/K$^+$-АТФазы.\src{moltbstructure}

\textbf{Почему различие модели принципиально.} Работа 2015 года связывает образование производного и устойчивость с антибактериальным действием; структурная работа 2024 года уточняет положение LPZS и функциональные последствия взаимодействия с дыхательным комплексом. Совокупность методов усиливает биологическое обоснование направления, но не устраняет межвидовых и фармакокинетических ограничений. В частности, исследование комплекса \textit{M. smegmatis} нельзя обозначать как клиническое испытание лечения туберкулеза или прямое доказательство эффективности обычной дозы лансопразола у пациента.\src{moltb,moltbstructure}

\textbf{Клинический статус, представленный в обзоре.} Описаны доклинические результаты, но не подтвержденная эффективность лечения туберкулеза у человека, клиническая дозировка LPZS или рекомендованная противотуберкулезная комбинация с лансопразолом. Наличие подходящей бактериальной мишени не доказывает достижения эффективной экспозиции после стандартной гастроэнтерологической дозы. Следовательно, обычный лансопразол нельзя считать клинически апробированным эквивалентом LPZS или самостоятельно добавлять к противотуберкулезной терапии на основании этих работ.\src{moltb,moltbstructure}

\textbf{Рифампицин: взаимодействие не равно терапевтической комбинации.} В разделе~\ref{sec:interactions} рифампицин указан как средство, способное снижать экспозицию ИПП. Это предупреждение не сообщает концентрацию LPZS в очаге инфекции и не доказывает эффективность либо безопасность противотуберкулезного сочетания. Режим компенсационного увеличения дозы лансопразола в обзоре не приведен. Поэтому из механистических и фармакокинетических наблюдений нельзя конструировать клиническую схему лечения.\src{uslabel,moltb,moltbstructure}

\section{Нефрологический профиль безопасности и маркеры интерстициального повреждения}
\label{sec:nephrology}

\subsection{Молекулярные маркеры интерстициального повреждения почек при терапии лансопразолом}
\label{sec:renal-interstitial-markers}
\textbf{Иммунный патогенез и причинная атрибуция.} ИПП-ассоциированный острый интерстициальный нефрит (ОИН) представляет воспалительное повреждение почечного интерстиция и канальцев, а не нефротический синдром и не прямое следствие блокады желудочной H$^+$/K$^+$-АТФазы. В патогенетической модели замедленной гиперчувствительности взаимодействие препарата или метаболита с белком может создавать иммуногенный комплекс, распознаваемый клеточным иммунитетом. Однако образование конкретного гаптенового аддукта лансопразола не доказано для каждого случая; временная связь назначения с ростом креатинина также не идентифицирует молекулярный механизм. Клиническое подозрение требует анализа всей лекарственной экспозиции, альтернативных причин и динамики после отмены, а при соответствующих показаниях --- морфологической верификации.\src{safeain,uslabel,renalniceaki} Лихорадка, сыпь и эозинофилия могут поддерживать предположение об иммунной реакции, но их отсутствие не исключает ОИН. Согласно инструкции, острый тубулоинтерстициальный нефрит возможен на любом этапе приема лансопразола, в том числе без характерных внепочечных проявлений. Отсутствие обычной коррекции дозы при почечной недостаточности не означает отсутствия иммунного риска.\src{uslabel}

\textbf{Функциональные показатели и исследовательские молекулярные маркеры.} Сывороточный креатинин и рСКФ характеризуют фильтрационную функцию, но не специфичны для интерстициального воспаления; лейкоцитурия, гематурия и протеинурия также требуют дифференциальной диагностики. Эозинофилурия не является самостоятельным подтверждающим или исключающим тестом.\src{renalniceaki,renaleosin2013} В биопсийно ориентированном исследовании Moledina и соавт. мочевые TNF-$\alpha$ и IL-9 улучшали различение ОИН и других причин повреждения почек по сравнению с клинической оценкой и доступными стандартными тестами. Эти результаты характеризуют воспалительный фенотип, но не доказывают гаптенизацию определенным метаболитом и не устанавливают специфический тест токсичности лансопразола. В настоящем алгоритме они рассматриваются как исследовательские показатели, а не обязательный скрининг всех получающих ИПП. Их определение не должно задерживать отмену подозреваемого препарата, оценку ОПП или консультацию нефролога; перенос диагностических характеристик на другую популяцию требует отдельной валидации.\src{renalmarkers2019}

\textbf{Сравнительные риски: клинические исходы и ограничения исследований 2025--2026 годов.} В ретроспективной когорте BIGAN, опубликованной онлайн 3 декабря 2024 года и включенной в выпуск 2025 года, скорректированный HR подтвержденного снижения рСКФ ниже 60 мл/мин/1,73 м$^2$ относительно ранитидина составлял 1,08 (95\% ДИ 0,69--1,68) для лансопразола, 0,99 (0,66--1,48) для омепразола и 0,98 (0,64--1,50) для пантопразола. Сопоставление этих величин не выявляет доказанного молекулоспецифичного преимущества, но также не является исследованием эквивалентности: каждый ИПП сравнивался с общим компаратором, а не непосредственно с другим ИПП. Биопсийно верифицированный ОИН и последующий переход этого ОИН в ХБП не составляли отдельной последовательной конечной точки.\src{renalbigan2025} Анализ TriNetX 2025 года характеризует ассоциации класса ИПП с почечными исходами, а FAERS 2026 года оценивает непропорциональность спонтанных сообщений посредством ROR. Наличие контрольной группы в ретроспективной когорте не превращает ее в рандомизированное испытание; ROR фармаконадзора нельзя подменять клиническим OR или HR. Эти источники не предоставляют сопоставимых оценок риска цепочки «ОИН --- ХБП» для трех названных молекул. Смешение по показанию, исходная функция почек, полиморбидность и сопутствующие нефротоксичные воздействия ограничивают причинную интерпретацию. Следовательно, отсутствие нужной сравнительной оценки означает неопределенность, а не отсутствие риска или основание объявлять лансопразол нефропротектором.\src{renaltrinetx2025,renalfaers2026}

\textbf{Риск-ориентированное лабораторное наблюдение.} При установленной ХБП либо дополнительных факторах почечного риска перед началом или пересмотром длительной терапии фиксируют исходные креатинин, рСКФ, калий и показатели мочи; частоту повторного обследования определяют клиническим состоянием и стадией ХБП, а не единым календарем для всех получающих ИПП. У пациентов без дополнительных факторов риска рутинный частый контроль креатинина исключительно по причине назначения ИПП не является универсальным требованием.\src{renalkdigo2024,acggerd} Основанием для незамедлительной оценки острого повреждения почек служат рост креатинина не менее чем на 26 мкмоль/л за 48 часов по критерию NICE, увеличение не менее чем в 1,5 раза за предшествующие 7 суток либо диурез менее 0,5 мл/кг/ч в течение более 6 часов у взрослого. Эти пороги относятся к ОПП любой этиологии, а не специфически к ОИН; достижение порога внутри лабораторного референсного интервала также клинически значимо.\src{renalniceaki} При быстро меняющемся креатинине расчетная СКФ не является надежным стационарным измерением фильтрации, поэтому приоритетны динамика креатинина, диурез и водно-электролитный баланс. У пациента со стабильной ХБП снижение рСКФ более чем на 20\% при последующем определении требует оценки причины, но не доказывает лекарственный нефрит. При интерпретации сопоставляют результат с документированным исходным уровнем, оценивают объемный статус и изменения сопутствующей терапии; повторное измерение не заменяет немедленных действий при выполнении критериев ОПП. Персистирование функциональных или структурных нарушений не менее 3 месяцев необходимо для установления хронического характера процесса; ожидание этого срока не должно задерживать диагностику и лечение острого эпизода.\src{renalniceckd,renalkdigo2024}

\textbf{Диагностический алгоритм при подозрении на замедленную гиперчувствительность.} На первом этапе устанавливают временную связь ухудшения с началом и изменениями лекарственной терапии, уточняют все препараты, включая НПВП и антибиотики, оценивают сыпь, лихорадку и другие системные проявления. При обоснованном подозрении на ИПП-индуцированный ОИН отменяют подозреваемый препарат, не ожидая появления полной клинической триады и не заменяя его автоматически другим ИПП.\src{uslabel} На втором этапе повторно определяют креатинин, калий и другие электролиты, исследуют мочу и осадок, выполняют общий анализ крови с лейкоцитарной формулой; исключают гиповолемию, инфекцию, обструкцию и иные причины ОПП. Ни эозинофилия крови, ни эозинофилурия не являются самостоятельным подтверждающим или исключающим тестом.\src{renalniceaki,renaleosin2013} На третьем этапе при подозрении на тубулоинтерстициальный нефрит обсуждают случай с нефрологом в пределах 24 часов; выраженные нарушения электролитного, кислотно-основного или водного баланса требуют немедленной специализированной помощи. Показания к биопсии определяют по диагностической неопределенности, тяжести и динамике нарушения функции, особенно если морфологический результат изменит тактику.\src{renalniceaki} Глюкокортикостероиды не назначают автоматически по факту сыпи или эозинофилурии: решение принимает нефролог с учетом вероятности лекарственного ОИН, инфекции, морфологических данных и ограничений доказательств лечения. После отмены документируют динамику восстановления функции и остаточное повреждение, обеспечивают последующее наблюдение и регистрацию подозреваемой нежелательной реакции; сохранение ХБП после ОИН оценивают индивидуально, без переноса групповых HR на прогноз конкретного пациента.\src{renalsteroids2018,renalkdigo2024}

\textbf{Сравнительная эпидемиология: необходимо различать ОИН и ОПП.} В популяционной когорте Antoniou и соавт. (2015) 290\,592 новых пользователей ИПП в возрасте не менее 66 лет сопоставлены с таким же числом лиц без назначения ИПП. За 120 дней наблюдения HR госпитализации с острым повреждением почек (ОПП) составил 2,52 (95\% ДИ 2,27--2,79), а для зарегистрированного ОИН --- 3,00 (1,47--6,14). Молекулоспецифичные оценки авторы представили для \emph{ОПП}, а не для биопсийно подтвержденного ОИН.\src{renalantoniou2015}
\begin{longtable}{P{0.25\textwidth}P{0.26\textwidth}P{0.40\textwidth}}
\toprule
ИПП & HR ОПП (95\% ДИ) & Основание сравнения \\
\midrule
\endhead
Лансопразол & 2,56 (1,85--3,55) & Соответствующая сопоставленная группа без ИПП \\
Омепразол & 2,94 (2,21--3,91) & Соответствующая сопоставленная группа без ИПП \\
Пантопразол & 2,43 (1,97--3,00) & Соответствующая сопоставленная группа без ИПП \\
\bottomrule
\end{longtable}
Эти HR не являются прямыми отношениями рисков лансопразола к омепразолу или пантопразолу. Их численное различие не устанавливает превосходства одной молекулы; наблюдательный дизайн и административная регистрация исходов ограничивают причинную интерпретацию. Нельзя переименовывать таблицу в «сравнительный риск ОИН» или выводить из нее доказанную нефропротекцию лансопразола.\src{renalantoniou2015}

В новозеландском исследовании Blank и соавт. (2014; когорта 572\,661 пользователя омепразола, пантопразола или лансопразола) для биопсийно подтвержденных случаев ОИН отношение шансов текущего приема к прошлому составило OR 5,16 (95\% ДИ 2,21--12,05). В анализе лиц, получавших только омепразол, соответствующее OR равнялось 4,00 (1,70--9,42). Для всей группы частота подтвержденных и вероятных случаев составляла 11,98 против 1,68 на 100\,000 человеко-лет при текущем и прошлом приеме. Отдельных сопоставимых OR для лансопразола и пантопразола здесь нет. Таким образом, доказательный материал поддерживает настороженность в отношении класса, но не полноценный сравнительный рейтинг трех препаратов по ОИН. OR, HR и абсолютную частоту нельзя использовать как взаимозаменяемые показатели.\src{renalblank2014}

\textbf{Что именно сравнивается в работах 2025--2026 годов.} Для практического чтения следует различать три вопроса: возникновение ОИН после начала ИПП, развитие или прогрессирование ХБП среди всех пользователей и сохранение ХБП после верифицированного ОИН. BIGAN и TriNetX характеризуют почечные исходы в группах экспозиции; FAERS выявляет непропорциональность сообщений. Ни один из этих подходов сам по себе не оценивает условную вероятность ХБП после ОИН для каждой молекулы. Поэтому отсутствие нужного OR/HR в таблице следует обозначать как отсутствие оценки, а не как отсутствие риска или равенство препаратов.\src{renalbigan2025,renaltrinetx2025,renalfaers2026}

\textbf{Когорта 2025 года: клинические HR, но не риск ОИН.} В исследовании BIGAN (González-Pérez и соавт.; выпуск 2025 года, онлайн 3 декабря 2024 года) включены 122\,606 новых пользователей кислотосупрессоров с исходной расчетной СКФ выше 60 мл/мин/1,73 м$^2$. Таблица воспроизводит скорректированные HR основного анализа во время лечения; каждый ИПП сравнивался с ранитидином, а не непосредственно с другим ИПП.\src{renalbigan2025}
\begin{longtable}{P{0.22\textwidth}P{0.34\textwidth}P{0.35\textwidth}}
\toprule
Препарат & HR подтвержденного снижения СКФ $<60$ (95\% ДИ) & HR ОПП по критериям Aberdeen (95\% ДИ) \\
\midrule
\endhead
Лансопразол & 1,08 (0,69--1,68) & 0,57 (0,42--0,78) \\
Омепразол & 0,99 (0,66--1,48) & 0,54 (0,42--0,70) \\
Пантопразол & 0,98 (0,64--1,50) & 0,44 (0,33--0,59) \\
\bottomrule
\end{longtable}
\textbf{Размер подгрупп и абсолютная частота уточняют смысл HR.} В BIGAN омепразол получали 105\,447 человек, пантопразол --- 4\,244, лансопразол --- 3\,784; ранитидин --- 3\,086. Для подтвержденного снижения СКФ ниже 60 мл/мин/1,73 м$^2$ нескорректированные частоты в основном анализе составляли соответственно 31,2; 44,4; 34,8 и 18,7 на 1\,000 человеко-лет. Эти показатели не являются процентами пациентов и не относятся к ОИН. Более высокая исходная частота при близком к единице скорректированном HR подчеркивает влияние состава групп: нельзя выводить сравнительную токсичность из одних абсолютных частот или делить независимые HR для получения прямого эффекта переключения.\src{renalbigan2025}

Средняя длительность основного анализа составляла лишь 0,7 года. Отсутствие значимого увеличения первого исхода не доказывает эквивалентности препаратов; HR ОПП ниже единицы не устанавливает нефропротекцию. Сравниваемые группы различались по клиническому составу, сохранялась возможность остаточного смешения. Биопсийный ОИН и последующая ХБП после него не являлись отдельной последовательной конечной точкой. Исторический компаратор не означает рекомендации назначать ранитидин.\src{renalbigan2025}

\textbf{Согласованность результатов 2025 года ограничена.} Исследование Chen и соавт. в TriNetX, опубликованное 13 июля 2025 года, после сопоставления по склонности к назначению, напротив, выявило более высокий риск неблагоприятных почечных исходов у пользователей ИПП относительно H$_2$-блокаторов. Для исхода со СКФ $<15$ мл/мин/1,73 м$^2$ HR составил 1,785 (95\% ДИ 1,718--1,854). Это оценка по классу, не специфическая для лансопразола и не условный риск ХБП среди перенесших ОИН. Различия дизайна, определения экспозиции и состава пациентов препятствуют механическому объединению этой оценки с BIGAN; статистическое сопоставление не устраняет все неизмеренные факторы.\src{renaltrinetx2025}

\textbf{Фармаконадзор 2026 года: ROR не равен клиническому OR или HR.} Wang и соавт. (2 февраля 2026 года) проанализировали 176\,680 случаев из FAERS за 2004--2024 годы. Ниже приведены сигналы из таблицы 2 публикации: отношение шансов сообщения (ROR) для конкретного препарата относительно остальных сообщений базы.\src{renalfaers2026}
\begin{longtable}{P{0.22\textwidth}P{0.34\textwidth}P{0.35\textwidth}}
\toprule
Препарат & ROR тубулоинтерстициального нефрита (95\% ДИ) & ROR ХБП (95\% ДИ) \\
\midrule
\endhead
Лансопразол & 18,22 (17,24--19,25) & 45,80 (44,86--46,75) \\
Омепразол & 11,01 (10,39--11,66) & 8,25 (7,99--8,52) \\
Пантопразол & 12,50 (11,76--13,29) & 22,39 (21,85--22,94) \\
\bottomrule
\end{longtable}
Это не частоты заболевания: отсутствует знаменатель всех экспонированных пациентов, а вероятность сообщения зависит от препарата и обстоятельств регистрации. Термин тубулоинтерстициального нефрита в базе не гарантирует биопсийную верификацию острого лекарственного процесса. Более высокий ROR лансопразола не доказывает большего причинного риска; отдельные сигналы нефрита и ХБП не устанавливают переход между ними.\src{renalfaers2026}

\textbf{Что можно заключить о переходе ОИН в ХБП.} В рассмотренных крупных работах 2025--2026 годов нет сопоставимых молекулоспецифичных OR/HR для последовательности «верифицированный ИПП-индуцированный ОИН --- сохраняющаяся ХБП». Для такой оценки нужны подтверждение причины ОИН, исходная функция почек, сроки отмены, последующее наблюдение и прямое сопоставление препаратов. Нельзя делить ROR ХБП на ROR нефрита или переносить общий HR снижения СКФ на пациентов после ОИН. Клиническая задача --- раннее выявление повреждения и оценка восстановления, а не выбор предположительно «нефробезопасного» ИПП по несопоставимым показателям.\src{renalbigan2025,renaltrinetx2025,renalfaers2026,renalniceaki}

\textbf{От ОИН к хроническому снижению функции: критерий исхода, а не сигнал базы.} При подозреваемом лекарственном ОИН необходимо связать в одной временной последовательности исходные показатели, начало повреждения, дату отмены и восстановление. Если нарушение сохраняется, сопоставляют расчетную СКФ и маркеры повреждения с докризисными значениями: уже существовавшая ХБП, неполное восстановление после ОПП и впервые установленная ХБП не являются одним исходом. При хроническом течении классификацию дополняют категорией альбуминурии. Документирование этой последовательности позволяет оценивать конкретного пациента, но не заменяет отсутствующего сравнительного исследования трех ИПП.\src{renalniceaki,renalniceckd}

\textbf{Ранние лабораторные ориентиры.} У коморбидного пациента с почечными факторами риска фиксируют исходный креатинин, расчетную СКФ и предшествующую динамику; последующие результаты сравнивают с этой базой, а не только с референсным диапазоном лаборатории. При быстро меняющемся креатинине расчетная СКФ менее надежна: ее снижение является сигналом пересмотра состояния, но не самостоятельным диагностическим критерием ОИН. Универсальная программа частых анализов для каждого здорового пользователя ИПП не обоснована; у пациента с ХБП, острым заболеванием или нефротоксичной полипрагмазией наблюдение индивидуализируют.\src{renalniceckd,acggerd}

\textbf{Расчетная СКФ: контроль динамики при установленной ХБП.} KDIGO 2024 рекомендует оценивать СКФ и альбуминурию не реже одного раза в год, чаще --- при повышенном риске прогрессирования, если результат изменит лечение. Изменение расчетной СКФ более чем на 20\% при последующем исследовании превышает ожидаемую вариабельность и требует оценки; удвоение ACR также требует выяснения причины. Это сигналы для обследования, не пороги диагностики ОИН и не основания автоматически отменять любой ИПП. При недостаточной надежности креатинина и значимости СКФ для решения используют совместную оценку по креатинину и цистатину C (eGFRcr-cys). Она уточняет фильтрацию, но не определяет лекарственную этиологию повреждения.\src{renalkdigo2024}

\textbf{Направление к нефрологу при хроническом ухудшении.} Среди критериев NICE NG203 --- устойчивое снижение расчетной СКФ не менее чем на 25\% с переходом в другую категорию СКФ за 12 месяцев либо устойчивое снижение не менее чем на 15 мл/мин/1,73 м$^2$ за 12 месяцев. Эти критерии длительного наблюдения не заменяют срочной оценки при подозрении на ОПП/ОИН: ожидание накопления годовой динамики недопустимо. У пациента без факторов риска рутинное назначение столь же интенсивного мониторинга только из-за приема лансопразола не следует из этих рекомендаций.\src{renalniceckd,renalniceaki,acggerd}

По NICE NG148 признаками ОПП служат увеличение креатинина не менее чем на 26 мкмоль/л за 48 часов, рост не менее чем на 50\% за 7 дней либо диурез менее 0,5 мл/кг/ч более 6 часов у взрослого. Эти пороги требуют оценки причины и тяжести, но не доказывают ОИН. В стационаре при ОПП или его риске креатинин обычно контролируют ежедневно; вне стационара частоту определяет клиническая ситуация. При подозрении исследуют мочу на кровь, белок, лейкоциты, нитриты и глюкозу, оценивают электролиты и водный баланс.\src{renalniceaki}

Эозинофилурия не должна служить тестом, исключающим ОИН: в исследовании Muriithi и соавт. с биопсийной верификацией ее диагностические характеристики были недостаточными. Нормальный результат не разрешает продолжать подозреваемый препарат при нарастающем повреждении почек; изменения мочевого осадка также требуют дифференциальной диагностики, а не автоматического приписывания ИПП.\src{renaleosin2013}

\textbf{Лабораторная настороженность: порог и действие.} Критерии ОПП, перечисленные выше, применяют к изменению относительно исходного креатинина даже при значении внутри лабораторного референсного интервала. Следующие ориентиры дополняют, но не заменяют клиническую оценку.
\begin{longtable}{P{0.31\textwidth}P{0.60\textwidth}}
\toprule
Сигнал & Интерпретация и действие интерниста \\
\midrule
\endhead
Достигнут любой критерий ОПП; нарастает олигурия & Срочно установить причину, оценить калий, кислотно-основное состояние и водный баланс; при подозрении на тубулоинтерстициальный нефрит обсудить случай с нефрологом в пределах 24 часов. Рефрактерные гиперкалиемия, ацидоз или отек легких требуют немедленной специализированной помощи.\src{renalniceaki} \\
Новая гематурия, протеинурия или лейкоцитурия на фоне роста креатинина & Исследовать мочевой осадок и исключить инфекцию и другие почечные заболевания. Ни наличие эозинофилов мочи, ни их отсутствие не подтверждает и не исключает ОИН самостоятельно.\src{renalniceaki,renaleosin2013} \\
Впервые выявлена СКФ $<60$ мл/мин/1,73 м$^2$ без признаков острого ухудшения & Повторить определение в пределах 2 недель; при подозрении на ОПП не ждать планового срока. Хронический характер устанавливают по сохранению нарушений не менее 3 месяцев, а не по одному анализу.\src{renalniceckd,renalniceaki} \\
Альбумин/креатинин мочи (ACR) $\geq3$ мг/ммоль & Значение 3--70 мг/ммоль подтвердить в последующей утренней пробе. Стойкая альбуминурия --- маркер повреждения почек, но не специфический тест ИПП-индуцированного ОИН. После острого эпизода оценивать динамику функции и маркеров повреждения; нельзя откладывать помощь на 3 месяца ради подтверждения ХБП.\src{renalniceckd,renalniceaki} \\
\bottomrule
\end{longtable}

\textbf{Действия врача при подозрении на лекарственный ОИН.} Следующий порядок представляет клинический синтез инструкции и нефрологических рекомендаций, а не схему самолечения.
\begin{enumerate}
\item Прекратить подозреваемый лансопразол, зафиксировать сроки экспозиции и нежелательную реакцию, пересмотреть все лекарства. Простое снижение дозы не заменяет отмену; автоматическая замена другим ИПП без нефрологической оценки не обоснована.\src{uslabel,renalblank2014}
\item Проверить динамику креатинина, диурез, калий и кислотно-основное состояние; оценить гиповолемию, сепсис, обструкцию и другие причины ОПП. Обсудить вероятный тубулоинтерстициальный нефрит с нефрологом как можно раньше, в пределах 24 часов. При рефрактерной гиперкалиемии, ацидозе или отеке легких требуется немедленная специализированная помощь.\src{renalniceaki}
\item Решить вопрос о биопсии при неясной причине, тяжелом или сохраняющемся нарушении функции, особенно если результат изменит лечение. Назначение глюкокортикоидов определяет нефролог с учетом морфологии, восстановления после отмены и риска инфекции. В ретроспективной серии 182 биопсийно подтвержденных лекарственных ОИН задержка стероидной терапии ассоциировалась с худшим восстановлением, но это не рандомизированное доказательство универсальной пользы пульс-терапии.\src{renaleosin2013,renalsteroids2018}
\item Контролировать восстановление функции после выписки и документировать подозреваемый препарат для предотвращения повторной экспозиции. Сохраняющееся снижение СКФ требует дальнейшего нефрологического наблюдения, а необходимость кислотосупрессии пересматривают отдельно.\src{renalniceaki,uslabel}
\end{enumerate}

\section{Гастропротекция в онкогематологии и интенсивной эндокринологии}
\label{sec:onco-endocrine}
\textbf{Показание определяется риском, а не интенсивностью лечения как таковой.} Высокодозная пульс-терапия метилпреднизолоном, онкогематологический диагноз и прием перорального ингибитора тирозинкиназ не образуют самостоятельного универсального показания к лансопразолу. В позиции SIF--AIGO--FIMMG рутинная профилактика ИПП при изолированном применении кортикостероидов не поддерживается; существенны язвенный анамнез и сочетание с НПВП. У коморбидного пациента дополнительно оценивают антитромботическую нагрузку и сохраняющийся риск верхнего желудочно-кишечного кровотечения. Если развивается критическое состояние, решение о стресс-профилактике принимают по отдельным критериям раздела~\ref{sec:stress-prophylaxis}.\src{oncoppiposition2016,cardioesc2024,sccm}

\textbf{Эндокринологическая конкретизация.} При активной орбитопатии Грейвса EUGOGO 2021 рассматривает внутривенный метилпреднизолон в специализированном центре; в разделе безопасности предусмотрено применение ИПП по необходимости, а не обязательный выбор лансопразола для каждого пациента. Следовательно, пульс-терапия требует оценки показания к гастропротекции наряду с контролем гликемии, артериального давления, печени и инфекционных рисков. Эти положения нельзя превращать в доказанную профилактическую схему «метилпреднизолон плюс ИТК плюс лансопразол»: подобная комбинация требует самостоятельной проверки каждого компонента.\src{oncoeugogo2021}

\textbf{Границы защиты слизистой.} Кислотосупрессия не устраняет саму миелосупрессию, тромбоцитопению или лекарственное повреждение других отделов ЖКТ. Ее не следует описывать как универсальную профилактику химиотерапевтического мукозита либо всех кровотечений. При установленном показании дозу и длительность лансопразола выбирают по соответствующей кислотозависимой патологии или риску: режимы 15 и 30 мг не являются двумя доказанными стандартами сопровождения любой пульс-терапии. После исчезновения показания назначение пересматривают; увеличение дозы «на всякий случай» особенно нежелательно при pH-чувствительной противоопухолевой терапии.\src{uslabel,oncoppiposition2016,oncocaltablet2026,agadepres}

\textbf{Два механизма взаимодействия нельзя смешивать.} У ряда пероральных ИТК растворимость и всасывание уменьшаются при повышении желудочного pH. Это фармацевтическое и абсорбционное взаимодействие, а не только конкуренция за CYP2C19. Поэтому предпочтительность лансопразола относительно омепразола при клопидогреле не переносится на онкологические молекулы. Короткий плазменный период полувыведения ИПП также не означает короткого изменения pH: разнесение приема на несколько часов не гарантирует устранения взаимодействия. Проверка CYP3A, сопутствующих ингибиторов и индукторов проводится дополнительно, независимо от проверки кислотосупрессии.\src{oncobosulif2026,oncodasatinib2026,oncocalcaps2026}

\textbf{Примеры, значимые для онкогематологии.} Таблица относится к указанным продуктам и лекарственным формам; ее нельзя распространять на все препараты с тем же действующим веществом без проверки инструкции.
\begin{longtable}{P{0.22\textwidth}P{0.34\textwidth}P{0.35\textwidth}}
\toprule
Препарат / форма & Проверенные данные & Следствие для выбора гастропротекции \\
\midrule
\endhead
Бозутиниб (Bosulif) & Лансопразол 60 мг при повторном приеме снижал $C_{\max}$ бозутиниба на 46\%, AUC --- на 26\%; изучалась однократная доза бозутиниба 400 мг натощак. & Инструкция предлагает вместо ИПП короткодействующий антацид или H$_2$-блокатор с интервалом более 2 часов от бозутиниба. Результат не определяет точную величину взаимодействия для 15/30 мг лансопразола.\src{oncobosulif2026} \\
Дазатиниб: стандартные таблетки в цитируемой инструкции & Даже через 22 часа после омепразола 40 мг AUC дазатиниба снижалась на 43\%, $C_{\max}$ --- на 42\%. & Не сочетать с ИПП или H$_2$-блокаторами. Антациды допустимы не менее чем за 2 часа до или после дозы; замена ИПП на фамотидин не является универсальным решением.\src{oncodasatinib2026} \\
Акалабрутиниб: капсулы Calquence & Омепразол 40 мг в течение 5 дней снижал AUC на 43\%. & Избегать ИПП; при H$_2$-блокаторе капсулу принимают за 2 часа до него, антацид разносят не менее чем на 2 часа.\src{oncocalcaps2026} \\
Акалабрутиниб: таблетки Calquence & При совместном приеме рабепразола клинически значимого изменения фармакокинетики не выявлено. & Ограничения для капсул нельзя автоматически переносить на таблетки. Форму и доступность проверяют; выбор или замена согласуются с онкогематологом.\src{oncocaltablet2026} \\
\bottomrule
\end{longtable}
Для отдельного продукта дазатиниба PHYRAGO инструкция сообщает отсутствие клинически значимого изменения фармакокинетики при омепразоле или фамотидине. Это дополнительный пример зависимости взаимодействия от лекарственной формы, а не разрешение сочетать любой дазатиниб с ИПП. Приведенные изменения AUC являются фармакокинетическими результатами; их нельзя численно приравнивать к снижению вероятности ремиссии или выживаемости.\src{oncophyrago2024}

\textbf{Совместное решение онкогематолога, эндокринолога и клинического фармаколога.} Сначала фиксируют конкретное показание к гастропротекции и точные наименование, форму и режим противоопухолевого препарата. Если показание отсутствует, исключают необоснованный ИПП. Если кислотосупрессия необходима, а инструкция запрещает сочетание, выбирают совместимую стратегию: допустимую альтернативу для желудочных симптомов либо согласованную коррекцию противоопухолевой формы/лечения. Антацид не следует считать равнозначной заменой ИПП для заживления язвы или профилактики высокорискового кровотечения. Нельзя самостоятельно повышать дозу ИТК в компенсацию сниженного всасывания или считать внутривенный ИПП способом обойти взаимодействие, обусловленное повышением желудочного pH. Эти решения требуют согласования целей обоих направлений лечения.\src{oncobosulif2026,oncodasatinib2026,oncocalcaps2026,oncocaltablet2026,oncophyrago2024,oncoppiposition2016}

До назначения и при пересмотре лечения сверяют весь список лекарств, включая безрецептурные антациды, функцию почек и печени, гемограмму и электролиты по риску. При высокодозном метотрексате учитывают отдельное предупреждение инструкции лансопразола о возможном замедлении выведения метотрексата; это не то же самое, что pH-зависимое уменьшение всасывания ИТК. При подозрении на лекарственный ОИН используют алгоритм раздела~\ref{sec:nephrology}. Таким образом, клиническая ценность лансопразола в этой нише определяется обоснованной гастропротекцией при совместимости с основным лечением, а не предполагаемой универсальной безопасностью молекулы.\src{uslabel,oncobosulif2026,oncocaltablet2026}

\chapter{Фармакогенетика, безопасность и персонализация антисекреторной терапии}
\label{chap:personalization}

\section*{Введение к главе}
\addcontentsline{toc}{section}{Введение к главе}
Персонализация антисекреторной терапии требует разграничивать физико-химические свойства молекулы и факторы, определяющие экспозицию у пациента. Липофильность и связанные с $pK_a$ равновесия ионизации описывают предпосылки распределения, накопления и кислотной активации; генотип CYP2C19, функция печени и взаимодействия изменяют доступность исходного вещества для этих процессов. Они не означают изменения аминокислотной последовательности желудочной помпы и не превращают печеночный метаболизм в обязательную стадию антисекреторной биоактивации. Для клинициста это различие важно при недостаточном ответе: до генотипической коррекции дозы необходимо проверить диагноз, приверженность, время приема и сопутствующую терапию.\src{molchem,molp450,cpic,uslabel,agagerd}

Ковалентная модификация мишени ограничивает возможность прямого переноса плазменной AUC на продолжительность эффекта и вероятность нежелательной реакции. Более высокая экспозиция или выраженный гастриновый ответ не равнозначны доказанной частоте атрофии, опухоли либо иного осложнения; липофильность также не служит самостоятельным показателем органной токсичности. Поэтому оценка безопасности должна разделять последствия кислотосупрессии, конкретные лекарственные реакции и наблюдательные ассоциации. В этой главе фармакогенетические результаты сопоставляются с клиническими ограничениями их переноса, а неодинаковая доказательная база разных ИПП не подменяется универсальным правилом пересчета доз по одному метаболическому признаку.\src{molcys,molrestore,leonova2015,cpic,uslabel,kidney}

При длительном лечении и отмене необходимо учитывать, что исчезновение свободного препарата и восстановление функциональной секреции происходят не как один процесс. Дисульфидный механизм ингибирования, возможная реактивация части помп и изменение доступного пула мишеней формируют биохимический контекст, но сами по себе не определяют индивидуальную схему снижения дозы. Практическое решение основывается на сохраняющемся показании, риске осложнений, особенностях пациента и оценке симптомов после прекращения приема. Последовательность главы поэтому включает проверку противопоказаний и взаимодействий, наблюдение, пересмотр дозы и различение временных проявлений отмены с рецидивом исходного заболевания.\src{molrestore,withdrawbell,withdrawniklasson,agadepres,uslabel}

\section{CYP2C19 и персонализация лечения}
\label{sec:pharmacogenetics}

\textbf{Согласование фенотипов статьи 2015 года и CPIC.} Леонова называет носителей CYP2C19 *1/*1 «экстенсивными, или быстрыми» метаболизаторами. В номенклатуре CPIC этот диплотип соответствует \textbf{нормальному} метаболизатору; быстрый фенотип соответствует *1/*17, ультрабыстрый --- *17/*17. Следовательно, слово «быстрый» в историческом исследовании нельзя автоматически переводить как современный rapid metabolizer. Примеры промежуточного фенотипа --- *1/*2 и *1/*3, медленного --- *2/*2, *2/*3 и *3/*3. Определение фенотипа требует сочетания аллелей, а не только наличия одного варианта. В таблицах статьи встречается обозначение CYP2C9, но обсуждаемые здесь данные относятся к CYP2C19.\src{leonova2015,cpic}

CPIC рассматривает связь CYP2C19 с экспозицией и эффектом лансопразола. Быстрый метаболизм может снижать кислотосупрессию, медленный --- повышать экспозицию. Эти рекомендации предназначены прежде всего для использования \textbf{уже известного генотипа}; они не устанавливают необходимость генотипирования каждого пациента с изжогой.\src{cpic}

\begin{longtable}{P{0.27\textwidth}P{0.66\textwidth}}
\toprule
Фенотип & Подход CPIC для лансопразола, омепразола и пантопразола \\
\midrule
\endhead
Ультрабыстрый (*17/*17) & Увеличение начальной суточной дозы на 100\%; допустимо разделение приема и необходим контроль результата. Категория CPIC: optional. \\
Быстрый (*1/*17) & Стандартная начальная доза; при \Hp\ и эрозивном эзофагите возможно увеличение на 50--100\% с контролем результата. \\
Нормальный (*1/*1) & Стандартная начальная доза; при \Hp\ и эрозивном эзофагите также можно рассмотреть увеличение на 50--100\%. \\
Промежуточный / медленный & Стандартная начальная доза; при хроническом лечении более 12 недель и достигнутом эффекте возможно снижение суточной дозы примерно на 50\% с наблюдением. \\
\bottomrule
\end{longtable}
Эти ориентиры не отменяют ограничений конкретной инструкции и коррекции при болезнях печени. Фенотип CYP2C19 не заменяет проверку приема, инфекции и альтернативного диагноза.\src{cpic}

Нельзя выбирать дозу только по этнической принадлежности: популяционные различия частот аллелей не определяют генотип конкретного человека. Если лечение не работает, разумнее сначала устранить более частые клинические причины, а фармакогенетику использовать как дополнительный инструмент.\src{cpic,agagerd}

\textbf{Уточнение ASGE 2025.} При недостаточном ответе на ИПП допускается определение полиморфизмов CYP2C19 с последующей коррекцией препарата или дозы; рекомендация условная, уверенность доказательств очень низкая. Это возможность для отобранных пациентов, а не обязательный тест до каждого назначения. В целом ASGE рекомендует минимальную необходимую интенсивность кислотосупрессии с обсуждением долгосрочной стратегии.\src{asge2025}

\textbf{Лансопразол, омепразол и рабепразол: зависимость экспозиции от CYP2C19.} Генотип влияет прежде всего на элиминацию исходного ИПП, а не на его химическую активацию у желудочной помпы. Лансопразол и омепразол существенно зависят от CYP2C19: снижение активности этого пути уменьшает клиренс и увеличивает AUC; сохраненная высокая активность, напротив, может ограничить экспозицию. Для лансопразола важны также CYP3A4 и неодинаковая скорость метаболизма R- и S-энантиомеров. У рабепразола значителен неферментативный переход в тиоэфир, поэтому изменение CYP2C19 затрагивает меньшую часть общей элиминации. Однако «меньшая зависимость» не означает отсутствия метаболизма через CYP2C19/CYP3A4 или невозможности лекарственных взаимодействий. Повышение AUC увеличивает доступность исходной молекулы для последующей кислотной активации; оно не является активацией ИПП печеночным ферментом.\src{leonova2015,molp450,molchem,cpic}

\textbf{Количественное сравнение: относительная экспозиция не равна относительной дозе.} В сводном тексте Леоновой отношения AUC медленных к экстенсивным метаболизаторам составляют около 4,3 для лансопразола, 6,3 для омепразола, 6 для пантопразола и 1,9 для рабепразола. Следовательно, этот материал описывает выраженную зависимость первых двух молекул от фенотипа и меньшую --- рабепразола, но не показывает большей генетической нестабильности лансопразола по сравнению с омепразолом. Это результаты разных работ, а не единое прямое сравнение: нельзя делить 6,3 на 4,3 для расчета относительной эффективности или умножать дозу лансопразола на 4,3. В частности, отношение PM/EM не измеряет дефицит AUC у носителя *17/*17. Стандартные категории NICE --- лансопразол 30 мг, омепразол 20 мг и рабепразол 20 мг/сут --- отражают дозовую классификацию, а не равенство концентраций или клиренса.\src{leonova2015,cpic,nice}

\textbf{Ультрабыстрый метаболизм и феноконверсия --- разные явления.} CYP2C19 *17 относится к аллелям повышенной функции: *17/*17 соответствует прогнозируемому ультрабыстрому фенотипу (UM), *1/*17 --- быстрому (RM), *1/*1 --- нормальному (NM). При UM исходная молекула лансопразола может быстрее удаляться до достижения достаточной доступности для кислотной активации. Это наследственная предпосылка низкой экспозиции, а не феноконверсия. Феноконверсия возникает, когда фактическая активность фермента перестает соответствовать прогнозу по ДНК из-за негенетического воздействия. У носителя *17/*17 сопутствующее ингибирование может функционально ослабить быстрый метаболизм, не изменив аллели; увеличение же дозы лишь меняет лекарственную нагрузку. Ни удвоение дозы, ни достижение клинического ответа не являются доказательством перехода UM в NM.\src{cpic,molfluvox}

\textbf{Механизм феноконверсии у прогнозируемого UM.} Последовательность включает появление ингибитора CYP2C19, уменьшение доступной ферментативной активности, снижение CYP2C19-зависимой части клиренса лансопразола и рост его экспозиции. Итог зависит от исходной активности, концентрации ингибитора, длительности сочетания и сохранности CYP3A4 и других путей; его нельзя определить только по наличию *17/*17. Возможен более медленный функциональный профиль, но автоматическое присвоение статуса PM носителю UM не обосновано. После прекращения ингибирования экспозиция может вновь снизиться по мере восстановления активности фермента, поэтому прежнее решение о дозе требует пересмотра. Индукция метаболизма действует в противоположном направлении и также способна нарушать соответствие между генетическим прогнозом и фактическим ответом. Эта схема объясняет взаимодействие, но не является валидированным коэффициентом коррекции для каждого ингибитора у UM.\src{cpic,molfluvox,uslabel}

\textbf{Количественный пример феноконверсии и граница его применимости.} В исследовании Miura и соавт. (2005) 18 здоровых участников распределялись по историческим группам homEM, hetEM и PM. После шести дней флувоксамина 25 мг дважды в сутки изучали разовый прием рацемического лансопразола 60 мг. У homEM AUC R- и S-энантиомеров увеличилась соответственно на 903\% и 1664\%, у hetEM --- на 462\% и 781\%. Это прирост относительно исходного значения, а не итоговая AUC в процентах. У homEM отношение $\mathrm{AUC}_R/\mathrm{AUC}_S$ уменьшилось с 12,7 до 6,4: S-энантиомер оказался чувствительнее к взаимодействию. Таким образом, генотипически сохраненный метаболизм не защищает от выраженного лекарственного повышения экспозиции. Однако homEM не означает специально исследованную группу *17/*17; эти числа нельзя переносить на UM, на обычные лечебные дозы или на иные ингибиторы. Работа не обосновывает назначение флувоксамина ради усиления кислотосупрессии.\src{molfluvox}

\textbf{Коррекция по CPIC: компенсировать низкую экспозицию, а не менять фенотип.} Для UM CPIC указывает увеличение \textbf{начальной суточной дозы на 100\%} с возможным разделением приема и контролем эффективности; для лансопразола, омепразола и пантопразола категория рекомендации --- \textit{optional}. Это генотипический ориентир уже на старте, а не правило, действующее только после неудачи стандартного курса. Для RM и NM предусмотрена стандартная начальная доза; увеличение на 50--100\% рассматривают при эрозивном эзофагите и эрадикации \Hp\ (категория \textit{moderate}). Ограниченность прямых данных по *17, особенно UM, не позволяет трактовать эти проценты как точное выравнивание AUC с NM. Смысл коррекции --- повысить вероятность достаточной кислотосупрессии, а не изменить генотип или гарантировать удвоение клинического эффекта. У IM/PM сохраняются условия возможного снижения дозы при эффективном длительном лечении, указанные в таблице.\src{cpic}

\textbf{Клинический смысл удвоения и ограничения арифметического переноса.} Если исходный лечебный режим при эрозивном эзофагите составляет 30 мг лансопразола в сутки, увеличение на 100\% означает 60 мг/сут; разделенный вариант --- 30 мг дважды в сутки. Это пример применения дозового ориентира, а не универсальное назначение при любой жалобе. До его использования проверяют диагноз, прием до еды, приверженность, печеночную функцию и сопутствующие ингибиторы либо индукторы. При эрадикации \Hp\ лансопразол 30 мг дважды в сутки уже может быть компонентом исходного протокола; нельзя автоматически удваивать этот режим повторно до 120 мг/сут только потому, что в генетической таблице указан процент. Необходимо определить, относительно какой исходной дозы применяется коррекция, и согласовать ее с конкретной схемой, инструкцией и результатом лечения. Генотип ИПП не устраняет устойчивость бактерии к антибиотикам.\src{cpic,uslabel,acggerd,acghp,maastricht}

При вероятной феноконверсии вследствие ингибитора механическое применение правила UM может привести уже не к компенсации низкой, а к избыточной экспозиции. Поэтому решение об увеличении дозы и решение о лекарственном взаимодействии принимают совместно; после изменения сопутствующей терапии показание к повышенной дозе оценивают заново. Самостоятельное назначение ингибитора CYP2C19, его отмена либо увеличение дозы ИПП по одному генетическому заключению не следуют из данного алгоритма.\src{molfluvox,cpic,uslabel}

\textbf{Что означает более слабая зависимость рабепразола.} По приведенному Леоновой исследованию 20 здоровых добровольцев на фоне рабепразола 20 мг/сут в течение восьми дней статистически значимых различий среднего внутрижелудочного pH между генотипическими группами не выявили, хотя AUC различалась. Это ограниченное наблюдение относительной устойчивости фармакодинамического эффекта, а не доказательство неизменной экспозиции или гарантированного ответа при *17/*17. В другом сопоставлении статьи, воспроизведенном ниже, средний pH на рабепразоле также численно менялся между фенотипами; результаты нельзя свести к абсолютной генотипической независимости.\src{leonova2015,cpic}

При недостаточной кислотосупрессии на правильно принимаемом лансопразоле и подтвержденном быстром метаболическом профиле можно обсуждать две разные стратегии: коррекцию дозы CYP2C19-зависимого ИПП либо переход на менее зависимый от этого пути рабепразол. Это клинический синтез, а не доказанное превосходство второй стратегии во всех ситуациях. CPIC не устанавливает для рабепразола и эзомепразола аналогичной формальной генотипической схемы коррекции из-за неоднородности доказательств (уровень C, отсутствие рекомендации), а не потому, что их доза никогда не требует изменения. Следовательно, рабепразол может уменьшать значимость вариабельности CYP2C19, но не освобождает от оценки заживления, эрадикации, взаимодействий и адекватности режима.\src{leonova2015,cpic,asge2025}

\textbf{Две стратегии персонализации: клинические ситуации, а не взаимозаменяемые назначения.} В следующей таблице разграничены дозовая компенсация для лансопразола и выбор менее CYP2C19-зависимой молекулы. Это авторское сопоставление приведенных данных; оно не является результатом прямого испытания двух стратегий.\src{leonova2015,cpic,asge2025}
\begin{longtable}{P{0.26\textwidth}P{0.34\textwidth}P{0.31\textwidth}}
\toprule
Ситуация & Лансопразол & Что можно заключить о рабепразоле \\
\midrule
\endhead
UM; начало лечения эрозивного эзофагита & Если исходный режим 30 мг/сут, пример удвоения --- 60 мг/сут, в том числе 30 мг дважды/сут; нужны проверка взаимодействий и контроль эффекта. & Менее зависимая от CYP2C19 альтернатива; превосходство замены над дозовой коррекцией лансопразола здесь не установлено. \\
RM или NM; эрозивный эзофагит либо эрадикация \Hp & Стандартная начальная доза; повышение на 50--100\% рассматривается в контексте конкретного протокола, а не любой жалобы. & Нельзя приписывать гарантированную эрадикацию или заживление одной лишь меньшей генотипической вариабельности. \\
UM по ДНК и сопутствующий ингибитор CYP2C19 & Сначала оценивают возможную феноконверсию; генотипическое удвоение без учета взаимодействия не обосновано. & Неферментативный путь не доказывает отсутствия всех взаимодействий; необходим пересмотр всей лекарственной схемы. \\
Уже достигнут объективный ответ на рабепразол & Переход на удвоенный лансопразол не следует из самого результата генотипирования. & Ответ поддерживает достаточность данного режима у данного пациента, но не универсальную стабильность при *17/*17. \\
\bottomrule
\end{longtable}

\textbf{Сравнение кислотосупрессии при конкретных режимах.} Таблица~5 статьи Леоновой содержит следующие средние значения 24-часового внутрижелудочного pH. Обозначение «экстенсивный» сохранено в историческом смысле, а не заменено современным быстрым фенотипом.\src{leonova2015}
\begin{longtable}{P{0.34\textwidth}P{0.18\textwidth}P{0.18\textwidth}P{0.18\textwidth}}
\toprule
Препарат и режим & Экстенсивные & Промежуточные & Медленные \\
\midrule
\endhead
Омепразол 20 мг дважды в сутки & 5,0 & 5,7 & 6,6 \\
Эзомепразол 20 мг дважды в сутки & 5,4 & 5,6 & 6,2 \\
Лансопразол 30 мг дважды в сутки & 4,7 & 5,4 & 6,4 \\
Рабепразол 10 мг дважды в сутки & 4,8 & 5,3 & 6,4 \\
\bottomrule
\end{longtable}
В этой таблице численно более высокий pH у экстенсивных метаболизаторов получен при эзомепразоле, чем при лансопразоле; это не равнозначно доказанному превосходству по заживлению или эрадикации. Дозы разных молекул неэквивалентны по массе, а средний pH не заменяет время удержания pH выше выбранного порога. Для рабепразола сама таблица также показывает различия между фенотипами: формулировку о меньшей зависимости нельзя превращать в утверждение об отсутствии любых различий.\src{leonova2015}

\textbf{Кратность приема и границы переноса исторических схем.} В статье сопоставлены рабепразол 40 мг один раз, 20 мг дважды и 10 мг четыре раза в сутки: общая суточная доза остается 40 мг, меняется ее распределение. Поэтому эффект дробления приема нельзя описывать как эффект увеличения суточной дозы. Авторская рекомендация 2015 года повышать дозу или кратность при недостаточной кислотосупрессии не заменяет дозовую таблицу CPIC и не устанавливает четырехкратный режим лансопразола для каждого носителя *1/*1. Сначала уточняют диагноз, приверженность, время приема, взаимодействия и, при эрадикации, антибиотикорезистентность; затем оценивают необходимость персонализации.\src{leonova2015,cpic,acghp}

\textbf{MDR1/P-гликопротеин: дополнительный, пока не дозоопределяющий фактор.} Обзор Леоновой рассматривает полиморфизм транспортера наряду с CYP2C19. В приведенном исследовании схемы с лансопразолом частота неудачи эрадикации составляла 18\% при генотипе C/C и 33\% при T/T. Эти наблюдения относятся к комбинированному лечению и не устанавливают самостоятельной причинности изменения транспорта лансопразола. Сама статья указывает на необходимость дальнейших исследований; она не содержит валидированного алгоритма коррекции дозы лансопразола по MDR1. Поэтому данные не служат основанием для обязательного тестирования MDR1 или для назначения определенной дозы по одному этому генотипу.\src{leonova2015}

\section{Противопоказания, осторожность и передозировка}

\textbf{Абсолютные ограничения.} По американской инструкции противопоказаны тяжелая гиперчувствительность к компонентам препарата и сочетание с рилпивирин-содержащими средствами. В других странах формулировки и дополнительные ограничения могут отличаться. Реакция на один ИПП требует осторожности при выборе другого, а не самостоятельного перебора препаратов.\src{uslabel}

\textbf{Ситуации для отдельного решения.} Тяжелая болезнь печени, предшествующая гипомагниемия, сложная лекарственная комбинация, беременность, лактация и детский возраст требуют проверки показания и формы. Не все эти состояния являются абсолютными противопоказаниями. Предполагаемая злокачественная опухоль не исключается ответом на ИПП и требует своевременной диагностики.\src{uslabel,emc}

\textbf{Передозировка.} Специфический антидот не описан; лечение симптоматическое, гемодиализ существенно не удаляет препарат. При случайном приеме избыточной дозы, особенно ребенком или одновременно с другими средствами, следует обратиться за медицинской / токсикологической помощью, сообщив препарат, дозу и время приема. Не следует самостоятельно вызывать рвоту или считать описанные в исследованиях высокие дозы доказательством безопасности любого количества.\src{emc,pharmalabel}

\section{Безопасность: что установлено, а что остается неопределенным}
\label{sec:safety}

\subsection{Обычные нежелательные реакции}
В инструкциях наиболее частыми называются диарея, боль в животе, тошнота и запор; возможны также головная боль и другие реакции. Частота зависит от исследования, возраста и комбинации. При эрадикации диарея или нарушение вкуса могут быть связаны с антибиотиками, поэтому нежелательное явление нельзя автоматически приписывать лансопразолу.\src{uslabel,emc}

\subsection{Редкие, но клинически значимые риски}
\begin{longtable}{P{0.30\textwidth}P{0.63\textwidth}}
\toprule
Реакция / предупреждение & Значение для наблюдения \\
\midrule
\endhead
Острый тубулоинтерстициальный нефрит & Возможен при приеме ИПП. При подозрении препарат отменяют и оценивают функцию почек. \\
Гипомагниемия & Иногда тяжелая, особенно при длительном лечении; возможны сопутствующие нарушения кальция/калия, судороги и аритмии. \\
Тяжелая гиперчувствительность, кожные реакции & Отек лица, нарушение дыхания, пузырная сыпь или поражение слизистых требуют срочной помощи и прекращения подозреваемого препарата. \\
Кожная / системная лекарственная волчанка & Новая характерная сыпь и системные проявления требуют оценки и пересмотра лечения. \\
Инфекционная диарея & Стойкая водянистая диарея, особенно после антибиотиков, требует исключения \textit{C. difficile} и других причин. \\
Дефицит B$_{12}$ & Риск возможен при многолетней кислотосупрессии; обследование определяется симптомами и факторами риска. \\
Полипы фундальных желез & Чаще описываются при длительном приеме; эндоскопическая оценка зависит от количества, размеров и морфологии. \\
\bottomrule
\end{longtable}
Это предупреждения инструкций, а не утверждение, что каждый пользователь столкнется с ними. Подозрение на конкретную тяжелую реакцию отличается от общего страха перед статистическими ассоциациями.\src{uslabel,hpra}

\subsection{Гипомагниемия: транспорт через эпителий и уязвимость пациента}
Магний пересекает кишечный эпителий как через клетки, так и между ними. Поэтому нарушение всасывания нельзя свести к единственному «магниевому насосу», который якобы ковалентно блокируется лансопразолом подобно ATP4A. В клеточной модели Caco-2 омепразол уменьшал перенос Mg$^{2+}$ и экспрессию клаудинов-7 и -12, участвующих в организации межклеточной проницаемости. Подкисление апикальной среды ослабляло эти изменения. Это аргумент в пользу зависимости транспорта от локальной среды и эпителиальных белков; препарат в данном опыте --- \textbf{омепразол}, не лансопразол.\src{safemgclaudin}

TRPM6 важен для магниевого гомеостаза, однако связь генетической вариабельности с осложнением не доказывает прямого связывания лекарства с каналом. В исследовании Hess и соавт. среди 133 длительно принимавших ИПП пациентов у 17 выявили гипомагниемию; варианты rs3750425 и rs2274924 в TRPM6 ассоциировались с повышенной уязвимостью. Доля 17/133 относится к выбранной исследовательской группе, а не к частоте осложнения у всех пользователей ИПП. Эти данные не являются валидированным алгоритмом подбора лансопразола по TRPM6.\src{safemggen}

Практически следует различать кишечное поступление, почечные потери, недостаточное питание и сочетания лекарств. Механистические гипотезы не отменяют зарегистрированного предупреждения о тяжелой гипомагниемии и возможных нарушениях кальция/калия. Вместе с тем они не обосновывают самостоятельное «подкисление кишечника» или профилактический прием больших доз магния: необходимость обследования и коррекции определяется клинической ситуацией.\src{uslabel,safemgclaudin}

\textbf{Клинические аргументы в пользу причинной связи.} В серии Mackay и Bladon описаны десять пациентов с гипомагниемией на фоне ИПП; нормализация после отмены и повторное снижение магния при возобновлении приема у части пациентов поддерживают лекарственную причинность сильнее, чем одномоментная ассоциация. Однако серия случаев не определяет популяционную частоту, сравнительный риск отдельных ИПП или неизбежность реакции. Сопутствующие диуретики и другие причины потерь магния необходимо учитывать отдельно. Клиническое подтверждение реакции не означает, что все звенья ее эпителиального механизма установлены.\src{safemgseries2010}

\subsection{Железо и витамин B12: разные пути возможной мальабсорбции}
\textbf{Средовой компонент.} Изменение желудочной кислотности может повлиять на доступность некоторых пищевых соединений и этапы высвобождения нутриентов. В инструкции отмечена возможность нарушения всасывания витамина B$_{12}$ при длительной гипо-/ахлоргидрии. Это не означает прямого разрушения витамина молекулой лансопразола; сам факт приема ИПП также не устанавливает причину анемии.\src{uslabel}

\textbf{Мальабсорбция не тождественна клиническому дефициту.} В экспериментальном исследовании десяти здоровых мужчин двухнедельный прием омепразола уменьшал всасывание связанного с пищевым белком витамина B$_{12}$. Это прямое подтверждение нарушения одного этапа усвоения, но не доказательство развития анемии или неврологического поражения после короткого курса. Истощение запасов зависит от исходного питания, продолжительности воздействия и иных причин дефицита. Исследование омепразола не дает количественной оценки риска для лансопразола; его результаты нельзя переносить на свободный витамин в препаратах как на идентичный пищевой субстрат.\src{safeb12trial1994}

\textbf{Регуляторный компонент обмена железа.} В работе Hamano и соавт. изучали цепь AhR--гепсидин--ферропортин. В клетках HepG2 омепразол повышал образование гепсидина; подавление AhR генетическим или фармакологическим методом уменьшало ответ. В мышиных опытах повышались печеночная экспрессия гепсидина и его уровень в крови, а количество ферропортина в двенадцатиперстной кишке и селезенке снижалось. Ферропортин обеспечивает выход железа из клетки, поэтому такая регуляция представляет отдельный от растворимости пищи механизм.\src{safeiron}

Авторы исследовали класс ИПП, но ключевые вмешательства для этой причинной цепи проводили с \textbf{омепразолом}. Наличие AhR-эффектов лансопразола в других моделях не позволяет автоматически объявить всю цепь доказанной у получающего его пациента. Железодефицит требует проверки кровопотери и других причин; приписывать его препарату только по биологически правдоподобной схеме нельзя.\src{safeiron,molahr}

\subsection{Микробиота: ослабление кислотного барьера не равно кишечной инфекции}
Кислая среда участвует в отборе микроорганизмов, проходящих из ротовой полости в кишечник. В рандомизированном исследовании Zhu и соавт. 49 здоровых добровольцев получали ИПП или H$_2$-блокатор семь дней. Метагеномный анализ слюны и кала выявил более выраженные изменения и перенос оральных видов в кишечник в группе ИПП. Это экспериментальное подтверждение эффекта на микробное сообщество, но не испытание профилактики или развития клинической инфекции и не отдельное доказательство эффекта лансопразола.\src{safemicrotrial}

В наблюдении Hojo и соавт. у 20 пациентов с рефлюкс-эзофагитом после восьми недель ИПП увеличивались количества отдельных Lactobacillus и Streptococcus. Однако лансопразол получали только двое, рабепразол --- двое, эзомепразол --- 16 человек; из этой выборки нельзя получить надежную отдельную оценку для лансопразола. Существенного увеличения частоты обнаружения бактерий в крови не было.\src{safemicrocohort}

Следует различать изменение относительной численности бактерии, ее устойчивую колонизацию, инфекцию и клинический исход. Присутствие вида, ассоциированного с заболеванием в других исследованиях, не означает, что препарат вызвал это заболевание. Эти работы сами по себе не обосновывают назначение антибиотика, пробиотика или «анализа на дисбактериоз» всем принимающим ИПП.\src{safemicrotrial,safemicrocohort}

\subsection{Микроскопический колит: важный сигнал именно для лансопразола}
В датском исследовании случай--контроль с 10\,652 гистологически подтвержденными случаями микроскопического колита наиболее сильная ассоциация наблюдалась при текущем приеме лансопразола. Скорректированное отношение шансов составляло 15,74 для коллагенового колита и 6,87 для лимфоцитарного; четкой зависимости от дозы не выявили. Это \textbf{отношения шансов в наблюдательном исследовании}, не абсолютная вероятность заболевания, не проценты и не доказанный коэффициент причинного риска у отдельного пациента.\src{safecolitis}

Европейское руководство UEG/EMCG 2021 года прямо разграничивает ассоциацию ИПП с микроскопическим колитом и доказанную причинность. Оно предлагает рассматривать отмену лекарства при подозреваемой временной связи его назначения с началом диареи; рекомендация слабая, качество доказательств очень низкое. Нормальная или почти нормальная эндоскопическая картина не исключает микроскопический колит: диагноз опирается на биопсии и гистологию.\src{safecolitisguide}

Для лансопразола нет установленной универсальной молекулярной цепи, которая позволяла бы объяснить эту ассоциацию прямой блокадой одного белка кишечного эпителия. Возможные изменения барьера и микробиоты нельзя выдавать за доказанный механизм всех случаев. Практический вывод --- не считать стойкую водянистую диарею автоматически безобидной реакцией и обсуждать с врачом причинную роль препарата и обследование, а не проверять ее самостоятельным повторным приемом.\src{safecolitisguide,safecolitis}

\subsection{Интерстициальный нефрит: не продолжение блокады желудочной помпы}
Механистический разбор ниже дополняется сравнительной эпидемиологией, лабораторными критериями и алгоритмом действий в разделе~\ref{sec:nephrology}; повторное изложение численных оценок здесь не приводится.

Отсутствие необходимости стандартной почечной коррекции дозы не означает отсутствия лекарственного повреждения почек. Для лансопразола описан биопсийно подтвержденный острый интерстициальный нефрит: в сообщении Jose и соавт. симптомы возникли после 45 дней приема, а после отмены и лечения наблюдалось восстановление. Это описание случая, которое не оценивает частоту реакции и не устанавливает конкретный антиген или сигнальный путь.\src{safeain}

Интерстициальное воспаление следует отличать от простого изменения транспорта H$^+$ или Mg$^{2+}$. Его рассматривают в контексте возможной лекарственной гиперчувствительности; ковалентное связывание с Cys813 желудочной ATPазы не является доказанным объяснением почечного поражения. При подозрении на такую реакцию ориентируются на клиническую оценку и отмену подозреваемого средства, а не на доказательство повышенной плазменной концентрации. Самостоятельная повторная проба приема недопустима.\src{safeain,uslabel}

\subsection{Сводное разграничение механизмов безопасности}
\begin{longtable}{P{0.22\textwidth}P{0.38\textwidth}P{0.31\textwidth}}
\toprule
Проблема & Что поддерживают источники & Чего они не устанавливают \\
\midrule
\endhead
Магний & Эпителиальные и генетические механизмы уязвимости, предупреждение инструкции.\src{safemgclaudin,safemggen,uslabel} & Прямую специфическую блокаду TRPM6 лансопразолом. \\
Железо & Доклиническую цепь AhR--гепсидин--ферропортин, преимущественно на омепразоле.\src{safeiron} & Причину анемии у каждого пользователя лансопразола. \\
Витамин B$_{12}$ & Экспериментальное уменьшение всасывания пищевого белково-связанного витамина при омепразоле.\src{safeb12trial1994} & Неизбежный клинический дефицит или деменцию при лансопразоле. \\
Гастрин / ECL-клетки & Адаптивный гормональный и трофический ответ на кислотосупрессию.\src{uslabel,vision2025} & Тождество гиперплазии и нейроэндокринной опухоли. \\
Кишечные инфекции & Рандомизированный сигнал для совокупности кишечных инфекций.\src{compass} & Столь же убедительное доказательство именно инфекции \textit{C. difficile}. \\

Микробиота & Изменения сообщества после кислотосупрессии.\src{safemicrotrial,safemicrocohort} & Неизбежное развитие инфекции или необходимость профилактического лечения. \\
Колит & Выраженную наблюдательную ассоциацию с лансопразолом.\src{safecolitis} & Абсолютную вероятность и единый доказанный молекулярный механизм. \\
Нефрит & Клинически значимую нежелательную реакцию и описания подтвержденных случаев.\src{safeain,uslabel} & Обязательную связь с накоплением лекарства или известным участком ATP4A. \\
\bottomrule
\end{longtable}

\subsection{Длительный прием: причинность и наблюдательные ассоциации}
\label{sec:longterm-causality}
\textbf{Методологическая рамка.} Смешение по показанию (\textit{confounding by indication}) возникает, когда причина назначения связана с изучаемым исходом; это не синоним любой ошибки отбора. Например, гастропротекция при антитромботической терапии выделяет популяцию с исходным сердечно-сосудистым риском. Обратная причинность возникает, если начальные проявления еще не распознанного заболевания приводят к назначению. Коррекция по измеренным ковариатам не устраняет автоматически остаточное смешение. Следует проверять компаратор, временную последовательность, определение исхода, абсолютную частоту и доверительный интервал.\src{acggerd}

\textbf{Не бинарное деление на «истинные» и «ложные» осложнения.} Физиологический эффект, признанная нежелательная реакция, правдоподобный механизм и наблюдательная ассоциация имеют разную доказательную силу. Отсутствие доказанной причинности не равнозначно доказательству отсутствия вреда; наличие механизма не определяет частоту клинического осложнения.
\begin{longtable}{P{0.22\textwidth}P{0.35\textwidth}P{0.34\textwidth}}
\toprule
Исход & Что обосновано & Ограничение причинного вывода \\
\midrule
\endhead
Гипергастринемия и ECL-клетки & Подавление кислоты вызывает компенсаторный гастриновый ответ; возможны трофические изменения ECL-клеток. & Повышенный гастрин не тождественен гиперплазии у каждого пациента, а гиперплазия не тождественна нейроэндокринной опухоли. Морфологические данные VISION рассмотрены отдельно.\src{uslabel,vision2025} \\
Гипомагниемия & Признанная редкая реакция при длительном приеме ИПП; тяжелое нарушение может требовать отмены наряду с восполнением магния. & Это не обязательный результат глубокой кислотосупрессии. Другие причины потерь магния и сопутствующие препараты оценивают отдельно.\src{uslabel} \\
Витамин B$_{12}$ & Гипо-/ахлоргидрия может ухудшать освобождение пищевого, связанного с белком B$_{12}$ и способствовать дефициту при длительной экспозиции. & Механизм мальабсорбции не доказывает клинический дефицит у каждого пользователя; важны питание, исходные запасы и другие причины.\src{uslabel} \\
\textit{Clostridioides difficile} & Биологически правдоподобный риск и регистрационное предупреждение; рандомизированные данные поддерживают сигнал кишечных инфекций в целом. & Нельзя приравнивать общий инфекционный исход к доказанному причинному эффекту именно на \textit{C. difficile}; учитывают антибиотики и госпитализации.\src{uslabel,compass} \\
Деменция, инфаркт миокарда, переломы & Наблюдательные ассоциации требуют проверки смешения; убедительный причинный эффект не установлен. & Нельзя объявить все ассоциации заведомо ложными. Для переломов сохраняется предупреждение инструкции; индивидуальный костный риск оценивают независимо.\src{acggerd,compass,uslabel} \\
\bottomrule
\end{longtable}

В литературе описаны связи ИПП с переломами, хронической болезнью почек, инфекциями, деменцией, сердечно-сосудистыми событиями и опухолями. Пациенты, получающие ИПП, нередко старше и имеют больше сопутствующих болезней; возможны смешение по показанию, обратная причинность и остаточные различия между группами. Поэтому относительный риск из базы наблюдений не равен доказанному лекарственному эффекту.\src{acggerd,compass}

\textbf{Рандомизированная проверка: величина эффекта и точность.} В COMPASS 17\,598 участников получали пантопразол 40 мг/сут или плацебо; медиана наблюдения составляла 3,01 года. Кишечные инфекции зарегистрированы у 1,4\% против 1,0\%: OR 1,33 (95\% ДИ 1,01--1,75), абсолютная разница по округленным долям --- 0,4 процентного пункта. Для \textit{C. difficile} было лишь 13 событий, и различие не достигло статистической значимости. Эти исходы нельзя объединять в утверждение о доказанном увеличении именно \textit{C. difficile}-инфекции. Большинство других предполагаемых осложнений значимо не различалось; однако испытание не исключает малых, редких и более поздних эффектов. Это данные \textbf{пантопразола}, а не прямое многолетнее испытание лансопразола.\src{compass}

\textbf{Важное уточнение 2024 года.} Вторичный анализ COMPASS у 8991 участника с парными измерениями функции почек обнаружил дополнительное снижение расчетной СКФ на 0,27 мл/мин/1,73 м$^2$ в год на пантопразоле. Анализ был post hoc, охватил около половины исходной рандомизированной популяции и не установил убедительного увеличения всех клинических почечных исходов. Он не доказывает идентичный эффект лансопразола, но делает формулировку «почечный риск полностью опровергнут» некорректной.\src{kidney}

\subsection{Рак желудка, атрофия и гипергастринемия}
Для оценки риска важны сама инфекция \Hp, исходная атрофия, кишечная метаплазия и показания, по которым назначался ИПП. Нельзя по наблюдательной ассоциации заключить, что лансопразол неизбежно вызывает рак, или что успешное подавление кислоты предотвращает его. Эрадикация инфекции и надлежащее наблюдение за предраковыми изменениями являются отдельными задачами.\src{maastricht,rfgastrit,jhp}

Повышение гастрина и хромогранина A на фоне кислотосупрессии может мешать диагностике. Это не самостоятельное доказательство нейроэндокринной опухоли. Если исследование необходимо, сроки отмены и повторного анализа согласуют с врачом и лабораторией.\src{uslabel}

\textbf{Фармакогенетический контекст гипергастринемии.} В таблице~9 обзора Леоновой при лансопразоле 30 мг в течение пяти дней отношение AUC гастрина на фоне лечения к исходному уровню составляло 1,6 у экстенсивных, 2,6 у промежуточных и 3,1 у медленных метаболизаторов. Это показатель гормонального ответа, а не частоты атрофии, опухоли или другого клинического осложнения. Краткосрочное увеличение экспозиции и гастринового ответа не позволяет рассчитать индивидуальный онкологический риск. В самой статье экспериментальные опухолевые наблюдения отделены от отсутствия клинического подтверждения; оснований превращать генотип CYP2C19 в самостоятельный маркер рака желудка эти сведения не дают.\src{leonova2015}

\subsection{Пятилетние данные VISION: биопсия важнее одного уровня гастрина}
В японском открытом рандомизированном исследовании VISION 202 из 208 пациентов достигли заживления эрозивного эзофагита и перешли на поддерживающее лечение. Поддерживающие режимы включали вонопразан 10 мг/сут или лансопразол 15 мг/сут; предусматривалось наблюдение до 260 недель с оценкой биопсий желудка. Число участников, достигших заживления, не означает, что все они завершили пятилетнее наблюдение.\src{vision2025}

В опубликованных результатах на 260-й неделе медиана гастрина составила 625 пг/мл при вонопразане и 200 пг/мл при лансопразоле; ECL-гиперплазия обнаруживалась у 4,9\% и 7,7\% соответственно. Гиперплазия париетальных клеток была чаще при вонопразане (97,1\% против 86,5\%). Эти признаки не равнозначны злокачественности: злокачественных эпителиальных изменений и желудочных нейроэндокринных опухолей в исследовании не наблюдали; зарегистрировано по одной аденоме в каждой группе.\src{vision2025}

Это прямые многолетние данные с группой лансопразола, но небольшая открытая выборка, выбывания и отсутствие группы без кислотосупрессии ограничивают выводы. Отсутствие редкой опухоли за пять лет не доказывает нулевого пожизненного риска. Результаты также не означают, что серийные биопсии нужны каждому принимающему ИПП вне исследовательского протокола.\src{vision2025,acggerd}

\section{Лекарственные взаимодействия}
\label{sec:interactions}

\subsection{Наиболее важные сочетания}
В таблице приведены не все взаимодействия, а те, которые особенно важно проверять перед назначением. Разнесение по времени помогает не всегда: повышение pH и изменение активности ферментов сохраняются дольше нескольких часов.\src{uslabel,sps}

\begin{longtable}{P{0.29\textwidth}P{0.33\textwidth}P{0.29\textwidth}}
\toprule
Препарат / группа & Проблема & Практический подход \\
\midrule
\endhead
Рилпивирин внутрь & Снижение всасывания и эффективности & Совместный прием противопоказан по инструкции США \\
Атазанавир, некоторые другие противовирусные средства & pH-зависимое изменение экспозиции & Проверить инструкцию всей схемы \\
Высокодозный метотрексат & Замедление выведения, токсичность & Возможна временная отмена ИПП по протоколу \\
Такролимус & Возможно увеличение концентрации & Контроль минимальной концентрации и функции почек \\
Варфарин & Возможное увеличение МНО & Контроль при начале/отмене и изменении схемы \\
Дигоксин & Возможное повышение экспозиции; роль магния & Оценка концентрации и электролитов по показаниям \\
Кетоконазол, итраконазол, некоторые противоопухолевые препараты & Снижение pH-зависимого всасывания & Проверить допустимость сочетания и форму \\
Рифампицин, зверобой & Снижение экспозиции ИПП & Избегать неподтвержденных сочетаний \\
Сукральфат & Может уменьшать всасывание лансопразола & Лансопразол раньше, интервал по инструкции \\
\bottomrule
\end{longtable}
Основа таблицы --- разделы взаимодействий американской инструкции. Для противовирусных и противоопухолевых средств ограничения зависят от конкретной молекулы, формы и дозы.\src{uslabel}

\subsubsection{Ревматология: метотрексат и границы предупреждения}
Для высокодозного метотрексата в обзоре указаны возможность замедления выведения и токсичности при сочетании с ИПП; допустимость временной отмены кислотосупрессии определяется соответствующим протоколом. Это лекарственное взаимодействие следует отличать от показания к гастропротекции при сопутствующем приеме НПВП: наличие желудочно-кишечного риска не делает сочетание автоматически безопасным, а наличие потенциального взаимодействия не задает универсальной схемы отмены.\src{uslabel,nice}

\textbf{Ограничение переноса на другие режимы.} Здесь не приведены порог дозы метотрексата, количественный риск при низкодозных режимах, срок прекращения лансопразола перед введением или срок его возобновления. Поэтому из предупреждения для высокодозной терапии нельзя выводить ни обязательную отмену ИПП у каждого ревматологического пациента, ни доказанную безопасность любого другого режима. Такие решения требуют проверки конкретной схемы и протокола; снижение дозы лансопразола с 30 до 15 мг в этом обзоре не представлено как доказанный способ устранения взаимодействия.\src{uslabel}

\subsubsection{Иммуносупрессоры и антитромботические средства: разные параметры контроля}
При возможном увеличении концентрации такролимуса рассматривают контроль минимальной концентрации и функции почек; при сочетании с варфарином --- контроль МНО при начале, отмене или изменении схемы. Для дигоксина отдельно учитывают возможное увеличение экспозиции и электролитные нарушения, включая магний; оценка концентрации и электролитов проводится по показаниям. Это разные задачи наблюдения, которые нельзя заменить одним общим заключением «лекарства совместимы».\src{uslabel}

Отсутствие необходимости стандартной почечной коррекции дозы лансопразола не отменяет мониторирования другого препарата и не исключает лекарственного интерстициального нефрита. Аналогично, нормальная переносимость ИПП не доказывает неизменности МНО или концентрации такролимуса. Приведенное изложение указывает предмет контроля, но не устанавливает универсальные интервалы анализов, целевые концентрации или алгоритмы изменения доз сопутствующей терапии.\src{uslabel,hpra,safeain}

\subsection{Почему эрадикационная комбинация требует отдельной проверки}
Кларитромицин, метронидазол и другие компоненты добавляют собственные ограничения и взаимодействия. Проверка только строки «лансопразол» не оценивает безопасность всего курса, особенно при приеме антиаритмических средств, антикоагулянтов, статинов и иммуносупрессоров. Необходимо сверить инструкции всех компонентов, а не считать фиксированный набор автоматически безопасным.\src{uslabel,jhp}

\section{Особые группы пациентов}

\subsection{Печеночная и почечная недостаточность}
\textbf{Печень.} В инструкции США при тяжелом нарушении, Child--Pugh C, указана доза 15 мг/сут. Европейские SmPC рекомендуют при умеренном или тяжелом поражении снижение суточной дозы на 50\% и наблюдение. Это реальное различие документов, а не два взаимозаменяемых правила; выбор зависит от локальной инструкции и показания.\src{uslabel,emc}

\textbf{Почки.} Стандартная коррекция дозы обычно не требуется, однако это не означает отсутствия риска интерстициального нефрита. При ХБП полезны оценка всей лекарственной нагрузки и индивидуальное наблюдение. При эрадикации ограничения антибиотиков по функции почек могут быть важнее фармакокинетики самого ИПП.\src{hpra,uslabel}

\subsection{Пожилые пациенты}
Возраст сам по себе не запрещает лечение; значение имеют язвенный анамнез, коморбидность и полипрагмазия. Европейская SmPC отмечает снижение клиренса и рекомендует не превышать 30 мг/сут без убедительного клинического основания. Одновременно именно у пожилого пациента может быть сильное показание к гастропротекции; необоснованная отмена не обязательно безопаснее продолжения.\src{emc,agadepres}

\subsection{Беременность}
Нельзя называть лансопразол «абсолютно безопасным» или переносить на современную практику старые буквенные категории FDA. По данным UKTIS/BUMPS, убедительной связи ИПП с основными неблагоприятными исходами беременности не установлено; омепразол чаще предпочитают из-за большего объема данных, другие ИПП возможны по врачебному решению.\src{uktis}

Регистрационные формулировки могут быть осторожнее: некоторые европейские инструкции не рекомендуют лансопразол при беременности из-за ограниченности данных. Поэтому учитывают тяжесть болезни, альтернативы и конкретную инструкцию. Данные о самом лансопразоле не означают допустимость комбинации антибиотиков для эрадикации во время беременности.\src{emc,uktis}

\subsection{Грудное вскармливание}
LactMed, обновление 15.02.2025, указывает на отсутствие опубликованных измерений содержания лансопразола в молоке и у ребенка и прямых данных об исходах у вскармливаемых детей. Предполагаемый низкий риск не равен доказанной безопасности. При необходимости лечения учитывают возраст и состояние ребенка; препараты с более изученной лактационной экспозицией могут быть предпочтительнее.\src{lactmed}

\subsection{Дети: регистрация и доказательства}
В США определенные формы зарегистрированы для ГЭРБ / эрозивного эзофагита с одного года. В возрасте 1--11 лет типовые дозы: при массе до 30 кг включительно --- 15 мг/сут, более 30 кг --- 30 мг/сут; курс до 12 недель. В 12--17 лет: 15 мг/сут при симптоматической неэрозивной ГЭРБ или 30 мг/сут при эрозивном эзофагите, до 8 недель.\src{uslabel}

Эти режимы приведены для сравнения регистрации, а не для самостоятельного лечения ребенка в РФ. Европейские инструкции отдельных препаратов могут не рекомендовать детское применение; российские ограничения следует проверять отдельно для конкретной формы.\src{emc,grls}

\textbf{Младенцы.} В рандомизированном исследовании 162 младенцев частота клинического ответа была одинаковой на лансопразоле и плацебо --- 54\%; серьезные нежелательные явления встречались чаще на лансопразоле. Следовательно, обычный плач и срыгивание не являются достаточным основанием для назначения ИПП. В маркировке США также присутствует предупреждение, основанное на доклинических данных, о риске утолщения клапанов сердца у пациентов младше года; это не оценка установленной частоты такого осложнения у детей.\src{infant,uslabel}

\section{Наблюдение при лечении}
\label{sec:monitoring}

\subsection{Что стоит оценить до назначения}
Ниже --- клинический контрольный список, синтезированный из рекомендаций и инструкций, а не обязательный одинаковый набор анализов для каждого пациента.\src{agagerd,agadepres,uslabel}
\begin{itemize}
\item Есть ли конкретное показание и ожидаемый срок лечения?
\item Нет ли тревожных признаков и показаний к эндоскопии?
\item Нужна ли диагностика \Hp\ до начала ИПП?
\item Какие лекарства, БАДы и безрецептурные средства уже принимаются?
\item Есть ли беременность, лактация, нарушение функции печени, ХБП, электролитные нарушения или тяжелая реакция на ИПП в прошлом?
\item Когда и по какому критерию будет оценен результат?
\end{itemize}

\subsection{Нужны ли всем магний, B12, креатинин и денситометрия?}
Нет оснований автоматически назначать всем пользователям ИПП регулярную денситометрию, расширенные витаминные панели и одинаковые лабораторные интервалы только из-за приема препарата. Подход зависит от исходного риска и клинических проявлений. При этом инструкция советует рассмотреть контроль магния до и во время длительного лечения, особенно при диуретиках или дигоксине.\src{acggerd,uslabel}

\begin{longtable}{P{0.34\textwidth}P{0.59\textwidth}}
\toprule
Ситуация & Что целесообразно обсудить \\
\midrule
\endhead
Диуретик, дигоксин, судороги, аритмия & Магний, калий, кальций; связь с лекарствами \\
Анемия, нейропатия, мальабсорбция & Общий анализ крови, B$_{12}$, железо и поиск причины \\
ХБП или подозрение на нефрит & Креатинин / СКФ, анализ мочи, оценка нефротоксичных лекарств \\
Перелом, высокий риск остеопороза & Стандартная оценка костного риска, а не только отмена ИПП \\
Длительная водянистая диарея & Инфекции, лекарственные причины, при необходимости иной диагностический поиск \\
Нет ответа на курс & Режим приема, эндоскопия / мониторирование рефлюкса по показаниям \\
\bottomrule
\end{longtable}
Это ориентиры индивидуальной оценки. Не вся анемия, диарея или почечная дисфункция у получающего ИПП пациента вызвана препаратом.\src{uslabel,agagerd,kidney}

\section{Снижение дозы и отмена}
\label{sec:withdrawal}

AGA рекомендует регулярно проверять показание; при его отсутствии рассматривать отмену, а при двукратном длительном приеме --- возможность перехода на однократный. Обычно не отменяют без специального основания терапию при осложненной ГЭРБ, пищеводе Баррета, контролируемом ИПП эозинофильном эзофагите или высоком риске желудочно-кишечного кровотечения.\src{agadepres}

После отмены возможно временное усиление кислотозависимых симптомов; одним из возможных объяснений является рикошетная гиперсекреция. Это физиологическое явление не равнозначно зависимости и не доказывает, что исходное заболевание требует пожизненного лечения.\src{agadepres}

\textbf{Единственной обязательной схемы постепенной отмены нет.} Допустимы как постепенное снижение, так и прекращение сразу у подходящих пациентов. Практический вариант --- сначала перейти с двукратного приема на однократный, затем обсудить 15 мг или прием по требованию при неосложненном заболевании; конкретные сроки индивидуальны. Альгинат или антацид могут использоваться для временных симптомов по согласованию.\src{agadepres,nice}

Не следует отменять эффективный ИПП только из-за неперсонализированного сообщения о возможном вреде. Но при подозрении на нефрит, анафилаксию или тяжелую кожную реакцию действует другая логика: прекращение подозреваемого препарата и медицинская оценка.\src{agadepres,uslabel}

\subsection{Патофизиология синдрома рикошета соляной кислоты при отмене ингибиторов протонной помпы}
\label{sec:rebound-mechanism}
\textbf{Определение феномена и нарушение отрицательной обратной связи.} Рикошетная гиперсекреция соляной кислоты (\textit{rebound acid hypersecretion}) означает превышение исходной, до лечения, кислотопродукции после прекращения антисекреторной терапии. Возврат секреции к исходному уровню и возобновление изжоги являются другими исходами: симптом может отражать рецидив ГЭРБ или повышенную чувствительность пищевода и сам по себе не доказывает гиперсекрецию.\src{withdrawbell,acggerd} Начальным звеном компенсаторного каскада служит подавление желудочной H$^+$/K$^+$-АТФазы и повышение внутрижелудочного pH. Уменьшение кислотной стимуляции антральных D-клеток ослабляет соматостатиновое торможение G-клеток; вследствие этого усиливаются выделение и синтез гастрина. Эксперимент Brand и Stone с омепразол-индуцированной ахлоргидрией на крысах связывает гастриновый ответ с изменением соматостатинового контроля, но не устанавливает его количественную величину у получающего лансопразол человека. Вторичная гипергастринемия может поддерживаться в период глубокой гипохлоргидрии; ее длительность на фоне лечения не означает обязательного сохранения после восстановления кислотности.\src{withdrawbrand1988,uslabel}

\textbf{Гастрин-зависимая перестройка ECL-клеточного звена.} Гастрин оказывает секреторное и трофическое действие через CCK$_2$-рецепторы энтерохромаффиноподобных (ECL) клеток оксинтной слизистой. Функциональная активация включает образование гистамина при участии гистидиндекарбоксилазы и его паракринное выделение; длительная стимуляция способна сопровождаться увеличением ECL-клеточного пула. Однако усиление функции отдельной клетки и морфологическая гиперплазия не тождественны. В мышиных и органоидных моделях Sheng и соавт. гипергастринемия приводила к расширению ECL-популяции преимущественно через CCK2R-положительные клетки-предшественники с участием MEK/ERK, а не через обязательное деление всех зрелых ECL-клеток. Эти результаты обосновывают клеточный механизм адаптации, но не позволяют считать гиперплазию неизбежной у каждого пациента или выводить ее выраженность из одной концентрации гастрина либо хромогранина A. Повышенный гормональный сигнал, секреторная активность и число клеток требуют раздельной оценки.\src{withdrawsheng2020,molcck,vision2025}

\textbf{Гистаминовая стимуляция и восстановление доступной мишени.} Гистамин активирует базолатеральные H$_2$-рецепторы париетальных клеток, сопряженные с G$_s$-белком, аденилатциклазой, cAMP и протеинкиназой A. Последующие изменения цитоскелета и мембранного транспорта обеспечивают функциональное включение секреторного аппарата. При сохраняющейся ковалентной модификации помпы этот рецепторный сигнал сам по себе не устраняет ингибирование фермента: на фоне приема ИПП возможны одновременно низкая кислотопродукция и усиленный гастрин-ECL-гистаминовый стимул.\src{molh2signal,molezrin,molcys} После последней дозы прекращается поступление нового ингибитора, но исчезновение препарата из плазмы не означает мгновенного разрыва всех дисульфидных аддуктов. Доступность активных помп возрастает по мере синтеза и включения новых молекул, а также восстановления части ранее ингибированных ферментов. Если регуляторная адаптация сохраняется дольше, чем функциональная блокада мишени, усиленный гистаминовый сигнал действует уже на восстановленный секреторный аппарат и может обеспечить кислотопродукцию выше исходной. Это временное несоответствие составляет патофизиологическое основание рикошета. Секреция соляной кислоты связана с переносом H$^+$ и Cl$^-$; речь не идет о высвобождении заранее накопленного запаса кислоты.\src{molrestore,molcys,molslc26,withdrawwaldum1996}

\textbf{Границы переноса механизма в клинический прогноз.} Формулировка об обязательном массивном гистаминовом всплеске после отмены чрезмерно категорична: сроки восстановления помп, нормализации гастрина и изменения ECL-клеточного пула неодинаковы, а клинический исход зависит от исходного заболевания, состояния слизистой и длительности лечения. В небольшом плацебо-контролируемом исследовании Bell и соавт. после 14 дней лансопразола 30 мг/сут кислотопродукция возвращалась к исходному уровню за 2--4 дня без превышения исходных значений. Этот результат не исключает рикошет после более продолжительной экспозиции, но опровергает его обязательность при любом курсе.\src{withdrawbell} Физиологические наблюдения после омепразола и исследования гастринового ответа при эзомепразоле относятся к иным режимам и не задают универсальной частоты синдрома для лансопразола. Поэтому в клиническом наблюдении различают измеренную гиперсекрецию, временные симптомы отмены и рецидив ГЭРБ; последние не должны автоматически интерпретироваться как подтверждение ECL-гиперплазии или основание для пожизненной терапии.\src{withdrawwaldum1996,withdrawboyce2017,agadepres}

\textbf{Интегральная схема рикошета.} Молекулярную последовательность следует читать в направлении «подавление секреции --- адаптация --- снятие блокады», а не как немедленный токсический эффект отмены:
\begin{center}
\begin{tabular}{c}
Ингибирование H$^+$/K$^+$-АТФазы $\longrightarrow$ гипохлоргидрия \\
$\Downarrow$ \\
Ослабление D--G-клеточной обратной связи $\longrightarrow$ гастрин $\uparrow$ \\
$\Downarrow$ \\
CCK$_2$-зависимая активация ECL; при длительном стимуле --- рост пула \\
$\Downarrow$ \\
Гистаминовый сигнал $\uparrow$ при еще ингибированных помпах \\
$\Downarrow$ \\
Отмена ИПП и восстановление доступных помп \\
$\Downarrow$ \\
Возможное превышение исходной кислотопродукции
\end{tabular}
\end{center}
Это авторская схема синтеза экспериментальных и физиологических данных. Ее последняя ступень условна: скорость восстановления помп и затухания гастринового/ECL-ответа различается, поэтому исчезновение лекарства из плазмы не задает момент максимальной секреции. Она не устанавливает обязательного рикошета у каждого пациента и не означает, что гиперплазия сохраняется у всех после лечения.\src{withdrawbrand1988,withdrawsheng2020,molrestore,withdrawwaldum1996,withdrawboyce2017}

\textbf{Почему нельзя обещать «резкий всплеск сразу после отмены».} Временная последовательность зависит от длительности лечения, состояния слизистой и изучаемого исхода. Гиперплазия и стойкая гипергастринемия после прекращения ИПП не обязательны; восстановление секреции и исчезновение гормонального ответа не следует считать синхронными. В исследовании Boyce и соавт. после 28 дней эзомепразола гастрин и CgA возвращались к исходным значениям за 2--3 дня, без рикошетной диспепсии; однако биомаркеры и симптомы не заменяют прямого измерения кислотопродукции. Механистическая возможность не устанавливает универсальный срок или обязательность клинического рикошета.\src{withdrawboyce2017,withdrawbell,vision2025}

Таким образом, гипохлоргидрия может формировать адаптивный гастриновый и ECL-клеточный ответ, но этот ответ не доказывает неизбежной гиперсекреции после каждого курса. Предупреждение о возможных временных симптомах и индивидуальный план отмены обоснованнее, чем утверждение о гарантированной гормональной «зависимости».\src{molrestore,uslabel,agadepres}

\subsection{Прямые данные об отмене: что именно измеряли}
\begin{longtable}{P{0.26\textwidth}P{0.37\textwidth}P{0.28\textwidth}}
\toprule
Исследование & Измеренный результат & Граница переноса \\
\midrule
\endhead
Waldum и соавт., 1996; 9 пациентов с рефлюкс-эзофагитом, \textbf{омепразол} 40 мг/сут 90 дней & Через 14 дней после отмены пентагастрин-стимулированная кислотопродукция была на 50\% выше исходной; возросла и базальная секреция.\src{withdrawwaldum1996} & Прямое измерение секреторного ответа, но малая выборка, другой ИПП и оценка не сразу после последней дозы. \\
Bell и соавт., 2001; 16 здоровых добровольцев, лансопразол 30 мг/сут или плацебо 14 дней & После отмены внутрижелудочный pH, базальная и стимулированная секреция возвращались к исходному уровню за 2--4 дня без превышения исходной кислотопродукции.\src{withdrawbell} & Небольшой короткий опыт не исключает рикошета после иной длительности лечения или у другой популяции. \\
Niklasson и соавт., 2010; 48 здоровых \Hp-отрицательных добровольцев, \textbf{пантопразол} 40 мг/сут или плацебо 28 дней & В первую неделю после отмены диспепсия отмечена у 11/25 (44\%) и 2/23 (9\%) соответственно. Различия среднего симптомного балла на третьей неделе не были значимыми.\src{withdrawniklasson} & Это не исследование лансопразола; симптомный балл и корреляция с гастрином не тождественны прямому измерению кислотопродукции. \\
\bottomrule
\end{longtable}

В работе Bell терминальный период полувыведения лансопразола составлял 1,11 часа, тогда как расчетные времена полувозврата среднего pH и гастрина к исходному уровню --- 22 и 19 часов. Это параметры конкретного исследования, а не сроки обязательной отмены перед диагностическим тестом. Авторы связывали восстановление преимущественно с синтезом новой помпы; более поздние биохимические данные о восстановлении дисульфидов показывают, почему это объяснение не следует считать единственно возможным.\src{withdrawbell,molrestore}

Таким образом, результаты двух исследований не образуют простого противоречия: различаются препарат, длительность приема и конечная точка. Цифру 44\% нельзя представлять как частоту рикошета после лансопразола. И отсутствие превышения секреции после 14 дней нельзя использовать как доказательство отсутствия симптомов после любого многолетнего курса.\src{withdrawbell,withdrawniklasson}

\subsection*{Сравнение способов отмены: успех депрескриптинга не равен отсутствию рикошета}
В небольшом открытом рандомизированном исследовании Hendricks и соавт. 2021 года анализировали 33 участника после выбываний. Значимого различия между постепенной и одномоментной отменой по сохранению отказа от ИПП через 12 месяцев не выявлено; при постепенном снижении симптомы были менее выражены в отдельные сроки наблюдения. Эти результаты не устанавливают обязательного преимущества одной схемы и не измеряют непосредственно превышение исходной кислотопродукции.\src{withdrawhendricks2021}

В исследовании Hojo и соавт., опубликованном 27 апреля 2026 года, 63 клинически стабильных пациента с ГЭРБ рандомизированы на немедленную отмену или уменьшение дозы наполовину на две недели с последующей отменой; включались также получавшие P-CAB. Успех через 24 недели составил 74,2\% и 75,0\% соответственно; разность рисков $-0{,}008$ (95\% ДИ от $-0{,}231$ до $0{,}215$). Критерий допускал до восьми дней спасательного приема за каждое 28-дневное окно: это не полный отказ от любой кислотосупрессии. Широкий ДИ не доказывает эквивалентности стратегий; результаты не являются отдельной проверкой лансопразола 30/15 мг и не подтверждают предотвращение физиологического рикошета.\src{withdrawhojo2026}

\subsection{Юго-Восточная Азия: консенсус по рациональному применению ИПП 2026 года}
Региональный документ содержит 20 согласованных положений и алгоритм пересмотра терапии. Он ориентирован на пациентов, у которых больше нет достаточного показания к ИПП; само достижение определенной длительности лечения не является основанием отменять гастропротекцию или поддержание при осложненной ГЭРБ.\src{seadepres2026,agadepres}

Для подходящих пациентов после менее четырех недель приема алгоритм допускает прекращение сразу; после четырех недель и более --- постепенное снижение или прекращение сразу. Предлагается наблюдать за возвратом симптомов около восьми недель, используя при необходимости согласованную симптоматическую помощь. Сохранение симптомов восемь недель и более требует переоценки исходного заболевания, а не автоматического объяснения всего периода «рикошетом». Это региональная стратегия наблюдения, \textbf{не установленная биологическая продолжительность гиперсекреции у каждого пациента}.\src{seadepres2026}

\subsection{Как применять эти сведения без необоснованной отмены}
Решение состоит из двух отдельных вопросов: сохраняется ли показание и как облегчить переход при допустимой отмене. История LA C/D-эзофагита, язвенной стриктуры, пищевода Баррета или высокого риска кровотечения может существенно изменить первый ответ. При сохраняющемся показании длительное лечение не следует объявлять «зависимостью» только потому, что симптомы возвращаются без него.\src{agadepres,acggerd,barrett2025}

У подходящего пациента полезно заранее согласовать способ снижения, допустимую временную помощь и момент повторной оценки. Такой план отличается от бесконечного самостоятельного чередования отмены и возобновления. Дисфагия, кровотечение, анемия или необъяснимое похудение требуют отдельной диагностики, а не ожидания окончания предполагаемого рикошета.\src{agadepres,acggerd}

Пошаговая реализация снижения дозы с критериями отбора, возврата к лечению и наблюдения приведена в завершающем алгоритме главы (раздел~\ref{sec:deprescribing-algorithm}).

\section{Практические ответы на частые вопросы}

\textbf{Можно ли пить постоянно?} Да, если сохраняется показание; нет, если постоянный прием стал привычкой без пересмотра диагноза. Длительность определяется риском рецидива и осложнений, а не универсальным «максимумом дней».\src{agadepres}

\textbf{Почему 15 мг иногда лечит, хотя стандартная доза 30 мг?} Потому что дозовая категория NICE и зарегистрированная доза при конкретном показании отвечают на разные вопросы.\src{nice,uslabel}

\textbf{Можно ли принимать после еды?} Для обычного лансопразола это менее оптимально; предпочтителен прием до еды. Свободный режим декслансопразола относится к другой форме препарата.\src{acggerd,dexlabel}

\textbf{Можно ли заменить капсулу двумя таблетками другого ИПП?} Не по арифметике миллиграммов. Нужны конкретная молекула, показание и эквивалентный клинический режим.\src{nice}

\textbf{Если изжога исчезла, значит ли это, что \Hp\ вылечена?} Нет. Нужен контроль эрадикации с правильными интервалами после антибиотиков, висмута и ИПП.\src{acghighlights}

\textbf{Если препарат не помогает, обязательно ли нужен более сильный?} Нет. Причиной могут быть неправильный прием, некислый рефлюкс или другое заболевание.\src{agagerd,lyon}

\textbf{Можно ли одновременно два разных ИПП?} Рассмотренные алгоритмы предлагают оптимизацию одного препарата, изменение кратности или смену терапии, а не рутинное дублирование класса. Это авторский практический вывод из рекомендаций.\src{agagerd,dgvs}

\textbf{Нужны ли всем пробиотики, витамины и кальций?} Сам прием лансопразола не создает такого универсального показания; вмешательства определяются конкретными рисками и клиническими задачами.\src{acggerd,acghp}

\section{Когда нельзя ограничиваться обзором или самолечением}
\textbf{Срочно обратиться за медицинской помощью} следует при рвоте кровью или «кофейной гущей», черном дегтеобразном стуле с ухудшением самочувствия, обмороке, внезапной сильной боли, затруднении дыхания, отеке лица или тяжелой пузырной сыпи. Боль или сдавление в груди, особенно с одышкой, потливостью или распространением в руку/челюсть, нельзя автоматически считать изжогой.\src{otc,uslabel,esge2026}

Прогрессирующее нарушение глотания, боль при глотании, непреднамеренное похудение, анемия, повторная рвота и устойчивые симптомы требуют очной диагностики. Временное улучшение на лансопразоле не исключает опухоль или другое серьезное заболевание.\src{uslabel,dgvs}

\section{Нормативный протокол пошагового депрескриптинга лансопразола}
\label{sec:deprescribing-algorithm}
\textbf{Область применения и нормативный статус алгоритма.} Перед депрескриптингом документируют исходное показание, завершенность необходимого курса, текущую дозу и риск осложнений. Рассматриваемая схема предназначена для взрослого с контролируемой неэрозивной ГЭРБ, получающего лансопразол 30 мг один раз в сутки, при отсутствии установленной потребности в постоянной кислотосупрессии. По принципам AGA 2022 прекращение ИПП обычно не рассматривают при осложненной ГЭРБ, пищеводе Баррета, эозинофильном эзофагите либо высоком риске желудочно-кишечного кровотечения; у получающих препарат дважды в сутки первоначально оценивают возможность перехода на однократный прием. AGA допускает как постепенное снижение дозы, так и одномоментную отмену. Следовательно, приведенные ниже интервалы 14--28 и 14 дней представляют согласуемую с пациентом практическую реализацию, а не обязательный календарь консенсуса или доказанный способ предотвращения гиперсекреции.\src{agadepres,withdrawfarrell2017}

\textbf{Первый этап: снижение дозы с 30 до 15 мг на 14--28 дней.} При подтвержденной возможности уменьшения терапии назначенный ежедневный режим 30 мг заменяют режимом 15 мг один раз в сутки до еды на протяжении двух--четырех недель. Используют подходящую лекарственную форму 15 мг; уменьшение дозы не должно сопровождаться разрушением гастрорезистентных гранул. Выбор 15 мг означает уменьшение лекарственной нагрузки при сохраняющемся контроле симптомов, но не гарантирует достаточности этой дозы для любого заболевания.\src{uslabel,nice} До начала этапа и к его завершению фиксируют частоту дневных и ночных симптомов, нарушения сна, влияние жалоб на повседневную активность и использование симптоматических средств. Решение о переходе к следующему этапу принимают по клиническому контролю; истечение календарного срока не требует автоматического дальнейшего снижения. При выраженном ухудшении уточняют соблюдение режима и диагноз, а не объясняют любой симптом предполагаемым рикошетом.\src{withdrawfarrell2017,acggerd}

\textbf{Второй этап: 15 мг через день в течение последующих 14 дней.} При приемлемом контроле после первого этапа ежедневный прием заменяют календарным альтернирующим режимом: 15 мг через день на протяжении двух недель. Такой режим является промежуточным вариантом уменьшения экспозиции; он не тождественен приему по требованию и не предполагает одинакового подавления секреции в каждый день интервала. В письменном плане отмечают дни приема, выраженность жалоб в дни без ИПП и потребность в дополнительном облегчении.\src{withdrawfarrell2017,molrestore} Сохранение ковалентного эффекта после отдельной дозы объясняет возможность более длительного антисекреторного действия, чем присутствие препарата в плазме, но не доказывает, что двухнедельный прием через день нормализует гастрин или устраняет ECL-гиперплазию. Поэтому этот этап не используют вместо постоянной высокорисковой гастропротекции; при выявлении сохраняющегося показания возвращаются к обоснованному поддерживающему лечению.\src{molrestore,agadepres}

\textbf{Третий этап: прекращение регулярного приема и согласованная помощь по требованию.} При удовлетворительном состоянии после второго этапа отменяют постоянное ежедневное или альтернирующее назначение. Если симптомов нет, повторные дозы ИПП не требуются. При эпизодической изжоге возможно кратковременное применение невсасывающегося антацида либо альгинатного препарата по инструкции с учетом противопоказаний, состава средства и интервалов с другими лекарствами; обязательное профилактическое назначение этих средств всем пациентам не предполагается.\src{nice} Если при неэрозивной ГЭРБ сохраняется потребность именно в ИПП, режим \textit{on-demand} означает согласованное возобновление приема до контроля симптомов с последующим прекращением, а не ожидание немедленного действия одной дозы по аналогии с антацидом.\src{acggerd,withdrawfarrell2017} Частые рецидивы, ночные жалобы или постоянная потребность в симптоматической помощи требуют пересмотра диагноза и минимальной эффективной терапии. Дисфагия, анемия и непреднамеренное похудение требуют обследования; кровавая рвота, мелена с ухудшением состояния или обморок --- неотложной оценки. Результат фиксируют раздельно как полную отмену ИПП, прекращение постоянного приема с использованием по требованию либо сохранение минимальной эффективной ежедневной дозы.\src{acggerd,withdrawfarrell2017}

\textbf{Календарь для письменного плана.} День 1 --- первый день приема 15 мг вместо 30 мг. Два варианта ниже конкретизируют авторский диапазон длительности первого этапа; они не являются отдельными протоколами AGA. При неудовлетворительном контроле переход не выполняют автоматически.
\begin{longtable}{P{0.27\textwidth}P{0.30\textwidth}P{0.34\textwidth}}
\toprule
Этап & Если шаг 1 длится 14 дней & Если шаг 1 длится 28 дней \\
\midrule
\endhead
Шаг 1: 15 мг ежедневно & Дни 1--14 & Дни 1--28 \\
Шаг 2: 15 мг через день & Дни 15--28; 7 приемов & Дни 29--42; 7 приемов \\
Шаг 3: без планового ИПП & С дня 29 & С дня 43 \\
\bottomrule
\end{longtable}
На втором этапе чередуют день приема 15 мг и день без ИПП, а не дозы 30 и 15 мг. При иной согласованной длительности первого этапа даты последующих переходов сдвигают, сохраняя 14-дневный второй этап. Режим по требованию начинается только при необходимости после прекращения календарного назначения; при отсутствии симптомов автоматически начинать новый курс не следует.

\textbf{Краткосрочная помощь на шаге 3.} В письменном плане указывают конкретный антацид, допустимую дозу по его инструкции, ограничения совместимости и срок пересмотра потребности; альгинат может быть согласованной альтернативой. Средство используют эпизодически при симптомах, а не как обязательный постоянный компонент отмены. NICE не рекомендует длительный частый прием антацидов как замену ведению заболевания. Регулярная потребность в них, ночные симптомы или ухудшение повседневной активности требуют внепланового пересмотра, а не бесконечного продления «прикрытия». Облегчение изжоги антацидом не доказывает, что нормализовались гастрин, ECL-клеточный пул или собственная кислотопродукция.\src{nice,withdrawfarrell2017}

\textbf{Что означает «полная отмена».} На шаге 3 прекращается плановое ежедневное или альтернирующее назначение. Полная отмена без повторных доз ИПП и эпизодическое возобновление по требованию --- разные возможные результаты, которые отдельно фиксируют в документации. При указанном календаре переход к шагу 3 приходится примерно на 29--43-й день от начала снижения; эти сроки не являются прогнозом нормализации гастрина или морфологии ECL-клеток.

\textbf{Наблюдение и основания для изменения плана.} Контакты примерно через 4 и 12 недель от начала депрескриптинга используют для оценки жалоб, сна, повседневной активности и частоты приема по требованию. Возврат симптомов не доказывает рикошетную гиперсекрецию: следует отличать временный эффект отмены от рецидива ГЭРБ (раздел~\ref{sec:rebound-mechanism}). При повторяющихся выраженных симптомах обсуждают возврат к минимальной эффективной ежедневной дозе и необходимое обследование; стремиться к нулевой дозе любой ценой не следует.\src{withdrawfarrell2017,acggerd}

\textbf{Тревожные признаки: не ждать планового визита.} Дисфагия, непреднамеренное похудение, стойкая рвота или анемия требуют обследования. Кровавая рвота, мелена с ухудшением состояния, обморок либо новая боль в груди требуют срочной оценки и не должны автоматически объясняться отменой ИПП.\src{acggerd}

\textbf{Документирование результата.} Итог пересмотра фиксируют одной из трех формулировок: «ИПП отменен, повторные дозы не требуются»; «регулярный ИПП отменен, используется согласованный режим по требованию»; «сохраняется минимальная эффективная ежедневная доза по указанному показанию». В каждой записи указывают фактическую дозу, частоту симптоматической помощи и дату следующего контакта. Такой учет отличает уменьшение лекарственной нагрузки от полного прекращения лечения и позволяет оценить успех без ошибочного требования нулевой дозы для каждого пациента. Это организационная реализация индивидуального пересмотра, а не дополнительная рекомендация AGA по длительности курса.\src{agadepres,withdrawfarrell2017}

\textbf{Итог алгоритма.} Цель постепенного перехода --- уменьшить необоснованную лекарственную нагрузку при сохранении контроля заболевания и понятном плане помощи при временных симптомах. Отмена, прием по требованию и сохранение минимальной эффективной постоянной дозы являются разными допустимыми исходами индивидуального пересмотра; ни один из них не следует выбирать вопреки сохраняющемуся показанию.\src{agadepres,withdrawfarrell2017}

\chapter{Сквозное междисциплинарное заключение}
\label{chap:conclusion}

\section{Сквозные клинико-фармакологические выводы}
\label{sec:integrated-conclusions}
Семь следующих выводов объединяют материалы четырех основных глав; завершающий раздел~\ref{sec:research-agenda} переводит выявленные ограничения в программу будущих исследований. Они представляют авторский синтез изложенной доказательной базы, а не самостоятельное клиническое руководство; сравнительные преимущества относятся только к тем исходам и условиям, которые рассмотрены в соответствующих разделах.

\begin{enumerate}[label=\textbf{\arabic*.},leftmargin=*,itemsep=0.7\baselineskip]
\item \textbf{Клиническая ценность определяется соответствием мишени, показания и исхода.} Лансопразол остается рациональным вариантом лечения кислотозависимых заболеваний и компонентом обоснованной гастропротекции. При эрадикации \Hp\ он обеспечивает кислотосупрессию в составе комбинации, но не заменяет антибиотики. Заживление язвы, контроль рефлюксных симптомов, предотвращение кровотечения и эрадикация инфекции являются разными задачами. Принадлежность к ранее разработанным ИПП не означает клинической неполноценности, как и появление P-CAB не устанавливает их предпочтительность при каждом показании. Выбор лансопразола определяется диагнозом, режимом, взаимодействиями и локальной инструкцией; ответ на лечение не заменяет разграничения доказанной ГЭРБ и функциональных симптомов. Именно связь между установленным заболеванием и измеримым результатом, а не универсальное ранжирование препаратов по «силе», определяет место молекулы в практике.\src{uslabel,agagerd,nice,acghp,maastricht,seoul,lyon}

\item \textbf{Общий механизм класса сочетается с индивидуальным фармакологическим профилем.} Кислотная активация и ковалентная модификация желудочной H$^+$/K$^+$-АТФазы объединяют лансопразол с другими ИПП; длительность действия определяется не только плазменной экспозицией, но и доступностью, блокадой и восстановлением мишени. Особенности кислотно-основных равновесий, мембранного переноса, стереоселективного метаболизма и участия CYP2C19 формируют профиль самой молекулы. Однако индивидуальный профиль не означает уникальной клинической эффективности: липофильность не доказывает преимущества в гортани или бронхах, а отдельный клеточный противовоспалительный эффект не устанавливает превосходства по исходам. Механистические различия объясняют возможную вариабельность ответа, но должны сопоставляться с данными у пациентов. Поэтому фармакологическая характеристика служит основанием для постановки клинического вопроса, а не его заранее готовым ответом.\src{molchem,molcys,molrestore,molstereokin,cpic,toppits,asthma}

\item \textbf{Дозы 15 и 30 мг обеспечивают гибкость, но не допускают произвольного пересчета режимов.} Лансопразол 15 мг один раз в сутки имеет профилактические применения при НПВП-ассоциированном риске и изучен во вторичной профилактике язв на аспирине; 30 мг один раз в сутки является лечебным режимом для ряда язвенных и эрозивных поражений, а 30 мг дважды в сутки используется в определенных эрадикационных схемах. Продолжительность и кратность приема определяются конкретной задачей. Испытания профилактики не измеряют скорость заживления исходной язвы; численное различие результатов 15 и 30 мг не доказывает превосходства большей дозы. Дробление неизменной суточной дозы, ее удвоение и переход к поддерживающему режиму представляют разные вмешательства. Дозовая гибкость имеет клинический смысл при правильном приеме, соблюдении ограничений инструкции и повторной оценке результата.\src{uslabel,hpra,aspirin15,misotrial,nsaidagrawal2000,acghp}

\item \textbf{Междисциплинарная гастропротекция требует раздельной оценки эффективности и взаимодействий.} В ревматологии и неврологии кислотосупрессия не устраняет все механизмы НПВП-повреждения и не заменяет пересмотр анальгетической терапии; предупреждение повторного язвенного осложнения отличается от лечения уже существующего дефекта. В кардиологии выбор лансопразола вместо омепразола или эзомепразола при клопидогреле обосновывается меньшим потенциалом клинического взаимодействия, а не доказанным снижением MACE именно этой молекулой; преимущества перед пантопразолом по сердечно-сосудистым исходам не установлены. В ОРИТ профилактика стрессового кровотечения определяется риском и не тождественна лечению состоявшегося кровотечения. Таким образом, общий антисекреторный механизм позволяет использовать класс ИПП в нескольких специальностях, но не делает одинаковыми показания, дозы, пути введения и доказательства для всех его представителей. Пошаговая реализация кардиологического решения приведена в подразделе~\ref{sec:dapt-protocol}.\src{nice,aspirintrial,sccm,esge2026,cardioplavix2025,sps,cardiocogent2010}

\item \textbf{Внепищеводные и экспериментальные эффекты требуют самостоятельной клинической верификации.} Отрицательные результаты лансопразола при персистирующих глоточных симптомах и плохо контролируемой астме у исследованных детей препятствуют эмпирическому расширению применения, но не исключают отдельного показания к лечению подтвержденной ГЭРБ. При ЭоЭ рекомендация ИПП относится к классу, а оценка лечения включает симптомы, эндоскопию и биопсию; успешный гистологический ответ не следует переименовывать в полное восстановление слизистой. Наблюдения Keap1--Nrf2--HO-1, CCL26 или IL-8 этих критериев не заменяют. Еще более строгая граница относится к LPZS: восстановленное производное и бактериальная дыхательная мишень образуют доклиническое направление, которое нельзя превращать в противотуберкулезное назначение исходного лансопразола. Биологическая убедительность механизма повышает ценность исследовательской гипотезы, но сама по себе не расширяет показаний.\src{toppits,asthma,eoe,eoeguide,eoestat6,molnrf2,molho1,molil8,moltb,moltbstructure}

\item \textbf{Фармакогенетическая персонализация учитывает фактический метаболический контекст.} Зависимость экспозиции лансопразола и омепразола от CYP2C19 выражена сильнее, чем у рабепразола, однако меньшая зависимость последнего не гарантирует неизменной экспозиции или одинакового ответа у всех пациентов. Генетически прогнозируемый ультрабыстрый метаболизм необходимо отличать от феноконверсии под действием сопутствующего лекарства. Ориентир CPIC по увеличению начальной суточной дозы у UM на 100\% направлен на компенсацию недостаточной экспозиции, а не на изменение фенотипа; применение требует учета инструкции, взаимодействий и исходной схемы. При эрозивном эзофагите и эрадикации повышение дозы, изменение кратности и замена молекулы являются разными решениями с неодинаковыми основаниями. Генетическое заключение не отменяет проверки приверженности, диагноза и антибиотикорезистентности и не служит самостоятельным основанием для бесконтрольного наращивания кислотосупрессии.\src{leonova2015,cpic,molfluvox,asge2025,acghp}

\item \textbf{Долгосрочная ценность определяется управляемым балансом пользы и риска.} Обоснованную терапию сохраняют, пока существует показание, а ее интенсивность и необходимость регулярно пересматривают; отмена не должна быть ни автоматической, ни неопределенно откладываемой. При оценке безопасности установленные лекарственные реакции отличают от наблюдательных ассоциаций; восстановление секреции после прекращения приема не совпадает с исчезновением препарата из плазмы. Региональные различия регистрации и лекарственных форм требуют проверки конкретного продукта, а не переноса иностранных допусков по аналогии. Итоговое место лансопразола --- не универсальное преимущество над всеми ИПП, а возможность обоснованно выбрать дозу, сочетание и продолжительность лечения при сохранении границ доказательности. Последовательность «показание --- режим --- объективный результат --- пересмотр» объединяет гастроэнтерологическую и междисциплинарную практику и составляет главный практический вывод монографии. Критерии нефрологической настороженности и ограничения интерпретации новых исследований систематизированы в разделе~\ref{sec:nephrology}.\src{agadepres,kidney,uslabel,molrestore,withdrawbell,emc,grls}
\end{enumerate}

\section{Перспективные научно-практические парадоксы лансопразола: векторы для будущих клинико-фармакологических и фундаментальных исследований}
\label{sec:research-agenda}

Настоящий раздел формулирует программу исследований, объединяющую биофизику мишени, клиническую гастроэнтерологию, междисциплинарную безопасность и фармакогенетику. Под парадоксом понимается несоответствие между наблюдением и предполагаемым объяснением либо между результатами, полученными на разных уровнях организации. Двадцать четыре направления сгруппированы в четыре взаимосвязанных блока. Каждый пункт содержит описание научного конфликта, условия проверки гипотезы и заключительную формулировку целевого направления теоретико-экспериментального анализа. Исследовательские направления не являются утверждениями о выполненных исследованиях или доказанных эффектах. Предполагаемое взаимодействие с белком не признается установленной мишенью без экспериментальной проверки; доклиническое наблюдение не расширяет клинических показаний.

Общая методология требует раздельного определения исходного рацемата, энантиомеров, лансопразол-сульфида и кислотно-активированных производных; измерения свободной, а не только общей концентрации; контроля pH, длительности экспозиции и жизнеспособности клеток. Для каждого направления задают компаратор, основной исход, критерий опровержения и условия переноса результатов. При сравнении клинических стратегий дополнительно учитывают генотип, сопутствующее лечение, чувствительность возбудителя, исходный риск и приверженность. Докинг, изменение экспрессии, ингибирование фермента и улучшение исхода пациента рассматриваются как разные ступени доказательства.

\subsection{Группа А. Фундаментальная биофизика и молекулярная кинетика мишени}
\begin{enumerate}[label=\textbf{\arabic*.},leftmargin=*,itemsep=0.9\baselineskip]

\item \textbf{Метаболизм лансопразол-сульфида при туберкулезе.} Восстановленный метаболит LPZS обладает доклинической антимикобактериальной активностью; установленная мишень --- QcrB, цитохром-$b$-содержащая субъединица комплекса $bc_1$. Обозначение DykB не следует использовать как синоним этой мишени. Конфликт заключается в различии внутриклеточной активности метаболита и достижимости эффективной свободной концентрации в человеческом легочном очаге при введении исходного препарата. Концентрация в гомогенате легкого не характеризует непосредственно экспозицию бактерии внутри макрофага.\src{moltb,moltbstructure}

Следует построить физиологически обоснованную модель раздельного образования и распределения лансопразола и LPZS, связав свободную тканевую экспозицию с экспериментальной фармакодинамической целью. Недостижимость цели при допустимой экспозиции исходного вещества будет основанием отклонить прямое перепрофилирование гастроэнтерологической схемы, даже при сохранении активности метаболита в клеточной модели.

\textit{Целевое направление теоретико-экспериментального анализа:} математическое и фармакокинетическое моделирование терапевтических концентраций лансопразол-сульфида в легочной ткани человека.

\item \textbf{Лизосомальный захват иммуноцитов.} Сопоставление $pK_{a1}\approx3{,}83$ с pH 4,5--5,0 задает задачу распределения слабого основания, но не доказывает достаточной кислотной активации и ковалентной блокады V-АТФазы. Лизосома, фагосома и секреторная гранула не являются взаимозаменяемыми компартментами; их pH необходимо измерять отдельно в макрофагах и нейтрофилах. Экспериментальные лизосомальные эффекты в опухолевых клетках не устанавливают тот же механизм в клетках крови при клинической экспозиции.\src{molchem,mollysosome}

Необходимы органелльная количественная фармакокинетика, измерение протонного транспорта и поиск белкового аддукта при контролируемом pH. Повышение органелльного pH без подтвержденного взаимодействия с V-АТФазой следует объяснять альтернативными механизмами; токсическое разрушение мембраны нельзя принимать за селективное ингибирование помпы.

\textit{Целевое направление теоретико-экспериментального анализа:} молекулярные механизмы внетаргетной аккумуляции лансопразола в лизосомальных компартментах клеток иммунной системы человека.

\item \textbf{Стереоселективная инверсия плазмы и ткани.} Преобладание R-энантиомера в плазменной экспозиции рацемата отражает стереоселективную фармакокинетику, но не устанавливает его пропорциональный вклад в модификацию желудочной помпы. Предположение о более быстром переносе S-формы остается самостоятельной гипотезой: плазменная AUC не измеряет канальцевую доставку, скорость активации и занятость цистеинов. Универсальную долю R-формы нельзя задавать без учета генотипа и режима.\src{futurestereo1996,molstereokin,molhumanstereo}

Следует сопоставить раздельные концентрации R- и S-форм, свободное связывание, мембранный перенос и образование продуктов активации. Модель должна предсказывать тканевой эффект при независимой проверке, а не воспроизводить только плазменные кривые. Отсутствие различий при равной свободной экспозиции будет аргументом против предполагаемого тканевого преимущества S-формы.

\textit{Целевое направление теоретико-экспериментального анализа:} энантиоселективная тканевая фармакокинетика лансопразола и математическое моделирование потоков энантиомеров.

\item \textbf{Восстановление дисульфидного аддукта и продолжительность блокады.} Ковалентная связь с помпой не тождественна химической невосстановимости. Эксперименты с восстановителями обосновывают возможность восстановления активности, но не определяют количественный вклад эндогенного глутатиона или тиоредоксина в живой слизистой. Конфликт состоит в необходимости отделить восстановление уже модифицированных молекул от синтеза и включения новых помп; доступ восстановителя к люминальному аддукту сам требует объяснения.\src{molcys,molrestore}

Предлагаются меченые пулы помп, количественная оценка аддуктов и раздельное изменение восстановительного потенциала и синтеза белка. Условия воспаления следует моделировать отдельно, сопоставляя доступность GSH/тиоредоксина и скорость восстановления аддукта. Если активность возвращается только за счет немодифицированного пула, гипотезу существенной физиологической регенерации аддукта следует ограничить, несмотря на ее воспроизводимость в бесклеточной системе.

\textit{Целевое направление теоретико-экспериментального анализа:} кинетика эндогенного восстановления дисульфидных связей H$^+$/K$^+$-АТФазы после ковалентной блокады лансопразолом.

\item \textbf{Аллостерическая конформационная подвижность E1--E2.} Каталитические переходы желудочной АТФазы связывают цитоплазматические домены с трансмембранным переносом ионов. Однако установленное действие кислотно-активированного ИПП на люминальные цистеины не доказывает предварительного связывания нейтрального лансопразола с цитоплазматическим участком в состоянии E1. Конфликт заключается в различении первичной аллостерии и вторичного изменения конформации после ковалентной модификации.\src{molstructure,molchem,molcys}

Молекулярную динамику необходимо сочетать с измерением связывания и каталитики в контролируемых состояниях фермента, отдельно исследуя исходное вещество и активированные продукты. Докинговая позиция без насыщаемого взаимодействия и функционального эффекта не подтверждает новую аллостерическую мишень.

\textit{Целевое направление теоретико-экспериментального анализа:} биофизическое моделирование предполагаемых аллостерических эффектов лансопразола на конформационные переходы H$^+$/K$^+$-АТФазы в каталитическом цикле.

\item \textbf{Рециклинг тубуловезикул и судьба модифицированных помп.} Участие Rab11a в стимулированном перемещении H$^+$/K$^+$-АТФазы подтверждено экспериментально. Из этого не следует, что лансопразол непосредственно нарушает Rab11a/SNARE-зависимый транспорт или направляет помпу к протеасомной деградации. Конфликт состоит в отделении нормальной перестройки секреторных мембран при снижении функции от специфического дефекта рециклинга ковалентно модифицированного белка.\src{futurerab1999,molstructure}

Следует одновременно измерять локализацию помп, взаимодействия транспортных белков, убиквитинирование и скорость деградации. Лизосомный и протеасомный пути должны проверяться раздельно. Сохраненный рециклинг при подавленной секреции опровергнет отождествление каталитической блокады с нарушением мембранного транспорта.

\textit{Целевое направление теоретико-экспериментального анализа:} влияние ковалентного ингибирования H$^+$/K$^+$-АТФазы лансопразолом на молекулярный трафик белковых комплексов париетальной клетки.

\end{enumerate}

\subsection{Группа Б. Клиническая гастроэнтерология и барьерные функции слизистой}
\begin{enumerate}[label=\textbf{\arabic*.},start=7,leftmargin=*,itemsep=0.9\baselineskip]

\item \textbf{Критическое плато внутрижелудочного pH при эрадикации \Hp.} Короткая экспозиция ИПП и обновление функционального пула помп могут ограничивать кислотосупрессию, но CYP2C19*17 не доказывает ускоренного синтеза помп. Генотип изменяет преимущественно метаболическую экспозицию; интенсивность секреции, включение резервных мишеней и чувствительность бактерии являются отдельными переменными. Высокая эрадикация в выбранной популяции не устанавливает паритета с P-CAB во всех генотипических группах.\src{cpic,hpphfuruta2001,hpphalcon2022}

Необходимы прямое рандомизированное сравнение оптимизированных полных схем, стратификация по CYP2C19 и антибиотикочувствительности, одинаковый контроль излечения и заранее заданная граница не меньшей эффективности. Суточная pH-метрия должна объяснять различия исходов, а не заменять их; отсутствие паритета должно сохраняться как допустимый результат.

\textit{Целевое направление теоретико-экспериментального анализа:} проверка терапевтического паритета лансопразола и калий-конкурентных блокаторов кислотопродукции в эрадикационных протоколах с учетом генетического полиморфизма.

\item \textbf{Ночной прорыв кислоты и автономная регуляция.} Ночной кислотный прорыв характеризует эпизод желудочного pH ниже 4 продолжительностью более часа; он не равнозначен доказанному патологическому рефлюксу в пищевод. Возможное рекрутирование немодифицированных помп при ночной стимуляции не означает разрушения ковалентной связи вагусным сигналом. Термин «вегетативная инверсия» требует операционального определения и не может использоваться как установленный механизм NAB.\src{futurenab1998,molrestore}

Следует синхронно исследовать время приема, питание, сон, кислотность и показатели автономной регуляции, дополняя клинические данные моделью рекрутирования помп. Различия ночного pH необходимо сопоставлять с симптомами и пищеводной экспозицией; изолированное изменение желудочного показателя не подтверждает клиническую необходимость усиления лечения.

\textit{Целевое направление теоретико-экспериментального анализа:} хронофармакология ночной секреции желудка и оценка роли вегетативной регуляции в возобновлении кислотопродукции после блокады АТФазы лансопразолом.

\item \textbf{Микробиомный перекрест верхних отделов ЖКТ.} Кислотосупрессия может сопровождаться усилением переноса оральных микроорганизмов в кишечник. Однако обнаружение стрептококковых последовательностей не доказывает жизнеспособной биопленки на эрозированной слизистой пищевода. Конфликт состоит в различении транзиторного переноса, колонизации и причинного участия в повреждении; антисекреторное лечение и состояние эпителия могут независимо изменять состав микробиоты.\src{safemicrotrial,safemicrocohort}

Продольное исследование должно включать парные оральные и слизистые образцы, пространственную визуализацию, оценку жизнеспособности и барьерной функции. Контроль антибиотиков, диеты и загрязнения при заборе необходим для причинного вывода. Динамику микробной колонизации следует сопоставлять со скоростью эпителизации при последовательной оценке одних и тех же участков слизистой. Отсутствие тканевой организации микроорганизмов будет основанием отказаться от термина «биопленка».

\textit{Целевое направление теоретико-экспериментального анализа:} трансформация локального микробиома гастроэзофагеальной зоны при глубокой кислотосупрессии лансопразолом и оценка кинетики эпителизации.

\item \textbf{Кальциевая сигнализация ECL-клеток при отмене.} Гипергастринемия и CCK2R-зависимое увеличение ECL-клеточного пула имеют экспериментальное обоснование. Независимое от CCK2-рецепторов изменение экспрессии потенциалзависимых кальциевых каналов Cav под действием лансопразола остается гипотезой. Увеличение суммарного кальциевого сигнала может отражать число клеток, их дифференцировку или усиление рецепторной стимуляции, а не прямую модуляцию канала.\src{withdrawsheng2020,withdrawbrand1988,withdrawboyce2017}

Необходимо сопоставить экспрессию Cav, токи в отдельных клетках, секрецию гистамина и динамику гастрина до отмены и после нее. Контроль CCK2R-сигнализации позволит проверить независимый механизм. Нормализация сигнала после учета числа ECL-клеток будет аргументом против прямого канального эффекта.

\textit{Целевое направление теоретико-экспериментального анализа:} молекулярные триггеры кальциевой сигнализации энтерохромаффиноподобных клеток при отмене лансопразола.

\item \textbf{Химиопревенция при пищеводе Барретта и трофические эффекты гастрина.} AGA 2025 рассматривает применение ИПП при пищеводе Барретта, но клиническое положение не доказывает молекулоспецифической онкопревенции лансопразола. Гастрин-зависимый рост клеток желудка нельзя автоматически переносить на метапластический эпителий пищевода. Конфликт заключается в оценке потенциально противоположных последствий уменьшения рефлюксного повреждения и гормональной адаптации при разных тканевых мишенях.\src{barrett2025,withdrawsheng2020}

Следует разделять желудочные и пищеводные модели, проверять экспрессию рецепторов гастрина и оценивать прогрессию дисплазии в клинической когорте. Пролиферативный маркер не равнозначен раку; отсутствие гастрин-зависимого эффекта в пищеводной ткани должно опровергать соответствующее звено гипотезы.

\textit{Целевое направление теоретико-экспериментального анализа:} сопоставление молекулярных последствий длительной кислотосупрессии лансопразолом и оценка рисков гастрин-индуцированной пролиферации эпителия.

\end{enumerate}

\subsection{Группа В. Междисциплинарные коморбидные риски и синергизм}
\begin{enumerate}[label=\textbf{\arabic*.},start=12,leftmargin=*,itemsep=0.9\baselineskip]

\item {\raggedright \textbf{ОИН: иммунное повреждение и смешение по показанию.} Лансопразол-ассоциированный ОИН является клинически значимой возможной реакцией, однако специфическая гаптенизация метаболитом не установлена для каждого случая. Эпидемиологическая ассоциация ИПП с почечным исходом может зависеть от полиморбидности и причин назначения. Эти объяснения не взаимоисключающие: наличие смешения в когорте не опровергает иммунный механизм у конкретного пациента, а биопсия не определяет популяционную частоту.\src{safeain,renalblank2014,renalbigan2025}\par}

Предлагается объединить клиническую верификацию, анализ лекарственных белковых аддуктов и клеточного иммунного ответа с сопоставленной контрольной группой. Диагностическое повторное введение подозреваемого препарата не требуется. Отсутствие специфического аддукта при подтвержденном ОИН должно вести к проверке альтернативного иммунного механизма, а не к отрицанию повреждения.

{\raggedright \textit{Целевое направление теоретико-экспериментального анализа:} иммунопатология лансопразол-индуцированного интерстициального нефрита и верификация предполагаемых гаптеновых механизмов.\par}

\item \textbf{ДАТТ-метаболический сдвиг при полипрагмазии.} Влияние лансопразола на активацию клопидогрела зависит от генотипа, экспозиции и состава сопутствующего лечения. Нельзя считать все статины, антикоагулянты и аспирин одинаковыми конкурентами CYP2C19. Константа ингибирования в микросомах, концентрация активного метаболита клопидогрела и частота MACE измеряют разные уровни взаимодействия; лабораторное различие может не иметь самостоятельного клинического эффекта.\src{cardioli2004,cardioogilvie2011,cardiofrelinger2012,cardiocpic2022}

Следует анализировать конкретные сочетания лекарств, раздельно моделировать CYP2C19 и CYP3A, измерять активный метаболит и тромбоцитарную функцию. Проспективная оценка ишемических и геморрагических исходов необходима для проверки клинической значимости; отсутствие эффекта на MACE должно ограничивать выводы даже при фармакокинетическом сигнале.

\textit{Целевое направление теоретико-экспериментального анализа:} молекулярные взаимодействия при двойной антитромбоцитарной терапии и оценка совокупного ингибирования цитохромов сопутствующими препаратами.

\item \textbf{Keap1--Nrf2--HO-1 и эотаксин-3 при эозинофильном эзофагите.} Экспериментальная активация Nrf2/HO-1 лансопразолом в желудочных клетках и подавление STAT6-зависимого эотаксина-3 в пищеводных моделях представляют разные механистические наблюдения. Их нельзя объединять в доказанную последовательную ось ремиссии ЭоЭ. «Лиганд-независимость» должна означать проверку независимости от определенного рецепторного сигнала, а не отсутствие какого-либо молекулярного взаимодействия.\src{molnrf2,eoestat6,eoe}

В пищеводных органоидах необходимо независимо изменять pH, экспозицию ИПП, Nrf2 и STAT6, измеряя CCL26 и барьерную функцию. Сохранение эффекта при блокаде Nrf2 опровергнет обязательность этой оси, но не исключит другого противовоспалительного механизма. Клинические биопсии должны проверять перенос результатов, а не подменять причинный эксперимент.

\textit{Целевое направление теоретико-экспериментального анализа:} внетаргетная модуляция интерлейкин-индуцированного воспаления и антиэозинофильный потенциал лансопразола.

\item \textbf{Кишечный транспорт железа и кислотозависимая мальабсорбция.} Попадание липофильной молекулы в энтероцит не доказывает ингибирования DMT1 или ферропортина. Снижение переноса железа может отражать его химическую доступность, изменение протонного градиента, регуляцию гепсидином либо прямое действие на транспортный белок. Для B12 отдельно проверяется захват комплекса внутреннего фактора с витамином через кубилин/амнионлесс; это иной субстрат и иной механизм.\src{safeiron,futurecubam2005,pharmalabel}

Нужны поляризованные кишечные модели с раздельным контролем апикального pH, доступного железа, гепсидина и свободной концентрации препарата. Экспрессию, локализацию и поток субстрата измеряют независимо. Исчезновение эффекта после выравнивания среды будет свидетельствовать против непосредственной блокады транспортера.

\textit{Целевое направление теоретико-экспериментального анализа:} прямое и опосредованное влияние длительной терапии лансопразолом на молекулярные транспортеры металлов в тонкой кишке.

\item \textbf{Гериатрический парадокс костной резорбции.} Предполагаемое ухудшение доступности кальция и прямые клеточные эффекты ИПП могут иметь разные направления. Нельзя переносить данные о V-АТФазе и омепразоле на лансопразол без проверки. В отдельной модели лансопразол при низких микромолярных концентрациях усиливал созревание остеокластов и резорбцию, тогда как более высокие концентрации подавляли образование клеток. Это требует отделения специфической антирезорбции от цитотоксичности.\src{futurebone2015,uslabel}

Следует сопоставить кишечное усвоение кальция, кислотность резорбционной лакуны, остеокластическую и остеобластическую функции при измеренных тканевых экспозициях. Возраст, питание, витамин D и сопутствующее лечение включают в клиническую модель. Изменение клеточного маркера без изменения механической прочности кости не устанавливает влияние на переломы.

\textit{Целевое направление теоретико-экспериментального анализа:} сопоставление разнонаправленных фармакодинамических эффектов лансопразола в костной ткани на системном и органном уровнях.

\item \textbf{Сосудистый эндотелиальный парадокс ADMA--DDAH.} Экспериментальные данные связывают ИПП с ингибированием диметиларгинина-диметиламиногидролазы, накоплением ADMA и уменьшением NO-зависимой вазодилатации. Но ингибирование eNOS эндогенным ADMA не означает непосредственного связывания лансопразола с eNOS. Доклиническая цепь также не доказывает причинного увеличения MACE при стандартной экспозиции у человека.\src{futureddah2013,cardiocogent2010}

Докинг следует использовать для выбора проверяемых позиций связывания, дополняя его биофизическим анализом, ферментной кинетикой и измерением ADMA/NO. Необходимы сравнение ИПП и контроль свободной концентрации. Несоответствие расчетного сродства экспериментальному ингибированию или клинически достижимой экспозиции будет ограничивать сосудистую гипотезу.

\textit{Целевое направление теоретико-экспериментального анализа:} сравнительный молекулярный докинг лансопразола в активный центр диметиларгининаминогидролазы сосудов человека с экспериментальной проверкой связывания.

\item \textbf{Транспортное взаимодействие с метотрексатом.} Замедленное выведение высокодозного метотрексата при сопутствующем ИПП и ингибирование ABCG2/BCRP в мембранных везикулах поддерживают исследование транспортного механизма. Однако экспериментальные ингибирующие концентрации превышали свободные плазменные концентрации ИПП; авторы не объясняли взаимодействие только BCRP. Результаты онкологических высокодозных режимов нельзя непосредственно переносить на еженедельное ревматологическое лечение.\src{futuremtx2009,uslabel}

Следует раздельно оценить вклад почечной функции, дозы метотрексата, локальной экспозиции и альтернативных транспортных путей. Везикулярные опыты должны сопоставляться с клинической фармакокинетикой. Сохранение взаимодействия при отсутствии достаточного ингибирования ABCG2 потребует пересмотра единственного транспортного объяснения, а не отрицания клинического риска.

\textit{Целевое направление теоретико-экспериментального анализа:} фармакокинетическое взаимодействие лансопразола с метотрексатом на уровне АТФ-связывающих транспортеров.

\item \textbf{Системный Th1/Th2-профиль и локальный ответ при ЭоЭ.} Уменьшение эозинофилии слизистой и продукции эпителиального эотаксина-3 не устанавливает системного перепрограммирования Т-хелперов. Ответ может быть обусловлен восстановлением барьера, уменьшением локальной антигенной стимуляции или прямым воздействием на иммунные клетки. Конфликт состоит в разграничении локальной эпителиальной реакции и изменения дифференцировки CD4$^+$-клеток; один сывороточный цитокин не характеризует весь Th2-профиль.\src{eoestat6,eoe,molnrf2}

Предлагаются парные исследования крови и биопсий до и после лечения, функциональные тесты Т-клеток и сопоставление с одинаковой кислотосупрессией альтернативным средством. Ремиссия слизистой без устойчивого изменения Т-клеточного профиля будет аргументом против обязательного системного механизма.

\textit{Целевое направление теоретико-экспериментального анализа:} влияние лансопразола на системный цитокиновый профиль и дифференцировку Т-хелперов 2-го типа при эозинофильном воспалении.

\item \textbf{ОРИТ: профилактика кровотечения и инфекционные исходы.} SCCM/ASHP 2024 связывают профилактику стрессовых кровотечений с клиническим риском и прекращением профилактики после его устранения. Гипохлоргидрия создает предпосылки изменения микробной экологии, но последовательность «СИБР --- транслокация --- сепсис» не является неизбежным следствием ИПП. Конфликт определяется балансом потенциальной пользы и инфекционного риска при влиянии тяжести болезни, питания и антибиотиков.\src{sccm,safemicrotrial}

Исследование должно раздельно регистрировать клинически значимое кровотечение, инфекцию, микробиом и признаки нарушения барьера; факт транслокации требует собственного подтверждения. Сопоставление групп по тяжести и экспозиции обязательно. Микробиомное изменение без увеличения инфекционных исходов не доказывает клинический вред, а классовые результаты не устанавливают особого риска лансопразола.

\textit{Целевое направление теоретико-экспериментального анализа:} оценка баланса пользы и риска антисекреторной терапии в ОРИТ и вероятности системной бактериальной транслокации.

\end{enumerate}

\subsection{Группа Г. Клиническая фармакогенетика и биотрансформация}
\begin{enumerate}[label=\textbf{\arabic*.},start=21,leftmargin=*,itemsep=0.9\baselineskip]

\item \textbf{Гериатрическая феноконверсия и возрастной клиренс.} Меньшая зависимость рабепразола от CYP2C19 вследствие неферментного превращения не означает неизменности его экспозиции при любой коморбидности. Доступность лансопразола 15 мг также не доказывает достаточного эффекта и безопасности для каждого пожилого пациента. Конфликт состоит в сравнении метаболической зависимости и управляемости дозы при неоднородных изменениях функции печени, воспаления, белкового связывания и полипрагмазии.\src{cpic,uslabel,leonova2015}

Продольное исследование должно учитывать биологическую хрупкость, генотип, свободную экспозицию и клинический ответ, сравнивая целые стратегии дозирования. Возрастная ассоциация без измеренного изменения метаболической активности не должна называться феноконверсией. Превосходство одной стратегии по безопасности или эффективности необходимо проверять, а не предполагать из пути биотрансформации.

\textit{Целевое направление теоретико-экспериментального анализа:} возрастная динамика фармакогенетического профиля ИПП в гериатрической практике интерниста.

\item \textbf{CYP2C19*17 и цитокин-индуцированная феноконверсия.} IL-6 подавлял экспрессию ряда CYP в культуре человеческих гепатоцитов, что обосновывает возможность расхождения генотипического прогноза и текущего метаболического состояния. Однако наличие *17 и системного воспаления не доказывает обязательного превращения UM в функционального PM: эффект зависит от степени подавления, исходной активности и сопутствующих воздействий. Генотип при этом не изменяется.\src{futureil62011,cpic}

Следует выполнять серийную оценку воспаления, исходного вещества и метаболитов на фоне восстановления пациента, учитывая питание, функцию печени и ингибиторы CYP. Обратимость изменения экспозиции при снижении воспаления усилит причинную модель; отсутствие связи после учета органной дисфункции ограничит самостоятельный вклад цитокинов.

\textit{Целевое направление теоретико-экспериментального анализа:} цитокин-индуцированная феноконверсия метаболизма лансопразола у критически больных пациентов.

\item \textbf{Интермедиат-зависимая инактивация CYP3A4.} Участие CYP3A в метаболизме лансопразола не доказывает, что промежуточный метаболит необратимо инактивирует этот фермент при увеличении метаболического потока. Обратимое ингибирование, индукция через регуляторные рецепторы и метаболизм-зависимая инактивация могут давать разные временные профили. Данные о взаимодействии ИПП с CYP2C19 нельзя автоматически переносить на CYP3A4.\src{molp450,molcypinhib,molpxr,cardioogilvie2011}

Необходимы опыты с предварительной инкубацией, контролем NADPH, времени и концентрации, последующим разведением и оценкой восстановления активности. Поиск аддукта должен дополнять кинетику. Сохранение активности после удаления свободного ингибитора будет свидетельствовать против необратимого механизма; положительный результат требует проверки при клинически достижимой экспозиции.

\textit{Целевое направление теоретико-экспериментального анализа:} оценка рисков интермедиат-индуцированной инактивации изофермента CYP3A4 при применении лансопразола.

\item \textbf{OATP1B1 и генетическая вариабельность печеночного захвата.} Полиморфизмы SLCO1B1 могут влиять на фармакокинетику доказанных субстратов OATP1B1, но принадлежность лансопразола к таким субстратам нельзя выводить из липофильности или печеночного метаболизма. Вначале необходимо установить вклад транспортера относительно пассивного поступления и CYP-зависимого удаления. Ассоциация варианта SLCO1B1 с плазменной концентрацией без функциональной проверки не идентифицирует механизм.\src{pharmalabel,molp450,cpic}

Следует сравнить захват в контрольных и экспрессирующих OATP1B1 клетках, затем исследовать функциональные варианты при одинаковой свободной концентрации и контроле метаболизма. Генотипическую когорту целесообразно анализировать после подтверждения транспортного вклада. Отсутствие насыщаемого специфического захвата ограничит гипотезу прямой роли SLCO1B1, даже при статистической ассоциации.

\textit{Целевое направление теоретико-экспериментального анализа:} роль полиморфизмов транспортера OATP1B1 и гена SLCO1B1 в межиндивидуальной вариабельности фармакокинетики лансопразола.

\end{enumerate}

\chapter*{Справочный аппарат}
\addcontentsline{toc}{chapter}{Справочный аппарат}
\section*{Сокращения}
\addcontentsline{toc}{section}{Сокращения}
\begin{description}[style=nextline]
\item[ИПП / PPI] Ингибитор протонной помпы / proton pump inhibitor.
\item[P-CAB] Калий-конкурентный блокатор кислотопродукции.
\item[ГЭРБ / GERD] Гастроэзофагеальная рефлюксная болезнь.
\item[ЭоЭ] Эозинофильный эзофагит.
\item[ЛОР; ОРИТ] Оториноларингология; отделение реанимации и интенсивной терапии.
\item[ОФВ$_1$] Объем форсированного выдоха за первую секунду.
\item[LPZS] Лансопразол-сульфид; восстановленное производное, отличное от исходного лансопразола-сульфоксида.
\item[ATP4A; ATP4B] Каталитическая $\alpha$- и вспомогательная $\beta$-субъединицы желудочной H$^+$/K$^+$-АТФазы.
\item[HO-1; ROS; IL-8] Гемоксигеназа-1; реактивные формы кислорода; интерлейкин-8.
\item[QcrB] Компонент микобактериального комплекса $bc_1$, рассматриваемый в экспериментальном направлении LPZS.
\item[ДПК; НПВП; ЖКТ] Двенадцатиперстная кишка; нестероидные противовоспалительные препараты; желудочно-кишечный тракт.
\item[LA] Лос-Анджелесская классификация эрозивного эзофагита.
\item[SmPC] Краткая характеристика лекарственного препарата.
\item[ACG; AGA; ESGE] Американская коллегия гастроэнтерологов; Американская гастроэнтерологическая ассоциация; Европейское общество гастроинтестинальной эндоскопии.
\item[CPIC] Консорциум по внедрению клинической фармакогенетики.
\item[ДАТТ; MACE] Двойная антитромбоцитарная терапия; большие неблагоприятные сердечно-сосудистые события (состав конечной точки определяется исследованием).
\item[ОИН; ОПП; ХБП] Острый интерстициальный нефрит; острое повреждение почек; хроническая болезнь почек.
\item[OR; HR; ROR] Отношение шансов; отношение мгновенных рисков; отношение шансов сообщения в фармаконадзоре.
\item[ACR] Отношение альбумина к креатинину мочи; в тексте приведено в мг/ммоль.
\item[FAERS] Система сообщений о нежелательных явлениях FDA (FDA Adverse Event Reporting System).
\item[ДИ; NNT; СКФ; МНО] Доверительный интервал; число пациентов, которых нужно лечить для предотвращения одного события за указанный срок; скорость клубочковой фильтрации; международное нормализованное отношение.
\end{description}

\section*{Источники и проверяемость}
\addcontentsline{toc}{section}{Источники и проверяемость}
Ссылки ниже ведут к исходным документам, их официальным страницам или библиографическим записям. Основная дата обращения --- 14.09.2026; источники целевых редакционных дополнений проверялись 15.09.2026. Индивидуальные даты и ограничения проверки указываются в описании соответствующего источника или раздела. Для динамических инструкций дата обновления страницы может отличаться от даты редакции медицинского текста. Регистрационные инструкции не являются независимыми сравнительными исследованиями; в основном тексте отдельно используются их разделы дозирования, безопасности, взаимодействий и применения у детей. Для китайского руководства 2026 года использована доступная аннотация; ограничения российского реестрового поиска указаны в начале обзора.

\begin{thebibliography}{999}
\addcontentsline{toc}{chapter}{\bibname}
\small
\raggedright

\bibitem{uslabel} DailyMed. \textit{Lansoprazole delayed-release capsules}. Medcore, LLC. Редакция 01/2026; страница обновлена 16.01.2026. Разделы 1--2, 4--8, Instructions for Use. \href{https://dailymed.nlm.nih.gov/dailymed/lookup.cfm?setid=18b4020f-d7e6-42f8-838d-2b5da152530e}{Официальный текст инструкции}.

\bibitem{hpra} Health Products Regulatory Authority, Ireland. \textit{Lansoprazole Teva Pharma 30 mg gastro-resistant capsules, hard}. PA1986/008/002. SmPC, 08.02.2024. \href{https://assets.hpra.ie/products/Human/29860/Licence_PA1986-008-002_08022024114836.pdf}{Документ регулятора ЕС}.

\bibitem{emc} Electronic Medicines Compendium. \textit{Lansoprazole 30 mg Orodispersible Tablets}; SmPC, product 4396; также Lupin, product 14713, обновление 01.10.2024. \href{https://www.medicines.org.uk/emc/product/4396/smpc}{SmPC 4396}; \href{https://www.medicines.org.uk/emc/product/14713/smpc}{SmPC 14713}.

\bibitem{pharmalabel} DailyMed. \textit{Lansoprazole delayed-release capsules}. Preferred Pharmaceuticals / исходный продукт Dr. Reddy's. Обновление страницы 12.01.2026; редакция медицинского текста 01/2023. Разделы 11--12: механизм действия, фармакодинамика, фармакокинетика. \href{https://dailymed.nlm.nih.gov/dailymed/lookup.cfm?setid=960d6601-aec5-4900-a400-385b6687a4bc}{Текст}.





\bibitem{otc} DailyMed. \textit{Lansoprazole USP 15 mg}, OTC Drug Facts. Dr. Reddy's Laboratories. Обновление 21.07.2026. \href{https://www.dailymed.nlm.nih.gov/dailymed/drugInfo.cfm?setid=4a701c0d-fef5-aa12-01e6-e3680a161a7d}{Маркировка OTC}.

\bibitem{dexlabel} DailyMed. \textit{Dexlansoprazole delayed-release capsules}. Редакция 02/2026. \href{https://www.dailymed.nlm.nih.gov/dailymed/drugInfo.cfm?setid=31d3d485-5f87-4405-b940-490fa4ccf9cc}{Инструкция}.

\bibitem{acggerd} Katz PO, Dunbar KB, Schnoll-Sussman FH, et al. ACG Clinical Guideline for the Diagnosis and Management of Gastroesophageal Reflux Disease. \textit{Am J Gastroenterol}. 2022;117:27--56. \href{https://doi.org/10.14309/ajg.0000000000001538}{doi:10.14309/ajg.0000000000001538}.

\bibitem{agagerd} Yadlapati R, Gyawali CP, Pandolfino JE. AGA Clinical Practice Update on the Personalized Approach to the Evaluation and Management of GERD: Expert Review. 2022. \href{https://doi.org/10.1016/j.cgh.2022.01.025}{doi:10.1016/j.cgh.2022.01.025}. \href{https://gastro.org/clinical-guidance/personalized-approach-to-the-evaluation-and-management-of-gastroesophageal-reflux-disease-gerd/}{Рекомендации AGA}.

\bibitem{agadepres} Targownik LE, Fisher DA, Saini SD. AGA Clinical Practice Update on De-Prescribing of Proton Pump Inhibitors: Expert Review. 2022. \href{https://doi.org/10.1053/j.gastro.2021.12.247}{doi:10.1053/j.gastro.2021.12.247}. \href{https://gastro.org/clinical-guidance/de-prescribing-proton-pump-inhibitors-ppis/}{Советы AGA}.

\bibitem{acghp} Chey WD, Howden CW, Moss SF, et al. ACG Clinical Guideline: Treatment of Helicobacter pylori Infection. \textit{Am J Gastroenterol}. 2024;119:1730--1753. \href{https://doi.org/10.14309/ajg.0000000000002968}{doi:10.14309/ajg.0000000000002968}. \href{https://gi.org/journals-publications/ebgi/schoenfeld_sep2024/}{Официальное резюме ACG}.

\bibitem{acghighlights} American College of Gastroenterology. \textit{ACG Guideline Highlights: Treatment of Helicobacter pylori Infection}. 2024. \href{https://gi.org/wp-content/uploads/2025/01/ACG-Hpylori-Guidelines-Highlights-2024-FINAL.pdf.pdf}{Таблицы режимов и контроль эрадикации}.

\bibitem{hpphblum1998} Blum RA, Hunt RH, Kidd SL, Shi H, Jennings DE, Greski-Rose PA. Dose-response relationship of lansoprazole to gastric acid antisecretory effects. \textit{Aliment Pharmacol Ther}. 1998;12:321--327. \href{https://doi.org/10.1046/j.1365-2036.1998.00306.x}{doi:10.1046/j.1365-2036.1998.00306.x}; \href{https://pubmed.ncbi.nlm.nih.gov/9690720/}{PMID:9690720}.

\bibitem{hpphfuruta2001} Furuta T, Shirai N, Xiao F, Ohashi K, Ishizaki T. Effect of high-dose lansoprazole on intragastic pH in subjects who are homozygous extensive metabolizers of cytochrome P4502C19. \textit{Clin Pharmacol Ther}. 2001;70:484--492. \href{https://doi.org/10.1067/mcp.2001.119721}{doi:10.1067/mcp.2001.119721}; \href{https://pubmed.ncbi.nlm.nih.gov/11719736/}{PMID:11719736}.

\bibitem{hpgrowth2012} Marcus EA, Inatomi N, Nagami GT, Sachs G, Scott DR. The effects of varying acidity on Helicobacter pylori growth and the bactericidal efficacy of ampicillin. \textit{Aliment Pharmacol Ther}. 2012;36:972--979. \href{https://doi.org/10.1111/apt.12059}{doi:10.1111/apt.12059}; \href{https://pubmed.ncbi.nlm.nih.gov/23009227/}{PMID:23009227}.

\bibitem{hphighdose2013} Prasertpetmanee S, Mahachai V, Vilaichone RK. Improved efficacy of proton pump inhibitor - amoxicillin - clarithromycin triple therapy for Helicobacter pylori eradication in low clarithromycin resistance areas or for tailored therapy. \textit{Helicobacter}. 2013;18:270--273. \href{https://doi.org/10.1111/hel.12041}{doi:10.1111/hel.12041}; \href{https://pubmed.ncbi.nlm.nih.gov/23356886/}{PMID:23356886}.

\bibitem{hptailored2007} Furuta T, Shirai N, Kodaira M, et al. Pharmacogenomics-based tailored versus standard therapeutic regimen for eradication of H. pylori. \textit{Clin Pharmacol Ther}. 2007;81:521--528. \href{https://doi.org/10.1038/sj.clpt.6100043}{doi:10.1038/sj.clpt.6100043}; \href{https://pubmed.ncbi.nlm.nih.gov/17215846/}{PMID:17215846}.

\bibitem{hpphalcon2022} Chey WD, Mégraud F, Laine L, López LJ, Hunt BJ, Howden CW. Vonoprazan Triple and Dual Therapy for Helicobacter pylori Infection in the United States and Europe: Randomized Clinical Trial. \textit{Gastroenterology}. 2022;163:608--619. \href{https://doi.org/10.1053/j.gastro.2022.05.055}{doi:10.1053/j.gastro.2022.05.055}; \href{https://pubmed.ncbi.nlm.nih.gov/35679950/}{PMID:35679950}.

\bibitem{hpcost2021} Tokunaga K, Suzuki C, Hasegawa M, Fujimori I. Cost Analysis in Helicobacter pylori Eradication Therapy Based on a Database of Health Insurance Claims in Japan. \textit{Clinicoecon Outcomes Res}. 2021;13:241--250. \href{https://doi.org/10.2147/CEOR.S297680}{doi:10.2147/CEOR.S297680}; \href{https://pubmed.ncbi.nlm.nih.gov/33889000/}{PMID:33889000}.

\bibitem{maastricht} Malfertheiner P, Megraud F, Rokkas T, et al. Management of Helicobacter pylori infection: the Maastricht VI/Florence consensus report. \textit{Gut}. 2022;71:1724--1762. \href{https://doi.org/10.1136/gutjnl-2022-327745}{doi:10.1136/gutjnl-2022-327745}.

\bibitem{dgvs} DGVS. \textit{S2k-Leitlinie Gastroösophageale Refluxkrankheit und eosinophile Ösophagitis}. AWMF 021-013, Version 3.0, 03/2023. \href{https://www.dgvs.de/leitlinien/oberer-gi-trakt/gastrooesophageale-refluxkrankheit-und-eosinophile-oesophagitis/?digital=1}{Официальная цифровая версия}.

\bibitem{nice} NICE. \textit{Gastro-oesophageal reflux disease and dyspepsia in adults: investigation and management}. CG184. 2014; обновление 18.10.2019. \href{https://www.nice.org.uk/guidance/cg184/chapter/Recommendations}{Рекомендации}; \href{https://www.nice.org.uk/guidance/cg184/chapter/appendix-a-dosage-information-on-proton-pump-inhibitors}{Appendix A: дозы ИПП}.

\bibitem{lyon} Gyawali CP, Yadlapati R, Fass R, et al. Updates to the modern diagnosis of GERD: Lyon consensus 2.0. \textit{Gut}. 2024;73:361--371. \href{https://doi.org/10.1136/gutjnl-2023-330616}{doi:10.1136/gutjnl-2023-330616}.

\bibitem{esge2026} Gralnek IM, Morris J, Laursen SB, et al. Endoscopic diagnosis and management of peptic ulcer bleeding: ESGE Guideline --- Update 2026. \textit{Endoscopy}. 2026;58:899--924. Онлайн 13.05.2026. \href{https://doi.org/10.1055/a-2863-8314}{doi:10.1055/a-2863-8314}. \href{https://www.esge.com/guidelines}{Реестр ESGE}.

\bibitem{jgerd} Japanese Society of Gastroenterology. Evidence-based clinical practice guidelines for gastroesophageal reflux disease 2021. Англоязычная публикация: \textit{J Gastroenterol}. 2022. \href{https://pmc.ncbi.nlm.nih.gov/articles/8938399/}{Полный текст, PMC8938399}.

\bibitem{jhp} Isomoto H, Shimoyama T, Ito M, et al. Guidelines for the management of Helicobacter pylori infection in Japan: 2024 revised edition. \textit{J Gastroenterol}. 2026;61:841--862. \href{https://doi.org/10.1007/s00535-026-02379-4}{doi:10.1007/s00535-026-02379-4}.

\bibitem{jpud2026} Japanese Society of Gastroenterology. Evidence-based clinical practice guidelines for peptic ulcer disease 2026. \textit{J Gastroenterol}. 2026. \href{https://doi.org/10.1007/s00535-026-02504-3}{doi:10.1007/s00535-026-02504-3}.

\bibitem{seoul} Huh CW, Chang JW, Park MI, et al. 2025 Focused Update of the Seoul Consensus on Gastroesophageal Reflux Disease: Evidence-based Recommendations on Acid Suppressive Therapy. \textit{J Neurogastroenterol Motil}. 2026;32:7--18; онлайн 29.09.2025. \href{https://doi.org/10.5056/jnm25128}{doi:10.5056/jnm25128}.

\bibitem{khp} Bang CS, Kang SJ, Lim H, et al. Evidence-Based Guidelines for the Diagnosis and Treatment of Helicobacter pylori Infection in Korea: 2025 Revised Edition. \textit{Helicobacter}. 2026;31:e70149. \href{https://doi.org/10.1111/hel.70149}{doi:10.1111/hel.70149}. \href{https://pubmed.ncbi.nlm.nih.gov/42439470/}{PMID:42439470}.

\bibitem{china2022} Zhou L, Lu H, Song Z, et al. 2022 Chinese national clinical practice guideline on Helicobacter pylori eradication treatment. \textit{Chin Med J}. 2022;135:2899--2910. \href{https://pmc.ncbi.nlm.nih.gov/articles/PMC10106216/}{Полный текст, PMC10106216}.

\bibitem{china2026} Helicobacter pylori Study Group, Chinese Society of Gastroenterology, Chinese Medical Association. Clinical practice guidelines for the diagnosis, treatment, and prevention of Helicobacter pylori infection in Chinese adults. \textit{Zhonghua Nei Ke Za Zhi}. 2026;65:651--680. На китайском языке; использована аннотация. \href{https://doi.org/10.3760/cma.j.cn112138-20260330-00183}{doi:10.3760/cma.j.cn112138-20260330-00183}. \href{https://pubmed.ncbi.nlm.nih.gov/42427040/}{PMID:42427040}.

\bibitem{rfgerd} Ивашкин В.Т. и соавт. Диагностика и лечение гастроэзофагеальной рефлюксной болезни: рекомендации российских профессиональных обществ. \textit{Рос журн гастроэнтерол гепатол колопроктол}. 2024;34(5):111--135. \href{https://doi.org/10.22416/1382-4376-2024-34-5-111-135}{doi:10.22416/1382-4376-2024-34-5-111-135}. \href{https://www.gastro.ru/userfiles/R_GERB2024.pdf}{Текст РГА}.

\bibitem{rfpud} Ивашкин В.Т. и соавт. Диагностика и лечение язвенной болезни у взрослых: клинические рекомендации российских профессиональных обществ. \textit{Рос журн гастроэнтерол гепатол колопроктол}. 2024;34(2):101--131. \href{https://doi.org/10.22416/1382-4376-2024-34-2-101-131}{doi:10.22416/1382-4376-2024-34-2-101-131}. \href{https://www.gastro.ru/userfiles/R_Yaz_2024.pdf}{Текст РГА}.

\bibitem{rfgastrit} Клинические рекомендации «Гастрит и дуоденит». 2024. Утвержденный текст в републикации справочной системы ГАРАНТ. \href{https://base.garant.ru/409334054/}{Доступный текст}. Републикация использована для чтения; статус последней редакции проверяется в Рубрикаторе отдельно.

\bibitem{registry} Минздрав России. \textit{Рубрикатор клинических рекомендаций}. Карточка «Гастрит и дуоденит», 708\_2. \href{https://cr.minzdrav.gov.ru/view-cr/708_2}{Карточка документа}. Полная автоматизированная проверка содержимого не выполнена.

\bibitem{grls} Минздрав России. \textit{Государственный реестр лекарственных средств}. \href{https://grls.minzdrav.gov.ru/}{ГРЛС}. Используется для проверки конкретного регистрационного удостоверения и инструкции; полный перечень действующих форм лансопразола в этом обзоре не верифицирован.

\bibitem{sccm} SCCM / ASHP. Guideline for the Prevention of Stress-Related Gastrointestinal Bleeding in Critically Ill Adults. 15.07.2024. \href{https://www.sccm.org/clinical-resources/guidelines/guidelines/sccm-ashp-guideline-prevention-of-ugib}{Рекомендации общества}.

\bibitem{eoe} Dellon ES, Muir AB, Katzka DA, et al. ACG Clinical Guideline: Diagnosis and Management of Eosinophilic Esophagitis. \textit{Am J Gastroenterol}. 2025;120:31--59. \href{https://doi.org/10.14309/ajg.0000000000003194}{doi:10.14309/ajg.0000000000003194}.


\bibitem{eoehighlights2025} American College of Gastroenterology. \textit{ACG Guideline Highlights: Diagnosis and Management of Eosinophilic Esophagitis}. 2025. \href{https://webfiles.gi.org/links/committees/DigitalCommsAndPubs/ACG-EOE-Guidelines-Highlights-2025.pdf}{Официальная памятка ACG}. Дата обращения: 15.09.2026.

\bibitem{eoestat6} Zhang X, Cheng E, Huo X, et al. Omeprazole blocks STAT6 binding to the eotaxin-3 promoter in eosinophilic esophagitis cells. \textit{PLoS One}. 2012;7:e50037. \href{https://doi.org/10.1371/journal.pone.0050037}{doi:10.1371/journal.pone.0050037}.

\bibitem{eoedupilumab2022} Dellon ES, Rothenberg ME, Collins MH, et al. Dupilumab in adults and adolescents with eosinophilic esophagitis. \textit{N Engl J Med}. 2022;387:2317--2330. \href{https://doi.org/10.1056/NEJMoa2205982}{doi:10.1056/NEJMoa2205982}. Источник более строгой исследовательской конечной точки, не испытание лансопразола.

\bibitem{eoeguide} ACG. A Look at the Updated ACG Eosinophilic Esophagitis Clinical Guidelines. 2025. \href{https://gi.org/journals-publications/ebgi/eluri_feb2025/}{Официальное резюме с практическими пояснениями}.

\bibitem{leonova2015} Леонова М.В. Фармакогенетика ингибиторов протонной помпы. \textit{Медицинский совет}. 2015;(17):96--102. Полный текст на сайте журнала: \href{https://www.med-sovet.pro/jour/article/view/433/433}{Медицинский совет}. Учтен при редакционном дополнении 15.09.2026. Численные результаты приведены по обзору, а не как повторная независимая проверка всех цитируемых в нем первичных исследований.

\bibitem{cpic} Lima JJ, Thomas CD, Barbarino J, et al. Clinical Pharmacogenetics Implementation Consortium Guideline for CYP2C19 and Proton Pump Inhibitor Dosing. \textit{Clin Pharmacol Ther}. 2021;109:1417--1423; онлайн 2020. \href{https://files.cpicpgx.org/data/guideline/publication/PPI/2020/32770672.pdf}{Документ CPIC}. PMID:32770672.

\bibitem{vontrial} Laine L, DeVault K, Katz P, et al. Vonoprazan Versus Lansoprazole for Healing and Maintenance of Healing of Erosive Esophagitis: A Randomized Trial. \textit{Gastroenterology}. 2023;164:61--71. \href{https://doi.org/10.1053/j.gastro.2022.09.041}{doi:10.1053/j.gastro.2022.09.041}.

\bibitem{fdctakeldalaunch} Takeda Pharmaceutical Company Limited. Takelda Combination Tablets, the Fixed-Dose Combination of Lansoprazole and Low-Dose Aspirin, is Now Available in Japan. 12.06.2014. \href{https://www.takeda.com/newsroom/newsreleases/2014/takelda-combination-tablets-the-fixed-dose-combination-of-lansoprazole-and-low-dose-aspirin-is-now-available-in-japan/}{Сообщение о начале продаж в Японии}. Ожидаемая польза для приверженности отделена от результатов сравнительного исследования.

\bibitem{fdctakeldalabel} PMDA. Takelda Combination Tablets (aspirin 100 mg / lansoprazole 15 mg). Японская инструкция, редакция 14.07.2026; разделы 4, 6, 14, 16--18. \href{https://www.pmda.go.jp/PmdaSearch/iyakuDetail/400061_3399102F1026_2_09}{Официальная инструкция}. Дата обращения: 15.09.2026.

\bibitem{fdctakeldareview} Pharmaceuticals and Medical Devices Agency. Takelda Combination Tablets: Review Report. 10.01.2014. Английская версия японского регистрационного отчета; разделы о составе, биоэквивалентности и клиническом обосновании. \href{https://www.pmda.go.jp/files/000208812.pdf}{Отчет PMDA}.

\bibitem{fdcnaprapac} FDA. Prevacid NapraPAC 500: lansoprazole delayed-release 15 mg capsules and naproxen 500 mg tablets kit. Инструкция, май 2009 года. \href{https://www.accessdata.fda.gov/drugsatfda_docs/label/2009/021507s009lbl.pdf}{Архивная инструкция FDA}. Исторический пример совместной упаковки, а не подтверждение текущей доступности.

\bibitem{fdcvimovo} Agencia Española de Medicamentos y Productos Sanitarios. Vimovo 500 mg/20 mg, comprimidos de liberación modificada. Ficha técnica; разделы 2--5. \href{https://cima.aemps.es/cima/dochtml/ft/73182/FT_73182.html}{Официальная инструкция AEMPS}. Дата обращения: 15.09.2026. Компонент ИПП --- эзомепразол, не лансопразол.

\bibitem{aspirintrial} Lai KC, Lam SK, Chu KM, et al. Lansoprazole for the Prevention of Recurrences of Ulcer Complications from Long-Term Low-Dose Aspirin Use. \textit{N Engl J Med}. 2002;346:2033--2038. \href{https://doi.org/10.1056/NEJMoa012877}{doi:10.1056/NEJMoa012877}.

\bibitem{toppits} O'Hara J, Stocken DD, Watson GC, et al. Use of proton pump inhibitors to treat persistent throat symptoms: multicentre, double blind, randomised, placebo controlled trial. \textit{BMJ}. 2021;372:m4903. \href{https://doi.org/10.1136/bmj.m4903}{doi:10.1136/bmj.m4903}.

\bibitem{asthma} Holbrook JT, et al. Lansoprazole for children with poorly controlled asthma: a randomized controlled trial. \textit{JAMA}. 2012. \href{https://pubmed.ncbi.nlm.nih.gov/22274684/}{PMID:22274684}.

\bibitem{infant} Orenstein SR, Hassall E, Furmaga-Jablonska W, et al. Multicenter, double-blind, randomized, placebo-controlled trial assessing the efficacy and safety of proton pump inhibitor lansoprazole in infants with symptoms of gastroesophageal reflux disease. \textit{J Pediatr}. 2009;154:514--520.e4. \href{https://doi.org/10.1016/j.jpeds.2008.09.054}{doi:10.1016/j.jpeds.2008.09.054}.

\bibitem{compass} Moayyedi P, Eikelboom JW, Bosch J, et al. Safety of Proton Pump Inhibitors Based on a Large, Multi-Year, Randomized Trial of Patients Receiving Rivaroxaban or Aspirin. \textit{Gastroenterology}. 2019;157:682--691.e2. \href{https://www.sciencedirect.com/science/article/pii/S0016508519409748}{Статья исследования}. PMID:31152740.

\bibitem{kidney} Pyne L, Smyth A, Molnar AO, et al. The Effects of Pantoprazole on Kidney Outcomes: Post Hoc Observational Analysis from the COMPASS Trial. \textit{J Am Soc Nephrol}. 2024;35:901--909. \href{https://doi.org/10.1681/ASN.0000000000000356}{doi:10.1681/ASN.0000000000000356}.

\bibitem{uktis} UKTIS / BUMPS. \textit{Proton pump inhibitors (PPIs)}. July 2025, Version 1.1. \href{https://www.medicinesinpregnancy.org/leaflets-a-z/proton-pump-inhibitors/}{Информация о применении при беременности}.

\bibitem{lactmed} National Institute of Child Health and Human Development. \textit{Lansoprazole}. Drugs and Lactation Database (LactMed). Обновление 15.02.2025. \href{https://www.ncbi.nlm.nih.gov/books/NBK501283/}{NBK501283}.


\bibitem{cardioli2004} Li XQ, Andersson TB, Ahlström M, Weidolf L. Comparison of inhibitory effects of the proton pump-inhibiting drugs omeprazole, esomeprazole, lansoprazole, pantoprazole, and rabeprazole on human cytochrome P450 activities. \textit{Drug Metab Dispos}. 2004;32:821--827. \href{https://doi.org/10.1124/dmd.32.8.821}{doi:10.1124/dmd.32.8.821}.

\bibitem{cardioogilvie2011} Ogilvie BW, Yerino P, Kazmi F, et al. The proton pump inhibitor, omeprazole, but not lansoprazole or pantoprazole, is a metabolism-dependent inhibitor of CYP2C19: implications for coadministration with clopidogrel. \textit{Drug Metab Dispos}. 2011;39:2020--2033. \href{https://doi.org/10.1124/dmd.111.041293}{doi:10.1124/dmd.111.041293}.

\bibitem{cardiofrelinger2012} Frelinger AL III, Lee RD, Mulford DJ, et al. A randomized, 2-period, crossover design study to assess the effects of dexlansoprazole, lansoprazole, esomeprazole, and omeprazole on the steady-state pharmacokinetics and pharmacodynamics of clopidogrel in healthy volunteers. \textit{J Am Coll Cardiol}. 2012;59:1304--1311. \href{https://doi.org/10.1016/j.jacc.2011.12.024}{doi:10.1016/j.jacc.2011.12.024}.

\bibitem{cardioplavix2025} Sanofi. \textit{Plavix (clopidogrel): U.S. prescribing information}. Редакция 05.2025; разделы 5.1 и 7.2. FDA/DailyMed. \href{https://dailymed.nlm.nih.gov/dailymed/fda/fdaDrugXsl.cfm?setid=de8b0b67-eb25-4684-83b5-7ad785314227}{Американская инструкция}. Дата обращения: 15.09.2026.

\bibitem{cardioesc2023} Byrne RA, Rossello X, Coughlan JJ, et al. 2023 ESC Guidelines for the management of acute coronary syndromes. \textit{Eur Heart J}. 2023;44:3720--3826. Раздел 13.3.7 и таблица 7. \href{https://doi.org/10.1093/eurheartj/ehad191}{doi:10.1093/eurheartj/ehad191}.

\bibitem{cardioesc2024} Vrints C, Andreotti F, Koskinas KC, et al. 2024 ESC Guidelines for the management of chronic coronary syndromes. \textit{Eur Heart J}. 2024;45:3415--3537. Раздел 4.3.1.4. \href{https://doi.org/10.1093/eurheartj/ehae177}{doi:10.1093/eurheartj/ehae177}.

\bibitem{cardiocogent2010} Bhatt DL, Cryer BL, Contant CF, et al. Clopidogrel with or without omeprazole in coronary artery disease. \textit{N Engl J Med}. 2010;363:1909--1917. \href{https://doi.org/10.1056/NEJMoa1007964}{doi:10.1056/NEJMoa1007964}.

\bibitem{sps} NHS Specialist Pharmacy Service. \textit{Using clopidogrel with proton pump inhibitors (PPIs)}. 14.09.2023. \href{https://www.sps.nhs.uk/articles/using-clopidogrel-with-proton-pump-inhibitors-ppis/}{Рекомендации NHS SPS}.

\bibitem{asge2025} ASGE Standards of Practice Committee; Desai M, et al. American Society for Gastrointestinal Endoscopy guideline on the diagnosis and management of GERD: summary and recommendations. \textit{Gastrointest Endosc}. 2025. \href{https://pubmed.ncbi.nlm.nih.gov/39692638/}{PMID:39692638}.

\bibitem{barrett2025} Wani S, et al. AGA Clinical Practice Guideline on Surveillance of Barrett's Esophagus. 20.10.2025. \href{https://doi.org/10.1053/j.gastro.2025.09.012}{doi:10.1053/j.gastro.2025.09.012}. \href{https://gastro.org/clinical-guidance/surveillance-of-barretts-esophagus/}{Рекомендации AGA}.

\bibitem{pedhp2024} Homan M, et al. Updated joint ESPGHAN/NASPGHAN guidelines for management of Helicobacter pylori infection in children and adolescents (2023). \textit{J Pediatr Gastroenterol Nutr}. 2024;79:758--785. \href{https://doi.org/10.1002/jpn3.12314}{doi:10.1002/jpn3.12314}. \href{https://naspghan.org/wp-content/uploads/2024/10/J-pediatr-gastroenterol-nutr-2024-Homan-Updated-joint-ESPGHAN-NASPGHAN-guidelines-for-management-of-Helicobacter.pdf}{Официальная копия NASPGHAN}.

\bibitem{rfpudmin} Клинические рекомендации «Язвенная болезнь». 2024. ID 277\_2; взрослые. Разработчики: Российское общество хирургов, Российская гастроэнтерологическая ассоциация и другие профильные общества. \href{https://base.garant.ru/411274125/}{Републикация утвержденного текста в ГАРАНТ}; \href{https://cr.minzdrav.gov.ru/view-cr/277_2}{Карточка Минздрава}; \href{https://apicr.minzdrav.gov.ru/api.ashx?id=277_2&op=GetClinrecPdf}{Прямая ссылка на официальный PDF}. Сверены доступные разделы; исчерпывающая проверка последующих редакций не выполнена.

\bibitem{indiahp} Singh SP, Ahuja V, Ghoshal UC, et al. Management of Helicobacter pylori infection: The Bhubaneswar Consensus Report of the Indian Society of Gastroenterology. \textit{Indian J Gastroenterol}. 2021;40:420--444. \href{https://doi.org/10.1007/s12664-021-01186-4}{doi:10.1007/s12664-021-01186-4}. \href{https://pubmed.ncbi.nlm.nih.gov/34219211/}{PMID:34219211}. Использованы положения 18, 20--24.

\bibitem{india2026} Sinha SK, Kochhar R, Rao KS, et al. Acid peptic disorder compass: An Indian expert consensus on comprehensive management. \textit{World J Gastrointest Pharmacol Ther}. 05.03.2026;17(1):112615. \href{https://doi.org/10.4292/wjgpt.v17.i1.112615}{doi:10.4292/wjgpt.v17.i1.112615}. \href{https://www.wjgnet.com/2150-5349/full/v17/i1/112615.htm}{Полный текст и раскрытие конфликтов интересов}.


\bibitem{nsaidagrawal2000} Agrawal NM, et al. Superiority of lansoprazole vs ranitidine in healing nonsteroidal anti-inflammatory drug-associated gastric ulcers: results of a double-blind, randomized, multicenter study. \textit{Arch Intern Med}. 2000. \href{https://pubmed.ncbi.nlm.nih.gov/10826458/}{PMID:10826458}.

\bibitem{nsaidomnium1998} Hawkey CJ, Karrasch JA, Szczepañski L, et al. Omeprazole compared with misoprostol for ulcers associated with nonsteroidal antiinflammatory drugs. \textit{N Engl J Med}. 1998;338:727--734. \href{https://doi.org/10.1056/NEJM199803123381105}{doi:10.1056/NEJM199803123381105}.

\bibitem{agaextra2023} Chen JW, Vela MF, Peterson KA, Carlson DA. AGA Clinical Practice Update on the Diagnosis and Management of Extraesophageal Gastroesophageal Reflux Disease: Expert Review. \textit{Clin Gastroenterol Hepatol}. 2023. \href{https://doi.org/10.1016/j.cgh.2023.01.040}{doi:10.1016/j.cgh.2023.01.040}.

\bibitem{aao2018} Stachler RJ, Francis DO, Schwartz SR, et al. Clinical Practice Guideline: Hoarseness (Dysphonia) (Update). \textit{Otolaryngol Head Neck Surg}. 2018;158(1 Suppl):S1--S42. \href{https://doi.org/10.1177/0194599817751030}{doi:10.1177/0194599817751030}.

\bibitem{chestcough2016} Kahrilas PJ, et al. Chronic Cough Due to Gastroesophageal Reflux in Adults: CHEST Guideline and Expert Panel Report. \textit{Chest}. 2016;150:1341--1360. \href{https://doi.org/10.1016/j.chest.2016.08.1458}{doi:10.1016/j.chest.2016.08.1458}.

\bibitem{gina2026} Global Initiative for Asthma. \textit{Global Strategy for Asthma Management and Prevention}. 2026. Раздел 6, Gastroesophageal reflux disease (GERD). \href{https://ginasthma.org/wp-content/uploads/2026/05/GINA-2026-Strategy-Report-WMS.pdf}{Отчет GINA 2026}. Дата обращения: 15.09.2026.

\bibitem{asthmaesome2009} American Lung Association Asthma Clinical Research Centers. Efficacy of esomeprazole for treatment of poorly controlled asthma. \textit{N Engl J Med}. 2009;360:1487--1499. \href{https://doi.org/10.1056/NEJMoa0806290}{doi:10.1056/NEJMoa0806290}.

\bibitem{misotrial} Graham DY, et al. Ulcer prevention in long-term users of nonsteroidal anti-inflammatory drugs: results of a double-blind, randomized, multicenter, active- and placebo-controlled study of misoprostol vs lansoprazole. \textit{Arch Intern Med}. 2002. \href{https://pubmed.ncbi.nlm.nih.gov/11802750/}{PMID:11802750}.

\bibitem{aspirin15} Sugano K, Matsumoto Y, Itabashi T, et al. Lansoprazole for secondary prevention of gastric or duodenal ulcers associated with long-term low-dose aspirin therapy: results of a prospective, multicenter, double-blind, randomized, double-dummy, active-controlled trial. \textit{J Gastroenterol}. 2011;46:724--735. \href{https://doi.org/10.1007/s00535-011-0397-7}{doi:10.1007/s00535-011-0397-7}. \href{https://pubmed.ncbi.nlm.nih.gov/21499703/}{PMID:21499703}.

\bibitem{nsaid15} Sugano K, Kontani T, Katsuo S, et al. Lansoprazole for secondary prevention of gastric or duodenal ulcers associated with long-term non-steroidal anti-inflammatory drug (NSAID) therapy: results of a prospective, multicenter, double-blind, randomized, double-dummy, active-controlled trial. \textit{J Gastroenterol}. 2012;47:540--552. \href{https://pubmed.ncbi.nlm.nih.gov/22388884/}{PMID:22388884}; \href{https://pmc.ncbi.nlm.nih.gov/articles/PMC3360874/}{Полный текст}.

\bibitem{bsgfd} Black CJ, et al. British Society of Gastroenterology guidelines on the management of functional dyspepsia. \textit{Gut}. 2022;71:1697--1723. \href{https://doi.org/10.1136/gutjnl-2022-327737}{doi:10.1136/gutjnl-2022-327737}. \href{https://pmc.ncbi.nlm.nih.gov/articles/PMC9380508/}{Полный текст}.

\bibitem{agapcab} Patel A, Laine L, Moayyedi P, Wu J. AGA Clinical Practice Update on Integrating Potassium-Competitive Acid Blockers Into Clinical Practice: Expert Review. \textit{Gastroenterology}. 2024;167:1228--1238. \href{https://doi.org/10.1053/j.gastro.2024.06.038}{doi:10.1053/j.gastro.2024.06.038}. \href{https://gastro.org/news/new-guidance-when-to-use-potassium-competitive-acid-blockers/}{Официальное резюме AGA}.

\bibitem{molchem} Shin JM, Cho YM, Sachs G. Chemistry of covalent inhibition of the gastric (H$^+$, K$^+$)-ATPase by proton pump inhibitors. \textit{J Am Chem Soc}. 2004;126:7800--7811. \href{https://doi.org/10.1021/ja049607w}{doi:10.1021/ja049607w}; \href{https://pubmed.ncbi.nlm.nih.gov/15212527/}{PMID:15212527}.

\bibitem{molphselect1998} Kromer W, Krüger U, Huber R, Hartmann M, Steinijans VW. Differences in pH-dependent activation rates of substituted benzimidazoles and biological in vitro correlates. \textit{Pharmacology}. 1998;56:57--70. \href{https://doi.org/10.1159/000028183}{doi:10.1159/000028183}; \href{https://pubmed.ncbi.nlm.nih.gov/9494064/}{PMID:9494064}.

\bibitem{molkinetics} Al-Matar AK, El-Eswed B, Tutunji MF. Kinetics of acid degradation of proton pump inhibitors in the presence of a thiol. \textit{Int J Chem Kinet}. 2009;41:498--506. \href{https://doi.org/10.1002/kin.20415}{doi:10.1002/kin.20415}. \href{https://eacademic.ju.edu.jo/maha.tutunji/Lists/Published%20Research/Attachments/37/Abstract.pdf}{Авторская копия в университетском репозитории}.

\bibitem{molnagaya89} Nagaya H, Satoh H, Kubo K, Maki Y. Possible mechanism for the inhibition of gastric (H$^+$ + K$^+$)-adenosine triphosphatase by the proton pump inhibitor AG-1749. \textit{J Pharmacol Exp Ther}. 1989;248:799--805. \href{https://pubmed.ncbi.nlm.nih.gov/2537417/}{PMID:2537417}.

\bibitem{molnagaya90} Nagaya H, Satoh H, Maki Y. Possible mechanism for the inhibition of acid formation by the proton pump inhibitor AG-1749 in isolated canine parietal cells. \textit{J Pharmacol Exp Ther}. 1990;252:1289--1295. \href{https://pubmed.ncbi.nlm.nih.gov/2156997/}{PMID:2156997}.

\bibitem{molcys} Besancon M, Simon A, Sachs G, Shin JM. Sites of reaction of the gastric H,K-ATPase with extracytoplasmic thiol reagents. \textit{J Biol Chem}. 1997;272:22438--22446. \href{https://doi.org/10.1074/jbc.272.36.22438}{doi:10.1074/jbc.272.36.22438}; \href{https://pubmed.ncbi.nlm.nih.gov/9278394/}{PMID:9278394}.

\bibitem{molrestore} Shin JM, Sachs G. Restoration of acid secretion following treatment with proton pump inhibitors. \textit{Gastroenterology}. 2002;123:1588--1597. \href{https://doi.org/10.1053/gast.2002.36593}{doi:10.1053/gast.2002.36593}.

\bibitem{molstructure} Abe K, Irie K, Nakanishi H, et al. Crystal structures of the gastric proton pump. \textit{Nature}. 2018;556:214--218. \href{https://doi.org/10.1038/s41586-018-0003-8}{doi:10.1038/s41586-018-0003-8}.

\bibitem{molk} Yamamoto K, Dubey V, Irie K, et al. A single K$^+$-binding site in the crystal structure of the gastric proton pump. \textit{eLife}. 2019;8:e47701. \href{https://doi.org/10.7554/eLife.47701}{doi:10.7554/eLife.47701}.

\bibitem{molselect} Madapally HV, Abe K, Dubey V, Khandelia H. Specific protonation of acidic residues confers K$^+$ selectivity to the gastric proton pump. \textit{J Biol Chem}. 2024;300:105542; онлайн 2023. \href{https://doi.org/10.1016/j.jbc.2023.105542}{doi:10.1016/j.jbc.2023.105542}.

\bibitem{molezrin} Ding X, Deng H, Wang D, et al. Phospho-regulated ACAP4-Ezrin interaction is essential for histamine-stimulated parietal cell secretion. \textit{J Biol Chem}. 2010;285:18769--18780. \href{https://doi.org/10.1074/jbc.M110.129007}{doi:10.1074/jbc.M110.129007}.

\bibitem{molp450} Pearce RE, Rodrigues AD, Goldstein JA, Parkinson A. Identification of the human P450 enzymes involved in lansoprazole metabolism. \textit{J Pharmacol Exp Ther}. 1996;277:805--816. \href{https://pubmed.ncbi.nlm.nih.gov/8627562/}{PMID:8627562}.

\bibitem{molurease} Nagata K, Satoh H, Iwahi T, et al. Potent inhibitory action of the gastric proton pump inhibitor lansoprazole against urease activity of Helicobacter pylori: unique action selective for H. pylori cells. \textit{Antimicrob Agents Chemother}. 1993;37:769--774. \href{https://doi.org/10.1128/AAC.37.4.769}{doi:10.1128/AAC.37.4.769}.

\bibitem{molgrowth} Nagata K, Takagi E, Tsuda M, et al. Inhibitory action of lansoprazole and its analogs against Helicobacter pylori: inhibition of growth is not related to inhibition of urease. \textit{Antimicrob Agents Chemother}. 1995;39:567--570. \href{https://doi.org/10.1128/AAC.39.2.567}{doi:10.1128/AAC.39.2.567}; \href{https://pubmed.ncbi.nlm.nih.gov/7726537/}{PMID:7726537}.

\bibitem{molresp} Nagata K, et al. Inhibitory activities of lansoprazole against respiration in Helicobacter pylori. \textit{Antimicrob Agents Chemother}. 2001;45:1522--1527. \href{https://doi.org/10.1128/AAC.45.5.1522-1527.2001}{doi:10.1128/AAC.45.5.1522-1527.2001}; \href{https://pmc.ncbi.nlm.nih.gov/articles/PMC90499/}{Полный текст}.

\bibitem{molnrf2} Takagi T, Naito Y, Okada H, et al. Lansoprazole, a proton pump inhibitor, mediates anti-inflammatory effect in gastric mucosal cells through the induction of heme oxygenase-1 via activation of NF-E2-related factor 2 and oxidation of kelch-like ECH-associating protein 1. \textit{J Pharmacol Exp Ther}. 2009;331:255--264. \href{https://doi.org/10.1124/jpet.109.152702}{doi:10.1124/jpet.109.152702}.

\bibitem{molho1} Schulz-Geske S, Erdmann K, Wong RJ, et al. Molecular mechanism and functional consequences of lansoprazole-mediated heme oxygenase-1 induction. \textit{World J Gastroenterol}. 2009;15:4392--4401. \href{https://doi.org/10.3748/wjg.15.4392}{doi:10.3748/wjg.15.4392}.

\bibitem{molil8} Isomoto H, Saenko VA, Kanazawa Y, et al. Enhanced expression of interleukin-8 and activation of nuclear factor kappa-B in endoscopy-negative gastroesophageal reflux disease. \textit{Am J Gastroenterol}. 2004;99:589--597. \href{https://doi.org/10.1111/j.1572-0241.2004.04110.x}{doi:10.1111/j.1572-0241.2004.04110.x}; \href{https://pubmed.ncbi.nlm.nih.gov/15089887/}{PMID:15089887}.

\bibitem{mollysosome} Takeda A, Takano N, Kokuba H, et al. Macrolide antibiotics enhance the antitumor effect of lansoprazole resulting in lysosomal membrane permeabilization-associated cell death. \textit{Int J Oncol}. 2020;57:1280--1292. \href{https://doi.org/10.3892/ijo.2020.5138}{doi:10.3892/ijo.2020.5138}.

\bibitem{moltb} Rybniker J, Vocat A, Sala C, et al. Lansoprazole is an antituberculous prodrug targeting cytochrome $bc_1$. \textit{Nat Commun}. 2015;6:7659. \href{https://doi.org/10.1038/ncomms8659}{doi:10.1038/ncomms8659}.

\bibitem{moltbstructure} Kovalova T, Król S, Gamiz-Hernandez AP, et al. Inhibition mechanism of potential antituberculosis compound lansoprazole sulfide. \textit{Proc Natl Acad Sci USA}. 2024;121:e2412780121. \href{https://doi.org/10.1073/pnas.2412780121}{doi:10.1073/pnas.2412780121}.

\bibitem{molkchannel} Lambrecht NWG, Yakubov I, Scott D, Sachs G. Identification of the K efflux channel coupled to the gastric H-K-ATPase during acid secretion. \textit{Physiol Genomics}. 2005;21:81--91. \href{https://doi.org/10.1152/physiolgenomics.00212.2004}{doi:10.1152/physiolgenomics.00212.2004}; \href{https://pubmed.ncbi.nlm.nih.gov/15613615/}{PMID:15613615}.

\bibitem{molh2signal} Wang L, Gantz I, DelValle J. Histamine H2 receptor activates adenylate cyclase and PLC via separate GTP-dependent pathways. \textit{Am J Physiol}. 1996;271:G613--G620. \href{https://doi.org/10.1152/ajpgi.1996.271.4.G613}{doi:10.1152/ajpgi.1996.271.4.G613}.

\bibitem{molprobes} Sugita Y, Nagao T, Urushidani T. Nonspecific effects of the pharmacological probes commonly used to analyze signal transduction in rabbit parietal cells. \textit{Eur J Pharmacol}. 1999;365:77--89. \href{https://doi.org/10.1016/S0014-2999(98)00855-3}{doi:10.1016/S0014-2999(98)00855-3}; \href{https://pubmed.ncbi.nlm.nih.gov/9988126/}{PMID:9988126}.

\bibitem{molm3} Aihara T, Fujishita T, Kanatani K, et al. Impaired gastric secretion and lack of trophic responses to hypergastrinemia in M3 muscarinic receptor knockout mice. \textit{Gastroenterology}. 2003;125:1774--1784. \href{https://doi.org/10.1053/j.gastro.2003.09.018}{doi:10.1053/j.gastro.2003.09.018}.

\bibitem{molcck} Kanai S, Hosoya H, Akimoto S, et al. Gastric acid secretion in cholecystokinin-1 receptor, -2 receptor, and -1, -2 receptor gene knockout mice. \textit{J Physiol Sci}. 2009;59:23--29. \href{https://doi.org/10.1007/s12576-008-0001-y}{doi:10.1007/s12576-008-0001-y}.

\bibitem{molsst} Zhao CM, Martinez V, Piqueras L, et al. Control of gastric acid secretion in somatostatin receptor 2 deficient mice: shift from endocrine/paracrine to neurocrine pathways. \textit{Endocrinology}. 2008;149:498--505. \href{https://doi.org/10.1210/en.2007-0238}{doi:10.1210/en.2007-0238}.

\bibitem{molca} Karhukorpi EK. Carbonic anhydrase II in rat acid secreting cells: comparison of osteoclasts with gastric parietal cells and kidney intercalated cells. \textit{Acta Histochem}. 1991;90:11--20. \href{https://doi.org/10.1016/S0065-1281(11)80151-0}{doi:10.1016/S0065-1281(11)80151-0}.

\bibitem{molae2} Gawenis LR, Ledoussal C, Judd LM, et al. Mice with a targeted disruption of the AE2 Cl$^{-}$/HCO$_3^{-}$ exchanger are achlorhydric. \textit{J Biol Chem}. 2004;279:30531--30539. \href{https://doi.org/10.1074/jbc.M403779200}{doi:10.1074/jbc.M403779200}.

\bibitem{molslc26} Xu J, Song P, Miller ML, et al. Deletion of the chloride transporter Slc26a9 causes loss of tubulovesicles in parietal cells and impairs acid secretion in the stomach. \textit{Proc Natl Acad Sci USA}. 2008;105:17955--17960. \href{https://doi.org/10.1073/pnas.0800616105}{doi:10.1073/pnas.0800616105}.

\bibitem{molclc2} Nighot MP, Nighot PK, Ma TY, et al. Genetic ablation of the ClC-2 Cl$^{-}$ channel disrupts mouse gastric parietal cell acid secretion. \textit{PLoS One}. 2015;10:e0138174. \href{https://doi.org/10.1371/journal.pone.0138174}{doi:10.1371/journal.pone.0138174}.

\bibitem{molstereokin} Kim KA, Kim MJ, Park JY, et al. Stereoselective metabolism of lansoprazole by human liver cytochrome P450 enzymes. \textit{Drug Metab Dispos}. 2003;31:1227--1234. \href{https://doi.org/10.1124/dmd.31.10.1227}{doi:10.1124/dmd.31.10.1227}; \href{https://pubmed.ncbi.nlm.nih.gov/12975331/}{PMID:12975331}.

\bibitem{molhumanstereo} Miura M, Tada H, Yasui-Furukori N, et al. Pharmacokinetic differences between the enantiomers of lansoprazole and its metabolite, 5-hydroxylansoprazole, in relation to CYP2C19 genotypes. \textit{Eur J Clin Pharmacol}. 2004;60:623--628. \href{https://doi.org/10.1007/s00228-004-0809-1}{doi:10.1007/s00228-004-0809-1}.

\bibitem{molfluvox} Miura M, Tada H, Yasui-Furukori N, et al. Enantioselective disposition of lansoprazole in relation to CYP2C19 genotypes in the presence of fluvoxamine. \textit{Br J Clin Pharmacol}. 2005;60:61--68. \href{https://doi.org/10.1111/j.1365-2125.2005.02381.x}{doi:10.1111/j.1365-2125.2005.02381.x}; \href{https://pmc.ncbi.nlm.nih.gov/articles/PMC1884909/}{Полный текст}.

\bibitem{molcypinhib} Liu KH, Kim MJ, Shon JH, et al. Stereoselective inhibition of cytochrome P450 forms by lansoprazole and omeprazole in vitro. \textit{Xenobiotica}. 2005;35:27--38. \href{https://doi.org/10.1080/00498250400026472}{doi:10.1080/00498250400026472}; \href{https://pubmed.ncbi.nlm.nih.gov/15788366/}{PMID:15788366}.

\bibitem{molahr} Novotna A, Srovnalova A, Svecarova M, et al. Differential effects of omeprazole and lansoprazole enantiomers on aryl hydrocarbon receptor in human hepatocytes and cell lines. \textit{PLoS One}. 2014;9:e98711. \href{https://doi.org/10.1371/journal.pone.0098711}{doi:10.1371/journal.pone.0098711}.

\bibitem{molpxr} Novotna A, Dvorak Z. Omeprazole and lansoprazole enantiomers induce CYP3A4 in human hepatocytes and cell lines via glucocorticoid receptor and pregnane X receptor axis. \textit{PLoS One}. 2014;9:e105580. \href{https://doi.org/10.1371/journal.pone.0105580}{doi:10.1371/journal.pone.0105580}.

\bibitem{safemgclaudin} Thongon N, Krishnamra N. Apical acidity decreases inhibitory effect of omeprazole on Mg$^{2+}$ absorption and claudin-7 and -12 expression in Caco-2 monolayers. \textit{Exp Mol Med}. 2012;44:684--693. \href{https://doi.org/10.3858/emm.2012.44.11.077}{doi:10.3858/emm.2012.44.11.077}.

\bibitem{safemggen} Hess MW, de Baaij JHF, Broekman MMTJ, et al. Common single nucleotide polymorphisms in transient receptor potential melastatin type 6 increase the risk for proton pump inhibitor-induced hypomagnesemia: a case-control study. \textit{Pharmacogenet Genomics}. 2017;27:83--88. \href{https://doi.org/10.1097/FPC.0000000000000259}{doi:10.1097/FPC.0000000000000259}.

\bibitem{safeiron} Hamano H, Niimura T, Horinouchi Y, et al. Proton pump inhibitors block iron absorption through direct regulation of hepcidin via the aryl hydrocarbon receptor-mediated pathway. \textit{Toxicol Lett}. 2020;318:86--91. \href{https://doi.org/10.1016/j.toxlet.2019.10.016}{doi:10.1016/j.toxlet.2019.10.016}.

\bibitem{safemicrotrial} Zhu J, Sun C, Li M, et al. Compared to histamine-2 receptor antagonist, proton pump inhibitor induces stronger oral-to-gut microbial transmission and gut microbiome alterations: a randomised controlled trial. \textit{Gut}. 2024;73:1087--1097. \href{https://doi.org/10.1136/gutjnl-2023-330168}{doi:10.1136/gutjnl-2023-330168}.

\bibitem{safemicrocohort} Hojo M, Asahara T, Nagahara A, et al. Gut microbiota composition before and after use of proton pump inhibitors. \textit{Dig Dis Sci}. 2018;63:2940--2949. \href{https://doi.org/10.1007/s10620-018-5122-4}{doi:10.1007/s10620-018-5122-4}.

\bibitem{safecolitis} Bonderup OK, Nielsen GL, Dall M, et al. Significant association between the use of different proton pump inhibitors and microscopic colitis: a nationwide Danish case-control study. \textit{Aliment Pharmacol Ther}. 2018;48:618--625. \href{https://doi.org/10.1111/apt.14916}{doi:10.1111/apt.14916}.

\bibitem{safecolitisguide} Miehlke S, Guagnozzi D, Zabana Y, et al. European guidelines on microscopic colitis: United European Gastroenterology and European Microscopic Colitis Group statements and recommendations. \textit{United European Gastroenterol J}. 2021;9:13--37. \href{https://doi.org/10.1177/2050640620951905}{doi:10.1177/2050640620951905}. \href{https://www.emcg-ibd.eu/pdf/Guidelines.pdf}{Документ EMCG}.

\bibitem{renalantoniou2015} Antoniou T, Macdonald EM, Hollands S, et al. Proton pump inhibitors and the risk of acute kidney injury in older patients: a population-based cohort study. \textit{CMAJ Open}. 2015;3:E166--E171. \href{https://doi.org/10.9778/cmajo.20140074}{doi:10.9778/cmajo.20140074}; \href{https://pmc.ncbi.nlm.nih.gov/articles/PMC4571830/}{Полный текст}.

\bibitem{renalblank2014} Blank ML, Parkin L, Paul C, Herbison P. A nationwide nested case-control study indicates an increased risk of acute interstitial nephritis with proton pump inhibitor use. \textit{Kidney Int}. 2014;86:837--844. \href{https://doi.org/10.1038/ki.2014.74}{doi:10.1038/ki.2014.74}; \href{https://pubmed.ncbi.nlm.nih.gov/24646856/}{PMID:24646856}.

\bibitem{renalniceaki} NICE. Acute kidney injury: prevention, detection and management. NG148. Обновление 16.10.2024. Разделы 1.3--1.5. \href{https://www.nice.org.uk/guidance/ng148/chapter/Recommendations}{Рекомендации NICE}.

\bibitem{renalniceckd} NICE. Chronic kidney disease: assessment and management. NG203. Раздел 1.1: измерение и интерпретация расчетной СКФ. \href{https://www.nice.org.uk/guidance/ng203/chapter/Recommendations}{Рекомендации NICE}. Дата обращения: 15.09.2026.

\bibitem{renaleosin2013} Muriithi AK, Nasr SH, Leung N. Utility of urine eosinophils in the diagnosis of acute interstitial nephritis. \textit{Clin J Am Soc Nephrol}. 2013;8:1857--1862. \href{https://doi.org/10.2215/CJN.01330213}{doi:10.2215/CJN.01330213}; \href{https://pubmed.ncbi.nlm.nih.gov/24052222/}{PMID:24052222}. Учтено наличие исправления 2018 года; численные диагностические характеристики здесь не воспроизводятся.

\bibitem{renalsteroids2018} Fernandez-Juarez G, Perez JV, Caravaca-Fontán F, et al. Duration of Treatment with Corticosteroids and Recovery of Kidney Function in Acute Interstitial Nephritis. \textit{Clin J Am Soc Nephrol}. 2018. \href{https://pubmed.ncbi.nlm.nih.gov/30397027/}{PMID:30397027}; \href{https://pmc.ncbi.nlm.nih.gov/articles/PMC6302327/}{Полный текст}.

\bibitem{oncoppiposition2016} Scarpignato C, Gatta L, Zullo A, Blandizzi C; SIF--AIGO--FIMMG Group. Effective and safe proton pump inhibitor therapy in acid-related diseases: a position paper addressing benefits and potential harms of acid suppression. \textit{BMC Med}. 2016;14:179. \href{https://pubmed.ncbi.nlm.nih.gov/27825371/}{PMID:27825371}; \href{https://www.webaigo.it/download/20180807092918.pdf}{Документ AIGO}.

\bibitem{oncoeugogo2021} Bartalena L, Kahaly GJ, Baldeschi L, et al. The 2021 European Group on Graves' orbitopathy (EUGOGO) clinical practice guidelines for the medical management of Graves' orbitopathy. \textit{Eur J Endocrinol}. 2021;185:G43--G67. \href{https://doi.org/10.1530/EJE-21-0479}{doi:10.1530/EJE-21-0479}; \href{https://orca.cardiff.ac.uk/id/eprint/143526/1/eje-21-0479.pdf}{Авторская версия, Cardiff University}.

\bibitem{oncobosulif2026} DailyMed. BOSULIF (bosutinib), tablets and capsules. Pfizer. Редакция 03/2026. Разделы 7.1 и 12.3. \href{https://dailymed.nlm.nih.gov/dailymed/lookup.cfm?setid=adc84ad5-a04d-4fee-9ba8-91f7abd928e3}{Официальная инструкция}.

\bibitem{oncodasatinib2026} DailyMed. DASATINIB tablets. Разделы взаимодействий и фармакокинетики; инструкция стандартной таблетированной формы. \href{https://dailymed.nlm.nih.gov/dailymed/drugInfo.cfm?setid=f60261dd-b8ea-4511-bed6-ad6d7f4dd828}{Официальная инструкция}. Дата обращения: 15.09.2026. Не переносить ограничения на иные формы без проверки.

\bibitem{oncocalcaps2026} DailyMed. CALQUENCE (acalabrutinib), capsules. AstraZeneca. Редакция 02/2026. Разделы 2, 7 и 12.3. \href{https://dailymed.nlm.nih.gov/dailymed/drugInfo.cfm?setid=dd4835ef-e1bc-4997-a399-1ffa2556fbfe}{Инструкция капсул}.

\bibitem{oncocaltablet2026} DailyMed. CALQUENCE (acalabrutinib), tablets. AstraZeneca. Редакция 02/2026. Разделы 5, 7 и 12.3. \href{https://dailymed.nlm.nih.gov/dailymed/drugInfo.cfm?setid=b1409184-3267-4d8c-8f17-4e01f69828b1}{Инструкция таблеток}.

\bibitem{oncophyrago2024} FDA. PHYRAGO (dasatinib), tablets. Инструкция, 2024; раздел 12.3. \href{https://www.accessdata.fda.gov/drugsatfda_docs/label/2024/216099s004lbl.pdf}{Регистрационный документ FDA}. Приведены данные взаимодействия именно этой формы с омепразолом и фамотидином.

\bibitem{safeain} Jose J, Saravu K, Khera K, et al. Acute interstitial nephritis related to lansoprazole administration. \textit{J Pak Med Assoc}. 2008;58:206--207. \href{https://pubmed.ncbi.nlm.nih.gov/18655432/}{PMID:18655432}.

\bibitem{safemgseries2010} Mackay JD, Bladon PT. Hypomagnesaemia due to proton-pump inhibitor therapy: a clinical case series. \textit{QJM}. 2010;103:387--395. \href{https://doi.org/10.1093/qjmed/hcq021}{doi:10.1093/qjmed/hcq021}.

\bibitem{safeb12trial1994} Marcuard SP, Albernaz L, Khazanie PG. Omeprazole therapy causes malabsorption of cyanocobalamin (vitamin B12). \textit{Ann Intern Med}. 1994;120:211--215. \href{https://doi.org/10.7326/0003-4819-120-3-199402010-00006}{doi:10.7326/0003-4819-120-3-199402010-00006}.

\bibitem{safedementia2023} Mehta RS, Kochar B, Zhou Z, et al. Association of Proton Pump Inhibitor Use With Incident Dementia and Cognitive Decline in Older Adults: A Prospective Cohort Study. \textit{Gastroenterology}. 2023;165:564--572.e1. \href{https://doi.org/10.1053/j.gastro.2023.05.052}{doi:10.1053/j.gastro.2023.05.052}.

\bibitem{withdrawhendricks2021} Hendricks E, Ajmeri AN, Singh MM, Mongalo M, Goebel LJ. A Randomized Open-Label Study of Two Methods of Proton Pump Inhibitors Discontinuation. \textit{Cureus}. 2021;13:e15022. \href{https://doi.org/10.7759/cureus.15022}{doi:10.7759/cureus.15022}.

\bibitem{withdrawhojo2026} Hojo M, Ueda K, Ueyama H, et al. Immediate discontinuation versus gradual tapering of proton pump inhibitors in long-term users: a randomized controlled trial. \textit{J Gastroenterol}. 2026;61:1083--1092; опубликовано онлайн 27.04.2026. \href{https://doi.org/10.1007/s00535-026-02419-z}{doi:10.1007/s00535-026-02419-z}.


\bibitem{withdrawbell} Bell N, Karol MD, Sachs G, et al. Duration of effect of lansoprazole on gastric pH and acid secretion in normal male volunteers. \textit{Aliment Pharmacol Ther}. 2001;15:105--113. \href{https://doi.org/10.1046/j.1365-2036.2001.00831.x}{doi:10.1046/j.1365-2036.2001.00831.x}.

\bibitem{withdrawniklasson} Niklasson A, Lindström L, Simrén M, Lindberg G, Björnsson E. Dyspeptic symptom development after discontinuation of a proton pump inhibitor: a double-blind placebo-controlled trial. \textit{Am J Gastroenterol}. 2010;105:1531--1537. \href{https://doi.org/10.1038/ajg.2010.81}{doi:10.1038/ajg.2010.81}. Изучался пантопразол, не лансопразол.

\bibitem{vision2025} Uemura N, Kinoshita Y, Haruma K, et al. Vonoprazan as a long-term maintenance treatment for erosive esophagitis: VISION, a 5-year, randomized, open-label study. \textit{Clin Gastroenterol Hepatol}. 2025;23:748--757.e5; онлайн 27.08.2024. \href{https://doi.org/10.1016/j.cgh.2024.08.004}{doi:10.1016/j.cgh.2024.08.004}; \href{https://pubmed.ncbi.nlm.nih.gov/39209187/}{PMID:39209187}.

\bibitem{seadepres2026} Leelakusolvong S, et al. Rationalizing proton pump inhibitor use: Southeast Asia consensus recommendations. \textit{J Gastroenterol Hepatol}. 2026. \href{https://doi.org/10.1111/jgh.70350}{doi:10.1111/jgh.70350}. Использованы положения о пересмотре лечения и алгоритм отмены.

\bibitem{priceusprev} Walmart. Prevacid 24HR, lansoprazole 15 mg, 42 capsules: онлайн-каталог и опубликованная цена. \href{https://www.walmart.com/browse/health-and-medicine/prevacid/976760_1396434_7324505_5124400}{Каталог Prevacid}. Дата обращения: 14.09.2026; цена зависит от витрины и обновления страницы.

\bibitem{priceusequate} Walmart. Equate Lansoprazole Delayed Release Capsules, 15 mg, 42 count. \href{https://www.walmart.com/ip/907145156}{Карточка Equate}. Дата обращения: 14.09.2026; использована доступная индексированная карточка с ценой 9,97 USD за 42 капсулы.

\bibitem{pricefrbio} Base de données publique des médicaments, France. LANSOPRAZOLE BIOGARAN 30 mg, comprimé orodispersible; présentation 28 comprimés, CIP 3400927871399. \href{https://base-donnees-publique.medicaments.gouv.fr/medicament/63062321/extrait}{Официальная карточка и цена}. Дата обращения: 14.09.2026.

\bibitem{pricejptake} Ministry of Health, Labour and Welfare, Japan. Перечень лекарственных цен, апрель 2026: Takepron OD 15 mg (2329023F1020), Takepron OD 30 mg (2329023F2027). \href{https://www.mhlw.go.jp/topics/2026/04/dl/tp20260401-01_01.pdf}{Перечень от 01.04.2026}; \href{https://www.mhlw.go.jp/topics/2026/04/dl/tp20260415-01_01.pdf}{Перечень от 15.04.2026}. Цены приведены за таблетку. Дата обращения: 14.09.2026.

\bibitem{priceinlanzol} Tata 1mg. Lanzol 30 Capsule (Cipla), 10 capsules; витрина Gurgaon, цена предложения и MRP. \href{https://www.1mg.com/drugs/lanzol-30-capsule-14655}{Карточка Lanzol 30}. Дата обращения: 14.09.2026; карточка содержит отметку обновления 14.09.2026. Дополнительные купоны не учитывались.

\bibitem{pricerulancid} Apteka.ru. Ланцид, капсулы 15 мг и 30 мг, по 30 шт.; московская витрина и варианты товара. \href{https://apteka.ru/product/lanczid-30-mg-30-sht-kapsuly-5e326aadca7bdc000192dfde/}{Карточка Ланцида}. Дата обращения: 14.09.2026; цены предложения, не сведения о государственной предельной цене.


\bibitem{cardioacc2025} Rao SV, et al. 2025 ACC/AHA/ACEP/NAEMSP/SCAI Guideline for the Management of Patients With Acute Coronary Syndromes. \textit{J Am Coll Cardiol}. 2025. Раздел 11.1: гастропротекция при высоком риске кровотечения. \href{https://doi.org/10.1016/j.jacc.2024.11.009}{doi:10.1016/j.jacc.2024.11.009}.

\bibitem{renalbigan2025} González-Pérez A, Martínez-Domínguez SJ, et al. Proton Pump Inhibitor Use and Worsening Kidney Function: A Retrospective Cohort Study Including 122,606 Acid-Suppressing Users. \textit{J Gen Intern Med}. 2025;40:818--827; онлайн 03.12.2024. \href{https://doi.org/10.1007/s11606-024-09213-8}{doi:10.1007/s11606-024-09213-8}. Таблица 2, основной анализ во время лечения.

\bibitem{renaltrinetx2025} Chen TP, Wu MJ, Chen CH, Tsai SF. Bone health, renal outcomes, and iron deficiency anaemia in proton pump inhibitor versus histamine-2 receptor antagonist users: a retrospective cohort study based on TriNetX global collaborative network data. \textit{BMJ Open Gastroenterol}. 2025;12:e001723. \href{https://doi.org/10.1136/bmjgast-2024-001723}{doi:10.1136/bmjgast-2024-001723}.

\bibitem{renalfaers2026} Wang Z, Jin J, Wang G, Zhou G, Jiang H. Systematic analysis of proton pump inhibitors-related adverse reactions using the FDA adverse event reporting system database. \textit{PLoS One}. 2026;21(2):e0340684. \href{https://doi.org/10.1371/journal.pone.0340684}{doi:10.1371/journal.pone.0340684}. Таблица 2; ROR сообщений, не клиническая частота.

\bibitem{withdrawfarrell2017} Farrell B, Pottie K, Thompson W, et al. Deprescribing proton pump inhibitors: Evidence-based clinical practice guideline. \textit{Can Fam Physician}. 2017;63:354--364. \href{https://pmc.ncbi.nlm.nih.gov/articles/PMC5429051/}{Полный текст}; \href{https://pubmed.ncbi.nlm.nih.gov/28500192/}{PMID:28500192}. Дополнительный источник по вариантам снижения дозы и наблюдению; не документ AGA.

\bibitem{withdrawbrand1988} Brand SJ, Stone D. Reciprocal regulation of antral gastrin and somatostatin gene expression by omeprazole-induced achlorhydria. \textit{J Clin Invest}. 1988;82:1059--1066. \href{https://doi.org/10.1172/JCI113662}{doi:10.1172/JCI113662}. Эксперимент на крысах.

\bibitem{withdrawsheng2020} Sheng W, Malagola E, et al. Hypergastrinemia Expands Gastric ECL Cells Through CCK2R$^+$ Progenitor Cells via ERK Activation. \textit{Cell Mol Gastroenterol Hepatol}. 2020;10:434--449. \href{https://doi.org/10.1016/j.jcmgh.2020.04.008}{doi:10.1016/j.jcmgh.2020.04.008}. Мышиные и органоидные модели.

\bibitem{withdrawwaldum1996} Waldum HL, Arnestad JS, Brenna E, et al. Marked increase in gastric acid secretory capacity after omeprazole treatment. \textit{Gut}. 1996;39:649--653. \href{https://doi.org/10.1136/gut.39.5.649}{doi:10.1136/gut.39.5.649}; \href{https://pubmed.ncbi.nlm.nih.gov/9026477/}{PMID:9026477}.

\bibitem{withdrawboyce2017} Boyce M, van den Berg F, Mitchell T, et al. Randomised trial of the effect of a gastrin/CCK$_2$ receptor antagonist on esomeprazole-induced hypergastrinaemia: evidence against rebound hyperacidity. \textit{Eur J Clin Pharmacol}. 2017;73:129--139; онлайн 29.10.2016. \href{https://doi.org/10.1007/s00228-016-2150-x}{doi:10.1007/s00228-016-2150-x}; \href{https://pubmed.ncbi.nlm.nih.gov/27796466/}{PMID:27796466}.

\bibitem{phthoring1999} Thoring M, Hedenström H, Eriksson LS. Rapid effect of lansoprazole on intragastric pH: a crossover comparison with omeprazole. \textit{Scand J Gastroenterol}. 1999;34:341--345. \href{https://doi.org/10.1080/003655299750026335}{doi:10.1080/003655299750026335}; \href{https://pubmed.ncbi.nlm.nih.gov/10365892/}{PMID:10365892}.

\bibitem{phbell1996} Bell NJ, Hunt RH. Time to maximum effect of lansoprazole on gastric pH in normal male volunteers. \textit{Aliment Pharmacol Ther}. 1996;10:897--904. \href{https://doi.org/10.1046/j.1365-2036.1996.103242000.x}{doi:10.1046/j.1365-2036.1996.103242000.x}; \href{https://pubmed.ncbi.nlm.nih.gov/8971286/}{PMID:8971286}.

\bibitem{hpsaito2005} Saito M, Yasui-Furukori N, Uno T, et al. Effects of clarithromycin on lansoprazole pharmacokinetics between CYP2C19 genotypes. \textit{Br J Clin Pharmacol}. 2005;59:302--309. \href{https://doi.org/10.1111/j.1365-2125.2004.02329.x}{doi:10.1111/j.1365-2125.2004.02329.x}; \href{https://pubmed.ncbi.nlm.nih.gov/15752376/}{PMID:15752376}. Фармакокинетическое исследование здоровых добровольцев; не испытание эрадикации у UM.

\bibitem{renalkdigo2024} Kidney Disease: Improving Global Outcomes (KDIGO) CKD Work Group. KDIGO 2024 Clinical Practice Guideline for the Evaluation and Management of Chronic Kidney Disease. \textit{Kidney Int}. 2024;105(Suppl 4S):S117--S314. Разделы 1.2.2 и 2.1. \href{https://kdigo.org/wp-content/uploads/2024/03/KDIGO-2024-CKD-Guideline.pdf}{Официальный текст KDIGO}. Критерии наблюдения за ХБП, не молекулоспецифичная программа мониторинга ИПП.

\bibitem{cardiocpic2022} Lee CR, Luzum JA, Sangkuhl K, et al. Clinical Pharmacogenetics Implementation Consortium Guideline for CYP2C19 Genotype and Clopidogrel Therapy: 2022 Update. \textit{Clin Pharmacol Ther}. 2022. \href{https://doi.org/10.1002/cpt.2526}{doi:10.1002/cpt.2526}; \href{https://files.cpicpgx.org/data/guideline/publication/clopidogrel/2022/35034351.pdf}{Документ CPIC}. Таблица 2: антиагрегантная терапия при сердечно-сосудистых показаниях.

\bibitem{futurestereo1996} Katsuki H, Yagi H, Arimori K, et al. Determination of R(+)- and S(-)-lansoprazole using chiral stationary-phase liquid chromatography and their enantioselective pharmacokinetics in humans. \textit{Pharm Res}. 1996;13:611--615. \href{https://doi.org/10.1023/A:1016062508580}{doi:10.1023/A:1016062508580}; \href{https://pubmed.ncbi.nlm.nih.gov/8710755/}{PMID:8710755}.

\bibitem{futurenab1998} Peghini PL, Katz PO, Bracy NA, Castell DO. Nocturnal recovery of gastric acid secretion with twice-daily dosing of proton pump inhibitors. \textit{Am J Gastroenterol}. 1998;93:763--767. \href{https://doi.org/10.1111/j.1572-0241.1998.221_a.x}{doi:10.1111/j.1572-0241.1998.221\_a.x}; \href{https://pubmed.ncbi.nlm.nih.gov/9625124/}{PMID:9625124}.

\bibitem{futurecubam2005} He Q, Madsen M, Kilkenney A, et al. Amnionless function is required for cubilin brush-border expression and intrinsic factor-cobalamin (vitamin B12) absorption in vivo. \textit{Blood}. 2005;106:1447--1453. \href{https://doi.org/10.1182/blood-2005-03-1197}{doi:10.1182/blood-2005-03-1197}; \href{https://pmc.ncbi.nlm.nih.gov/articles/PMC1895201/}{Полный текст}.

\bibitem{methodema2010} European Medicines Agency. Public statement: Interaction between clopidogrel and proton-pump inhibitors. EMA/174948/2010, 17.03.2010. \href{https://www.ema.europa.eu/en/documents/public-statement/public-statement-interaction-between-clopidogrel-and-proton-pump-inhibitors_en.pdf}{Официальное заявление EMA}. Регуляторная основа разграничения взаимодействий отдельных ИПП с клопидогрелом.

\bibitem{hpstability1997} Erah PO, Goddard AF, Barrett DA, Shaw PN, Spiller RC. The stability of amoxycillin, clarithromycin and metronidazole in gastric juice: relevance to the treatment of Helicobacter pylori infection. \textit{J Antimicrob Chemother}. 1997;39:5--12. \href{https://doi.org/10.1093/jac/39.1.5}{doi:10.1093/jac/39.1.5}; \href{https://pubmed.ncbi.nlm.nih.gov/9044021/}{PMID:9044021}. Экспериментальная оценка химической стабильности, не испытание эрадикационной схемы с лансопразолом.

\bibitem{cardiorko2024} Аверков О.В., Арутюнян Г.К., Дупляков Д.В. и др. Острый коронарный синдром без подъема сегмента ST электрокардиограммы. Клинические рекомендации 2024. Российское кардиологическое общество. \textit{Российский кардиологический журнал}. 2025;30(5):6319. \href{https://doi.org/10.15829/1560-4071-2025-6319}{doi:10.15829/1560-4071-2025-6319}. Национальный контекст ведения ОКС; не источник доказанного превосходства лансопразола по MACE.

\bibitem{futurerab1999} Duman JG, Tyagarajan K, Kolsi MS, Moore HP, Forte JG. Expression of rab11a N124I in gastric parietal cells inhibits stimulatory recruitment of the H$^+$-K$^+$-ATPase. \textit{Am J Physiol}. 1999;277:C361--C372. \href{https://doi.org/10.1152/ajpcell.1999.277.3.C361}{doi:10.1152/ajpcell.1999.277.3.C361}; \href{https://pubmed.ncbi.nlm.nih.gov/10484323/}{PMID:10484323}. Исследование мембранного транспорта, не доказательство его нарушения лансопразолом.

\bibitem{futurebone2015} Mishima K, Kitoh H, Ohkawara B, et al. Lansoprazole Upregulates Polyubiquitination of the TNF Receptor-Associated Factor 6 and Facilitates Runx2-mediated Osteoblastogenesis. \textit{EBioMedicine}. 2015;2:2046--2061. \href{https://doi.org/10.1016/j.ebiom.2015.11.024}{doi:10.1016/j.ebiom.2015.11.024}; \href{https://pubmed.ncbi.nlm.nih.gov/26844285/}{PMID:26844285}. Доклинические данные, включая концентрационно-зависимые эффекты на остеокласты.

\bibitem{futureddah2013} Ghebremariam YT, LePendu P, Lee JC, et al. Unexpected effect of proton pump inhibitors: elevation of the cardiovascular risk factor asymmetric dimethylarginine. \textit{Circulation}. 2013;128:845--853. \href{https://doi.org/10.1161/CIRCULATIONAHA.113.003602}{doi:10.1161/CIRCULATIONAHA.113.003602}; \href{https://pubmed.ncbi.nlm.nih.gov/23825361/}{PMID:23825361}. Экспериментальная гипотеза DDAH--ADMA--NO, не доказательство клинического риска конкретного ИПП.

\bibitem{futuremtx2009} Suzuki K, Doki K, Homma M, et al. Co-administration of proton pump inhibitors delays elimination of plasma methotrexate in high-dose methotrexate therapy. \textit{Br J Clin Pharmacol}. 2009;67:44--49. \href{https://doi.org/10.1111/j.1365-2125.2008.03303.x}{doi:10.1111/j.1365-2125.2008.03303.x}; \href{https://pubmed.ncbi.nlm.nih.gov/19076159/}{PMID:19076159}. Клинический анализ и везикулярное исследование BCRP; высокодозные онкологические режимы.

\bibitem{futureil62011} Dickmann LJ, Patel SK, Rock DA, Wienkers LC, Slatter JG. Effects of interleukin-6 (IL-6) and an anti-IL-6 monoclonal antibody on drug-metabolizing enzymes in human hepatocyte culture. \textit{Drug Metab Dispos}. 2011;39:1415--1422. \href{https://doi.org/10.1124/dmd.111.038679}{doi:10.1124/dmd.111.038679}; \href{https://pubmed.ncbi.nlm.nih.gov/21555507/}{PMID:21555507}. Культура человеческих гепатоцитов; не клиническая верификация феноконверсии UM при лансопразоле.

\bibitem{renalmarkers2019} Moledina DG, Wilson FP, Pober JS, et al. Urine TNF-$\alpha$ and IL-9 for clinical diagnosis of acute interstitial nephritis. \textit{JCI Insight}. 2019;4(10):e127456. \href{https://doi.org/10.1172/jci.insight.127456}{doi:10.1172/jci.insight.127456}. Биопсийно ориентированное исследование диагностических маркеров ОИН; не молекулоспецифичный тест лансопразола.

\end{thebibliography}
\end{document}
